\documentclass[9pt]{extarticle}

\usepackage[a3paper,landscape,margin=12mm]{geometry}
\usepackage[T1]{fontenc}
\usepackage{lmodern}
\usepackage{microtype}
\usepackage{amsmath,amssymb,amsthm,mathtools}
\usepackage{booktabs,array}
\usepackage{enumitem}
\usepackage[hidelinks,hypertexnames=false]{hyperref}
\usepackage[nameinlink,noabbrev]{cleveref}

\setlength{\parindent}{0pt}
\setlength{\parskip}{5pt}
\setlength{\emergencystretch}{2em}
\allowdisplaybreaks
\setlist{itemsep=2pt,topsep=4pt}

% A multiline monospaced proof term that remains a mathematical atom.
% This is used only where an ordinary \texttt block would become one
% unbreakable line extending beyond the page.
\newenvironment{mathleanproof}[1]{%
  \vcenter\bgroup
  \hbox\bgroup
  \begin{minipage}{#1}%
  \ttfamily\raggedright\parindent=0pt\parskip=0pt\obeylines\obeyspaces
}{%
  \end{minipage}%
  \egroup
  \egroup
}

\newtheorem{lemma}{Lemma}[section]
\newtheorem{proposition}[lemma]{Proposition}
\newtheorem{theorem}[lemma]{Theorem}
\newtheorem{corollary}[lemma]{Corollary}
\theoremstyle{definition}
\newtheorem{definition}[lemma]{Definition}
\theoremstyle{remark}
\newtheorem{remark}[lemma]{Remark}

\newcommand{\R}{\mathbb R}
\newcommand{\N}{\mathbb N}
\newcommand{\supp}{\operatorname{supp}}
\newcommand{\dd}{\,\mathrm d}
\newcommand{\one}{\mathbf 1}
\newcommand{\col}{\operatorname{col}}
\newcommand{\Dbulk}{\mathcal D_{\mathrm{bulk}}}
\newcommand{\Dminus}{\mathcal D_{-\infty}}
\newcommand{\A}{\mathcal A}
\newcommand{\eps}{\varepsilon}
\newcommand{\norm}[1]{\left\lVert #1\right\rVert}
\newcommand{\abs}[1]{\left\lvert #1\right\rvert}
\newcommand{\Var}{\operatorname{Var}}
\newcommand{\card}{\operatorname{card}}

\title{Lean Formalization Blueprint for Sections 4 and 5\\
of the Classical Conditional R\'enyi-Entropy Classification\\
\large Fully dependency-closed mathematical specification}
\author{}
\date{20 August 2026}

\begin{document}
\maketitle

\begin{abstract}
This document is a fully dependency-closed internal specification for the Lean
formalization accompanying the complete-proof Sections~4--5 paper.  It proves the exact parameter
ranges for which conditional entropy functionals obtained as exponential
averages of R\'enyi entropies are monotone under conditionally mixing
channels.  No definition or theorem from the preceding sections of the paper
is assumed.  Before the mathematical proof, Sections~0--2 fix total Lean-ready
data types and state every local finite-sum, support, normalization, compactified
order, measure, R\'enyi, convexity, channel, limit, and differentiation lemma
used later.  General results supplied by Lean and mathlib are isolated in an
explicit import boundary.  The parameter measure is an arbitrary Borel
probability measure on $[0,+\infty]$ and may accumulate at the Shannon point.
The singular weighted moment is defined through nonnegative one-sided
integrals with a total guarded real-valued version, zero conditioning columns
are treated explicitly, and every
limiting or differentiation argument is accompanied by a domination or
uniformity statement.  The predicates for stochastic conversion,
monotonicity, embeddings, curvature, and admissible parameters are named
explicitly rather than left in prose.  Sections~3--5 then give the setup,
sufficiency proof, and necessity proof in theorem-dependency order.  The
classification theorem itself is declared only after all of its dependencies,
so the displayed order can be copied into Lean without introducing an axiom.
\end{abstract}

\tableofcontents

The target Lean module begins with \texttt{noncomputable section}.  This is
required by the canonical Jordan decomposition, total suprema, and Bochner
integrals used below; it introduces no axiom.
Immediately after it, declare \texttt{universe u v}; these are the two
universe levels used by the polymorphic carriers and uniform-limit lemmas.

\setcounter{section}{-1}
\section{Formalization contract and total data model}
\label{sec:formalization-contract}

\subsection{Logical boundary}

The intended Lean development may import the ordinary foundations already
available in Lean and mathlib: classical logic, finite types and finite sums,
the real and nonnegative real numbers, elementary topology, derivatives of
real elementary functions, finite-dimensional Euclidean spaces, Borel
measures, finite and probability measures, signed measures or Jordan
decompositions, pushforwards, monotone convergence, dominated convergence,
and differentiation under a finite measure.  These foundational results are
not claims of the paper.

No result from another paper section may be imported.  In
particular, the formalization must not assume the semiring classification,
the earlier characterization of homomorphisms or derivations, the appendix on
convexity, or an unproved statement that a candidate functional is monotone.
All paper-specific facts used in Sections~4 and~5 are stated below.

One possible list of Lean imports is
\begin{verbatim}
Mathlib.Analysis.SpecialFunctions.Pow.Real
Mathlib.Analysis.SpecialFunctions.Log.Basic
Mathlib.Analysis.Calculus.Deriv.Basic
Mathlib.Analysis.Convex.SpecificFunctions.Basic
Mathlib.Data.ENNReal.Basic
Mathlib.MeasureTheory.Integral.Bochner.Basic
Mathlib.MeasureTheory.Integral.DominatedConvergence
Mathlib.MeasureTheory.Measure.Support
Mathlib.MeasureTheory.Measure.WithDensityFinite
Mathlib.MeasureTheory.VectorMeasure.Decomposition.Jordan
Mathlib.Topology.Instances.ENNReal
Mathlib.Tactic.FinCases
Mathlib.Tactic.NormNum
\end{verbatim}
Immediately after the imports, use
\begin{verbatim}
open MeasureTheory Filter
open scoped BigOperators MeasureTheory ENNReal NNReal Topology
\end{verbatim}
so that the measure bundles, integrals, finite sums, filters, and
neighbourhood notation in the signatures below resolve without additional
paper-level aliases.
The exact import names may change with mathlib.  The mathematical interface
specified here is independent of those names.

Every object below is total.  Expressions such as $x/\norm{x}_1$ are used
only under a proof that $x\neq0$.  A conditional distribution is never formed
from a zero column.  A minimum or maximum over conditioning outcomes is taken
only after proving that the finite set of positive-mass outcomes is nonempty.

\subsection{Recommended carriers and coercions}

For a finite alphabet of size $d\geq1$, a Lean implementation may use
$\operatorname{Fin}(d)$.  More generally, every alphabet in this document is
a finite nonempty type.  In Lean this means that the declarations carry both
\texttt{Fintype} and \texttt{Nonempty} instances.  The block constructions
use $\operatorname{Fin}(d)$ only after a proof of $1\leq d$.  A real vector is
a function $x:I\to\R$.

Define
\begin{align*}
 \operatorname{Nonneg}(x)&:\Longleftrightarrow
       \forall i:I,\ 0\leq x_i,\\
 \operatorname{l1Mass}(x)&:=\sum_{i:I}x(i),
 &\norm{x}_1&:=\operatorname{l1Mass}(x),\\
 \supp(x)&:=\operatorname{Finset.univ}.\operatorname{filter}
       (i\mapsto x_i\neq0),\\
 \operatorname{PosMass}(x)&:\Longleftrightarrow
       \operatorname{Nonneg}(x)\ \wedge\ 0<\norm{x}_1.
\end{align*}
First define the closed additive submonoid
\[
 \operatorname{NonnegAddSubmonoid}(I):=
 \left\{
 \begin{array}{l}
  \operatorname{carrier}:=\{x:I\to\R\mid\operatorname{Nonneg}(x)\},\\
  \operatorname{zero\_mem}':=\texttt{by intro i; exact le\_rfl},\\
  \operatorname{add\_mem}':=\texttt{by intro x hx z hz i;
    exact add\_nonneg (hx i) (hz i)}
 \end{array}\right\}
 \in\operatorname{AddSubmonoid}(I\to\R).
\]
Fix the following subtype carriers (rather than leaving a choice of
representations to the implementation):
\begin{align*}
 \operatorname{ConeVec}(I)
 &:=\operatorname{NonnegAddSubmonoid}(I),\\
 \operatorname{PosConeVec}(I)
 &:=\{x:\operatorname{ConeVec}(I)\mathbin{//}x.1\neq(0:I\to\R)\},\\
 \operatorname{ProbVec}(I)
 &:=\{p:I\to\R\mathbin{//}
       \operatorname{Nonneg}(p)\ \wedge\ \operatorname{l1Mass}(p)=1\}.
\end{align*}
These are Lean subtypes, with their usual coercions to functions.  Equality
of probability vectors is extensional equality of the underlying functions.
Whenever $\log\card\supp(x)$ is written, the natural cardinal is coerced to
$\R$ before applying $\log$.

The cone operations are proof-carrying constructors, defined immediately
with the carriers:
\begin{align*}
 \operatorname{coneZero}_I
 &:=(0:\operatorname{ConeVec}(I)),\\
 \operatorname{coneAdd}(x,z)
 &:=x+z,\\
 \operatorname{coneScale}(c,h_c,x)
 &:=\left\langle(i\mapsto c\,x.1(i)),
      \texttt{by intro i; exact mul\_nonneg h\_c (x.2 i)}
      \right\rangle,\\
 \operatorname{coneMix}(\lambda,h_\lambda,x,z)
 &:=\left\langle
       (i\mapsto\lambda x.1(i)+(1-\lambda)z.1(i)),\
       \texttt{by intro i; exact add\_nonneg
        (mul\_nonneg h\_lambda.1 (x.2 i))
        (mul\_nonneg (sub\_nonneg.mpr h\_lambda.2) (z.2 i))}
       \right\rangle,
 \quad h_\lambda:0\leq\lambda\ \wedge\ \lambda\leq1.
\end{align*}
The subtype instance inherited from \texttt{NonnegAddSubmonoid} is exactly
\[
 \texttt{instance : AddCommMonoid (ConeVec I) := inferInstance}.
\]
Its zero and addition are $\operatorname{coneZero}_I$ and
$\operatorname{coneAdd}$.  This makes finite cone sums literal Lean sums;
no negation or real-module instance is installed.  For a finite type $J$,
define the exact total function
\[
 \operatorname{coneSum}(z):=\sum_{j:J}z(j),
 \qquad z:J\to\operatorname{ConeVec}(I).
\]
If
$c:J\to\R$, $h_c:\forall j,0\leq c_j$, and
$z:J\to\operatorname{ConeVec}(I)$, define
\[
 \operatorname{weightedConeSum}(c,h_c,z):=
 \operatorname{coneSum}\!\left(
  j\mapsto\operatorname{coneScale}(c(j),h_c(j),z(j))\right).
\]
For probability vectors define the separately bundled convex combination
\[
 \operatorname{mixProbVec}(\lambda,h_\lambda,p,q):=
 \left\langle
  (i\mapsto\lambda p.1(i)+(1-\lambda)q.1(i)),
  \left\langle
   \texttt{by intro i; exact add\_nonneg
    (mul\_nonneg h\_lambda.1 (p.2.1 i))
    (mul\_nonneg (sub\_nonneg.mpr h\_lambda.2) (q.2.1 i))},
   \texttt{by simp [l1Mass, p.2.2, q.2.2]}
  \right\rangle\right\rangle
 \in\operatorname{ProbVec}(I),
 \qquad h_\lambda:0\leq\lambda\ \wedge\ \lambda\leq1.
\]
This name is intentionally distinct from the later constructor
\texttt{probMix} for probability \emph{measures}.  For a function
$f:\operatorname{ProbVec}(I)\to\R$, define
\begin{align*}
 \operatorname{SimplexConcave}_I(f)
 &: \Longleftrightarrow
 \forall p,q:\operatorname{ProbVec}(I),\ \forall\lambda:\R,\
 \forall h_\lambda:0\leq\lambda\ \wedge\ \lambda\leq1,\\[-2pt]
 &\hspace{15mm}
 f(\operatorname{mixProbVec}(\lambda,h_\lambda,p,q))
 \geq\lambda f(p)+(1-\lambda)f(q),\\
 \operatorname{SimplexConvex}_I(f)
 &: \Longleftrightarrow
 \forall p,q:\operatorname{ProbVec}(I),\ \forall\lambda:\R,\
 \forall h_\lambda:0\leq\lambda\ \wedge\ \lambda\leq1,\\[-2pt]
 &\hspace{15mm}
 f(\operatorname{mixProbVec}(\lambda,h_\lambda,p,q))
 \leq\lambda f(p)+(1-\lambda)f(q).
\end{align*}
For $r:\R$ with $h_0:0\leq r$ and $h_1:r\leq1$, define
\[
 \operatorname{binaryProb}(r,h_0,h_1):=
 \left\langle
   \bigl(i\mapsto\operatorname{Fin.cases}(r,(\_\mapsto1-r),i)\bigr),
   \left\langle
    \texttt{by
      intro i
      exact Fin.cases h\_0
        (fun \_ => sub\_nonneg.mpr h\_1) i},
    \texttt{by simp [l1Mass]}
   \right\rangle\right\rangle
 \in\operatorname{ProbVec}(\operatorname{Fin}2).
\]
For a finite type $I$, define the following three separately typed
constructors:
\begin{align*}
 \operatorname{posScale}(c,h_c,x)&:\operatorname{PosConeVec}(I),
 &c&:\R,&h_c&:0<c,&x&:\operatorname{PosConeVec}(I),\\[-2pt]
 \operatorname{posScale}(c,h_c,x)
 &:=\left\langle
   \operatorname{coneScale}(c,\operatorname{le_of_lt}h_c,x.1),
   \begin{mathleanproof}{0.46\linewidth}
by
  intro hs
  apply x.2
  funext i
  have hi : c * x.1.1 i = 0 := by
    have hcoord :=
      congrArg (fun y : ConeVec I => y.1 i) hs
    simpa [coneScale] using hcoord
  exact (mul\_eq\_zero.mp hi).resolve\_left (ne\_of\_gt h\_c)
   \end{mathleanproof}
   \right\rangle,\\
 \operatorname{posMix}(\lambda,h_\lambda,x,z)&:\operatorname{PosConeVec}(I),
 &\lambda&:\R,&h_\lambda&:0<\lambda\ \wedge\ \lambda<1,
 &x,z&:\operatorname{PosConeVec}(I),\\[-2pt]
 \operatorname{posMix}(\lambda,h_\lambda,x,z)
 &:=\left\langle
   \operatorname{coneMix}(\lambda,
      \langle\operatorname{le_of_lt}h_\lambda.1,
             \operatorname{le_of_lt}h_\lambda.2\rangle,x.1,z.1),
   \begin{mathleanproof}{0.46\linewidth}
by
  intro hs
  apply x.2
  funext i
  have hi :
      lambda * x.1.1 i + (1 - lambda) * z.1.1 i = 0 := by
    have hcoord :=
      congrArg (fun y : ConeVec I => y.1 i) hs
    simpa [coneMix] using hcoord
  have hz : 0 <= (1 - lambda) * z.1.1 i :=
    mul\_nonneg
      (sub\_nonneg.mpr (le\_of\_lt h\_lambda.2)) (z.1.2 i)
  have hxle : lambda * x.1.1 i <= 0 := by
    calc
      lambda * x.1.1 i <=
          lambda * x.1.1 i + (1 - lambda) * z.1.1 i :=
        le\_add\_of\_nonneg\_right hz
      \_ = 0 := hi
  have hx0 : lambda * x.1.1 i = 0 :=
    le\_antisymm hxle
      (mul\_nonneg (le\_of\_lt h\_lambda.1) (x.1.2 i))
  exact (mul\_eq\_zero.mp hx0).resolve\_left
    (ne\_of\_gt h\_lambda.1)
   \end{mathleanproof}
   \right\rangle .
\end{align*}
Also define
\[
 \operatorname{toPosCone}(x,h):=
 \left\langle x,\texttt{by
   intro h0
   apply h
   exact Subtype.ext h0}\right\rangle
 \in\operatorname{PosConeVec}(I),
 \qquad x:\operatorname{ConeVec}(I),\quad h:x\neq0,
\]
where $h_0:x.1=(0:I\to\R)$.  Thus subtype inequality against the installed
cone zero is converted explicitly into the underlying-function inequality
stored by \texttt{PosConeVec}.  No additive-group or real-module instance is
asserted for the cone subtype; nonnegative scaling and convex mixing always
use the proof-carrying constructors above.

For a nonnegative vector of positive mass, define the normalization
\begin{equation}
 \widehat x_i:=x_i/\norm{x}_1.
 \label{eq:foundation-normalize}
\end{equation}
This fixes one total Lean function
\[
 \operatorname{normalize}:
 \forall\{I:\texttt{Type u}\},
 [\operatorname{Fintype}I][\operatorname{Nonempty}I],\
 \operatorname{PosConeVec}(I)\longrightarrow\operatorname{ProbVec}(I).
\]
Its implementation first proves that a finite sum of nonnegative
coordinates of a nonzero vector is positive and then divides by that sum.
Whenever a later statement has $x:\operatorname{ConeVec}(I)$ and
$h:x\neq0$, the notation $\widehat x$ means exactly
\texttt{normalize (toPosCone x h)}; the proof argument is not left for
elaboration to infer.  It is never implemented by
unguarded division followed by the convention $0^{-1}=0$.

For finite nonempty types $I,J$, define the tensor product and the uniform
probability vector by
\[
 (p\otimes q)_{(i,j)}:=p_iq_j,
 \qquad (u_I)_i:=\frac1{\card I}.
\]
The exact bundled names are
\[
 \operatorname{probTensor}:
 \operatorname{ProbVec}(I)\to\operatorname{ProbVec}(J)
 \to\operatorname{ProbVec}(I\times J),
 \qquad
 \operatorname{uniformProb}_I:\operatorname{ProbVec}(I).
\]
If $e:I\hookrightarrow J$ is an injection, define the zero extension
$\operatorname{zeroExtendRaw}(e,x):J\to\R$ by
\[
 (e_*x)_j:=
 \begin{cases}
  x_i,&\text{if }j=e(i)\text{ for the necessarily unique }i,\\
  0,&\text{if }j\notin\operatorname{range}(e).
 \end{cases}
\]
The definition is total in Lean by a decidable membership test on the finite
range and the unique preimage supplied by injectivity.  A bijective relabeling
is the special case in which $e$ is an equivalence.
Package this raw map into
\begin{align*}
 \operatorname{zeroExtendCone}(e)&:
  \operatorname{ConeVec}(I)\to\operatorname{ConeVec}(J),\\
 \operatorname{zeroExtendProb}(e)&:
  \operatorname{ProbVec}(I)\to\operatorname{ProbVec}(J),
\end{align*}
using the coordinatewise nonnegativity and mass-preservation proofs.  Paper
notation $e_*x$ denotes the appropriate bundled function according to the
statically known input type.

For a finite type $I$, let $\one_I:I\to\R$ be the constant-one vector.  If
$J$ is finite, $n:J\to\N$ has $1\leq n_j$, and
$z_j:\operatorname{Fin}(n_j)\to\R$, define the vertical concatenation
\[
 \operatorname{concat}\{J\}(n,z):
 \bigl(\mathop{\Sigma}_{j:J}\operatorname{Fin}(n(j))\bigr)\to\R,
 \qquad \operatorname{concat}(n,z)\langle j,r\rangle:=z(j)(r).
\]
The exact dependent-function input is
$z:(j:J)\to\operatorname{Fin}(n(j))\to\R$; the lower bound on $n(j)$ is
needed only when a nonempty instance for the sigma type is requested, not by
the definition of \texttt{concat} itself.
The later notation $\col(z^{(1)},\ldots,z^{(m)})$ is this function on the
finite sigma type; $\one_N$ abbreviates $\one_{\operatorname{Fin}(N)}$.

For every finite nonempty type $J$ and $a:J\to\R$, define total finite
extrema
\[
 \operatorname{finMin}(a):=
 \operatorname{Finset.univ}.\operatorname{inf}'
   (\operatorname{Finset.univ\_nonempty})(a),
 \qquad
 \operatorname{finMax}(a):=
 \operatorname{Finset.univ}.\operatorname{sup}'
   (\operatorname{Finset.univ\_nonempty})(a).
\]
The notation $\min_{j:J}a_j$ and $\max_{j:J}a_j$ abbreviates these named
functions.  Minimizer and maximizer lemmas return witnesses in $J$ together
with the corresponding equality.

\begin{lemma}[A zero extension of a nonzero cone vector is nonzero]
\label{lem:fnd-zero-extend-nonzero}
Let $I,J$ be finite types, let $e:I\hookrightarrow J$, let
$x:\operatorname{ConeVec}(I)$, and let $h_x:x\neq0$.  The named conclusion
is
\[
 \operatorname{zeroExtend\_ne\_zero}(e,h_x):
 \operatorname{zeroExtendCone}(e,x)\neq0.
\]
\end{lemma}

\begin{proof}
If the extension were zero, evaluation at every $e(i)$ and injectivity of
$e$ would show $x(i)=0$ for every $i$, contradicting $h_x$.
\end{proof}

\begin{lemma}[Tensor and zero-extension package]
\label{lem:fnd-tensor-extension-package}
Let $I,J$ be finite nonempty types, let $x:\operatorname{ConeVec}(I)$, and
let $e:I\hookrightarrow J$ be an injection.  Then the following literal
conjunction holds:
\begin{align*}
 &\norm{\operatorname{zeroExtendCone}(e,x)}_1=\norm{x}_1\\
 &\quad\wedge\quad
 \supp(\operatorname{zeroExtendCone}(e,x))
 =\operatorname{Finset.map}\langle e,e.\operatorname{injective}\rangle
    (\supp(x))\\
 &\quad\wedge\quad
 \card\supp(\operatorname{zeroExtendCone}(e,x))=\card\supp(x)\\
 &\quad\wedge\quad
 \Bigl(\forall h_x:x\neq0,\quad
 \operatorname{normalize}\!\left(
   \operatorname{toPosCone}
    (\operatorname{zeroExtendCone}(e,x),
     \operatorname{zeroExtend\_ne\_zero}(e,h_x))\right)
 =\operatorname{zeroExtendProb}\!\left(e,
   \operatorname{normalize}(\operatorname{toPosCone}(x,h_x))\right)
 \Bigr).
\end{align*}
where the final equality is between bundled probability vectors.
\end{lemma}

\begin{proof}
All sums over $J$ split into the range of $e$ and its complement.  The
complement contributes zero, and the sum over the range is reindexed by the
bijection from $I$ to $\operatorname{range}(e)$.  This proves the mass and
support assertions.  The normalization identity follows coordinatewise from
mass preservation.
\end{proof}

A joint nonnegative matrix on finite types $X,Y$ is a function
$P:X\to Y\to\R$ with nonnegative entries.  Fix the subtype
\[
 \operatorname{JointProb}(X,Y):=
 \left\{P:X\to Y\to\R\ \mathbin{//}\
 (\forall x:X,\forall y:Y,\ 0\leq P(x,y))
 \ \wedge\ \sum_{x:X}\sum_{y:Y}P(x,y)=1\right\}.
\]
For finite nonempty $X,Y$, introduce the following declarations immediately
after the carrier:
\begin{align*}
 \operatorname{column}(P,y)&:\operatorname{ConeVec}(X),
 &\operatorname{column}(P,y)(x)&:=P(x,y),\\
 \operatorname{colMass}(P,y)&:\R,
 &\operatorname{colMass}(P,y)&:=\norm{\operatorname{column}(P,y)}_1,\\
 \operatorname{activeFinset}(P)&:\operatorname{Finset}(Y),
 &\operatorname{activeFinset}(P)&:=
   \operatorname{Finset.univ}.\operatorname{filter}
      (y\mapsto0<\operatorname{colMass}(P,y)),\\
 \operatorname{Active}(P)&:=
   \{y:Y\mathbin{//}y\in\operatorname{activeFinset}(P)\}.
\end{align*}
The bundled column uses the entrywise nonnegativity field of $P$.  For
$y:\operatorname{Active}(P)$, define
\begin{equation}
 \operatorname{conditional}(P,y):=
 \operatorname{normalize}
 \bigl(\operatorname{toPosCone}(\operatorname{column}(P,y),h_y)\bigr)
 \in\operatorname{ProbVec}(X),
 \label{eq:foundation-columns}
\end{equation}
where $h_y:\operatorname{column}(P,y)\neq0$ follows from the strict
column-mass inequality stored in $y$.  Define also
\begin{equation}
 Y_+(P):=\operatorname{activeFinset}(P).
 \label{eq:positive-column-finset}
\end{equation}
The paper notation $x_y$, $m_y$, and $P_{X\mid y}$ abbreviates respectively
\texttt{column P y}, \texttt{colMass P y}, and
\texttt{conditional P y}.  The subsequently declared theorem
\texttt{positiveColumnsNonempty} supplies a
\texttt{Nonempty (Active P)} instance locally wherever an extremum is formed.

For injections $e_X:X\hookrightarrow\widetilde X$ and
$e_Y:Y\hookrightarrow\widetilde Y$, define the proof-carrying joint zero
extension
\[
 \operatorname{jointZeroExtend}(e_X,e_Y):
 \operatorname{JointProb}(X,Y)\to
 \operatorname{JointProb}(\widetilde X,\widetilde Y)
\]
by putting $P(x,y)$ at $(e_X(x),e_Y(y))$ and zero elsewhere.  Define also
\begin{align*}
 \operatorname{jointTensor}&:
 \operatorname{JointProb}(X,Y)\to\operatorname{JointProb}(X',Y')
 \to\operatorname{JointProb}(X\times X',Y\times Y'),\\
 \operatorname{independentJoint}&:
 \operatorname{ProbVec}(X)\to\operatorname{ProbVec}(Y)
 \to\operatorname{JointProb}(X,Y),
\end{align*}
by
\[
 \operatorname{jointTensor}(P,Q)((x,x'),(y,y')):=P(x,y)Q(x',y'),
 \qquad
 \operatorname{independentJoint}(p,q)(x,y):=p(x)q(y).
\]
The notation $E_{e_X,e_Y}P$ and $P\otimes Q$ always denotes these bundled
constructors, never a raw function whose probability proof is inferred.

\begin{definition}[Active outcomes under zero extension]
\label{def:fnd-active-zero-extension}
Let $X,Y,\widetilde X,\widetilde Y$ be finite nonempty types, let
$P:\operatorname{JointProb}(X,Y)$, and let
$e_X:X\hookrightarrow\widetilde X$ and
$e_Y:Y\hookrightarrow\widetilde Y$ be injections.  Define
\[
 \operatorname{activeZeroExtendEquiv}(e_X,e_Y,P):
   \operatorname{Active}(P)\simeq
   \operatorname{Active}(\operatorname{jointZeroExtend}(e_X,e_Y,P)).
\]
Its forward map sends $y$ to $e_Y(y)$; its inverse uses the unique preimage
forced by positivity of an extended column.  The inverse laws are proved by
subtype extensionality.
\end{definition}

\begin{definition}[Active outcomes under tensor products]
\label{def:fnd-active-tensor}
Let $X,Y,X',Y'$ be finite nonempty types, let
$P:\operatorname{JointProb}(X,Y)$, and let
$Q:\operatorname{JointProb}(X',Y')$.  Define
\[
 \operatorname{activeTensorEquiv}(P,Q):
   \operatorname{Active}(\operatorname{jointTensor}(P,Q))\simeq
   \operatorname{Active}(P)\times\operatorname{Active}(Q).
\]
The forward map separates an active pair into its two active coordinates;
the inverse is subtype pairing.  Positivity of a product of nonnegative
column masses is equivalent to positivity of each factor.
\end{definition}

\begin{lemma}[Independent column mass]
\label{lem:fnd-independent-column-mass}
For finite nonempty $X,Y$, $p:\operatorname{ProbVec}(X)$,
$q:\operatorname{ProbVec}(Y)$, and $y:Y$,
\[
 \operatorname{colMass}(\operatorname{independentJoint}(p,q),y)=q(y).
\]
\end{lemma}

\begin{proof}
Factor $q(y)$ out of the finite sum and use $\sum_xp(x)=1$.
\end{proof}

\begin{definition}[Active independent outcome]
\label{def:fnd-independent-active}
For finite nonempty $X,Y$, $p:\operatorname{ProbVec}(X)$,
$q:\operatorname{ProbVec}(Y)$, $y:Y$, and $h_y:0<q(y)$, define
\[
 \operatorname{independentActive}(p,q,y,h_y):
 \operatorname{Active}(\operatorname{independentJoint}(p,q))
\]
by bundling $y$ with the membership proof obtained from
\cref{lem:fnd-independent-column-mass} and $h_y$.
\end{definition}

\begin{lemma}[Joint extension and tensor package]
\label{lem:fnd-joint-extension-package}
Let $X,Y,X',Y',\widetilde X,\widetilde Y$ be finite nonempty types, let
$P:\operatorname{JointProb}(X,Y)$,
$Q:\operatorname{JointProb}(X',Y')$,
$p:\operatorname{ProbVec}(X)$, and $q:\operatorname{ProbVec}(Y)$, and let
$e_X:X\hookrightarrow\widetilde X$ and
$e_Y:Y\hookrightarrow\widetilde Y$ be injections.  Then the following
literal conjunction holds:
\begin{align*}
 &\Bigl(\forall y:\operatorname{Active}(P),\
  \texttt{let }\widetilde y:=
   \operatorname{activeZeroExtendEquiv}(e_X,e_Y,P)(y);\\[-2pt]
 &\qquad
 \operatorname{colMass}
   (\operatorname{jointZeroExtend}(e_X,e_Y,P),\widetilde y.1)
 =\operatorname{colMass}(P,y.1)\ \wedge\\[-2pt]
 &\qquad
 \operatorname{conditional}
   (\operatorname{jointZeroExtend}(e_X,e_Y,P),\widetilde y)
 &=\operatorname{zeroExtendProb}
   (e_X,\operatorname{conditional}(P,y))\Bigr)\\
 &\quad\wedge\Bigl(
 \forall w:\operatorname{Active}(\operatorname{jointTensor}(P,Q)),\
 \texttt{let }r:=\operatorname{activeTensorEquiv}(P,Q)(w);\\[-2pt]
 &\qquad
 \operatorname{colMass}(\operatorname{jointTensor}(P,Q),w.1)
 =\operatorname{colMass}(P,r.1.1)\operatorname{colMass}(Q,r.2.1)
 \ \wedge\\[-2pt]
 &\qquad
 \operatorname{conditional}(\operatorname{jointTensor}(P,Q),w)
 &=\operatorname{probTensor}
   (\operatorname{conditional}(P,r.1),
    \operatorname{conditional}(Q,r.2))\Bigr)\\
 &\quad\wedge\Bigl(\forall y:Y,\
  \operatorname{colMass}(\operatorname{independentJoint}(p,q),y)=q(y)\Bigr)\\
 &\quad\wedge\Bigl(\forall(y:Y)(h_y:0<q(y)),\
  \operatorname{conditional}
   (\operatorname{independentJoint}(p,q),
    \operatorname{independentActive}(p,q,y,h_y))=p\Bigr).
\end{align*}
\end{lemma}

\begin{proof}
For zero extension, reindex the two finite sums as in
\cref{lem:fnd-tensor-extension-package}.  A column outside the range of
$e_Y$ is zero.  A column in the range has the same mass as its unique
preimage and is the row-zero-extension of that preimage column; normalization
commutes with the extension by
\cref{lem:fnd-tensor-extension-package}.  For tensor products, distribute the
finite double sum.  The column mass is $m_yn_{y'}$, which is positive exactly
when both factors are positive, and division by that product gives the tensor
of the two normalized columns.
\end{proof}

\begin{lemma}[Nonzero nonnegative mass]
\label{lem:fnd-nonzero-mass}
Let $I:\texttt{Type u}$ carry $[\operatorname{Fintype}I]$, let
$x:I\to\R$, and let $h_x:\forall i:I,0\leq x(i)$.  Then the literal
conjunction
\[
 \left(\sum_{i:I}x(i)=0\Longleftrightarrow x=0\right)
 \quad\wedge\quad
 \left(x\neq0\Longleftrightarrow0<\sum_{i:I}x(i)\right)
\]
holds.
\end{lemma}

\begin{proof}
If every coordinate is zero, the displayed finite sum is zero.  Conversely,
if the finite sum of nonnegative terms is zero, then for each index $i$ one
has $0\leq x_i\leq\sum_jx_j=0$, and antisymmetry gives $x_i=0$.
\end{proof}

\begin{corollary}[Nonzero cone vector has positive mass]
\label{cor:fnd-cone-nonzero-mass}
Let $I$ be a finite type and let $x:\operatorname{ConeVec}(I)$.  Then
\[
 x\neq0\Longleftrightarrow0<\norm{x}_1.
\]
\end{corollary}

\begin{proof}
Apply \cref{lem:fnd-nonzero-mass} to the underlying function of $x$ and its
stored coordinatewise nonnegativity proof.
\end{proof}

\begin{lemma}[Normalization package]
\label{lem:fnd-normalization-package}
Let $I$ be a finite nonempty type, let $x:\operatorname{ConeVec}(I)$ satisfy
$h_x:x\neq0$, and let $c:\R$ satisfy $h_c:0<c$.  Put
$\widehat x:=\operatorname{normalize}(\operatorname{toPosCone}(x,h_x))$.
Then the literal conjunction
\begin{align*}
 &
 \operatorname{normalize}(\operatorname{posScale}(c,h_c,
 \operatorname{toPosCone}(x,h_x)))=\widehat x
 \quad\wedge\quad
 \supp(\widehat x)=\supp(x)\\
 &\quad\wedge\quad
 (\forall i:I,x(i)=\norm{x}_1\widehat x(i))\\
 &\quad\wedge\quad
 \Bigl(\forall(p:\operatorname{ProbVec}(I))
  (h_p:\langle p,p.2.1\rangle\neq(0:\operatorname{ConeVec}(I))),
 \operatorname{normalize}
   (\operatorname{toPosCone}(\langle p,p.2.1\rangle,h_p))=p
 \Bigr)
\end{align*}
holds; all equalities involving normalizations are equalities of bundled
probability vectors.
\end{lemma}

\begin{proof}
All statements follow by finite-sum distributivity, positivity of
$\norm{x}_1$, and cancellation of positive real scalars.  The support claim
uses $\norm{x}_1\neq0$.
\end{proof}

\begin{lemma}[Positive conditioning outcomes are nonempty]
\label{lem:fnd-positive-columns-nonempty}
Let $X,Y$ be finite nonempty types and let
$P:\operatorname{JointProb}(X,Y)$.  The theorem named
\texttt{positiveColumnsNonempty} has the literal conjunctive conclusion
\[
 \operatorname{Finset.Nonempty}(\operatorname{activeFinset}(P))
 \quad\wedge\quad
 \sum_{y\in\operatorname{activeFinset}(P)}
   \operatorname{colMass}(P,y)=1.
\]
\end{lemma}

\begin{proof}
Every $m_y$ is nonnegative and their sum is one.  Thus at least one is
positive.  A zero-mass column is identically zero by
\cref{lem:fnd-nonzero-mass}, and removing its zero contribution does not
change the sum.
\end{proof}

\begin{definition}[Nonempty active-outcome instance]
\label{def:fnd-active-nonempty}
For finite nonempty $X,Y$ and $P:\operatorname{JointProb}(X,Y)$, define
\[
 \operatorname{activeNonempty}(P):\operatorname{Nonempty}
   (\operatorname{Active}(P))
 :=\operatorname{match}
   (\operatorname{positiveColumnsNonempty}(P)).1
   \operatorname{with}
   \mid\langle y,h_y\rangle\Rightarrow
       \langle\langle y,h_y\rangle\rangle.
\]
Every later definition using \texttt{finMin} or \texttt{finMax} installs
\texttt{letI := activeNonempty P} locally rather than asserting a global
instance that could cause instance-search loops.
\end{definition}

\subsection{Stochastic matrices and conditionally mixing channels}

For finite types $I,J$, define the exact predicates on raw matrices
\begin{align*}
 \operatorname{NonnegMatrix}_{I,J}(D)
 &:\Longleftrightarrow\forall i:I,\forall j:J,\ 0\leq D(i,j),\\
 \operatorname{RowStochastic}_{I,J}(D)
 &:\Longleftrightarrow\operatorname{NonnegMatrix}_{I,J}(D)\ \wedge\
   \forall i:I,\ \sum_{j:J}D(i,j)=1,\\
 \operatorname{SubStochastic}_{I,J}(D)
 &:\Longleftrightarrow\operatorname{NonnegMatrix}_{I,J}(D)\ \wedge\
   \forall i:I,\ \sum_{j:J}D(i,j)\leq1,\\
 \operatorname{DoublyStochastic}_I(S)
 &:\Longleftrightarrow\operatorname{NonnegMatrix}_{I,I}(S)\ \wedge\
   (\forall i:I,\ \sum_{j:I}S(i,j)=1)\ \wedge\
   (\forall j:I,\ \sum_{i:I}S(i,j)=1).
\end{align*}
The matrix action is the total function
\begin{equation}
 \operatorname{matrixAction}(S,x)(i):=\sum_{j:I}S(i,j)x(j).
 \label{eq:foundation-matrix-action}
\end{equation}
When $S$ is nonnegative, package this raw function with its coordinatewise
nonnegativity proof as
$\operatorname{coneMatrixAction}(S,h_S,x):\operatorname{ConeVec}(I)$.
When $S$ is doubly stochastic, package the same action on a probability vector
as
$\operatorname{probMatrixAction}(S,h_S,p):\operatorname{ProbVec}(I)$.

A finite conditional mixing datum from $X\times Y$ to $X\times Y'$ is the
following structure, whose auxiliary index is deliberately
$\operatorname{Fin}(n)$ so that no universe-valued existential is required:
\[
 \operatorname{CMData}(X,Y,Y'):=
 \left\{
 \begin{array}{l}
 n:\N,\quad h_n:0<n,\\
 S:\operatorname{Fin}(n)\to X\to X\to\R,\quad
 h_S:\forall k,\operatorname{DoublyStochastic}_X(S_k),\\
 D:\operatorname{Fin}(n)\to Y\to Y'\to\R,\quad
 h_D:\forall k,\operatorname{NonnegMatrix}_{Y,Y'}(D_k),\\
 h_{\rm norm}:\forall y:Y,\
  \sum_{k:\operatorname{Fin}(n)}\sum_{y':Y'}D_k(y,y')=1
 \end{array}\right\}.
\]
Thus the stored normalization proof is exactly
\begin{equation}
 \sum_{k:\operatorname{Fin}(n)}\sum_{y':Y'}D_k(y,y')=1
 \qquad(y\in Y).
 \label{eq:foundation-conditional-normalization}
\end{equation}
For $C:\operatorname{CMData}(X,Y,Y')$ and
$P:\operatorname{JointProb}(X,Y)$, define the raw output
\begin{equation}
 \operatorname{cmOutputRaw}(C,P)(x,y')
 :=\sum_{k:\operatorname{Fin}(C.n)}\sum_{y:Y}\sum_{z:X}
 C.S(k)(x,z)P(z,y)C.D(k)(y,y'),
 \label{eq:foundation-channel-output}
\end{equation}
whose entrywise nonnegativity and mass-one conclusions are stated in the
subsequently declared theorem \texttt{fndCmOutputRaw}.  Define the relation directly
against the raw function:
\[
 \operatorname{CMRel}_{X,Y,Y'}:
 \operatorname{JointProb}(X,Y)\to\operatorname{JointProb}(X,Y')
 \to\operatorname{Prop},
\]
and, for $P:\operatorname{JointProb}(X,Y)$ and
$Q:\operatorname{JointProb}(X,Y')$, define
\[
 \operatorname{CMRel}(P,Q):\Longleftrightarrow
 \exists C:\operatorname{CMData}(X,Y,Y'),\quad
 \forall x:X,\forall y':Y',\quad
 Q(x,y')=\operatorname{cmOutputRaw}(C,P)(x,y').
\]
The notation $P\xrightarrow{\mathrm{cm}}Q$ abbreviates this proposition.
This presentation avoids using a bundled output before its mass proof has
been declared.  It is equivalent to requiring each $D_k$ to be
sub-stochastic and $\sum_kD_k$ to be row-stochastic, but no later proof
depends on that equivalent presentation.

Fix a universe $u$ and define the genuine higher-rank carrier
\begin{equation}
 \texttt{PolyJointFunctional.\{u\}}:=
 \forall\{X,Y:\texttt{Type u}\},\
 [\operatorname{Fintype}X][\operatorname{Nonempty}X]
 [\operatorname{Fintype}Y][\operatorname{Nonempty}Y],\
 \operatorname{JointProb}(X,Y)\to\R.
 \label{eq:poly-joint-functional}
\end{equation}
The actual declaration header is
\[
 \texttt{abbrev PolyJointFunctional.\{u\} : Type (u+1) :=}
\]
followed by the displayed dependent function type.  Thus the suffix
\texttt{.\{u\}} is universe application, not part of the identifier.
For $F:\texttt{PolyJointFunctional.\{u\}}$, define
\begin{align*}
 \operatorname{CMMonotone}(F)
 &: \Longleftrightarrow
 \forall\{X,Y,Y':\texttt{Type u}\},\
 \forall[\operatorname{Fintype}X][\operatorname{Nonempty}X]
 [\operatorname{Fintype}Y][\operatorname{Nonempty}Y]
 [\operatorname{Fintype}Y'][\operatorname{Nonempty}Y'],\
 \forall P:\operatorname{JointProb}(X,Y),\
 \forall Q:\operatorname{JointProb}(X,Y'),\\[-2pt]
 &\hspace{48mm}\operatorname{CMRel}(P,Q)\Longrightarrow F(P)\leq F(Q).
\end{align*}
Separately, let $X:\texttt{Type u}$ already carry
$[\operatorname{Fintype}X][\operatorname{Nonempty}X]$ and let
$F:\texttt{PolyJointFunctional.\{u\}}$.  The exact declaration header is
\[
 \operatorname{CMMonotoneAtRow}
 (X:\texttt{Type u})[\operatorname{Fintype}X]
 [\operatorname{Nonempty}X](F:\texttt{PolyJointFunctional.\{u\}}):
 \operatorname{Prop},
\]
with body
\begin{align*}
 \operatorname{CMMonotoneAtRow}(X,F)
 &: \Longleftrightarrow
 \forall\{Y,Y':\texttt{Type u}\},\
 \forall[\operatorname{Fintype}Y][\operatorname{Nonempty}Y]
 [\operatorname{Fintype}Y'][\operatorname{Nonempty}Y'],\
 \forall P:\operatorname{JointProb}(X,Y),\
 \forall Q:\operatorname{JointProb}(X,Y'),\\[-2pt]
 &\hspace{48mm}\operatorname{CMRel}(P,Q)\Longrightarrow F(P)\leq F(Q).
\end{align*}
Thus the row instances are parameters of the definition and are not rebound
inside its proposition.  The notation
$\operatorname{CMMonotone}_X(F)$ abbreviates
\texttt{CMMonotoneAtRow X F}.  Every later
occurrence of ``monotone under conditionally mixing channels'' means
$\operatorname{CMMonotone}$ with this inequality direction.

To make the embedding relation independent of the main manuscript, let
$X,Y,X',Y'$ be finite nonempty types, let
$P:\operatorname{JointProb}(X,Y)$, and let
$Q:\operatorname{JointProb}(X',Y')$.  Define the exact proposition
\begin{equation}
 \operatorname{CEmbeds}(P,Q)
 \label{eq:foundation-embedded-preorder}
\end{equation}
to mean that there exist $n_Z,n_U,n_V:\N$ and injections
$e_X:X\hookrightarrow\operatorname{Fin}(n_Z+1)$,
$e'_X:X'\hookrightarrow\operatorname{Fin}(n_Z+1)$,
$e_Y:Y\hookrightarrow\operatorname{Fin}(n_U+1)$, and
$e'_Y:Y'\hookrightarrow\operatorname{Fin}(n_V+1)$, such that
\[
 \operatorname{CMRel}
 \bigl(\operatorname{jointZeroExtend}(e_X,e_Y,P),
       \operatorname{jointZeroExtend}(e'_X,e'_Y,Q)\bigr).
\]
The notation $P\preccurlyeq_{\mathrm c}Q$ abbreviates
\texttt{CEmbeds P Q}.
This is the full conditional-majorization relation used at the end of the
document.  The definition includes relabelings because equivalences are
injections and includes arbitrary finite zero embeddings because the common
types may be enlarged further.

\begin{lemma}[Doubly stochastic mass preservation]
\label{lem:fnd-ds-mass}
Let $I$ be a finite nonempty type, let $S:I\to I\to\R$ satisfy
$h_S:\operatorname{DoublyStochastic}_I(S)$, and let
$x:\operatorname{ConeVec}(I)$.  Put
$Sx:=\operatorname{coneMatrixAction}(S,h_S.1,x)$.  Then
$\norm{Sx}_1=\norm{x}_1$.
\end{lemma}

\begin{proof}
Nonnegativity follows termwise.  Interchange the two finite sums and use the
column-sum identity of $S$.
\end{proof}

\begin{lemma}[A doubly stochastic action preserves nonzeroness]
\label{lem:fnd-ds-action-nonzero}
Let $I$ be a finite nonempty type, let $S:I\to I\to\R$ have proof
$h_S:\operatorname{DoublyStochastic}_I(S)$, let
$x:\operatorname{ConeVec}(I)$, and let $h_x:x\neq0$.  Then the named proof
term
\[
 \operatorname{dsAction\_ne\_zero}(S,h_S,x,h_x):
 \operatorname{coneMatrixAction}(S,h_S.1,x)\neq0
\]
exists.
\end{lemma}

\begin{proof}
Otherwise \cref{lem:fnd-ds-mass} and
\cref{lem:fnd-nonzero-mass} would give $\norm{x}_1=0$ and hence $x=0$.
\end{proof}

\begin{lemma}[Doubly stochastic normalization]
\label{lem:fnd-ds-normalization}
Let $I$ be a finite nonempty type, let $S:I\to I\to\R$ have proof
$h_S:\operatorname{DoublyStochastic}_I(S)$, let
$x:\operatorname{ConeVec}(I)$, and let $h_x:x\neq0$.  Put
$Sx:=\operatorname{coneMatrixAction}(S,h_S.1,x)$.  Then
\[
 \operatorname{normalize}\!\left(
  \operatorname{toPosCone}
   (Sx,\operatorname{dsAction\_ne\_zero}(S,h_S,x,h_x))\right)
 =\operatorname{probMatrixAction}
   \left(S,h_S,
    \operatorname{normalize}(\operatorname{toPosCone}(x,h_x))\right)
\]
as an equality of bundled probability vectors.
\end{lemma}

\begin{proof}
Divide the linear identity $Sx=\norm{x}_1S\widehat x$ by the positive mass,
using \cref{lem:fnd-ds-mass}.
\end{proof}

\begin{lemma}[Raw conditional output is stochastic]
\label{lem:fnd-cm-output-raw}
Let $X,Y,Y'$ be finite nonempty types, let
$C:\operatorname{CMData}(X,Y,Y')$, and let
$P:\operatorname{JointProb}(X,Y)$.  The theorem declaration is named
\texttt{fndCmOutputRaw} and has the literal conjunctive conclusion
\[
 \left(\forall x:X,\forall y':Y',\
  0\leq\operatorname{cmOutputRaw}(C,P)(x,y')\right)
 \quad\wedge\quad
 \left(\sum_{x:X}\sum_{y':Y'}
  \operatorname{cmOutputRaw}(C,P)(x,y')=1\right).
\]
\end{lemma}

\begin{proof}
Every summand in \cref{eq:foundation-channel-output} is a product of
nonnegative numbers.  For the total mass, interchange the finite sums, use
the column-sum field of each \texttt{C.S k}, and finally use
\cref{eq:foundation-conditional-normalization}.  The result is the total
mass of $P$, namely one.
\end{proof}

\begin{definition}[Bundled conditional output]
\label{def:fnd-cm-output}
For finite nonempty $X,Y,Y'$, $C:\operatorname{CMData}(X,Y,Y')$, and
$P:\operatorname{JointProb}(X,Y)$, define
\[
 \operatorname{cmOutput}(C,P):=
 \left\langle\operatorname{cmOutputRaw}(C,P),
   \operatorname{fndCmOutputRaw}(C,P)\right\rangle
 \in\operatorname{JointProb}(X,Y'),
\]
where the second component is exactly the nonnegativity-and-mass proof in
\cref{lem:fnd-cm-output-raw}.
\end{definition}

\begin{lemma}[Conditional output column]
\label{lem:fnd-conditional-channel-stochastic}
Let $X,Y,Y'$ be finite nonempty types, let
$C:\operatorname{CMData}(X,Y,Y')$, and let
$P:\operatorname{JointProb}(X,Y)$, and let $y':Y'$.  Its $y'$-th bundled
column is
\begin{equation}
 \operatorname{column}(\operatorname{cmOutput}(C,P),y')
 =\operatorname{coneSum}\!\left((k,y)\mapsto
 \operatorname{coneScale}
 \bigl(C.D\,k\,y\,y',C.h_D\,k\,y\,y',
       \operatorname{coneMatrixAction}
       (C.S\,k,C.h_S(k).1,\operatorname{column}(P,y))\bigr)\right),
 \label{eq:foundation-output-column}
\end{equation}
where each displayed scalar multiple is the proof-carrying
\texttt{coneScale} using $C.h_D$.
\end{lemma}

\begin{proof}
The identity is finite-sum associativity and extensionality of the bundled
cone subtype.  Its nonnegativity proof fields are definitionally the ones
used by \texttt{coneScale} and \texttt{coneMatrixAction}.
\end{proof}

\begin{lemma}[Nonzero witness in a nonnegative vector sum]
\label{lem:fnd-nonzero-sum-witness}
Let $I,J$ be finite types, let $c:J\to\R$ have proof
$h_c:\forall j,0\leq c_j$, and let
$z:J\to\operatorname{ConeVec}(I)$.  Then the following literal conjunction
holds:
\begin{align*}
&\Bigl[
 \operatorname{weightedConeSum}(c,h_c,z)\neq0
 \quad\Longleftrightarrow\quad
 \exists j:J,\ c_j>0\ \wedge\ z_j\neq0
 \Bigr]\\
&\quad\wedge\quad
 \Bigl[\forall j_0:J,\quad
 c_{j_0}>0\to z_{j_0}\neq0\to
 \operatorname{weightedConeSum}(c,h_c,z)\neq0\Bigr].
\end{align*}
\end{lemma}

\begin{proof}
The reverse implication follows because a nonzero nonnegative vector has a
strictly positive coordinate by \cref{lem:fnd-nonzero-mass}; the same
coordinate of the sum is then positive.  Conversely, if no such index
exists, every weighted vector is zero and so is their finite sum.
\end{proof}

\begin{definition}[Two-column normalization]
\label{def:fnd-merge-norm}
For a finite type $I$, $x,z:\operatorname{ConeVec}(I)$, and $\lambda:\R$,
define
\[
 \operatorname{mergeNorm}(x,z,\lambda):=
 \lambda\norm{x}_1+(1-\lambda)\norm{z}_1.
\]
\end{definition}

\begin{lemma}[Positivity of the merge normalization]
\label{lem:fnd-merge-norm-pos}
Let $I$ be a finite nonempty type, let $x,z:\operatorname{ConeVec}(I)$ have
proofs $h_x:x\neq0$ and $h_z:z\neq0$, and let $\lambda:\R$ have proof
$h_\lambda:0<\lambda\ \wedge\ \lambda<1$.  Then the named conclusion is
\[
 \operatorname{mergeNorm\_pos}(h_x,h_z,h_\lambda):
 0<\operatorname{mergeNorm}(x,z,\lambda)
\]
.
\end{lemma}

\begin{proof}
Both masses are positive by \cref{lem:fnd-nonzero-mass}, and both
coefficients are positive.
\end{proof}

\begin{definition}[Raw merge input]
\label{def:fnd-merge-input-raw}
For a finite type $I$, $x,z:\operatorname{ConeVec}(I)$, and $\lambda:\R$,
put $Z:=\operatorname{mergeNorm}(x,z,\lambda)$ and define
\begin{align*}
 \operatorname{mergeInputRaw}(x,z,\lambda)(i,j)
 &:={\operatorname{Fin.cases}}
   \left(\frac{\lambda x(i)}Z\right)
   \left(r\mapsto\frac{(1-\lambda)z(i)}Z\right)(j),
 \qquad j:\operatorname{Fin}(2).
\end{align*}
\end{definition}

\begin{definition}[Raw merge output]
\label{def:fnd-merge-output-raw}
For a finite type $I$, $x,z:\operatorname{ConeVec}(I)$, and $\lambda:\R$,
put $Z:=\operatorname{mergeNorm}(x,z,\lambda)$ and define
\begin{align*}
 \operatorname{mergeOutputRaw}(x,z,\lambda)(i,j)
 &:=\frac{\lambda x(i)+(1-\lambda)z(i)}Z,
 \qquad j:\operatorname{Fin}(1).
\end{align*}
\end{definition}

\begin{definition}[Bundled merge input]
\label{def:fnd-merge-input}
Let $I$ be a finite nonempty type, let $x,z:\operatorname{ConeVec}(I)$ have
proofs $h_x:x\neq0$ and $h_z:z\neq0$, and let $\lambda:\R$ have proof
$h_\lambda:0<\lambda\ \wedge\ \lambda<1$.  Define
\[
 \operatorname{mergeInput}(x,z,h_x,h_z,\lambda,h_\lambda)
 :=\left\langle\operatorname{mergeInputRaw}(x,z,\lambda),
 h_{\rm input}\right\rangle
 :\operatorname{JointProb}(I,\operatorname{Fin}(2)).
\]
The proof $h_{\rm input}$ supplies entrywise nonnegativity and mass one using
\texttt{mergeNorm\_pos}.
\end{definition}

\begin{definition}[Bundled merge output]
\label{def:fnd-column-merge-data}
Let $I$ be a finite nonempty type, let $x,z:\operatorname{ConeVec}(I)$ have
proofs $h_x:x\neq0$ and $h_z:z\neq0$, and let $\lambda:\R$ have proof
$h_\lambda:0<\lambda\ \wedge\ \lambda<1$.  Define
\[
 \operatorname{mergeOutput}(x,z,h_x,h_z,\lambda,h_\lambda)
 :=\left\langle\operatorname{mergeOutputRaw}(x,z,\lambda),
 h_{\rm output}\right\rangle
 :\operatorname{JointProb}(I,\operatorname{Fin}(1)).
\]
The proof $h_{\rm output}$ supplies entrywise nonnegativity and mass one
using \texttt{mergeNorm\_pos}.
\end{definition}

\begin{lemma}[Column merging is a conditional channel]
\label{lem:fnd-column-merge-channel}
Let $I$ be a finite nonempty type, let $x,z:\operatorname{ConeVec}(I)$ have
proofs $h_x:x\neq0$ and $h_z:z\neq0$, and let $\lambda:\R$ have proof
$h_\lambda:0<\lambda\ \wedge\ \lambda<1$.  Then
\[
 \operatorname{CMRel}
 \bigl(\operatorname{mergeInput}(x,z,h_x,h_z,\lambda,h_\lambda),
       \operatorname{mergeOutput}(x,z,h_x,h_z,\lambda,h_\lambda)\bigr).
\]
A witness has $C.n=1$, the identity doubly stochastic matrix $C.S_0$, and
$C.D_0(j,0)=1$ for both $j:\operatorname{Fin}(2)$.
\end{lemma}

\begin{proof}
Positivity of the denominator follows from
\cref{lem:fnd-merge-norm-pos}.  Take one identity doubly stochastic matrix and
the unique row-stochastic map from a two-element conditioning alphabet to a
singleton.  Division by $Z$ makes the input a probability distribution and
is preserved by the channel.
\end{proof}

\subsection{The compactified parameter space and measures}

The parameter type is the ordered one-point compactification
\begin{equation}
 \overline{\R}_+:=[0,+\infty].
\end{equation}
A direct Lean carrier is $\operatorname{WithTop}(\R_{\geq0})$, fixed by the
literal declaration
\[
 \texttt{abbrev Param := WithTop NNReal}.
\]
Thus $\overline\R_+$ and \texttt{Param} are synonyms, equivalently
$\R_{\geq0}\cup\{\top\}$.  The symbols $0,1,+\infty$, the order intervals,
and the Borel $\sigma$-algebra always refer to this type.  Coercion of a
finite parameter to a real number is performed only after a proof that the
parameter is not $+\infty$.

All positive measures below are Borel measures on $\overline\R_+$.  A
probability measure is a positive measure of total mass one.  The
implementation uses mathlib's bundled \texttt{ProbabilityMeasure} and
\texttt{FiniteMeasure} types and fixes these coercions:
\begin{align*}
 \operatorname{probMeasure}(\tau)&:=(\uparrow\tau:
       \operatorname{Measure}(\overline\R_+)),\\
 \operatorname{finiteMeasure}(\sigma)&:=(\uparrow\sigma:
       \operatorname{Measure}(\overline\R_+)),\\
 \operatorname{atom}(\tau,a)&:\operatorname{ENNReal}:=
       \operatorname{probMeasure}(\tau)(\{a\}).
\end{align*}
Every analytic API---integration, \texttt{lintegral}, restriction, support,
scalar multiplication, and equality of measures---is applied to these
coerced \texttt{Measure} values.  Paper expressions such as $\tau(E)$ and
$\int f\dd\tau$ suppress only this coercion.  In particular, a finite Dirac
decomposition is an equality in \texttt{Measure}, never an equality between
a bundled probability measure and an unbundled measure.

For a Borel map $T$ with a supplied \texttt{AEMeasurable} proof $h_T$
relative to $\operatorname{probMeasure}(\tau)$, define
\[
 \operatorname{probMap}(\tau,T,h_T)
 :=\texttt{tau.map (f := T) h\_T}.
\]
Its coerced underlying measure is the ordinary pushforward
$T_\#\operatorname{probMeasure}(\tau)$.  Every later bundled probability
pushforward uses \texttt{probMap}; statements about its moments use the coerced
underlying measure.

For a finite positive measure $\sigma$ and a map $T$, define the total
bundled pushforward
\[
 \operatorname{finiteMap}(\sigma,T):=\texttt{sigma.map T}.
\]
Unlike \texttt{ProbabilityMeasure.map}, the \texttt{FiniteMeasure.map}
constructor is total; an \texttt{AEMeasurable} proof is supplied to each
evaluation or integration theorem.  Raw statements may equivalently use
\texttt{Measure.map} on $\operatorname{finiteMeasure}(\sigma)$.

Define the bundled Dirac constructor with its full signature by
\[
 \operatorname{diracProb}:
   \overline\R_+\to\operatorname{ProbabilityMeasure}(\overline\R_+),
 \qquad
 \operatorname{diracProb}(a):=
  \left\langle\operatorname{Measure.dirac}(a),
    \texttt{by infer\_instance}\right\rangle.
\]
Its underlying measure is $\operatorname{Measure.dirac}(a)$.  Define
\[
 \operatorname{diracFinite}(a):=
   \operatorname{diracProb}(a).\operatorname{toFiniteMeasure},
 \qquad
 \operatorname{diracRaw}(a):=
   \operatorname{probMeasure}(\operatorname{diracProb}(a)).
\]
Thus later notation $\delta_a$ must be expanded to one of these three names
according to its statically required carrier.  For probability measures
$\tau,\rho$, a real $\eps$, and
$h_\eps:0\leq\eps\ \wedge\ \eps\leq1$, define
$\operatorname{probMix}(\tau,\rho,\eps,h_\eps)$ by the literal constructor
\[
 \left\langle
  \operatorname{ENNReal.ofReal}(1-\eps)\mathbin{\bullet}
    \operatorname{probMeasure}(\tau)
  +\operatorname{ENNReal.ofReal}(\eps)\mathbin{\bullet}
    \operatorname{probMeasure}(\rho),
  \ \texttt{by
    have h1 : 0 <= 1 - eps := sub\_nonneg.mpr h\_eps.2
    simp [h1, h\_eps.1]}
 \right\rangle.
\]
The proof simplifies the value on \texttt{Set.univ}, both probability-mass
fields, and the identity $(1-\eps)+\eps=1$.  This is the only probability
mixture constructor used below; it does not assume a convex-space instance
on \texttt{ProbabilityMeasure}.

\begin{definition}[Nonnegative scaling of a finite measure]
\label{def:fnd-finite-scale}
For $c:\R$, $h_c:0\leq c$, and
$\sigma:\operatorname{FiniteMeasure}(\overline\R_+)$, define
\[
 \operatorname{finiteScale}(c,h_c,\sigma):=
 \left\langle
  \operatorname{ENNReal.ofReal}(c)\mathbin{\bullet}
    \operatorname{finiteMeasure}(\sigma),
  \texttt{by infer\_instance}\right\rangle
 \in\operatorname{FiniteMeasure}(\overline\R_+),
\]
where typeclass inference uses finiteness of an \texttt{ENNReal.ofReal}
scalar multiple.  The proof argument $h_c$ records the intended
real scaling and permits \texttt{ENNReal.toReal\_ofReal} simplification.
\end{definition}

A finite signed measure is mathlib's \texttt{SignedMeasure}.  The following
one-declaration-at-a-time aliases expose its canonical API.

\begin{definition}[Positive measure as a signed measure]
\label{def:fnd-signed-lift}
Let $E:\texttt{Type u}$ carry $[\operatorname{MeasurableSpace}E]$ and
$\sigma:\operatorname{FiniteMeasure}(E)$, define
\[
 \operatorname{signedLift}(\sigma):=
  ((\uparrow\sigma:\operatorname{Measure}(E))).\operatorname{toSignedMeasure}
  \in\operatorname{SignedMeasure}(E).
\]
\end{definition}

\begin{definition}[Positive Jordan part]
\label{def:fnd-signed-pos}
Let $E:\texttt{Type u}$ carry $[\operatorname{MeasurableSpace}E]$.
For $\mu:\operatorname{SignedMeasure}(E)$, define
\[
 \operatorname{signedPos}(\mu):=
 \mu.\operatorname{toJordanDecomposition}.\operatorname{posPart}
 \in\operatorname{Measure}(E).
\]
\end{definition}

\begin{definition}[Negative Jordan part]
\label{def:fnd-signed-neg}
Let $E:\texttt{Type u}$ carry $[\operatorname{MeasurableSpace}E]$.
For $\mu:\operatorname{SignedMeasure}(E)$, define
\[
 \operatorname{signedNeg}(\mu):=
 \mu.\operatorname{toJordanDecomposition}.\operatorname{negPart}
 \in\operatorname{Measure}(E).
\]
\end{definition}

\begin{definition}[Total variation]
\label{def:fnd-signed-tv}
Let $E:\texttt{Type u}$ carry $[\operatorname{MeasurableSpace}E]$.
For $\mu:\operatorname{SignedMeasure}(E)$, define
\[
 \operatorname{signedTV}(\mu):=\mu.\operatorname{totalVariation}.
\]
The paper aliases $\mu^+$, $\mu^-$, and $\abs\mu$ mean respectively
\texttt{signedPos mu}, \texttt{signedNeg mu}, and \texttt{signedTV mu}.
\end{definition}

\begin{lemma}[Total variation as the Jordan sum]
\label{lem:fnd-signed-tv-jordan}
Let $E:\texttt{Type u}$ carry $[\operatorname{MeasurableSpace}E]$.
For $\mu:\operatorname{SignedMeasure}(E)$,
\[
 \operatorname{signedTV}(\mu)=
 \operatorname{signedPos}(\mu)+\operatorname{signedNeg}(\mu).
\]
\end{lemma}

\begin{proof}
This is the defining theorem for
\texttt{SignedMeasure.totalVariation} after unfolding the three aliases.
\end{proof}

\begin{definition}[Canonical signed integral]
\label{eq:foundation-signed-integral}
Let $E:\texttt{Type u}$ carry $[\operatorname{MeasurableSpace}E]$.
For $\mu:\operatorname{SignedMeasure}(E)$ and $f:E\to\R$, define
\[
 \operatorname{signedIntegral}(\mu,f):=
 \int f\dd\operatorname{signedPos}(\mu)
 -\int f\dd\operatorname{signedNeg}(\mu).
\]
This is a total real function; every algebraic theorem about it below carries
integrability against both displayed parts.
\end{definition}

\begin{lemma}[Signed-integral congruence under total variation]
\label{lem:fnd-signed-integral-congr-ae}
Let $E:\texttt{Type u}$ carry $[\operatorname{MeasurableSpace}E]$, let
$\mu:\operatorname{SignedMeasure}(E)$, and let $f,g:E\to\R$.  Then the
named theorem has the literal implication type
\[
 \operatorname{signedIntegral\_congr\_ae}:\quad
 \operatorname{Filter.EventuallyEq}
  (\operatorname{ae}(\operatorname{signedTV}(\mu)),f,g)
 \Longrightarrow
 \operatorname{signedIntegral}(\mu,f)
 =\operatorname{signedIntegral}(\mu,g).
\]
\end{lemma}

\begin{proof}
By \cref{lem:fnd-signed-tv-jordan}, each Jordan part is bounded above by
\texttt{signedTV $\mu$}.  Filter monotonicity therefore turns the supplied
almost-everywhere equality into one for \texttt{signedPos $\mu$} and one for
\texttt{signedNeg $\mu$}.  Unfold \texttt{signedIntegral}, apply mathlib's
\texttt{integral\_congr\_ae} to the two ordinary integrals, and subtract.
\end{proof}

\begin{definition}[Finite negative measure data]
\label{def:fnd-negative-measure}
Define the exact one-field structure
\[
 \operatorname{NegativeMeasure}:=
 \left\{\operatorname{pos}:
   \operatorname{FiniteMeasure}(\overline\R_+)\right\}.
\]
\end{definition}

\begin{definition}[Signed value of negative measure data]
\label{def:fnd-negative-measure-signed}
For $m:\operatorname{NegativeMeasure}$, define
\[
 \operatorname{NegativeMeasure.signed}(m):=
 -\operatorname{signedLift}(m.\operatorname{pos}).
\]
Thus later negative-measure propositions are parameterized by the positive
field and never coerce $-\mu$ to a positive measure.
\end{definition}

\begin{definition}[One-signedness]
\label{def:fnd-one-signed}
Let $E:\texttt{Type u}$ carry $[\operatorname{MeasurableSpace}E]$.
For $\mu:\operatorname{SignedMeasure}(E)$, define
\[
 \operatorname{OneSigned}(\mu):\Longleftrightarrow
 \operatorname{signedPos}(\mu)=0\ \lor\
 \operatorname{signedNeg}(\mu)=0.
\]
For a finite positive measure $\nu$, the notation $\nu^{\mathrm s}$ means
\texttt{signedLift nu}.
\end{definition}

For a raw positive Borel measure $\nu:\operatorname{Measure}(\overline\R_+)$,
define
\[
 \operatorname{suppMeasure}:
 \operatorname{Measure}(\overline\R_+)\to\operatorname{Set}(\overline\R_+),
 \qquad
 \operatorname{suppMeasure}(\nu)
 :=\{a\mid\forall U:\operatorname{Set}(\overline\R_+),\
      \operatorname{IsOpen}(U)\to a\in U\to\nu(U)\neq0\}.
\]
For $\mu:\operatorname{SignedMeasure}(\overline\R_+)$, define
\begin{equation}
 \operatorname{suppSigned}(\mu)
 :=\operatorname{suppMeasure}(\operatorname{signedTV}(\mu)).
 \label{eq:foundation-measure-support}
\end{equation}
The overloaded paper notation $\supp(\nu)$ and $\supp(\mu)$ denotes the
first or second function according to the statically known measure type; a
Lean implementation uses the two displayed names and no ambiguous coercion.

\begin{lemma}[Positive-measure support complement]
\label{lem:fnd-support-complement}
Let $\nu:\operatorname{Measure}(\overline\R_+)$ carry an
$[\operatorname{IsFiniteMeasure}(\nu)]$ instance.  Then the following
literal conjunction holds:
\begin{align*}
 &\operatorname{IsOpen}
   (\operatorname{suppMeasure}(\nu)^{\mathrm c})\ \wedge\
 \nu(\operatorname{suppMeasure}(\nu)^{\mathrm c})=0\\
 &\quad\wedge\ \forall C:\operatorname{Set}(\overline\R_+),\
 \operatorname{IsClosed}(C)\Longrightarrow\\[-2pt]
 &\qquad\Bigl[
  \bigl(\operatorname{suppMeasure}(\nu)\subseteq C
       \Longleftrightarrow\nu(C^{\mathrm c})=0\bigr)\ \wedge\\[-2pt]
 &\qquad\quad
  \bigl(\operatorname{suppMeasure}(\nu)\subseteq C\Longrightarrow
    \operatorname{Measure.restrict}(\nu,C)=\nu\ \wedge\\[-2pt]
 &\qquad\qquad
    (\forall f:\overline\R_+\to\R,\
      \operatorname{Integrable}(f,\nu)\Longrightarrow
      \int_C f\dd\nu=\int f\dd\nu)\ \wedge\\[-2pt]
 &\qquad\qquad
    (\forall g:\overline\R_+\to\operatorname{ENNReal},\
      \operatorname{Measurable}(g)\Longrightarrow
      \int^-_C g\dd\nu=\int^-g\dd\nu)\bigr)
 \Bigr].
\end{align*}
\end{lemma}

\begin{proof}
Every point outside the support has an open null neighbourhood.  Choose a
countable subcover using second countability; the union remains null.  The
forward implication follows.  Conversely, if $\nu(C^c)=0$ and $x\notin C$,
then the open set $C^c$ is a null neighbourhood of $x$, so $x$ is outside
the support.  If $\nu(C^c)=0$, the defining formula for measure restriction
gives $\operatorname{Measure.restrict}(\nu,C)=\nu$ on every measurable set; the two
integration identities follow by congruence of measures.
\end{proof}

\begin{lemma}[Signed-measure support complement]
\label{lem:fnd-signed-support-complement}
Let $\mu:\operatorname{SignedMeasure}(\overline\R_+)$.  Then
\begin{align*}
 &\operatorname{signedTV}(\mu)
   (\operatorname{suppSigned}(\mu)^{\mathrm c})=0\\
 &\quad\wedge\ \forall C:\operatorname{Set}(\overline\R_+),\
 \operatorname{IsClosed}(C)\Longrightarrow\\[-2pt]
 &\qquad\Bigl[
  \bigl(\operatorname{suppSigned}(\mu)\subseteq C
   \Longleftrightarrow\operatorname{signedTV}(\mu)(C^{\mathrm c})=0\bigr)
   \ \wedge\\[-2pt]
 &\qquad\quad
  \bigl(\operatorname{suppSigned}(\mu)\subseteq C\Longrightarrow
   \operatorname{Measure.restrict}(\operatorname{signedPos}(\mu),C)
      =\operatorname{signedPos}(\mu)\ \wedge\
   \operatorname{Measure.restrict}(\operatorname{signedNeg}(\mu),C)
      =\operatorname{signedNeg}(\mu)\ \wedge\\[-2pt]
 &\qquad\qquad
   \forall f:\overline\R_+\to\R,\
   \operatorname{Integrable}(f,\operatorname{signedPos}(\mu))\Longrightarrow
   \operatorname{Integrable}(f,\operatorname{signedNeg}(\mu))\Longrightarrow\\[-2pt]
 &\hspace{62mm}
   \operatorname{signedIntegral}
    (\mu.\operatorname{restrict}(C),f)
   =\operatorname{signedIntegral}(\mu,f)\bigr)
 \Bigr].
\end{align*}
\end{lemma}

\begin{proof}
Apply \cref{lem:fnd-support-complement} to the finite measure
$\operatorname{signedTV}(\mu)$.  Both Jordan parts are absolutely continuous
with respect to $\operatorname{signedTV}(\mu)$;
their restricted-measure identities and the signed-integral equality follow.
\end{proof}

\begin{lemma}[Pushforward integration]
\label{lem:fnd-pushforward-integration}
Let $\nu:\operatorname{Measure}(\overline\R_+)$ carry an
$[\operatorname{IsFiniteMeasure}(\nu)]$ instance, let
$T:\overline\R_+\to\overline\R_+$ have proof
$h_T:\operatorname{AEMeasurable}(T,\nu)$, and let
$f:\overline\R_+\to\R$ and $C:\R$ have proofs
$h_f:\operatorname{Measurable}(f)$, $h_C:0\leq C$, and
$h_{\rm bnd}:\forall a,\abs{f(a)}\leq C$.  Then
\[
 \int f(\beta)\dd(\operatorname{Measure.map}(T,\nu))(\beta)
 =\int f(T(\alpha))\dd\nu(\alpha).
\]
\end{lemma}

\begin{proof}
First prove the identity for indicator functions from the definition of the
pushforward, then for nonnegative simple functions by linearity, for
nonnegative measurable functions by monotone approximation, and for bounded
real functions by positive and negative parts.
\end{proof}

\begin{lemma}[Pushforward nonnegative integration]
\label{lem:fnd-pushforward-lintegral}
Let $\nu:\operatorname{Measure}(\overline\R_+)$, let
$T:\overline\R_+\to\overline\R_+$ have proof
$h_T:\operatorname{AEMeasurable}(T,\nu)$, and let
$g:\overline\R_+\to\operatorname{ENNReal}$ have proof
$h_g:\operatorname{Measurable}(g)$.  Then
\[
 \int^- g(\beta)\dd(\operatorname{Measure.map}(T,\nu))(\beta)
 =\int^- g(T(\alpha))\dd\nu(\alpha).
\]
\end{lemma}

\begin{proof}
Apply the \texttt{lintegral\_map} theorem to $h_T$ and the measurability
proof of $g$.
\end{proof}

\begin{lemma}[Bounded finite-measure integration]
\label{lem:fnd-bounded-integral}
Let $\mu:\operatorname{SignedMeasure}(\overline\R_+)$, let $C:\R$ have
$h_C:0\leq C$, and let $f:\overline\R_+\to\R$ have proofs
$h_f:\operatorname{Measurable}(f)$ and
$h_{\rm bnd}:\forall a,\abs{f(a)}\leq C$.  Then the literal conjunction
\begin{align*}
 &\operatorname{Integrable}(f,\operatorname{signedPos}(\mu))\ \wedge\
  \operatorname{Integrable}(f,\operatorname{signedNeg}(\mu))\\
 &\quad\wedge\quad
 \abs{\operatorname{signedIntegral}(\mu,f)}\leq
 C\,\operatorname{ENNReal.toReal}\!\bigl(
   \operatorname{signedTV}(\mu)(\operatorname{Set.univ})\bigr)
\end{align*}
holds.
\end{lemma}

\begin{proof}
Apply the corresponding positive-measure estimate to both parts of the
Jordan decomposition.
\end{proof}

\subsection{Exact target and dependency discipline}

The formal target is the final theorem \texttt{mainClassification}. Its four clauses are
to be represented as four Lean theorems and then bundled into one final
theorem.  Each theorem is universally quantified over finite nonempty
alphabets and over all conditional mixing data of compatible sizes.
"Monotone" means the inequality for every such datum, not only for column
mergers or maps on the conditioning register.

The recommended implementation order is:
\begin{enumerate}
 \item the total carriers and elementary finite-sum facts of this section;
 \item the R\'enyi and measure library of Section~1;
 \item the curvature, perspective, channel, and limit library of Section~2;
 \item the definitions and reductions of Section~3;
 \item sufficiency in Section~4;
 \item the perturbation and necessity library of Section~5;
 \item the four clauses of the final theorem.
\end{enumerate}
Every later proof refers only backward in this order.  The few results that
are already available in mathlib may be discharged by an existing theorem;
the displayed statements still specify the exact specialization and all side
conditions that the Lean proof must expose.

\subsection{Logarithm and normalization convention}

Throughout this document, $\log$ is the natural logarithm and $\exp$ is its
inverse. The resulting fair-bit value is $\log2$; this is proved only after
the total R\'enyi definition, by the subsequently declared uniform-distribution
lemma and fair-bit corollary. This is exactly the convention of the
complete-proof Sections~4--5 paper, which is the specification for this Lean
development.

\section{Closed R\'enyi and parameter-measure library}
\label{sec:renyi-foundations}

\subsection{Total definition of the extended-order entropy}

Let $p$ be a probability vector on a nonempty finite type $I$.  Define
$s(p):=\card\supp(p)$ and $m(p):=\max_i p_i$.  The support is nonempty, so
$s(p)\geq1$ and $m(p)>0$.  For a finite real order $a>0$, $a\neq1$, set
\begin{equation}
 H_a(p):=\frac{\log\bigl(\sum_i p_i^a\bigr)}{1-a}.
 \label{eq:fnd-renyi-finite}
\end{equation}
Here and throughout, powers of nonnegative real coordinates by real orders
are \texttt{Real.rpow}; in particular the zero-base value at positive order
is total and equals zero.
At the remaining orders set
\begin{align}
 H_0(p)&:=\log s(p),\label{eq:fnd-renyi-zero}\\
 H_1(p)&:=-\sum_i p_i\log p_i,
 \quad\text{where }0\log0:=0,\label{eq:fnd-renyi-one}\\
 H_{+\infty}(p)&:=-\log m(p).
 \label{eq:fnd-renyi-infty}
\end{align}
For every fixed finite nonempty type $I$, define the polymorphic total
function
\[
 H_I:\overline\R_+\to\operatorname{ProbVec}(I)\to\R
\]
by the following exhaustive body (the proof arguments of the conditionals
are used only to refine the branches):
\begin{align*}
 H_I(\alpha,p):={}&
 \texttt{if }h_\top:\alpha=\top\texttt{ then }
   -\log\!\left(\operatorname{finMax}(i\mapsto p.1(i))\right)\\[-2pt]
 &\texttt{else let }a:\R:=
   \bigl((\alpha.\operatorname{untop}(h_\top):\operatorname{NNReal}):\R\bigr)
   \texttt{; }\\[-2pt]
 &\texttt{if }h_0:a=0\texttt{ then }
   \log\!\left((\operatorname{Finset.card}(\supp(p)):\N):\R\right)\\[-2pt]
 &\texttt{else if }h_1:a=1\texttt{ then }
   -\sum_{i:I}\left(\texttt{if }p.1(i)=0\texttt{ then }0
      \texttt{ else }p.1(i)\log(p.1(i))\right)\\[-2pt]
 &\texttt{else }
   \frac{\log\!\left(\sum_{i:I}
      \operatorname{Real.rpow}(p.1(i),a)\right)}{1-a}.
 \tag{H}\label{eq:fnd-renyi-total}
\end{align*}
The four preceding displayed formulas are the corresponding simplification
lemmas for this body, not additional partial definitions.  In Lean the
alphabet is an implicit parameter:
\texttt{H \{I : Type u\} [Fintype I] [Nonempty I] :
Param $\to$ ProbVec I $\to$ Real}.  There is no single ill-typed carrier called ``all finite probability
vectors.''

\begin{lemma}[Power sum is positive]
\label{lem:fnd-power-sum-positive}
Let $I$ be a finite nonempty type, let $p:\operatorname{ProbVec}(I)$, and
let $a:\R$ have proof $h_a:0<a$.  Then
$\sum_i p_i^a>0$.
\end{lemma}

\begin{proof}
At least one coordinate of $p$ is positive.  Its positive real power is a
strictly positive summand; all other summands are nonnegative.
\end{proof}

\begin{lemma}[Uniform distribution]
\label{lem:fnd-uniform-renyi}
Let $d:\N$ have proof $h_d:0<d$ and let
$\alpha:\overline\R_+$.  The exact conclusion is the dependent let
\[
 \texttt{letI : Nonempty (Fin d) := }
 \langle\langle0,h_d\rangle\rangle\texttt{; }\quad
 H_\alpha(\operatorname{uniformProb}_{\operatorname{Fin}(d)})
 =\log(d:\R).
\]
\end{lemma}

\begin{proof}
For finite $\alpha\neq1$, the power sum is $d^{1-\alpha}$.  The three
endpoint clauses follow directly from their definitions.
\end{proof}

\begin{corollary}[Fair-bit normalization]
\label{cor:fnd-uniform-bit}
For every $\alpha:\overline\R_+$,
\begin{equation}
 H_\alpha(\operatorname{uniformProb}_{\operatorname{Fin}(2)})=\log2.
 \label{eq:natural-log-normalization}
\end{equation}
\end{corollary}

\begin{proof}
Apply \cref{lem:fnd-uniform-renyi} with $d=2$ and its canonical positivity
proof.
\end{proof}

\begin{lemma}[Entropy bounds]
\label{lem:fnd-renyi-bounds}
Let $I$ be a finite nonempty type, let $p:\operatorname{ProbVec}(I)$, let
$d:\N$ have proof $h_d:\operatorname{Fintype.card}(I)\leq d$, and let
$\alpha:\overline\R_+$.  Then the exact conclusion is
\[
 0\leq H_\alpha(p)
 \quad\wedge\quad H_\alpha(p)\leq\log(d:\R)
 \quad\wedge\quad
 H_\alpha(p)\leq\log\bigl((\card(\supp(p)):\N):\R\bigr).
\]
\end{lemma}

\begin{proof}
Delete the zero coordinates.  For $0<a<1$, concavity of $r\mapsto r^a$
and Jensen's inequality give
$1\leq\sum p_i^a\leq s^{1-a}$.  For $a>1$, convexity gives
$s^{1-a}\leq\sum p_i^a\leq1$.  Taking logarithms and accounting for the sign
of $1-a$ gives the claim.  The Shannon bound follows from Jensen applied to
$-r\log r$, and the endpoint bounds follow from
$1/s\leq\max_i p_i\leq1$.
\end{proof}

\begin{lemma}[Continuity at order zero]
\label{lem:fnd-renyi-cont-zero}
Let $I$ be a finite nonempty type and let $p:\operatorname{ProbVec}(I)$.
Then
\[
 \operatorname{Tendsto}(a\mapsto H_a(p))
   (\mathcal N(0:\overline\R_+))(\mathcal N(H_0(p))).
\]
\end{lemma}

\begin{proof}
After restricting to the $s$ positive coordinates, each $p_i^a$ tends to
one, so the finite sum tends to $s$.  Also $1-a\to1$.  Apply continuity of
the logarithm at the positive number $s$.
\end{proof}

\begin{lemma}[Continuity at the Shannon order]
\label{lem:fnd-renyi-cont-one}
Let $I$ be a finite nonempty type and let $p:\operatorname{ProbVec}(I)$.
Then
\[
 \operatorname{Tendsto}(a\mapsto H_a(p))
   (\mathcal N(1:\overline\R_+))(\mathcal N(H_1(p))).
\]
\end{lemma}

\begin{proof}
Delete zero coordinates and put $Z(a)=\sum_i p_i^a$.  Then $Z$ is smooth,
$Z(1)=1$, and $Z'(1)=\sum_i p_i\log p_i$.  The fundamental theorem of
calculus gives
\[
 \frac{\log Z(a)-\log Z(1)}{a-1}
 =\int_0^1\frac{Z'(1+s(a-1))}{Z(1+s(a-1))}\dd s.
\]
Choose $0<\delta<1/2$.  After shrinking $\delta$, continuity and $Z(1)=1$
give $Z(r)\geq1/2$ on $[1-\delta,1+\delta]$.  The quotient $Z'/Z$ is then
continuous on that compact interval, hence uniformly continuous.  Therefore
the displayed integrand converges uniformly in $s\in[0,1]$ to $Z'(1)$ as
$a\to1$.  Negating the quotient proves the claim.  This quantitative
denominator bound is the side condition needed by Lean's division and
integral-continuity lemmas and avoids an unformalized appeal to
l'H\^opital's rule.
\end{proof}

\begin{lemma}[Continuity at infinite order]
\label{lem:fnd-renyi-cont-infty}
Let $I$ be a finite nonempty type and let $p:\operatorname{ProbVec}(I)$.
Then
\[
 \operatorname{Tendsto}(a\mapsto H_a(p))
   (\mathcal N(+\infty:\overline\R_+))
   (\mathcal N(H_{+\infty}(p))).
\]
\end{lemma}

\begin{proof}
Let $m=\max_i p_i$ and let $s$ be the number of positive coordinates.  Then
\[
 m^a\leq\sum_i p_i^a\leq s m^a.
\]
For $a>1$, division by $1-a<0$ gives a two-sided estimate whose error from
$-\log m$ is at most $\log s/(a-1)$.  This tends to zero.
In the $\operatorname{WithTop}(\R_{\geq0})$ topology, a net of finite orders
converges to $+\infty$ exactly when its extracted real values eventually
exceed every real bound; hence the real estimate is precisely the endpoint
continuity statement on $\overline\R_+$.
\end{proof}

\begin{proposition}[Continuity and measurability in the order]
\label{prop:fnd-renyi-order-continuity}
For every finite nonempty $I$ and $p:\operatorname{ProbVec}(I)$, the exact
return proposition is
\begin{align*}
 &\operatorname{Continuous}(\alpha\mapsto H_\alpha(p))
 \quad\wedge\quad
 \operatorname{Measurable}(\alpha\mapsto H_\alpha(p))\\
 &\quad\wedge\quad
 \exists C:\R,\quad0\leq C\ \wedge\
 \forall\alpha:\overline\R_+,\abs{H_\alpha(p)}\leq C.
\end{align*}
\end{proposition}

\begin{proof}
Continuity away from $0,1,+\infty$ follows from continuity of finite sums,
real powers, and logarithms.  Use
\cref{lem:fnd-renyi-cont-zero,lem:fnd-renyi-cont-one,lem:fnd-renyi-cont-infty}
at the three exceptional orders.
\end{proof}

\subsection{Algebraic, order, and curvature facts}

\begin{lemma}[Relabeling and zero-embedding invariance]
\label{lem:fnd-renyi-embedding}
Let $I,J$ be finite nonempty types, let $p:\operatorname{ProbVec}(I)$, let
$e:I\hookrightarrow J$, let $r:I\simeq J$, and let
$\alpha:\overline\R_+$.  Then the following literal conjunction holds:
\begin{align*}
&H_\alpha(\operatorname{zeroExtendProb}(e,p))=H_\alpha(p)\\
&\quad\wedge\quad
 H_\alpha(\operatorname{zeroExtendProb}(r.\operatorname{toEmbedding},p))
 =H_\alpha(p).
\end{align*}
\end{lemma}

\begin{proof}
Finite sums are invariant under bijective reindexing.  At positive finite
orders, the added zeros contribute zero to the power sum; at order one they
contribute $0\log0=0$.  They do not change the support cardinality or the
maximum, which proves the two other endpoint cases.
\end{proof}

\begin{lemma}[Tensor-product additivity]
\label{lem:fnd-renyi-tensor}
Let $I,J$ be finite nonempty types, let $p:\operatorname{ProbVec}(I)$,
$q:\operatorname{ProbVec}(J)$, and $\alpha:\overline\R_+$.  Then
\[
 H_\alpha(\operatorname{probTensor}(p,q))=H_\alpha(p)+H_\alpha(q).
\]
\end{lemma}

\begin{proof}
For finite $\alpha\neq1$, factor the double power sum.  At $0$, supports
multiply; at $1$, expand the double sum and use that both marginal sums are
one; at $+\infty$, maxima multiply.
\end{proof}

\begin{definition}[Finite-order power norm]
\label{def:fnd-lp-norm}
For a finite type $I$, define
\[
 \operatorname{lpNorm}_I(a,x):=
 \operatorname{Real.rpow}
 \left(\sum_{i:I}\operatorname{Real.rpow}(x(i),a)\right)(1/a)
\]
for $a:\R$ and $x:\operatorname{ConeVec}(I)$.
\end{definition}

\begin{lemma}[Scale and norm formula]
\label{lem:fnd-renyi-norm-formula}
Let $I$ be finite and
nonempty, let $x:\operatorname{ConeVec}(I)$ satisfy $h_x:x\neq0$, and let
$a:\R$ have proofs $h_a:0<a$ and $h_{a1}:a\neq1$.  Put
$\widehat x:=\operatorname{normalize}(\operatorname{toPosCone}(x,h_x))$.
Then the following literal conjunction holds:
\begin{align}
&
 H_{\operatorname{ENNReal.ofReal}(a)}(\widehat x)=\frac{a}{1-a}
   \bigl(\log\operatorname{lpNorm}_I(a,x)-\log\norm{x}_1\bigr)
 \label{eq:fnd-renyi-norm-formula}
\\[-2pt]
&\quad\wedge\quad
 H_{\top}(\widehat x)
 =-\bigl(\log(\operatorname{finMax}(i\mapsto x(i)))
          -\log\norm{x}_1\bigr).
\end{align}
\end{lemma}

\begin{proof}
Use
$\sum_i(x_i/\norm{x}_1)^a=\norm{x}_a^a/\norm{x}_1^a$.
At infinity use
$H_\infty(\widehat x)=\log\norm{x}_1-\log\max_i x_i$.
\end{proof}

\begin{lemma}[Schur concavity under a doubly stochastic matrix]
\label{lem:fnd-renyi-schur}
Let $I$ be a finite nonempty type, let $S:I\to I\to\R$ satisfy
$h_S:\operatorname{DoublyStochastic}_I(S)$, let
$p:\operatorname{ProbVec}(I)$, and let $\alpha:\overline\R_+$.  Then
\[
 H_\alpha(\operatorname{probMatrixAction}(S,h_S,p))\geq H_\alpha(p).
\]
\end{lemma}

\begin{proof}
For $0<a<1$, concavity of $r^a$ applied to each row gives
$(Sp)_i^a\geq\sum_jS(i,j)p_j^a$.  Summing and using the column sums gives
$\sum_i(Sp)_i^a\geq\sum_jp_j^a$; the factor $1/(1-a)$ is positive.  For
$a>1$, both inequalities reverse, including the negative factor
$1/(1-a)$, so the entropy inequality has the same direction.  At $a=1$ use
concavity of $-r\log r$.  At $+\infty$, every output coordinate is a convex
combination of input coordinates, so its maximum cannot exceed the input
maximum.

At $a=0$, let $A=\supp(p)$ and $B=\supp(Sp)$.  For $i\notin B$ and
$j\in A$, nonnegativity and $(Sp)_i=0$ imply $S(i,j)=0$.  Hence
\[
 \card A=\sum_{j\in A}\sum_iS(i,j)
 =\sum_{i\in B}\sum_{j\in A}S(i,j)
 \leq\sum_{i\in B}1=\card B.
\]
Taking logarithms proves the endpoint claim.
\end{proof}

\begin{lemma}[Concavity for orders at most one]
\label{lem:fnd-renyi-concavity}
Let $I$ be a finite nonempty type and let
$\alpha:\overline\R_+$ have proofs $h_{\alpha,0}:0\leq\alpha$ and
$h_{\alpha,1}:\alpha\leq1$.  Then
\[
 \operatorname{SimplexConcave}_I(p\mapsto H_\alpha(p)).
\]
\end{lemma}

\begin{proof}
For $0<a<1$, the power sum $p\mapsto\sum_i p_i^a$ is concave and positive.
The map $r\mapsto\log r/(1-a)$ is increasing and concave, so their
composition is concave.  At $a=1$, sum the coordinatewise concavity of
$-r\log r$.

At $a=0$, for $0<\lambda<1$,
$\supp(\lambda p+(1-\lambda)q)=\supp(p)\cup\supp(q)$.  Therefore
\[
 H_0(\lambda p+(1-\lambda)q)
 \geq\max\{H_0(p),H_0(q)\}
 \geq\lambda H_0(p)+(1-\lambda)H_0(q).
\]
The cases $\lambda=0,1$ are equalities.
\end{proof}

\begin{lemma}[Strict binary entropy gap]
\label{lem:fnd-strict-binary-gap}
Let $r:\R$ satisfy $h_r:0<r\ \wedge\ r<1$ and
$h_{\rm ne}:r\neq1/2$, and put
$p:=\operatorname{binaryProb}
\,r\,\operatorname{le_of_lt}(h_r.1)\,\operatorname{le_of_lt}(h_r.2)$.  Then
\[
 H_{(0:\overline\R_+)}(p)=\log2
 \quad\wedge\quad
 \forall\alpha:\overline\R_+,\quad0<\alpha\to H_\alpha(p)<\log2.
\]
\end{lemma}

\begin{proof}
The support claim follows from positivity of both entries.  For
$0<a<1$, strict concavity of $s^a$ gives the strict comparison between
$r^a+(1-r)^a$ and $2^{1-a}$ required by $1/(1-a)>0$.  For $a>1$, strict
convexity gives the reverse power-sum comparison and division by
$1-a<0$ again gives a strict entropy upper bound.  At $a=1$ use strict
concavity of Shannon entropy.  At infinity,
$\max\{r,1-r\}>1/2$.
\end{proof}

\subsection{Integrated entropy and signed integration}

For every finite nonempty type $I$, define two differently named functions:
\begin{align}
 \operatorname{integratedEntropyPos}_I
 &: \operatorname{Measure}(\overline\R_+)
    \to\operatorname{ProbVec}(I)\to\R,
 &\operatorname{integratedEntropyPos}_I(\nu,p)
 &:=\int H_\alpha(p)\dd\nu(\alpha),
 \label{eq:fnd-integrated-entropy}\\
 \operatorname{integratedEntropySigned}_I
 &: \operatorname{SignedMeasure}(\overline\R_+)
    \to\operatorname{ProbVec}(I)\to\R,
 &\operatorname{integratedEntropySigned}_I(\mu,p)
 &:=\operatorname{signedIntegral}
    (\mu,\alpha\mapsto H_\alpha(p)).
\end{align}
The positive function is total on raw measures; all analytic properties below
assume finiteness.  Paper notation is only an abbreviation with the following
mandatory expansions:
\begin{align*}
 \A_\tau(p)&:=\operatorname{integratedEntropyPos}_I
   (\operatorname{probMeasure}(\tau),p)
 &&(\tau:\operatorname{ProbabilityMeasure}),\\
 \A_\sigma(p)&:=\operatorname{integratedEntropyPos}_I
   (\operatorname{finiteMeasure}(\sigma),p)
 &&(\sigma:\operatorname{FiniteMeasure}),\\
 \A_\mu(p)&:=\operatorname{integratedEntropySigned}_I(\mu,p)
 &&(\mu:\operatorname{SignedMeasure}).
\end{align*}
No Lean declaration overloads one constant solely by these three carriers.

\begin{lemma}[Dirac integration and support]
\label{lem:fnd-dirac-package}
For every finite nonempty $I$, every $\alpha:\overline\R_+$, and every
$p:\operatorname{ProbVec}(I)$, the following literal conjunction holds:
\begin{align*}
&
 \operatorname{integratedEntropyPos}_I
  (\operatorname{diracRaw}(\alpha),p)=H_\alpha(p)\\
&\quad\wedge\quad
 \operatorname{suppMeasure}(\operatorname{diracRaw}(\alpha))=\{\alpha\}.
\end{align*}
The bundled value $\operatorname{diracProb}(\alpha)$ already has the type of
a probability measure, so this is not an additional theorem conclusion.
\end{lemma}

\begin{proof}
The integral identity is the defining integral rule for a Dirac measure.
Every open neighbourhood of $\alpha$ has Dirac mass one, whereas a point
different from $\alpha$ has an open neighbourhood not containing $\alpha$
because the parameter space is Hausdorff.  This proves the support identity;
the mass-one claim is immediate.
\end{proof}

\begin{lemma}[Integrated entropy is finite and linear in the measure]
\label{lem:fnd-integrated-entropy-finite}
Let $I$ be a finite nonempty type, let
$p:\operatorname{ProbVec}(I)$, let $d:\N$ have proof
$h_d:\operatorname{Fintype.card}(I)\leq d$, let
$\mu:\operatorname{SignedMeasure}(\overline\R_+)$, let
$\sigma:\operatorname{FiniteMeasure}(\overline\R_+)$, and let
$\tau:\operatorname{ProbabilityMeasure}(\overline\R_+)$.  Then the following
single proposition holds:
\begin{align*}
&\abs{\operatorname{integratedEntropySigned}_I(\mu,p)}
 \leq(\log(d:\R))\,
 \operatorname{ENNReal.toReal}\!\bigl(
   \operatorname{signedTV}(\mu)(\operatorname{Set.univ})\bigr)\\
&\quad\wedge\quad
 0\leq\operatorname{integratedEntropyPos}_I
       (\operatorname{finiteMeasure}(\sigma),p)\\[-2pt]
&\quad\wedge\quad
 \operatorname{integratedEntropyPos}_I
       (\operatorname{finiteMeasure}(\sigma),p)
 \leq\operatorname{ENNReal.toReal}
       (\operatorname{finiteMeasure}(\sigma)(\operatorname{Set.univ}))\log(d:\R)\\
&\quad\wedge\quad
 \operatorname{integratedEntropyPos}_I
       (\operatorname{probMeasure}(\tau),p)\leq\log(d:\R)\\
&\quad\wedge\quad
 \forall t:\R,\quad
 \operatorname{integratedEntropySigned}_I
 (t\mathbin{\bullet}\operatorname{signedLift}(\sigma),p)
 =t\operatorname{integratedEntropyPos}_I
   (\operatorname{finiteMeasure}(\sigma),p).
\end{align*}
Here the finite total-variation mass is coerced to $\R$.  The last conjunct
is the precise linearity instance needed later; unrestricted linearity in
two signed measures is supplied by the standard signed-integral API.
\end{lemma}

\begin{proof}
Combine \cref{lem:fnd-renyi-bounds,lem:fnd-bounded-integral}.  Positivity of
the integral gives the one-sided bound for a positive measure; its total
mass is finite, so its extended value has the displayed real coercion.
Linearity is the standard linearity theorem for the canonical signed
integral.
\end{proof}

\begin{lemma}[Integrated structural properties]
\label{lem:fnd-integrated-properties}
Let $I,J,K$ be finite nonempty types, let
$\sigma:\operatorname{FiniteMeasure}(\overline\R_+)$, let
$p:\operatorname{ProbVec}(I)$ and $q:\operatorname{ProbVec}(J)$, let
$S:I\to I\to\R$ satisfy $h_S:\operatorname{DoublyStochastic}_I(S)$, and let
$e:I\hookrightarrow K$.  Put
\begin{align*}
A_I&:=p\mapsto\operatorname{integratedEntropyPos}_I
 (\operatorname{finiteMeasure}(\sigma),p),\\
A_J&:=q\mapsto\operatorname{integratedEntropyPos}_J
 (\operatorname{finiteMeasure}(\sigma),q),\\
A_{I\times J}&:=z\mapsto\operatorname{integratedEntropyPos}_{I\times J}
 (\operatorname{finiteMeasure}(\sigma),z),\\
A_K&:=z\mapsto\operatorname{integratedEntropyPos}_K
 (\operatorname{finiteMeasure}(\sigma),z).
\end{align*}
Then the following literal conjunction holds:
\begin{align*}
&A_I(\operatorname{probMatrixAction}(S,h_S,p))\geq A_I(p)\\
&\quad\wedge\quad
 A_{I\times J}(\operatorname{probTensor}(p,q))=A_I(p)+A_J(q)\\
&\quad\wedge\quad
 A_K(\operatorname{zeroExtendProb}(e,p))=A_I(p)\\
&\quad\wedge\quad
 \Bigl(\forall r:I\simeq K,\quad
 A_K(\operatorname{zeroExtendProb}(r.\operatorname{toEmbedding},p))=A_I(p)\Bigr)\\
&\quad\wedge\quad
 \Bigl(\operatorname{suppMeasure}(\operatorname{finiteMeasure}(\sigma))
 \subseteq\operatorname{Set.Icc}(0,1)
 \to\operatorname{SimplexConcave}_I(A_I)\Bigr).
\end{align*}
\end{lemma}

\begin{proof}
Integrate
\cref{lem:fnd-renyi-schur,lem:fnd-renyi-tensor,lem:fnd-renyi-embedding,lem:fnd-renyi-concavity},
respectively.  The support
condition permits restriction to $[0,1]$ by
\cref{lem:fnd-support-complement}.
\end{proof}

\begin{lemma}[Finite sum and integral]
\label{lem:fnd-finite-sum-integral}
Let $Y:\texttt{Type u}$ carry $[\operatorname{Fintype}Y]$, let
$\mu:\operatorname{SignedMeasure}(\overline\R_+)$, let $c:Y\to\R$, and let
$f:Y\to\overline\R_+\to\R$.  Assume
\[
 \forall y:Y,\quad
 \operatorname{Measurable}(f(y))\ \wedge\
 \operatorname{Integrable}(f(y),\operatorname{signedPos}(\mu))\ \wedge\
 \operatorname{Integrable}(f(y),\operatorname{signedNeg}(\mu)).
\]
Then
\[
 \operatorname{signedIntegral}
  \left(\mu,\alpha\mapsto\sum_{y:Y}c(y)f(y,\alpha)\right)
 =\sum_{y:Y}c(y)\operatorname{signedIntegral}(\mu,f(y)).
\]
\end{lemma}

\begin{proof}
Induct on the finite set and use linearity of the finite signed integral.
\end{proof}

\begin{lemma}[Zero integral of a strictly positive kernel]
\label{lem:fnd-zero-integral-positive}
Let $\tau:\operatorname{Measure}(\overline\R_+)$ carry an
$[\operatorname{IsFiniteMeasure}(\tau)]$ instance.  The following two
implications form one literal conjunction:
\begin{align*}
 &\Bigl(\forall f:\overline\R_+\to\R,\
   \operatorname{Measurable}(f)\Longrightarrow
   (\forall a,0\leq f(a))\Longrightarrow
   \operatorname{Integrable}(f,\tau)\Longrightarrow
   \int f\dd\tau=0\Longrightarrow
   \tau(\{a\mid0<f(a)\})=0\Bigr)\\
 &\quad\wedge\Bigl(
   \forall(E:\operatorname{Set}(\overline\R_+))(C:\R)
          (g:\overline\R_+\to\R),\
   \operatorname{MeasurableSet}(E)\Longrightarrow0\leq C\Longrightarrow\\[-2pt]
 &\hspace{22mm}
   \operatorname{Measurable}(g)\Longrightarrow
   (\forall a\in E,0\leq g(a)\wedge g(a)\leq C)\Longrightarrow
   \int_E g\dd\tau=C\,\operatorname{ENNReal.toReal}(\tau(E))\Longrightarrow\\[-2pt]
 &\hspace{52mm}
   \forall^{\mathrm{ae}}a\ \partial
      (\operatorname{Measure.restrict}(\tau,E)),\
   g(a)=C\Bigr).
\end{align*}
\end{lemma}

\begin{proof}
The set $\{f>0\}$ is the countable union, indexed by $n:\N$, of
$E_n=\{f\geq1/(n+1:\R)\}$.  If one $E_n$ had positive measure, monotonicity of the
integral would give
$\int f\dd\tau\geq
\operatorname{ENNReal.toReal}(\tau(E_n))/(n+1:\R)>0$; finiteness of
$\tau$ makes the coercion strictly positive.  Thus every $E_n$ is
null.  For the second claim apply the first to
$\mathbf1_E(C-g)$; its integral is
$C\operatorname{ENNReal.toReal}(\tau(E))-\int_Eg\dd\tau=0$.  This function is
integrable because it is Borel and bounded by $C$ against a finite measure.
\end{proof}

\subsection{The singular weight and one-sided moments}

Define the total function $\omega:\overline\R_+\to\R$ by the literal case
split
\[
 \omega(a):=
 \texttt{if htop : a = $\top$ then -1 else}
 \quad\texttt{let r : Real := ((a.untop htop : NNReal) : Real);}
 \quad\texttt{if r = 1 then 0 else r / (1-r)}.
\]
Equivalently, in paper notation,
\begin{equation}
 \omega(a)=\frac a{1-a}\quad(a<+\infty,\ a\neq1),
 \qquad \omega(1)=0,\qquad \omega(+\infty)=-1.
 \label{eq:fnd-omega}
\end{equation}
The arbitrary value at $1$ makes the Lean definition total; every integral
of $\omega$ is restricted away from $\{1\}$.

\begin{lemma}[Elementary properties of the singular weight]
\label{lem:fnd-omega-properties}
The following single literal conjunction holds:
\begin{align*}
&\omega(0)=0\\
&\quad\wedge\quad
 (\forall a\in\operatorname{Set.Ioo}(0,1),\ 0<\omega(a))\\
&\quad\wedge\quad
 \operatorname{StrictMonoOn}(\omega,\operatorname{Set.Ico}(0,1))\\
&\quad\wedge\quad
 (\forall a:\overline\R_+,\ 1<a\to\omega(a)<0)\\
&\quad\wedge\quad
 \Bigl(\forall(a:\overline\R_+)(h_\top:a\neq\top),\quad
 1<a\to
 -\omega(a)=
 \frac{((a.\operatorname{untop}(h_\top):\R_{\geq0}):\R)}
      {((a.\operatorname{untop}(h_\top):\R_{\geq0}):\R)-1}\Bigr)\\
&\quad\wedge\quad -\omega(\top)=1\\
&\quad\wedge\quad
 \Bigl(\forall r:\overline\R_+,\ r<1\to
  \exists C:\R,\ 0\leq C\ \wedge\
  \forall a\in\operatorname{Set.Icc}(0,r),\abs{\omega(a)}\leq C\Bigr)\\
&\quad\wedge\quad
 \Bigl(\forall(r:\overline\R_+)(h_\top:r\neq\top),\ 1<r\to
  \exists C,\eps:\R,\ 0\leq C\ \wedge\ 0<\eps\ \wedge\
  \forall a\in\operatorname{Set.Ioi}(r),\quad
   \abs{\omega(a)}\leq C\ \wedge\ \eps\leq\abs{\omega(a)}\Bigr).
\end{align*}
\end{lemma}

\begin{proof}
For finite $a\neq1$, direct algebra gives
$\omega'(a)=1/(1-a)^2>0$.  The signs follow from the two denominator signs.
The bounds follow from the explicit rational formula on a set separated from
$1$ and the limit $\omega(a)\to-1$ as $a\to+\infty$.
\end{proof}

Define the nonnegative extended-real kernels with Lean-safe names
\[
 \operatorname{omegaLower}(a):=
 \begin{cases}
  \operatorname{ENNReal.ofReal}(\omega(a)),
    &a\in\operatorname{Set.Ico}((0:\operatorname{Param}),1),\\
  0,&\text{otherwise},
 \end{cases}
\]
\[
 \operatorname{omegaUpper}(a):=
 \begin{cases}
  \operatorname{ENNReal.ofReal}(-\omega(a)),
    &a\in\operatorname{Set.Ioi}((1:\operatorname{Param})),\\
  0,&\text{otherwise}.
 \end{cases}
\]
The symbols $\omega_<$ and $\omega_>$ are notation for
\texttt{omegaLower} and \texttt{omegaUpper}; they do not introduce further
constants.
\begin{lemma}[Measurability of the singular kernels]
\label{lem:fnd-omega-measurable}
The following literal conjunction holds:
\[
 \operatorname{Measurable}(\omega)
 \quad\wedge\quad
 \operatorname{Measurable}(\operatorname{omegaLower})
 \quad\wedge\quad
 \operatorname{Measurable}(\operatorname{omegaUpper}).
\]
\end{lemma}

\begin{proof}
On each of the Borel pieces $[0,1)$, $\{1\}$,
$(1,+\infty)$, and $\{+\infty\}$, the functions are either constant or the
composition of the continuous finite-value extraction with the rational
function $a/(1-a)$, whose denominator is nonzero there.  Measurability follows
by the finite Borel pasting lemma and measurability of
\texttt{ENNReal.ofReal}.
\end{proof}

For a raw positive Borel measure
$\nu:\operatorname{Measure}(\overline\R_+)$, define the nonnegative extended
integrals, using Lean's \texttt{lintegral},
\begin{align}
 \operatorname{MLower}(\nu)&:=\int^-\operatorname{omegaLower}\dd\nu,
 &M_<(\nu)&:=\operatorname{MLower}(\nu),\label{eq:fnd-M-lower}\\
 \operatorname{MUpper}(\nu)&:=\int^-\operatorname{omegaUpper}\dd\nu,
 &M_>(\nu)&:=\operatorname{MUpper}(\nu).
 \label{eq:fnd-M-upper}
\end{align}
These are first defined in $[0,+\infty]$.  Only after proofs
$M_<(\nu)<+\infty$ and $M_>(\nu)<+\infty$ may they be interpreted as the
positive and negative parts of an ordinary signed integral.  For a literally
total Lean definition, set
\begin{align}
 \operatorname{MomFin}(\nu)
 &:\Longleftrightarrow
 \operatorname{MLower}(\nu)<\top\ \wedge\
 \operatorname{MUpper}(\nu)<\top,
 \label{eq:fnd-moment-finite}\\
 \operatorname{MReal}(\nu)
 &:=\operatorname{ENNReal.toReal}(M_<(\nu))
   -\operatorname{ENNReal.toReal}(M_>(\nu)).
 \label{eq:fnd-M-real}
\end{align}
The notation $M_{\R}(\nu)$ is a pure notation for
$\operatorname{MReal}(\nu)$.
Here \texttt{ENNReal.toReal} is total and has the library value zero at
$+\infty$.  No theorem is allowed to use that arbitrary value: every use of
$M_{\R}$ below carries a proof of $\operatorname{MomFin}(\nu)$.  For a
probability measure $\tau$, every paper expression $M_<(\tau)$,
$M_>(\tau)$, $M_{\R}(\tau)$, or $\operatorname{MomFin}(\tau)$ expands by
first applying $\operatorname{probMeasure}(\tau)$; for a bundled finite
measure it expands using $\operatorname{finiteMeasure}$.  We use the
short notation
\begin{equation}
 M(\nu):=M_{\R}(\nu).
 \label{eq:fnd-M-signed}
\end{equation}
Thus no subtraction in the extended nonnegative reals is used.  If $c\geq0$,
the notation $M_<(\nu)\leq c$ means the typed inequality
$M_<(\nu)\leq\operatorname{ENNReal.ofReal}(c)$; in particular it implies
$M_<(\nu)<+\infty$.

\begin{lemma}[Moment restriction and scaling]
\label{lem:fnd-moment-algebra}
Let $\nu:\operatorname{Measure}(\overline\R_+)$ carry an
$[\operatorname{IsFiniteMeasure}(\nu)]$ instance, let $c:\R$ have
proof $h_c:0\leq c$, and define the raw scaled measure
$\nu_c:=\operatorname{ENNReal.ofReal}(c)\mathbin{\bullet}\nu$.  Then the
following single nested conjunction holds:
\begin{align*}
&\Bigl(
 \operatorname{MLower}(\nu_c)
  =\operatorname{ENNReal.ofReal}(c)\operatorname{MLower}(\nu)
 \ \wedge\ 
 \operatorname{MUpper}(\nu_c)
  =\operatorname{ENNReal.ofReal}(c)\operatorname{MUpper}(\nu)\Bigr)\\
&\quad\wedge\Bigl(0<c\to
  (\operatorname{MomFin}(\nu_c)\Longleftrightarrow
   \operatorname{MomFin}(\nu))\ \wedge\
  (\operatorname{MomFin}(\nu)\to
   \operatorname{MReal}(\nu_c)=c\operatorname{MReal}(\nu))\Bigr)\\
&\quad\wedge\Bigl(
 \operatorname{suppMeasure}(\nu)\subseteq
   \operatorname{Set.Icc}(0,1)\to
 \operatorname{MUpper}(\nu)=0\Bigr)\\
&\quad\wedge\Bigl(\forall a_*:\overline\R_+,\ 1<a_*\to
 \operatorname{suppMeasure}(\nu)
  \subseteq\operatorname{Set.Icc}(0,1)\cup\{a_*\}\to\\[-2pt]
&\hspace{22mm}
 \operatorname{MUpper}(\nu)
 =\operatorname{ENNReal.ofReal}(-\omega(a_*))\nu(\{a_*\})\ 
 \wedge\ \operatorname{MUpper}(\nu)<\top\Bigr)\\
&\quad\wedge\Bigl(\operatorname{MomFin}(\nu)\to
 \operatorname{Integrable}(\omega,\nu)\ \wedge\\[-2pt]
&\hspace{12mm}
 \int a,\operatorname{ENNReal.toReal}(\operatorname{omegaLower}(a))
   \dd\nu
 =\operatorname{ENNReal.toReal}(\operatorname{MLower}(\nu))\ \wedge\\[-2pt]
&\hspace{12mm}
 \int a,\operatorname{ENNReal.toReal}(\operatorname{omegaUpper}(a))
   \dd\nu
 =\operatorname{ENNReal.toReal}(\operatorname{MUpper}(\nu))\ \wedge\\[-2pt]
&\hspace{12mm}
 \int a,\omega(a)\dd\nu=\operatorname{MReal}(\nu)\Bigr)\\
&\quad\wedge\Bigl(\nu(\{1\})=0\to\operatorname{MomFin}(\nu)\to
 \operatorname{MReal}(\nu)
 =\int_{\{1\}^{\mathrm c}}\omega\dd\nu\Bigr)\\
&\quad\wedge\Bigl(
 \forall(\sigma:\operatorname{Measure}(\overline\R_+))
 [\operatorname{IsFiniteMeasure}(\sigma)],\
 \forall a,b:\R,\quad0\leq a\to0\leq b\to
 \operatorname{MomFin}(\nu)\to\operatorname{MomFin}(\sigma)\to\\[-2pt]
&\hspace{12mm}
 \texttt{let }\rho:=
   \operatorname{ENNReal.ofReal}(a)\mathbin{\bullet}\nu
   +\operatorname{ENNReal.ofReal}(b)\mathbin{\bullet}\sigma\texttt{; }\
 \operatorname{MomFin}(\rho)\ \wedge\
 \operatorname{MReal}(\rho)
  =a\operatorname{MReal}(\nu)+b\operatorname{MReal}(\sigma)\Bigr)\\
&\quad\wedge\Bigl(\forall a:\overline\R_+,\ a\neq1\to
 \operatorname{MomFin}(\operatorname{diracRaw}(a))\ \wedge\
 \operatorname{MReal}(\operatorname{diracRaw}(a))=\omega(a)\Bigr).
\end{align*}
\end{lemma}

\begin{proof}
Use homogeneity of nonnegative integration and
\cref{lem:fnd-support-complement}.  In the exceptional-support case, the
restriction above one is concentrated on the singleton $\{a_*\}$.  Under
finiteness, the ordinary real integrals of $\omega_<$ and $\omega_>$ equal
$\operatorname{ENNReal.toReal}(M_<)$ and $\operatorname{ENNReal.toReal}(M_>)$,
respectively.  Subtracting those two real identities is exactly the
positive-part-minus-negative-part definition of the signed integral.
When $c>0$, multiplication by the finite nonzero value
$\operatorname{ENNReal.ofReal}(c)$ preserves and reflects non-topness.
Applying \texttt{ENNReal.toReal\_mul} to the two identities in the first
clause and subtracting proves the guarded real scaling formula.
For the final clause, linearity of \texttt{lintegral} for sums and
nonnegative scalar multiples gives the two one-sided identities in
$[0,+\infty]$.  Their terms are finite by hypothesis.  Apply
\texttt{ENNReal.toReal\_add}, \texttt{ENNReal.toReal\_mul}, and
\texttt{ENNReal.toReal\_ofReal}, and subtract the upper identity from the
lower one.
For a Dirac measure, evaluate both nonnegative kernels at $a$ and inspect the
two cases $a<1$ and $a>1$ (including $a=+\infty$); the two finite values
subtract to $\omega(a)$.
\end{proof}

\begin{lemma}[Support and strict moment witness]
\label{lem:fnd-exceptional-witness-positive}
Let $\nu:\operatorname{Measure}(\overline\R_+)$ carry an
$[\operatorname{IsFiniteMeasure}(\nu)]$ instance, let $t:\R$
have proof $h_t:t<0$, and let $a_*:\overline\R_+$ have proof $h_*:1<a_*$.  Assume
$h_{\rm supp}:\operatorname{suppMeasure}(\nu)
 \subseteq\operatorname{Set.Icc}(0,1)\cup\{a_*\}$,
$h_<:\operatorname{MLower}(\nu)<\top$, and
$h_M:\operatorname{MReal}(\nu)\leq1/t$.  Then the exact conclusion is
\[
 0<\nu(\{a_*\})
 \quad\wedge\quad a_*\in\operatorname{suppMeasure}(\nu).
\]
\end{lemma}

\begin{proof}
If the singleton had mass zero, then $M_>(\nu)=0$ and, after coercing the
finite lower moment, $M_{\R}(\nu)=\operatorname{ENNReal.toReal}(M_<(\nu))\geq0$,
whereas $1/t<0$.  This contradicts the assumed
inequality.
\end{proof}

\subsection{Total finite-range discretizations}

Define the proof-carrying upper-order subtype
\[
 \operatorname{UpperParam}:=\{a:\overline\R_+\mathbin{//}1<a\}.
\]
For $k:\N$, put $q_k:=k+2$.  Define the proof-parameterized helper
\[
 \operatorname{grid}:(k:\N)\to(a:\overline\R_+)\to(a<1)\to
 \overline\R_+.
\]
For $h_a:a<1$, let $r_a:\R_{\geq0}$ be the finite value extracted from $a$
using $h_a$, and set
\[
 \operatorname{grid}(k,a,h_a):=
 \left(\operatorname{NNReal.ofReal}
 \frac{\operatorname{Nat.floor}
   (((q_k:\N):\R)\mathbin{*}(r_a:\R))}{((q_k:\N):\R)}
 :\overline\R_+\right).
\]
Now define three total functions with literal signatures
\begin{align*}
 \operatorname{TPos},\operatorname{TLow}
 &: \N\to\overline\R_+\to\overline\R_+,\\
 \operatorname{TExc}
 &: \operatorname{UpperParam}\to\N\to
    \overline\R_+\to\overline\R_+,
\end{align*}
by
\begin{equation}
 \operatorname{TPos}(k,a)=
 \begin{cases}
  \operatorname{grid}(k,a,h_a),&h_a:a<1,\\
  \operatorname{ENNReal.ofReal}(1-1/(q_k:\R)),&a=1,\\
  0,&1<a,
 \end{cases}
 \label{eq:fnd-total-discretization}
\end{equation}
\[
 \operatorname{TLow}(k,a)=
 \begin{cases}
  \operatorname{grid}(k,a,h_a),&h_a:a<1,\\
  1,&a=1,\\
  0,&1<a,
 \end{cases}
\]
and
\begin{equation}
 \operatorname{TExc}(a_*,k,a)=
 \begin{cases}
  \operatorname{grid}(k,a,h_a),&h_a:a<1,\\
  \operatorname{ENNReal.ofReal}(1-1/(q_k:\R)),&a=1,\\
  a_*.1,&a=a_*.1,\\
  0,&\text{otherwise}.
 \end{cases}
 \label{eq:fnd-exceptional-discretization}
\end{equation}
The ordering of these decidable branches is part of the definitions.  Define
named Borel proofs
\[
 h_{\rm pos}(k):\operatorname{Measurable}(\operatorname{TPos}(k,\cdot)),
 \quad h_{\rm low}(k):\operatorname{Measurable}(\operatorname{TLow}(k,\cdot)),
 \quad h_{\rm exc}(a_*,k):
 \operatorname{Measurable}(\operatorname{TExc}(a_*,k,\cdot)).
\]
For $\tau:\operatorname{ProbabilityMeasure}(\overline\R_+)$, define the
bundled sequences
\begin{align*}
 \operatorname{discPos}(\tau,k)&:=
  \operatorname{probMap}(\tau,\operatorname{TPos}(k,\cdot),
                         h_{\rm pos}(k).\operatorname{aemeasurable}),\\
 \operatorname{discLow}(\tau,k)&:=
  \operatorname{probMap}(\tau,\operatorname{TLow}(k,\cdot),
                         h_{\rm low}(k).\operatorname{aemeasurable}),\\
 \operatorname{discExc}(\tau,a_*,k)&:=
  \operatorname{probMap}(\tau,\operatorname{TExc}(a_*,k,\cdot),
                         h_{\rm exc}(a_*,k).\operatorname{aemeasurable}).
\end{align*}
No term of any sequence carries an arithmetic side condition.

\begin{lemma}[Basic discretization facts]
\label{lem:fnd-discretization-basic}
Let $\tau:\operatorname{ProbabilityMeasure}(\overline\R_+)$, let
$a_*:\operatorname{UpperParam}$, and let $k:\N$.  The exact return
proposition is the following nested conjunction:
\begin{align*}
&\Bigl[
 \operatorname{Set.Finite}
   (\operatorname{Set.range}(\operatorname{TPos}(k,\cdot)))
 \wedge
 \operatorname{Set.Finite}
   (\operatorname{Set.range}(\operatorname{TLow}(k,\cdot)))\\[-2pt]
&\hspace{31mm}\wedge
 \operatorname{Set.Finite}
   (\operatorname{Set.range}(\operatorname{TExc}(a_*,k,\cdot)))
 \Bigr]\\
&\quad\wedge\Bigl[
 \operatorname{Set.Finite}
  (\operatorname{suppMeasure}(\operatorname{probMeasure}
    (\operatorname{discPos}(\tau,k))))
 \wedge
 \operatorname{Set.Finite}
  (\operatorname{suppMeasure}(\operatorname{probMeasure}
    (\operatorname{discLow}(\tau,k))))\\[-2pt]
&\hspace{31mm}\wedge
 \operatorname{Set.Finite}
  (\operatorname{suppMeasure}(\operatorname{probMeasure}
    (\operatorname{discExc}(\tau,a_*,k))))
 \Bigr]\\
&\quad\wedge\Bigl[
 \forall a\in\operatorname{Set.Ico}(0,1),\quad
 0\leq\operatorname{TPos}(k,a)
 \wedge\operatorname{TPos}(k,a)\leq a
 \wedge\omega(\operatorname{TPos}(k,a))\leq\omega(a)
 \Bigr]\\
&\quad\wedge\Bigl[
 \bigl(\forall a:\overline\R_+,\ a\leq1\to
   \operatorname{Tendsto}
   (m\mapsto\operatorname{TPos}(m,a))\texttt{ atTop }\mathcal N(a)\bigr)\\[-2pt]
&\hspace{14mm}\wedge
 \bigl(\forall a:\overline\R_+,\ a\leq1\to
   \operatorname{Tendsto}
   (m\mapsto\operatorname{TLow}(m,a))\texttt{ atTop }\mathcal N(a)\bigr)\\[-2pt]
&\hspace{14mm}\wedge
 \bigl(\forall a:\overline\R_+,\ (a\leq1\vee a=a_*.1)\to
   \operatorname{Tendsto}
   (m\mapsto\operatorname{TExc}(a_*,m,a))\texttt{ atTop }\mathcal N(a)\bigr)
 \Bigr].
\end{align*}
\end{lemma}

\begin{proof}
Use the floor inequalities with $q_k=k+2$ and the monotonicity of $\omega$
on $[0,1)$.  Each range is contained in a finite grid together with at most
$0,1-1/q_k,1$, and $a_*.1$.  The support of a pushforward is contained in
the closure of its range, hence is finite here.
\end{proof}

\begin{lemma}[Discretization package]
\label{lem:fnd-discretization-package}
Let $\tau:\operatorname{ProbabilityMeasure}(\overline\R_+)$, let
$a_*:\operatorname{UpperParam}$, and let $k:\N$.  Abbreviate the raw
measures
\[
 \nu:=\operatorname{probMeasure}(\tau),\quad
 \nu_{\rm pos}:=\operatorname{probMeasure}(\operatorname{discPos}(\tau,k)),
 \quad
 \nu_{\rm low}:=\operatorname{probMeasure}(\operatorname{discLow}(\tau,k)),
\]
\[
 \nu_{\rm exc}:=\operatorname{probMeasure}
   (\operatorname{discExc}(\tau,a_*,k)).
\]
The theorem's exact return type is
\begin{align*}
&\Bigl[
 \operatorname{suppMeasure}(\nu)\subseteq\operatorname{Set.Icc}(0,1)
 \wedge\nu(\{1\})=0\ \to\\[-2pt]
&\hspace{9mm}
 \operatorname{suppMeasure}(\nu_{\rm pos})
   \subseteq\operatorname{Set.Ico}(0,1)
 \wedge\nu_{\rm pos}(\{1\})=0
 \wedge M_<(\nu_{\rm pos})\leq M_<(\nu)\\[-2pt]
&\hspace{9mm}{}\wedge
 M_>(\nu_{\rm pos})=M_>(\nu)
 \wedge M_>(\nu)=0
 \Bigr]\\
&\quad\wedge\Bigl[
 \operatorname{suppMeasure}(\nu)\subseteq\operatorname{Set.Icc}(0,1)
 \to
 \operatorname{suppMeasure}(\nu_{\rm low})
   \subseteq\operatorname{Set.Icc}(0,1)
 \Bigr]\\
&\quad\wedge\Bigl[
 \operatorname{suppMeasure}(\nu)
   \subseteq\operatorname{Set.Icc}(0,1)\cup\{a_*.1\}
 \wedge\nu(\{1\})=0\ \to\\[-2pt]
&\hspace{9mm}
 \operatorname{suppMeasure}(\nu_{\rm exc})
   \subseteq\operatorname{Set.Ico}(0,1)\cup\{a_*.1\}
 \wedge\nu_{\rm exc}(\{1\})=0
 \wedge M_<(\nu_{\rm exc})\leq M_<(\nu)\\[-2pt]
&\hspace{9mm}{}\wedge M_>(\nu_{\rm exc})=M_>(\nu)
 \wedge\bigl(\operatorname{MomFin}(\nu)\to
   \operatorname{MomFin}(\nu_{\rm exc})
   \wedge M_{\R}(\nu_{\rm exc})\leq M_{\R}(\nu)\bigr)
 \Bigr].
\end{align*}
Every moment function is therefore applied to an explicitly typed raw
measure.  No moment inequality is asserted in the middle implication,
because that branch fixes a possible atom at $1$.
\end{lemma}

\begin{proof}
The floor function is Borel, and the definitions have finitely many values on
$[0,1]$ plus at most the value $a_*.1$.  The pointwise facts follow from
$\lfloor q_k(r_a:\R)\rfloor\leq q_k(r_a:\R)
 <\lfloor q_k(r_a:\R)\rfloor+1$ and
\cref{lem:fnd-omega-properties}.  Pushforwards preserve total mass and a
finite-range pushforward has finite support.  Apply
\cref{lem:fnd-pushforward-lintegral} to the two nonnegative moment kernels.
In the first and third clauses the exceptional value at $1$ is ignored only
after using $\tau(\{1\})=0$; this hypothesis is necessary.  The exceptional
map fixes the only possible upper support point $a_*.1$.
\end{proof}

\begin{lemma}[Common-measure dominated convergence]
\label{lem:fnd-discretization-dct}
Let $I$ be a finite nonempty type, let $p:\operatorname{ProbVec}(I)$, let
$\nu:\operatorname{Measure}(\overline\R_+)$ carry an
$[\operatorname{IsFiniteMeasure}(\nu)]$ instance, and let
$a_*:\operatorname{UpperParam}$.  Define
\begin{align*}
 L_{\rm pos}(k)&:=\int H_\beta(p)\dd
   \operatorname{Measure.map}(\operatorname{TPos}(k,\cdot),\nu)(\beta),\\
 L_{\rm low}(k)&:=\int H_\beta(p)\dd
   \operatorname{Measure.map}(\operatorname{TLow}(k,\cdot),\nu)(\beta),\\
 L_{\rm exc}(k)&:=\int H_\beta(p)\dd
   \operatorname{Measure.map}(\operatorname{TExc}(a_*,k,\cdot),\nu)(\beta),\\
 L&:=\int H_\alpha(p)\dd\nu(\alpha).
\end{align*}
Then the following literal conjunction holds:
\begin{align*}
 &\Bigl(
   \operatorname{suppMeasure}(\nu)
      \subseteq\operatorname{Set.Icc}(0,1)\ \wedge\ \nu(\{1\})=0
   \Longrightarrow
   \operatorname{Tendsto}(L_{\rm pos},\texttt{atTop},\mathcal N(L))
  \Bigr)\\
 &\quad\wedge
  \Bigl(
   \operatorname{suppMeasure}(\nu)
      \subseteq\operatorname{Set.Icc}(0,1)
   \Longrightarrow
   \operatorname{Tendsto}(L_{\rm low},\texttt{atTop},\mathcal N(L))
  \Bigr)\\
 &\quad\wedge
  \Bigl(
   \operatorname{suppMeasure}(\nu)
      \subseteq\operatorname{Set.Icc}(0,1)\cup\{a_*.1\}
   \ \wedge\ \nu(\{1\})=0
   \Longrightarrow
   \operatorname{Tendsto}(L_{\rm exc},\texttt{atTop},\mathcal N(L))
  \Bigr).
\end{align*}
\end{lemma}

\begin{proof}
Use \cref{lem:fnd-pushforward-integration}.  The exceptional sets on which
pointwise convergence can fail are, respectively,
\[
 \{1\}\cup\operatorname{Set.Ioi}(1),\qquad
 \operatorname{Set.Ioi}(1),\qquad
 \{1\}\cup\bigl(\operatorname{Set.univ}\setminus
   (\operatorname{Set.Icc}(0,1)\cup\{a_*.1\})\bigr).
\]
Each is $\nu$-null by the corresponding hypotheses and
\cref{lem:fnd-support-complement}.  Thus
$\operatorname{TPos}(k,\alpha)$,
$\operatorname{TLow}(k,\alpha)$, or
$\operatorname{TExc}(a_*,k,\alpha)$ tends to $\alpha$ almost everywhere in
the corresponding branch.  Apply
\cref{prop:fnd-renyi-order-continuity}.  The bound
$\log(\operatorname{Fintype.card}(I):\R)$ from
\cref{lem:fnd-renyi-bounds} is integrable against the finite measure.  One
common pushforward sequence works for a finite list because it is independent
of the probability vector.
\end{proof}

\section{Conditional functionals and elementary curvature closure}
\label{sec:conditional-foundations}

\subsection{Total conditional candidates}

For every universe $u$, finite nonempty types $X,Y$, and
$P:\operatorname{JointProb}(X,Y)$, use the nonempty finite subtype
$\operatorname{Active}(P)$ from \cref{eq:positive-column-finset}.  For a
probability measure $\tau$ and $t\neq0$, define the temperate and tropical
candidates below.  The derivation candidate $H_{0,\sigma}$ is defined by the
same displayed formula for every finite positive measure $\sigma$:
\begin{align}
 \operatorname{HTemp}(t,h_t,\tau)(P)
 &:=\frac1t\log\sum_{y:\operatorname{Active}(P)}
   \operatorname{colMass}(P,y.1)
   \exp\!\left(t\operatorname{integratedEntropyPos}_X
     (\operatorname{probMeasure}(\tau),
      \operatorname{conditional}(P,y))\right),
   \label{eq:fnd-H-temperate}\\
 \operatorname{HZero}(\sigma)(P)
 &:=\sum_{y:\operatorname{Active}(P)}
   \operatorname{colMass}(P,y.1)
   \operatorname{integratedEntropyPos}_X
    (\operatorname{finiteMeasure}(\sigma),
     \operatorname{conditional}(P,y)),\label{eq:fnd-H-zero}\\
 \operatorname{HMinus}(\tau)(P)
 &:=\left(\texttt{letI : Nonempty (Active P) := activeNonempty P; }
   \operatorname{finMin}
   \left(y\mapsto\operatorname{integratedEntropyPos}_X
    (\operatorname{probMeasure}(\tau),
     \operatorname{conditional}(P,y))\right)\right),
   \label{eq:fnd-H-minus}\\
 \operatorname{HPlus}(\tau)(P)
 &:=\left(\texttt{letI : Nonempty (Active P) := activeNonempty P; }
   \operatorname{finMax}
   \left(y\mapsto\operatorname{integratedEntropyPos}_X
    (\operatorname{probMeasure}(\tau),
     \operatorname{conditional}(P,y))\right)\right).
 \label{eq:fnd-H-plus}
\end{align}
The exact Lean signatures return the higher-rank carrier from
\cref{eq:poly-joint-functional}:
\begin{align*}
 \operatorname{HTemp}&:(t:\R)\to(t\neq0)\to
   \operatorname{ProbabilityMeasure}(\overline\R_+)\to
   \texttt{PolyJointFunctional.\{u\}},\\
 \operatorname{HZero}&:
   \operatorname{FiniteMeasure}(\overline\R_+)\to
   \texttt{PolyJointFunctional.\{u\}},\\
 \operatorname{HMinus},\operatorname{HPlus}&:
   \operatorname{ProbabilityMeasure}(\overline\R_+)\to
   \texttt{PolyJointFunctional.\{u\}}.
\end{align*}
Here \texttt{FiniteMeasure} means a finite positive Borel measure.  The
paper notation suppresses the proof $t\neq0$; proof irrelevance makes the
result independent of that proof, and every use below obtains it from a sign
hypothesis or states it explicitly.  For a probability measure $\tau$, the
notation $H_{0,\tau}$ means
\texttt{HZero tau.toFiniteMeasure}; it does not rely on a coercion chosen by
elaboration.
The symbols $H_{t,\tau}$, $H_{-\infty,\tau}$, and $H_{+\infty,\tau}$ are
pure paper notation for the three corresponding named declarations above.
The tropical expressions are auxiliary extensions of the notation in the
main body.  The named main-body candidates are
$H_{0,\alpha}=H_{0,\operatorname{diracFinite}(\alpha)}$ and
$H_{+\infty,0}=H_{+\infty,\operatorname{diracProb}(0)}$.

\begin{lemma}[Well-definedness and positivity]
\label{lem:fnd-conditional-well-defined}
Let $X,Y$ be finite nonempty types, let
$P:\operatorname{JointProb}(X,Y)$, let
$\tau:\operatorname{ProbabilityMeasure}(\overline\R_+)$, let
$\sigma:\operatorname{FiniteMeasure}(\overline\R_+)$, and let $t:\R$ have
proof $h_t:t\neq0$.  Define
\[
 S_{t,\tau}(P):=
 \sum_{y:\operatorname{Active}(P)}\operatorname{colMass}(P,y.1)
 \exp\!\left(t\operatorname{integratedEntropyPos}_X
   (\operatorname{probMeasure}(\tau),
    \operatorname{conditional}(P,y))\right).
\]
Then the literal conjunction
\begin{align*}
 0<S_{t,\tau}(P)\quad&\wedge\quad
 0\leq\operatorname{HTemp}(t,h_t,\tau)(P)\\
 &\wedge\quad0\leq\operatorname{HZero}(\sigma)(P)\\
 &\wedge\quad0\leq\operatorname{HMinus}(\tau)(P)\\
 &\wedge\quad0\leq\operatorname{HPlus}(\tau)(P)
\end{align*}
holds.  The minimum and maximum are the already total
\texttt{finMin}/\texttt{finMax} on the locally nonempty type
\texttt{Active P}; no ``well-definedness'' predicate is needed.
\end{lemma}

\begin{proof}
The active set is finite and nonempty by
\cref{lem:fnd-positive-columns-nonempty}.  Each integrated entropy is finite
and nonnegative.  The exponential sum has strictly positive terms and hence
is positive.  To prove nonnegativity for finite $t$, let
$a_y=\A_\tau(P_{X\mid y})\geq0$.  If $t>0$, then
$\sum_ym_ye^{ta_y}\geq\sum_ym_y=1$.  If $t<0$, then the sum is at most one;
division of its nonpositive logarithm by $t<0$ is nonnegative.
\end{proof}

\begin{lemma}[Finite extrema algebra]
\label{lem:fnd-finite-extrema}
Let $J,K$ be finite nonempty types.  For real families $a,b:J\to\R$,
and independently typed real families $c:J\to\R$ and $d:K\to\R$,
the theorem returned to Lean is the literal four-way conjunction
\begin{align*}
 &\abs{\operatorname{finMin}(a)-\operatorname{finMin}(b)}
   \leq \operatorname{finMax}
      \bigl(j\mapsto\abs{a(j)-b(j)}\bigr)\\
 {}\wedge{}\;&
 \abs{\operatorname{finMax}(a)-\operatorname{finMax}(b)}
   \leq \operatorname{finMax}
      \bigl(j\mapsto\abs{a(j)-b(j)}\bigr)\\
 {}\wedge{}\;&
 \operatorname{finMin}
   \bigl(jk:J\times K\mapsto c(jk.1)+d(jk.2)\bigr)
 =\operatorname{finMin}(c)+\operatorname{finMin}(d)\\
 {}\wedge{}\;&
 \operatorname{finMax}
   \bigl(jk:J\times K\mapsto c(jk.1)+d(jk.2)\bigr)
 =\operatorname{finMax}(c)+\operatorname{finMax}(d).
\end{align*}
\end{lemma}

\begin{proof}
For every $j$, $a_j\geq b_j-\max\abs{a-b}$ and
$a_j\leq b_j+\max\abs{a-b}$; take minima or maxima and reverse the roles.
For the product identities, the right-hand side is a lower, respectively
upper, bound for all pairs and is attained by a pair of finite-set
minimizers, respectively maximizers.
\end{proof}

\begin{definition}[Finite log-sum-exp]
\label{def:fnd-log-sum-exp}
Let $J$ be a finite type and let $p,a:J\to\R$.  Define the total function
\[
 \operatorname{logSumExp}(p,a,t):=
 \frac1t\log\sum_{j:J}p_j\exp(t a_j).
\]
Its value at $t=0$ is the library's total-division value and is never used;
the limit below is taken within the punctured neighbourhood.
\end{definition}

\begin{lemma}[Log-sum-exp bounds and limits]
\label{lem:fnd-log-sum-exp}
Let $J$ be finite and nonempty, let $p,a:J\to\R$ satisfy
$h_p:\forall j,0<p_j$ and $h_{\rm sum}:\sum_jp_j=1$.
Let $j_-,j_+:J$ satisfy
$h_-:a_{j_-}=\operatorname{finMin}(a)$ and
$h_+:a_{j_+}=\operatorname{finMax}(a)$.  Its exact proposition is
\begin{align*}
 &\Bigl(\forall t:\R,\ t<0\to
   \bigl(\operatorname{finMin}(a)
      \leq\operatorname{logSumExp}(p,a,t)\\[-2pt]
 &\hspace{115pt}{}\wedge
      \operatorname{logSumExp}(p,a,t)
      \leq\operatorname{finMin}(a)+\tfrac{\log(p(j_-))}{t}\bigr)\Bigr)\\
 {}\wedge{}\;&\Bigl(\forall t:\R,\ 0<t\to
   \bigl(\operatorname{finMax}(a)+\tfrac{\log(p(j_+))}{t}
      \leq\operatorname{logSumExp}(p,a,t)\\[-2pt]
 &\hspace{115pt}{}\wedge
      \operatorname{logSumExp}(p,a,t)\leq\operatorname{finMax}(a)\bigr)\Bigr)\\
 {}\wedge{}\;&
 \operatorname{Tendsto}
   \bigl(t\mapsto\operatorname{logSumExp}(p,a,t)\bigr)
   \texttt{ atBot }\,\mathcal N\bigl(\operatorname{finMin}(a)\bigr)\\
 {}\wedge{}\;&
 \operatorname{Tendsto}
   \bigl(t\mapsto\operatorname{logSumExp}(p,a,t)\bigr)
   \texttt{ atTop }\,\mathcal N\bigl(\operatorname{finMax}(a)\bigr)\\
 {}\wedge{}\;&
 \operatorname{Tendsto}
   \bigl(t\mapsto\operatorname{logSumExp}(p,a,t)\bigr)
   \bigl(\operatorname{nhdsWithin}\,(0:\R)\,
      ((\{0\}:\operatorname{Set}(\R))^{\mathrm c})\bigr)
   \mathcal N\!\left(\sum_{j:J}p(j)a(j)\right).
\end{align*}
\end{lemma}

\begin{proof}
For $t<0$, use
$p_{j_-}e^{tm}\leq\sum_jp_je^{ta_j}\leq e^{tm}$, take logarithms, and divide
by the negative number $t$.  The positive case uses
$p_{j_+}e^{tM}\leq\sum_jp_je^{ta_j}\leq e^{tM}$ and division by $t>0$.
The error terms tend to zero.  For the limit at zero put
$Z(t)=\sum_jp_je^{ta_j}$.  Then $Z$ is differentiable,
$Z(0)=1$, and $Z'(0)=\sum_jp_ja_j$.  The fundamental theorem of calculus
gives, for $t\neq0$,
\[
 \frac{\log Z(t)}t
 =\int_0^1\frac{Z'(st)}{Z(st)}\dd s.
\]
The integrand converges uniformly on the compact $s$-interval to $Z'(0)$,
because $Z$ and $Z'$ are continuous and $Z$ stays positive.  Taking the
limit proves the last formula without invoking l'H\^opital's rule.
\end{proof}

\begin{definition}[Joint embedding invariance]
\label{def:fnd-joint-embedding-invariant}
For $F:\texttt{PolyJointFunctional.\{u\}}$, define
\begin{align*}
 \operatorname{JointEmbeddingInvariant}(F):\Longleftrightarrow{}
 &\forall\{X,Y,X',Y':\texttt{Type u}\},\
 \forall[\operatorname{Fintype}X][\operatorname{Nonempty}X]
 [\operatorname{Fintype}Y][\operatorname{Nonempty}Y]\
 [\operatorname{Fintype}X'][\operatorname{Nonempty}X']
 [\operatorname{Fintype}Y'][\operatorname{Nonempty}Y'],\\[-2pt]
 &\forall(e_X:X\hookrightarrow X')(e_Y:Y\hookrightarrow Y')
 (P:\operatorname{JointProb}(X,Y)),\quad
 F(\operatorname{jointZeroExtend}(e_X,e_Y,P))=F(P).
\end{align*}
\end{definition}

\begin{definition}[Joint tensor additivity]
\label{def:fnd-joint-tensor-additive}
For $F:\texttt{PolyJointFunctional.\{u\}}$, define
\begin{align*}
 \operatorname{JointTensorAdditive}(F):\Longleftrightarrow{}
 &\forall\{X,Y,X',Y':\texttt{Type u}\},\
 \forall[\operatorname{Fintype}X][\operatorname{Nonempty}X]
 [\operatorname{Fintype}Y][\operatorname{Nonempty}Y]\
 [\operatorname{Fintype}X'][\operatorname{Nonempty}X']
 [\operatorname{Fintype}Y'][\operatorname{Nonempty}Y'],\\[-2pt]
 &\forall(P:\operatorname{JointProb}(X,Y))
 (Q:\operatorname{JointProb}(X',Y')),\quad
 F(\operatorname{jointTensor}(P,Q))=F(P)+F(Q).
\end{align*}
\end{definition}

\begin{definition}[Uniform normalization]
\label{def:fnd-uniform-normalized}
For $F:\texttt{PolyJointFunctional.\{u\}}$, define
\begin{align*}
 \operatorname{UniformNormalized}(F):\Longleftrightarrow{}
 &\forall(d:\N)(h_d:0<d),\
 \forall\{Y:\texttt{Type u}\}
 [\operatorname{Fintype}Y][\operatorname{Nonempty}Y],\
 \forall q:\operatorname{ProbVec}(Y),\\[-2pt]
 &\texttt{letI : Nonempty (Fin d) := }
 \langle\langle0,h_d\rangle\rangle\texttt{; }\quad
 F\bigl(\operatorname{independentJoint}
   (\operatorname{uniformProb}_{\operatorname{Fin}(d)},q)\bigr)
 =\log(d:\R).
\end{align*}
\end{definition}

\begin{definition}[Elementary polymorphic functional predicates]
\label{def:fnd-easy-functional-predicates}
For $F:\texttt{PolyJointFunctional.\{u\}}$, define the single proposition
\[
 \operatorname{EasyConditionalAxioms}(F):\Longleftrightarrow
 \operatorname{JointEmbeddingInvariant}(F)\ \wedge\
 \operatorname{JointTensorAdditive}(F)\ \wedge\
 \operatorname{UniformNormalized}(F).
\]
Every binder of the three component predicates is literal in
\cref{def:fnd-joint-embedding-invariant,def:fnd-joint-tensor-additive,def:fnd-uniform-normalized};
no alphabet or instance is supplied by prose.
\end{definition}

\begin{proposition}[Easy conditional-entropy axioms]
\label{prop:fnd-easy-axioms}
Let $\tau:\operatorname{ProbabilityMeasure}(\overline\R_+)$, let
$\sigma:\operatorname{FiniteMeasure}(\overline\R_+)$, and let $t:\R$ have
proof $h_t:t\neq0$.  Then the following single literal conjunction holds:
\begin{align*}
 &\operatorname{EasyConditionalAxioms}(\operatorname{HTemp}(t,h_t,\tau))\\
 &\quad\wedge
 \operatorname{EasyConditionalAxioms}
   (\operatorname{HZero}(\tau.\operatorname{toFiniteMeasure}))\\
 &\quad\wedge
 \operatorname{EasyConditionalAxioms}(\operatorname{HMinus}(\tau))\\
 &\quad\wedge
 \operatorname{EasyConditionalAxioms}(\operatorname{HPlus}(\tau))\\
 &\quad\wedge
 \operatorname{JointEmbeddingInvariant}(\operatorname{HZero}(\sigma))\\
 &\quad\wedge
 \operatorname{JointTensorAdditive}(\operatorname{HZero}(\sigma))\\
 &\quad\wedge
 \Bigl(\forall(d:\N)(h_d:0<d),\
 \forall\{Y:\texttt{Type u}\}
 [\operatorname{Fintype}Y][\operatorname{Nonempty}Y],\
 \forall q:\operatorname{ProbVec}(Y),\\[-2pt]
 &\hspace{17mm}
 \texttt{letI : Nonempty (Fin d) := }
 \langle\langle0,h_d\rangle\rangle\texttt{; }\quad
 \operatorname{HZero}(\sigma)
 \bigl(\operatorname{independentJoint}
   (\operatorname{uniformProb}_{\operatorname{Fin}(d)},q)\bigr)\\[-2pt]
 &\hspace{49mm}
 =\operatorname{ENNReal.toReal}
   (\operatorname{finiteMeasure}(\sigma)(\operatorname{Set.univ}))\log(d:\R)
 \Bigr).
\end{align*}
\end{proposition}

\begin{proof}
Relabeling and embedding preserve the active column masses and conditionals
up to relabeling and zero embedding; apply
\cref{lem:fnd-renyi-embedding}.  For tensor products, active outcomes are
pairs, their masses multiply, and their conditional distributions tensor.
Use \cref{lem:fnd-renyi-tensor}.  In the finite-$t$ case the exponential sum
factors into a product and the logarithm turns it into a sum.  The derivation
case uses finite distributivity.  The tropical cases use
\cref{lem:fnd-finite-extrema}.  The uniform normalization follows from
\cref{lem:fnd-uniform-renyi} and the fact that the column masses sum to one.
For finite positive $\sigma$, integration of the constant function
$\alpha\mapsto\log d$ gives its finite real total mass times $\log d$.
\end{proof}

\subsection{Column functions and finite mixing rules}

For a finite nonempty type $I$, and using the proof-carrying constructors
from Section~\ref{sec:formalization-contract}, define
\begin{align}
 \operatorname{PosHomOne}(F)
 &:\Longleftrightarrow
 \forall x:\operatorname{ConeVec}(I),\ \forall c:\R,\
 \forall h_c:0\leq c,\quad
 F(\operatorname{coneScale}(c,h_c,x))=cF(x),
 \label{eq:fnd-pos-hom-one}\\
 \operatorname{ScaleInvariant}(f)
 &:\Longleftrightarrow
 \forall x:\operatorname{PosConeVec}(I),\ \forall c:\R,\
 \forall h_c:0<c,\quad
 f(\operatorname{posScale}(c,h_c,x))=f(x),
 \label{eq:fnd-scale-invariant}\\
 \operatorname{QCvx}(f)
 &:\Longleftrightarrow
 \forall x,z:\operatorname{PosConeVec}(I),\ \forall\lambda:\R,\
 \forall h_\lambda:0<\lambda\ \wedge\ \lambda<1,\quad
 f(\operatorname{posMix}(\lambda,h_\lambda,x,z))
 \leq\max\{f(x),f(z)\},
 \label{eq:fnd-quasiconvex}\\
 \operatorname{MaxQCave}(f)
 &:\Longleftrightarrow
 \forall x,z:\operatorname{PosConeVec}(I),\ \forall\lambda:\R,\
 \forall h_\lambda:0<\lambda\ \wedge\ \lambda<1,\quad
 f(\operatorname{posMix}(\lambda,h_\lambda,x,z))
 \geq\max\{f(x),f(z)\}.
 \label{eq:fnd-max-qcave}
\end{align}
Here $F:\operatorname{ConeVec}(I)\to\R$ and
$f:\operatorname{PosConeVec}(I)\to\R$.  The names are intentionally
distinct from any library predicate named \emph{quasi-concave}.
For $F:\operatorname{ConeVec}(I)\to\R$, also define
\begin{align}
 \operatorname{ConcaveCone}_I(F)
 &: \Longleftrightarrow
 \forall x,z:\operatorname{ConeVec}(I),\
 \forall\lambda:\R,\ \forall h_\lambda:
    0\leq\lambda\ \wedge\ \lambda\leq1,\quad
 F(\operatorname{coneMix}(\lambda,h_\lambda,x,z))
 \geq\lambda F(x)+(1-\lambda)F(z),
 \label{eq:fnd-concave-cone}\\
 \operatorname{ConvexCone}_I(F)
 &: \Longleftrightarrow
 \forall x,z:\operatorname{ConeVec}(I),\
 \forall\lambda:\R,\ \forall h_\lambda:
    0\leq\lambda\ \wedge\ \lambda\leq1,\quad
 F(\operatorname{coneMix}(\lambda,h_\lambda,x,z))
 \leq\lambda F(x)+(1-\lambda)F(z).
 \label{eq:fnd-convex-cone}
\end{align}

For each finite nonempty type $I$, introduce four genuine Lean identifiers;
their exact signatures are split according to their measure domains:
\begin{align*}
 \operatorname{columnPhi}_I
   &:\R\to\operatorname{ProbabilityMeasure}(\overline\R_+)
       \to\operatorname{ConeVec}(I)\to\R,\\
 \operatorname{aTrop}_I,\operatorname{gTrop}_I
   &:\operatorname{ProbabilityMeasure}(\overline\R_+)
       \to\operatorname{PosConeVec}(I)\to\R,\\
 \operatorname{columnDeriv}_I
   &:\operatorname{FiniteMeasure}(\overline\R_+)
       \to\operatorname{ConeVec}(I)\to\R.
\end{align*}
Here \texttt{FiniteMeasure} means a finite positive Borel measure, as above.
For $t:\R$, a Borel probability measure $\tau$, a finite positive Borel
measure $\sigma$, and $x:\operatorname{ConeVec}(I)$, define
\begin{align}
 \operatorname{columnPhi}_I(t,\tau,x)
 &:=
 \begin{cases}
  \norm{x}_1\exp(t\A_\tau(\widehat x)),&x\neq0,\\
  0,&x=0,
 \end{cases}\label{eq:fnd-column-Phi}\\
 \operatorname{columnDeriv}_I(\sigma,x)
 &:=
 \begin{cases}
  \norm{x}_1\A_\sigma(\widehat x),&x\neq0,\\
  0,&x=0,
 \end{cases}\label{eq:fnd-column-D}\\
 \operatorname{aTrop}_I(\tau,z)&:=\A_\tau(\widehat z),\qquad
 \operatorname{gTrop}_I(\tau,z):=-\A_\tau(\widehat z)
 \qquad(z:\operatorname{PosConeVec}(I)).
 \label{eq:fnd-column-tropical}
\end{align}
For readability, install only the following local notations after these
definitions:
\[
 \Phi_{t,\tau}(x):=\operatorname{columnPhi}_I(t,\tau,x),\quad
 D_\sigma(x):=\operatorname{columnDeriv}_I(\sigma,x),\quad
 a_\tau(z):=\operatorname{aTrop}_I(\tau,z),\quad
 g_\tau(z):=\operatorname{gTrop}_I(\tau,z).
\]
Every later occurrence of the four paper symbols expands to these declared
identifiers; no unindexed constant named $\Phi$, $D$, $a$, or $g$ is used.
In either nonzero branch of $\operatorname{columnPhi}_I$ or
$\operatorname{columnDeriv}_I$, the displayed normalization is
literally \texttt{normalize (toPosCone x h)}, where $h:x\neq0$ is the
branch proof.  Likewise, later notation $a_\tau(x)$ or $g_\tau(x)$ for an
unbundled nonzero cone vector explicitly means application to
\texttt{toPosCone x h}; no coercion is inferred from a proposition.

\begin{lemma}[Column identities and homogeneity]
\label{lem:fnd-column-identities}
Let $I,J$ be finite nonempty types.
The exact theorem is the following conjunction, with all measure and proof
arguments exposed:
\begin{align*}
 &\Bigl(\forall(\tau:\operatorname{ProbabilityMeasure}(\overline\R_+))
                 (t:\R),\quad
    \operatorname{PosHomOne}(\Phi_{t,\tau})
    \wedge\operatorname{ScaleInvariant}(a_\tau)
    \wedge\operatorname{ScaleInvariant}(g_\tau)\\[-2pt]
 &\hspace{45pt}{}\wedge
    \forall(h_t:t\neq0)(P:\operatorname{JointProb}(I,J)),\quad
    \sum_{j:J}\Phi_{t,\tau}(\operatorname{column}(P,j))
     =\exp\bigl(t\operatorname{HTemp}(t,h_t,\tau)(P)\bigr)\Bigr)\\
 {}\wedge{}\;&
 \Bigl(\forall(\sigma:\operatorname{FiniteMeasure}(\overline\R_+)),\quad
    \operatorname{PosHomOne}(D_\sigma)
    \wedge
    \forall P:\operatorname{JointProb}(I,J),\quad
    \sum_{j:J}D_\sigma(\operatorname{column}(P,j))
      =\operatorname{HZero}(\sigma)(P)\Bigr).
\end{align*}
Zero columns contribute zero to both displayed sums; the tropical extrema
elsewhere range over the already bundled active subtype.
\end{lemma}

\begin{proof}
Use \cref{lem:fnd-normalization-package} and the definitions.  In the column
sum, a nonzero column has mass $m_y$ and normalized vector
$P_{X\mid y}$.  A zero column is identically zero by
\cref{lem:fnd-nonzero-mass}.  Define
$S:=\sum_{y\in Y_+(P)}m_y\exp(t\A_\tau(P_{X\mid y}))$.
The active set is nonempty, every active mass and every exponential are
strictly positive, so obtain the named proof $h_S:0<S$.
The definition gives $tH_{t,\tau}(P)=\log S$ using $h_t:t\neq0$; applying
\texttt{Real.exp\_log} with the proof $S>0$ yields the first identity.
The second identity is finite distributivity of the column masses through
the defining integral for $\A_\sigma$.
\end{proof}

\begin{lemma}[Homogeneous curvature and finite sums]
\label{lem:fnd-homogeneous-finite-sums}
Let $I$ be a finite nonempty type, let $J$ be a finite type, and let
$F:\operatorname{ConeVec}(I)\to\R$ have proof
$h_F:\operatorname{PosHomOne}(F)$, and let
$z:J\to\operatorname{ConeVec}(I)$.
For $c:J\to\R$ and $h_c:\forall j,0\leq c(j)$, the exact return proposition
is the four-way conjunction
\begin{align*}
 &\Bigl(\operatorname{ConcaveCone}_I(F)\to
   \sum_{j:J}F(z(j))\leq F(\operatorname{coneSum}(z))\Bigr)\\
 {}\wedge{}\;&
 \Bigl(\operatorname{ConvexCone}_I(F)\wedge F(0)=0\to
   F(\operatorname{coneSum}(z))\leq\sum_{j:J}F(z(j))\Bigr)\\
 {}\wedge{}\;&
 \Bigl(\operatorname{ConcaveCone}_I(F)\to
   \sum_{j:J}c(j)F(z(j))
    \leq F(\operatorname{weightedConeSum}(c,h_c,z))\Bigr)\\
 {}\wedge{}\;&
 \Bigl(\operatorname{ConvexCone}_I(F)\wedge F(0)=0\to
   F(\operatorname{weightedConeSum}(c,h_c,z))
    \leq\sum_{j:J}c(j)F(z(j))\Bigr).
\end{align*}
\end{lemma}

\begin{proof}
For two vectors, apply curvature at their midpoint and multiply by two.
Induct on $n$.  Zero terms follow from $F(0)=0$, which follows from
homogeneity in the concave case.
\end{proof}

\begin{lemma}[Finite quasi mixing]
\label{lem:fnd-finite-quasi-mixing}
Let $I,J$ be finite nonempty types, let
$z:J\to\operatorname{PosConeVec}(I)$, let $c:J\to\R$ satisfy
$h_c:\forall j,0<c_j$, and put
\[
 s:=\operatorname{weightedConeSum}
 \bigl(c,j\mapsto\operatorname{le_of_lt}(h_c(j)),j\mapsto z_j.1\bigr).
\]
Let $h_s:s\neq0$ be an explicit proof argument (it can be constructed from
\cref{lem:fnd-nonzero-sum-witness}), and put
$s^+:=\operatorname{toPosCone}(s,h_s)$.  Let
$f:\operatorname{PosConeVec}(I)\to\R$.
Then the exact theorem is
\begin{align*}
 &\Bigl(\operatorname{ScaleInvariant}(f)\wedge\operatorname{QCvx}(f)
   \to f(s^+)\leq\operatorname{finMax}(j\mapsto f(z(j)))\Bigr)\\
 {}\wedge{}\;&
 \Bigl(\operatorname{ScaleInvariant}(f)\wedge\operatorname{MaxQCave}(f)
   \to \operatorname{finMax}(j\mapsto f(z(j)))\leq f(s^+)\Bigr).
\end{align*}
\end{lemma}

\begin{proof}
Normalize the positive coefficients by their sum; scale invariance deletes
the common scalar.  Induct on the number of summands using the binary
definition.
\end{proof}

\begin{lemma}[Nonempty weighted support]
\label{lem:fnd-weighted-support-nonempty}
Let $I,J$ be finite nonempty types, let
$z:J\to\operatorname{ConeVec}(I)$, let $c:J\to\R$ satisfy
$h_c:\forall j,0\leq c(j)$, and put
\[
 s:=\operatorname{weightedConeSum}(c,h_c,z).
\]
Let $h_s:s\neq0$ be an explicit proof argument, and put
\[
 K:=\{j:J\mathbin{//}0<c(j)\wedge z(j)\neq0\}.
\]
Then the exact conclusion is
\[
 \operatorname{weightedSupportNonempty}(c,h_c,z,h_s):
 \operatorname{Nonempty}(K).
\]
\end{lemma}

\begin{proof}
If $K$ were empty, every weighted summand would be zero, contradicting
$h_s$ and \cref{lem:fnd-nonzero-sum-witness}.
\end{proof}

\begin{lemma}[Finite quasi mixing after deletion of zero summands]
\label{lem:fnd-finite-quasi-mixing-delete-zero}
Let $I,J$ be finite nonempty types, let
$z:J\to\operatorname{ConeVec}(I)$, let $c:J\to\R$ satisfy
$h_c:\forall j,0\leq c(j)$, and put
\[
 s:=\operatorname{weightedConeSum}(c,h_c,z).
\]
Let $h_s:s\neq0$ be an explicit proof argument and put
\[
 K:=\{j:J\mathbin{//}0<c(j)\wedge z(j)\neq0\}.
\]
For every $f:\operatorname{PosConeVec}(I)\to\R$, the exact conclusion is
\begin{align*}
&\texttt{letI : Nonempty K :=
 weightedSupportNonempty c h\_c z h\_s; }\\[-2pt]
&\texttt{let }z^+:=
 (k\mapsto\operatorname{toPosCone}(z(k.1),k.2.2))\texttt{; }\quad
 \texttt{let }s^+:=\operatorname{toPosCone}(s,h_s)\texttt{; }\\[-2pt]
 &\Bigl(\operatorname{ScaleInvariant}(f)\wedge\operatorname{QCvx}(f)
   \to f(s^+)\leq\operatorname{finMax}(k\mapsto f(z^+(k)))\Bigr)\\
 {}\wedge{}\;&
 \Bigl(\operatorname{ScaleInvariant}(f)\wedge\operatorname{MaxQCave}(f)
   \to \operatorname{finMax}(k\mapsto f(z^+(k)))\leq f(s^+)\Bigr).
\end{align*}
\end{lemma}

\begin{proof}
Delete the zero summands, use the resulting positive coefficients, and apply
\cref{lem:fnd-finite-quasi-mixing}.  The prior named lemma supplies the
instance required to elaborate both finite maxima.
\end{proof}

\subsection{Perspective and preservation of curvature}

\begin{definition}[Total simplex perspective]
\label{def:fnd-perspective}
Let $I$ be finite and nonempty.  For
$f:\operatorname{ProbVec}(I)\to\R$, define the total function
\[
 \operatorname{perspective}_I(f):\operatorname{ConeVec}(I)\to\R,
 \qquad
 \operatorname{perspective}_I(f)(x):=
 \begin{cases}
  \norm{x}_1 f\bigl(\operatorname{normalize}
    (\operatorname{toPosCone}(x,h_x))\bigr),&h_x:x\neq0,\\
  0,&x=0.
 \end{cases}
\]
\end{definition}

\begin{lemma}[Perspective of a concave simplex function]
\label{lem:fnd-perspective}
Let $I$ be finite and nonempty and let
$f:\operatorname{ProbVec}(I)\to\R$.  If
$\operatorname{SimplexConcave}_I(f)$, then
\[
 \operatorname{ConcaveCone}_I(\operatorname{perspective}_I(f)).
\]
\end{lemma}

\begin{proof}
Let $x,z\neq0$, $0<\lambda<1$, and
$s=\lambda\norm{x}_1+(1-\lambda)\norm{z}_1>0$.  Put
$\theta=\lambda\norm{x}_1/s$.  Coordinatewise algebra gives
\[
 \widehat{\lambda x+(1-\lambda)z}
 =\theta\widehat x+(1-\theta)\widehat z.
\]
Apply concavity of $f$ and multiply by $s$.  If either vector is zero, the
claim follows from degree-one homogeneity.
\end{proof}

\begin{definition}[Integrated cone family]
\label{def:fnd-integrate-cone-family}
Let $I$ be finite and nonempty, let
$\nu:\operatorname{Measure}(\overline\R_+)$, and let
$F:\overline\R_+\to\operatorname{ConeVec}(I)\to\R$.  Define
\[
 \operatorname{integrateConeFamily}(\nu,F):
 \operatorname{ConeVec}(I)\to\R,
 \qquad
 \operatorname{integrateConeFamily}(\nu,F)(x):=
 \int F(a,x)\dd\nu(a).
\]
\end{definition}

\begin{lemma}[Positive integration preserves curvature]
\label{lem:fnd-integration-curvature}
Let $I$ be a finite nonempty type, let
$\nu:\operatorname{Measure}(\overline\R_+)$ carry an
$[\operatorname{IsFiniteMeasure}(\nu)]$ instance, and let
$F:\overline\R_+\to\operatorname{ConeVec}(I)\to\R$.  Assume
\[
 h_F:\forall x:\operatorname{ConeVec}(I),\quad
 \operatorname{Measurable}(a\mapsto F(a,x))\ \wedge\
 \operatorname{Integrable}(a\mapsto F(a,x),\nu).
\]
Then the following literal conjunction holds:
\begin{align*}
&\Bigl[
  (\forall^{\mathrm{ae}}a\ \partial\nu,\operatorname{ConcaveCone}_I(F(a)))
  \Longrightarrow
  \operatorname{ConcaveCone}_I(\operatorname{integrateConeFamily}(\nu,F))
 \Bigr]\\
&\qquad\wedge\
 \Bigl[
  (\forall^{\mathrm{ae}}a\ \partial\nu,\operatorname{ConvexCone}_I(F(a)))
  \Longrightarrow
  \operatorname{ConvexCone}_I(\operatorname{integrateConeFamily}(\nu,F))
 \Bigr].
\end{align*}
\end{lemma}

\begin{proof}
Integrate the pointwise curvature inequality.  Positivity of the measure
preserves its direction.
\end{proof}

\begin{definition}[Uniform pointwise error]
\label{def:fnd-uniform-point-error}
Let $E:\texttt{Type u}$, $K:\texttt{Type v}$, let
$f_n:\N\to E\to K\to\R$, and let $f:E\to K\to\R$.  Define
\[
 \operatorname{uniformPointError}(f_n,f,n,e):=
 \operatorname{sSup}
 \{r:\R\mid\exists u:K,\ r=\abs{f_n(n,e,u)-f(e,u)}\}.
\]
\end{definition}

\begin{definition}[Uniform positive-integral error]
\label{def:fnd-uniform-integral-error-pos}
Let $E:\texttt{Type u}$ carry $[\operatorname{MeasurableSpace}E]$,
$K:\texttt{Type v}$, let
$\nu:\operatorname{Measure}(E)$, let $f_n:\N\to E\to K\to\R$, and let
$f:E\to K\to\R$.  Define
\[
 \operatorname{uniformIntegralErrorPos}(\nu,f_n,f,n):=
 \operatorname{sSup}\left\{r:\R\ \middle|\ \exists u:K,\
 r=\abs{\int f_n(n,e,u)\dd\nu(e)-\int f(e,u)\dd\nu(e)}\right\}.
\]
\end{definition}

\begin{definition}[Uniform signed-integral error]
\label{def:fnd-uniform-integral-error-signed}
Let $E:\texttt{Type u}$ carry $[\operatorname{MeasurableSpace}E]$,
$K:\texttt{Type v}$, let
$\mu:\operatorname{SignedMeasure}(E)$, let
$f_n:\N\to E\to K\to\R$, and let $f:E\to K\to\R$.  Define
\[
 \operatorname{uniformIntegralErrorSigned}(\mu,f_n,f,n):=
 \operatorname{sSup}\left\{r:\R\ \middle|\ \exists u:K,\
 r=\abs{\operatorname{signedIntegral}(\mu,e\mapsto f_n(n,e,u))
       -\operatorname{signedIntegral}(\mu,e\mapsto f(e,u))}\right\}.
\]
\end{definition}

\begin{lemma}[Uniform dominated convergence for a positive measure]
\label{lem:fnd-uniform-dct}
Let $E:\texttt{Type u}$ carry $[\operatorname{MeasurableSpace}E]$, let
$K:\texttt{Type v}$ carry $[\operatorname{Nonempty}K]$, let
$\nu:\operatorname{Measure}(E)$ carry an
$[\operatorname{IsFiniteMeasure}(\nu)]$ instance, let
$f_n:\N\to E\to K\to\R$, let $f:E\to K\to\R$, and let $G:E\to\R$.
Assume
\begin{align*}
 &\forall n\,u,\quad
   \operatorname{Measurable}(e\mapsto f_n(n,e,u))\ \wedge\
   \operatorname{Integrable}(e\mapsto f_n(n,e,u),\nu),\\
 &\forall u,\quad
   \operatorname{Measurable}(e\mapsto f(e,u))\ \wedge\
   \operatorname{Integrable}(e\mapsto f(e,u),\nu),\\
 &\forall n\,e,\quad
   \operatorname{BddAbove}
    \{r:\R\mid\exists u:K,r=\abs{f_n(n,e,u)-f(e,u)}\},\\
 &\forall n,\quad
   \operatorname{Measurable}(\operatorname{uniformPointError}(f_n,f,n)),\\
 &\forall^{\mathrm{ae}}e\ \partial\nu,\quad
   \operatorname{Tendsto}
    (n\mapsto\operatorname{uniformPointError}(f_n,f,n,e))
    (\texttt{atTop})(\mathcal N(0)),\\
 &\operatorname{Integrable}(G,\nu),\qquad
   \forall n\,e,\quad
   0\leq\operatorname{uniformPointError}(f_n,f,n,e)\ \wedge\
   \operatorname{uniformPointError}(f_n,f,n,e)\leq G(e).
\end{align*}
Then
\begin{align*}
 &\Bigl(\forall n,\quad
  0\leq\operatorname{uniformIntegralErrorPos}(\nu,f_n,f,n)\ \wedge\
  \operatorname{uniformIntegralErrorPos}(\nu,f_n,f,n)
   \leq\int\operatorname{uniformPointError}(f_n,f,n,e)\dd\nu(e)\Bigr)\\
 &\quad\wedge\quad\operatorname{Tendsto}
   (n\mapsto\operatorname{uniformIntegralErrorPos}(\nu,f_n,f,n))
   (\texttt{atTop})(\mathcal N(0)).
\end{align*}
\end{lemma}

\begin{proof}
The triangle inequality gives the pointwise upper bound before taking the
supremum.  Dominated convergence applied to the measurable error functions
proves the limit.
\end{proof}

\begin{lemma}[Uniform dominated convergence for a signed measure]
\label{lem:fnd-uniform-dct-signed}
Let $E:\texttt{Type u}$ carry $[\operatorname{MeasurableSpace}E]$, let
$K:\texttt{Type v}$ carry $[\operatorname{Nonempty}K]$, let
$\mu:\operatorname{SignedMeasure}(E)$, let
$f_n:\N\to E\to K\to\R$, let $f:E\to K\to\R$, and let $G:E\to\R$.
Assume literally
\begin{align*}
 &\forall n\,u,\quad
   \operatorname{Measurable}(e\mapsto f_n(n,e,u))\ \wedge\
   \operatorname{Integrable}
    (e\mapsto f_n(n,e,u),\operatorname{signedPos}(\mu))\ \wedge\
   \operatorname{Integrable}
    (e\mapsto f_n(n,e,u),\operatorname{signedNeg}(\mu)),\\
 &\forall u,\quad
   \operatorname{Measurable}(e\mapsto f(e,u))\ \wedge\
   \operatorname{Integrable}
    (e\mapsto f(e,u),\operatorname{signedPos}(\mu))\ \wedge\
   \operatorname{Integrable}
    (e\mapsto f(e,u),\operatorname{signedNeg}(\mu)),\\
 &\forall n\,e,\quad
   \operatorname{BddAbove}
    \{r:\R\mid\exists u:K,r=\abs{f_n(n,e,u)-f(e,u)}\},\\
 &\forall n,\quad
   \operatorname{Measurable}(\operatorname{uniformPointError}(f_n,f,n)),\\
 &\forall^{\mathrm{ae}}e\ \partial\operatorname{signedTV}(\mu),\quad
   \operatorname{Tendsto}
    (n\mapsto\operatorname{uniformPointError}(f_n,f,n,e))
    (\texttt{atTop})(\mathcal N(0)),\\
 &\operatorname{Integrable}(G,\operatorname{signedTV}(\mu)),\qquad
   \forall n\,e,\quad
   0\leq\operatorname{uniformPointError}(f_n,f,n,e)\ \wedge\
   \operatorname{uniformPointError}(f_n,f,n,e)\leq G(e).
\end{align*}
Then
\begin{align*}
 &\Bigl(\forall n,\quad
  0\leq\operatorname{uniformIntegralErrorSigned}(\mu,f_n,f,n)\ \wedge\
  \operatorname{uniformIntegralErrorSigned}(\mu,f_n,f,n)
   \leq\int\operatorname{uniformPointError}(f_n,f,n,e)
          \dd\operatorname{signedTV}(\mu)(e)\Bigr)\\
 &\quad\wedge\quad\operatorname{Tendsto}
   (n\mapsto\operatorname{uniformIntegralErrorSigned}(\mu,f_n,f,n))
   (\texttt{atTop})(\mathcal N(0)).
\end{align*}
\end{lemma}

\begin{proof}
Apply the total-variation bound to each difference kernel and then use
dominated convergence with the finite positive measure
$\operatorname{signedTV}(\mu)$.
\end{proof}

\begin{lemma}[Measurability of a compact-parameter supremum]
\label{lem:fnd-compact-sup-measurable}
Let $E:\texttt{Type u}$ carry $[\operatorname{MeasurableSpace}E]$,
and let $K:\texttt{Type v}$ carry
$[\operatorname{TopologicalSpace}K][\operatorname{CompactSpace}K]
[\operatorname{PseudoMetrizableSpace}K]
[\operatorname{SecondCountableTopology}K][\operatorname{Nonempty}K]$.
Let $h:E\to K\to\R$ satisfy
\begin{align*}
 &\forall u:K,\quad\operatorname{Measurable}(e\mapsto h(e,u)),\\
 &\forall e:E,\quad\operatorname{Continuous}(h(e)),\\
 &\forall e:E,\quad
   \operatorname{BddAbove}(\operatorname{Set.range}(h(e))).
\end{align*}
Then the literal conjunction
\begin{align*}
 &\operatorname{Measurable}\!\left(
  e\mapsto\operatorname{sSup}(\operatorname{Set.range}(h(e)))\right)\\
 &\quad\wedge\quad
 \operatorname{Measurable}\!\left(
  e\mapsto\operatorname{sSup}
   (\operatorname{Set.range}(u\mapsto\abs{h(e,u)}))\right)
\end{align*}
holds.
\end{lemma}

\begin{proof}
Choose a countable dense subset $D\subseteq K$.  Continuity and compactness
give
\[
 \sup_{u\in K}h(e,u)=\sup_{u\in D}h(e,u).
\]
The right-hand side is a countable supremum of measurable real functions and
is therefore measurable.  Apply the same argument to the continuous function
$u\mapsto\abs{h(e,u)}$.
\end{proof}

\begin{lemma}[Finite pointwise limits preserve curvature]
\label{lem:fnd-pointwise-curvature-limit}
Let $I$ be a finite nonempty type, let
$F_n:\operatorname{ConeVec}(I)\to\R$ for $n:\N$, and let
$F:\operatorname{ConeVec}(I)\to\R$.  Assume the pointwise filter
statement
\[
 \forall x:\operatorname{ConeVec}(I),\qquad
 \operatorname{Tendsto}(n\mapsto F_n(x),\texttt{atTop},
                         \mathcal N(F(x))).
\]
Then the literal conjunction
\[
 \Bigl((\forall n:\N,\operatorname{ConcaveCone}_I(F_n))
       \Longrightarrow\operatorname{ConcaveCone}_I(F)\Bigr)
 \quad\wedge\quad
 \Bigl((\forall n:\N,\operatorname{ConvexCone}_I(F_n))
       \Longrightarrow\operatorname{ConvexCone}_I(F)\Bigr)
\]
holds.
\end{lemma}

\begin{proof}
Write the curvature inequality for each function and pass to the limit using
finite limit algebra and closedness of the real order.
\end{proof}

\begin{definition}[Second derivative operator]
\label{def:fnd-second-deriv}
$\operatorname{secondDeriv}(f,\lambda):=
\operatorname{deriv}(\operatorname{deriv}f)(\lambda)$ for
$f:\R\to\R$ and $\lambda:\R$.
\end{definition}

\begin{lemma}[Line curvature tests]
\label{lem:fnd-line-curvature}
Let $I:\operatorname{Set}(\R)$ and $f:\R\to\R$ have proofs
$h_I^{\rm open}:\operatorname{IsOpen}(I)$,
$h_I^{\rm cvx}:\operatorname{Convex}(\R,I)$, and
$h_f:\operatorname{ContDiffOn}(\R,2,f,I)$.
Then the literal conjunction
\begin{align*}
 &\Bigl(\operatorname{ConcaveOn}(\R,I,f)\Longrightarrow
   \forall\lambda\in I,\operatorname{secondDeriv}(f,\lambda)\leq0\Bigr)\\
 &\quad\wedge\quad
 \Bigl(\operatorname{ConvexOn}(\R,I,f)\Longrightarrow
   \forall\lambda\in I,0\leq\operatorname{secondDeriv}(f,\lambda)\Bigr)
\end{align*}
holds.
\end{lemma}

\begin{proof}
Apply the midpoint curvature inequality at $\lambda\pm h$, divide by $h^2$,
and take $h\to0$.
\end{proof}

\begin{lemma}[Eventual obstruction from a limiting second derivative]
\label{lem:fnd-line-curvature-eventual}
Let $I:\operatorname{Set}(\R)$ have proofs
$h_I^{\rm open}:\operatorname{IsOpen}(I)$,
$h_I^{\rm cvx}:\operatorname{Convex}(\R,I)$, and $h_0:0\in I$.  Let
$f_{\rm seq}:\N\to\R\to\R$ and $c:\R$ have proof
\[
 \forall n:\N,\quad
 \operatorname{ContDiffOn}(\R,2,f_{\rm seq}(n),I).
\]
Then
\begin{align*}
 &\Bigl(
  \operatorname{Tendsto}
   (n\mapsto\operatorname{secondDeriv}(f_{\rm seq}(n),0))
   (\texttt{atTop})(\mathcal N(c))\ \wedge\ 0<c
  \Longrightarrow
  \operatorname{Filter.Eventually}
   (n\mapsto\neg\operatorname{ConcaveOn}(\R,I,f_{\rm seq}(n)))
   (\texttt{Filter.atTop})\Bigr)\\
 &\quad\wedge\quad
 \Bigl(
  \operatorname{Tendsto}
   (n\mapsto\operatorname{secondDeriv}(f_{\rm seq}(n),0))
   (\texttt{atTop})(\mathcal N(c))\ \wedge\ c<0
  \Longrightarrow
  \operatorname{Filter.Eventually}
   (n\mapsto\neg\operatorname{ConvexOn}(\R,I,f_{\rm seq}(n)))
   (\texttt{Filter.atTop})\Bigr).
\end{align*}
\end{lemma}

\begin{proof}
Convergence and the strict sign of $c$ give the corresponding strict sign of
the second derivative eventually; apply
\cref{lem:fnd-line-curvature} to rule out the incompatible curvature.
\end{proof}

\setcounter{section}{2}
\section{Setup, conventions, and final statement}
\label{sec:setup}

This section begins with a notation-level synopsis of Sections~0--2 for the
reader of the classification proof.  It creates no second declaration for
any previously defined constant: every recalled symbol is an alias of the
uniquely labelled total definition above, and repeated formulas are
restatement or simplification lemmas.  The shape-reduction predicates,
admissibility predicates, and target proposition introduced later in this
section are new declarations and are labelled at their first occurrence.

\subsection{Basic spaces, measures, and asymptotic notation}

All alphabets are finite, and logarithms are natural. Changing the logarithm
base only rescales the parameter $t$ and does not change the classification.
We write $[0,+\infty]$ for the one-point compactification of $[0,+\infty)$,
equipped with its Borel $\sigma$-algebra.

For $d\in\N$ with $1\leq d$, the \emph{positive cone} is $\R^d_{>0}$, the
\emph{nonnegative cone} is $\R^d_{\geq0}$, and the punctured nonnegative cone
is
\[
 \R^d_{\geq0,\times}:=\R^d_{\geq0}\setminus\{0\}.
\]
Recall that $\one_N$ is the constant-one vector on $\operatorname{Fin}(N)$
and that $\col(z^{(1)},\ldots,z^{(m)})$ is the sigma-type concatenation fixed
in Section~\ref{sec:formalization-contract}.

A finite positive measure is a nonnegative countably additive Borel measure
with finite total mass. A finite negative measure is the negative of a finite
positive measure. A finite signed measure $\mu$ has a Jordan decomposition
$\mu=\mu^+-\mu^-$ and total variation
$\abs{\mu}:=\mu^++\mu^-$. Throughout this note,
\[
 \supp(\mu):=\supp\abs{\mu}
\]
for a signed measure. A point $r$ is called \emph{$\mu$-null} when
$\abs{\mu}(\{r\})=0$. Thus a null point need not lie outside the topological
support.  The term \emph{one-signed} has the definition fixed in
Section~\ref{sec:formalization-contract}.

If $T$ is a Borel map and $\tau$ is a Borel measure, the pushforward
$T_\#\tau$ is defined by
\[
 (T_\#\tau)(E):=\tau(T^{-1}(E))
\]
for every Borel set $E$.

Every convergence statement below is either pointwise or is written with an
explicit supremum over its parameter set.  In particular, no asymptotic
Landau symbol is needed in the formal theorem statements.

\subsection{R\'enyi entropies and conditional functionals}

For a probability vector $p=(p_i)_{i=1}^d$, recall
\begin{equation}
 H_\alpha(p)
 :=\frac{1}{1-\alpha}\log\sum_{i=1}^d p_i^\alpha,
 \qquad \alpha\in(0,+\infty)\setminus\{1\},
 \label{eq:renyi-definition}
\end{equation}
with the continuous endpoint values
\begin{equation}
 H_0(p)=\log\card\supp(p),\qquad
 H_1(p)=-\sum_i p_i\log p_i,\qquad
 H_\infty(p)=-\log\max_i p_i.
 \label{eq:renyi-endpoints}
\end{equation}
For every fixed finite-dimensional $p$, the map
$\alpha\mapsto H_\alpha(p)$ is continuous on $[0,+\infty]$ and
$0\leq H_\alpha(p)\leq\log d$.

Let $\tau$ be a Borel probability measure on $[0,+\infty]$. Recall
\begin{equation}
 \A_\tau(p):=\int_{[0,+\infty]}H_\alpha(p)\dd\tau(\alpha).
 \label{eq:A-tau}
\end{equation}
This integral is always finite for a finite-dimensional probability vector.
In the derivation discussion, the same notation is used when $\tau$ is an
arbitrary finite positive measure; normalization of that measure is irrelevant
to curvature and monotonicity.

For a joint probability distribution $P_{XY}$, the conditional distribution
$P_{X\mid Y=y}$ is used only when $P_Y(y)>0$. This convention is adopted in
every formula below. For $t\in\R\setminus\{0\}$, the earlier definition is
\begin{equation}
 H_{t,\tau}(X\mid Y)_P
 :=\frac1t\log\sum_{y:P_Y(y)>0}P_Y(y)
       \exp\!\left(t\A_\tau(P_{X\mid Y=y})\right).
 \label{eq:H-temperate}
\end{equation}
Its pointwise limits are
\begin{align}
 H_{0,\tau}(X\mid Y)_P
 &:=\sum_{y:P_Y(y)>0}P_Y(y)\A_\tau(P_{X\mid Y=y}),
 \label{eq:H-zero}\\
 H_{-\infty,\tau}(X\mid Y)_P
 &:=\min_{y:P_Y(y)>0}\A_\tau(P_{X\mid Y=y}),
 \label{eq:H-minus-infty}\\
 H_{+\infty,\tau}(X\mid Y)_P
 &:=\max_{y:P_Y(y)>0}\A_\tau(P_{X\mid Y=y}).
 \label{eq:H-plus-infty}
\end{align}
For a point measure $\tau=\delta_\alpha$, we also write
$H_{0,\alpha}:=H_{0,\delta_\alpha}$.

Every functional above is invariant under relabelling and additive under
tensor products. The issue studied here is monotonicity under conditionally
mixing channels.

\subsection{Shape properties and conditionally mixing channels}

A function $F$ on a convex set is quasi-convex if
\[
 F(\lambda x+(1-\lambda)z)\leq\max\{F(x),F(z)\}
 \qquad(0<\lambda<1),
\]
and is \emph{max-quasi-concave} if the reverse inequality holds with the same
maximum on the right-hand side.  This second property is exactly
$\operatorname{MaxQCave}(F)$ from \cref{eq:fnd-max-qcave}; it is sometimes
called strongly quasi-concave in the manuscript, but it is kept as a distinct
formal predicate to avoid confusion with mathlib's pre-existing notion.

We regard $P_{XY}$ as a nonnegative matrix whose columns are
\[
 x_y:=P_{X,Y=y}=P_Y(y)P_{X\mid Y=y}.
\]
A conditionally mixing channel has the form
\begin{equation}
 Q=\sum_{i=1}^m S_iPD_i,
 \label{eq:conditional-channel}
\end{equation}
where each $S_i$ is doubly stochastic on the $X$-system, each $D_i$ has
nonnegative entries, and
\begin{equation}
 \sum_{i=1}^m\sum_{y'}D_i(y,y')=1
 \qquad\text{for every input column }y.
 \label{eq:D-normalisation}
\end{equation}
The output column indexed by $y'$ is therefore
\[
 q_{y'}=\sum_{i,y}D_i(y,y')S_i x_y.
\]

For $x\in\R^d_{\geq0,\times}$, recall
\begin{equation}
 \widehat x:=\frac{x}{\norm{x}_1},\qquad
 \Phi_{t,\tau}(x):=\norm{x}_1
 \exp\!\left(t\A_\tau(\widehat x)\right),
 \label{eq:column-functional}
\end{equation}
and recall $\Phi_{t,\tau}(0)=0$. Also recall
\begin{equation}
 D_\tau(x):=\norm{x}_1\A_\tau(\widehat x)
 \quad(x\neq0),\qquad D_\tau(0):=0.
 \label{eq:derivation-column}
\end{equation}
Then
\begin{equation}
 \sum_y\Phi_{t,\tau}(x_y)
 =\exp\!\left(tH_{t,\tau}(X\mid Y)_P\right),
 \label{eq:matrix-column-sum}
\end{equation}
and
\[
H_{0,\tau}(X\mid Y)_P=\sum_{y:P_Y(y)>0}D_\tau(x_y).
\]

The predicates $\operatorname{ConcaveCone}_X$ and
$\operatorname{ConvexCone}_X$ used below are the uniquely labelled
definitions \cref{eq:fnd-concave-cone,eq:fnd-convex-cone}; they are not
redeclared here.

\begin{proposition}[Temperate and derivation shape reduction]
\label{prop:temperate-shape-reduction}
Let $X:\texttt{Type u}$ carry
$[\operatorname{Fintype}X][\operatorname{Nonempty}X]$, let
$\tau:\operatorname{ProbabilityMeasure}(\overline\R_+)$, let
$\sigma:\operatorname{FiniteMeasure}(\overline\R_+)$, and let $t:\R$.
Then the following literal conjunction holds:
\begin{align*}
 &\Bigl(\forall h_t:0<t,\quad
 \operatorname{CMMonotoneAtRow}
   (X,\operatorname{HTemp}(t,\operatorname{ne\_of\_gt}(h_t),\tau))
 \Longleftrightarrow
 \operatorname{ConcaveCone}_X(\Phi_{t,\tau})\Bigr)\\
 &\quad\wedge
 \Bigl(\forall h_t:t<0,\quad
 \operatorname{CMMonotoneAtRow}
   (X,\operatorname{HTemp}(t,\operatorname{ne\_of\_lt}(h_t),\tau))
 \Longleftrightarrow
 \operatorname{ConvexCone}_X(\Phi_{t,\tau})\Bigr)\\
 &\quad\wedge
 \Bigl(
 \operatorname{CMMonotoneAtRow}(X,\operatorname{HZero}(\sigma))
 \Longleftrightarrow
 \operatorname{ConcaveCone}_X(D_\sigma)\Bigr).
\end{align*}
\end{proposition}

\begin{proof}
The functions $\Phi_{t,\tau}$ and, independently, $D_\sigma$ are positively
homogeneous of degree one. Hence concavity implies superadditivity and convexity implies
subadditivity. Every R\'enyi entropy is Schur concave, so for a doubly
stochastic matrix $S$,
\[
 \Phi_{t,\tau}(Sx)\geq\Phi_{t,\tau}(x)\quad(t>0),
 \qquad
 \Phi_{t,\tau}(Sx)\leq\Phi_{t,\tau}(x)\quad(t<0),
\]
and $D_\sigma(Sx)\geq D_\sigma(x)$.

For $t>0$, concavity and \cref{eq:conditional-channel,eq:D-normalisation}
give
\begin{align*}
 \sum_{y'}\Phi_{t,\tau}(q_{y'})
 &\geq\sum_{i,y,y'}D_i(y,y')\Phi_{t,\tau}(S_i x_y)\\
 &\geq\sum_{i,y,y'}D_i(y,y')\Phi_{t,\tau}(x_y)
 =\sum_y\Phi_{t,\tau}(x_y).
\end{align*}
For $t<0$, convexity gives the reverse inequality. Combining this with
\cref{eq:matrix-column-sum}, and remembering that multiplication by $t<0$
reverses the logarithmic inequality, proves monotonicity. For the derivation,
concavity and \cref{eq:conditional-channel,eq:D-normalisation} give
\[
 \sum_{y'}D_\sigma(q_{y'})
 \geq\sum_{i,y,y'}D_i(y,y')D_\sigma(S_i x_y)
 \geq\sum_{i,y,y'}D_i(y,y')D_\sigma(x_y)
 =\sum_yD_\sigma(x_y),
\]
which is the required monotonicity inequality.

Conversely, start with two columns $\lambda x$ and $(1-\lambda)z$ and merge
them into $\lambda x+(1-\lambda)z$. If necessary, multiply the whole
matrix by one common scalar so that its total mass is one. Monotonicity gives
the desired inequality for the normalized matrix, and positive homogeneity
cancels the common scalar. Cases in which one vector vanishes follow from
positive homogeneity and the definitions at the origin.
\end{proof}

\begin{corollary}[Global temperate and derivation shape reduction]
\label{cor:global-temperate-shape-reduction}
Let $\tau:\operatorname{ProbabilityMeasure}(\overline\R_+)$, let
$\sigma:\operatorname{FiniteMeasure}(\overline\R_+)$, and let $t:\R$.
The following literal conjunction holds:
\begin{align*}
 &\Bigl(\forall h_t:0<t,\quad
 \operatorname{CMMonotone}
  (\operatorname{HTemp}(t,\operatorname{ne\_of\_gt}(h_t),\tau))
 \Longleftrightarrow
 \forall\{X:\texttt{Type u}\}
 [\operatorname{Fintype}X][\operatorname{Nonempty}X],\
 \operatorname{ConcaveCone}_X(\Phi_{t,\tau})\Bigr)\\
 &\quad\wedge
 \Bigl(\forall h_t:t<0,\quad
 \operatorname{CMMonotone}
  (\operatorname{HTemp}(t,\operatorname{ne\_of\_lt}(h_t),\tau))
 \Longleftrightarrow
 \forall\{X:\texttt{Type u}\}
 [\operatorname{Fintype}X][\operatorname{Nonempty}X],\
 \operatorname{ConvexCone}_X(\Phi_{t,\tau})\Bigr)\\
 &\quad\wedge
 \Bigl(\operatorname{CMMonotone}(\operatorname{HZero}(\sigma))
 \Longleftrightarrow
 \forall\{X:\texttt{Type u}\}
 [\operatorname{Fintype}X][\operatorname{Nonempty}X],\
 \operatorname{ConcaveCone}_X(D_\sigma)\Bigr).
\end{align*}
\end{corollary}

\begin{proof}
Unfold \texttt{CMMonotone}, commute its universal row-type binders with the
equivalence, and apply \cref{prop:temperate-shape-reduction} rowwise.
\end{proof}

For the tropical cases, recall the scale-invariant functions on the punctured
nonnegative cone
\begin{equation}
 g_\tau(x):=-\A_\tau(\widehat x),\qquad
 a_\tau(x):=\A_\tau(\widehat x).
 \label{eq:tropical-column-functions}
\end{equation}
They are not positively homogeneous of degree one; they are invariant under
multiplication of $x$ by a positive scalar.

\begin{proposition}[Tropical shape reduction]
\label{prop:tropical-shape-reduction}
Let $X:\texttt{Type u}$ carry
$[\operatorname{Fintype}X][\operatorname{Nonempty}X]$, and let
$\tau:\operatorname{ProbabilityMeasure}(\overline\R_+)$.  Then the literal
conjunction
\begin{align*}
 &\Bigl(
 \operatorname{CMMonotoneAtRow}(X,\operatorname{HMinus}(\tau))
 \Longleftrightarrow \operatorname{QCvx}(g_\tau)\Bigr)\\
 &\quad\wedge
 \Bigl(
 \operatorname{CMMonotoneAtRow}(X,\operatorname{HPlus}(\tau))
 \Longleftrightarrow \operatorname{MaxQCave}(a_\tau)\Bigr)
\end{align*}
holds.  Both predicates on the right are the total predicates on
$\operatorname{PosConeVec}(X)$ from
\cref{eq:fnd-quasiconvex,eq:fnd-max-qcave}; no zero-vector convention occurs
in the proposition.
\end{proposition}

\begin{proof}
Because $\A_\tau$ is Schur concave,
$g_\tau(Sx)\leq g_\tau(x)$ and $a_\tau(Sx)\geq a_\tau(x)$ for every doubly
stochastic $S$.

Assume first that $g_\tau$ is quasi-convex. Iterating quasi-convexity over the
nonzero summands of each nonzero output column gives
\[
 g_\tau(q_{y'})
 \leq\max_{\substack{i,y:D_i(y,y')>0\\x_y\neq0}}g_\tau(S_i x_y)
 \leq\max_{y:P_Y(y)>0}g_\tau(x_y).
\]
Taking the maximum over the nonzero output columns yields
$-H_{-\infty,\tau}(Q)\leq-H_{-\infty,\tau}(P)$, which is the desired
monotonicity. Conversely, merging two columns gives quasi-convexity; a common
normalization of the two-column matrix is harmless because $g_\tau$ is
scale-invariant.

Now assume that $a_\tau$ is max-quasi-concave and choose an input column
$x_{y_0}$ maximizing $a_\tau$. By \cref{eq:D-normalisation}, some pair
$(i,y')$ has $D_i(y_0,y')>0$.  The vector $S_ix_{y_0}$ is nonzero by mass
preservation, and all terms in $q_{y'}$ are nonnegative.  Therefore
$q_{y'}\neq0$ by \cref{lem:fnd-nonzero-sum-witness}.
Apply the second clause of \cref{lem:fnd-finite-quasi-mixing} to its nonzero
summands; scale invariance removes their positive coefficients.
Max-quasi-concavity then gives
\[
 a_\tau(q_{y'})\geq a_\tau(S_i x_{y_0})\geq a_\tau(x_{y_0}).
\]
Thus the maximum cannot decrease. The converse again follows by merging two
columns, with any common normalization removed by scale invariance.
\end{proof}

\begin{corollary}[Global tropical shape reduction]
\label{cor:global-tropical-shape-reduction}
For $\tau:\operatorname{ProbabilityMeasure}(\overline\R_+)$, the literal
conjunction
\begin{align*}
 &\Bigl(\operatorname{CMMonotone}(\operatorname{HMinus}(\tau))
 \Longleftrightarrow
 \forall\{X:\texttt{Type u}\}
 [\operatorname{Fintype}X][\operatorname{Nonempty}X],\
 \operatorname{QCvx}(g_\tau)\Bigr)\\
 &\quad\wedge
 \Bigl(\operatorname{CMMonotone}(\operatorname{HPlus}(\tau))
 \Longleftrightarrow
 \forall\{X:\texttt{Type u}\}
 [\operatorname{Fintype}X][\operatorname{Nonempty}X],\
 \operatorname{MaxQCave}(a_\tau)\Bigr)
\end{align*}
holds.
\end{corollary}

\begin{proof}
Unfold \texttt{CMMonotone} and apply
\cref{prop:tropical-shape-reduction} to each row type.
\end{proof}

\subsection{The singular weight and the exact parameter sets}

No constants are redeclared in this subsection.  We use the total Borel
function $\omega$ from \cref{eq:fnd-omega}, the nonnegative extended kernels
$\omega_<,\omega_>$ and their \texttt{lintegrals}
$M_<,M_>$ from
\cref{eq:fnd-M-lower,eq:fnd-M-upper}, the finiteness predicate
$\operatorname{MomFin}$ from \cref{eq:fnd-moment-finite}, and the total real
value
\[
 M_{\R}(\tau)=\operatorname{ENNReal.toReal}(M_<(\tau))
              -\operatorname{ENNReal.toReal}(M_>(\tau))
\]
from \cref{eq:fnd-M-real}.  We continue to abbreviate the last quantity by
$M(\tau)$.  Every substantive use of it is guarded by
$\operatorname{MomFin}(\tau)$.  Under that guard and
$\tau(\{1\})=0$, it equals the ordinary signed integral of $\omega$ away
from $1$ by \cref{lem:fnd-moment-algebra}.

\begin{definition}[Positive temperate admissibility]
\label{def:named-admissibility}
For $t:\R$ and
$\tau:\operatorname{ProbabilityMeasure}(\overline\R_+)$, define the exact
proposition
\begin{align}
 \operatorname{PosAdm}(t,\tau):\Longleftrightarrow{}
 &0<t\ \wedge\
 \operatorname{suppMeasure}(\operatorname{probMeasure}(\tau))
   \subseteq\operatorname{Set.Icc}(0,1)\ \wedge\
 \operatorname{probMeasure}(\tau)(\{1\})=0\ \wedge\nonumber\\[-2pt]
 &\operatorname{MLower}(\operatorname{probMeasure}(\tau))
   \leq\operatorname{ENNReal.ofReal}(1/t).
 \label{eq:Dbulk-positive}
\end{align}
\end{definition}

\begin{definition}[Negative lower-support admissibility]
\label{def:negative-lower-admissibility}
For $t:\R$ and
$\tau:\operatorname{ProbabilityMeasure}(\overline\R_+)$, define
\begin{equation}
 \operatorname{NegLowerAdm}(t,\tau):\Longleftrightarrow
 t<0\ \wedge\
 \operatorname{suppMeasure}(\operatorname{probMeasure}(\tau))
   \subseteq\operatorname{Set.Icc}(0,1).
 \label{eq:Dbulk-negative-lower}
\end{equation}
\end{definition}

\begin{definition}[Negative exceptional admissibility]
\label{def:negative-exceptional-admissibility}
For $t:\R$, $\tau:\operatorname{ProbabilityMeasure}(\overline\R_+)$, and
$a_*:\overline\R_+$, define
\begin{align}
 \operatorname{NegExcAdm}(t,\tau,a_*):\Longleftrightarrow{}
 &t<0\ \wedge\ 1<a_*\ \wedge\
 \operatorname{suppMeasure}(\operatorname{probMeasure}(\tau))
 \subseteq\operatorname{Set.Icc}(0,1)\cup\{a_*\}\ \wedge\nonumber\\[-2pt]
 &\operatorname{probMeasure}(\tau)(\{1\})=0\ \wedge\
 \operatorname{MLower}(\operatorname{probMeasure}(\tau))<\top\ \wedge\
 \operatorname{MomFin}(\operatorname{probMeasure}(\tau))\ \wedge\nonumber\\[-2pt]
 &\operatorname{MReal}(\operatorname{probMeasure}(\tau))\leq1/t.
 \label{eq:Dbulk-negative-exceptional}
\end{align}
The explicit \texttt{MomFin} conjunct is retained even though it follows
from the support and lower-moment conjuncts; it makes the guard on
\texttt{MReal} syntactically available after destructuring the predicate.
\end{definition}

\begin{definition}[Negative tropical exceptional admissibility]
\label{def:negative-tropical-exceptional-admissibility}
For $\tau:\operatorname{ProbabilityMeasure}(\overline\R_+)$ and
$a_*:\overline\R_+$, define
\begin{align}
 \operatorname{NegTropExcAdm}(\tau,a_*):\Longleftrightarrow{}
 &1<a_*\ \wedge\
 \operatorname{suppMeasure}(\operatorname{probMeasure}(\tau))
 \subseteq\operatorname{Set.Icc}(0,1)\cup\{a_*\}\ \wedge\nonumber\\[-2pt]
 &\operatorname{probMeasure}(\tau)(\{1\})=0\ \wedge\
 \operatorname{MLower}(\operatorname{probMeasure}(\tau))<\top\ \wedge\
 \operatorname{MomFin}(\operatorname{probMeasure}(\tau))\ \wedge\nonumber\\[-2pt]
 &\operatorname{MReal}(\operatorname{probMeasure}(\tau))\leq0.
 \label{eq:Dminus-exceptional}
\end{align}
\end{definition}

\begin{definition}[Temperate parameter set]\label{def:Dbulk}
Define the curried predicate
\[
 \operatorname{DBulk}:
 \R\to\operatorname{ProbabilityMeasure}(\overline\R_+)\to
 \operatorname{Prop}
\]
by the literal right-hand side
\[
 \operatorname{DBulk}(t,\tau):\Longleftrightarrow
 \operatorname{PosAdm}(t,\tau)\ \lor\
 \operatorname{NegLowerAdm}(t,\tau)\ \lor\
 \exists a_*:\overline\R_+,\operatorname{NegExcAdm}(t,\tau,a_*).
\]
There is no redundant outer $t\neq0$ conjunct: nonzeroness is a derived
fact from the sign field of whichever disjunct holds.  The notation
$(t,\tau)\in\Dbulk$ abbreviates \texttt{DBulk t tau}.
\end{definition}

\begin{lemma}[A temperate admissible parameter is nonzero]
\label{lem:Dbulk-ne}
For every $t:\R$ and
$\tau:\operatorname{ProbabilityMeasure}(\overline\R_+)$,
\[
 \operatorname{DBulk}(t,\tau)\Longrightarrow t\neq0.
\]
This declaration is named \texttt{DBulk.ne}.
\end{lemma}

\begin{proof}
Unfold \texttt{DBulk}; cases on the three defining disjuncts and use the
stored strict sign field in the selected admissibility predicate.
\end{proof}

\begin{definition}[Negative tropical parameter set]\label{def:Dminus}
Define
\[
 \operatorname{DMinus}:
 \operatorname{ProbabilityMeasure}(\overline\R_+)\to\operatorname{Prop}
\]
by
\[
 \operatorname{DMinus}(\tau):\Longleftrightarrow
 \operatorname{suppMeasure}(\operatorname{probMeasure}(\tau))
   \subseteq\operatorname{Set.Icc}(0,1)
 \ \lor\
 \exists a_*:\overline\R_+,\operatorname{NegTropExcAdm}(\tau,a_*).
\]
The notation $\tau\in\Dminus$ abbreviates \texttt{DMinus tau}.
\end{definition}

\begin{remark}[Atom versus accumulation at the Shannon point]
Whenever it is imposed, the condition $\tau(\{1\})=0$ excludes an atom at
$\alpha=1$ while still allowing the topological support to accumulate at
$1$. This is strictly weaker than $\supp(\tau)\subseteq[0,1)$, which would
also exclude such accumulation because the support is closed. In the positive
branch, $M_{>}(\tau)=0$, so \cref{eq:Dbulk-positive} is precisely the finite
moment condition on $[0,1)$.
\end{remark}

\begin{definition}[Target classification proposition]
\label{def:classification-target}
For universe $u$ and
$\tau:\operatorname{ProbabilityMeasure}(\overline\R_+)$, define the
single proposition.  The actual declaration header is
\[
 \texttt{def Classifies.\{u\}
  (tau : ProbabilityMeasure Param) : Prop :=}
\]
and its body is
\begin{align*}
 \texttt{Classifies.\{u\}}(\tau):\Longleftrightarrow{}
 &\Bigl(\forall(t:\R)(h_t:t\neq0),\quad
   \operatorname{CMMonotone}(\operatorname{HTemp}(t,h_t,\tau))
   \Longleftrightarrow\operatorname{DBulk}(t,\tau)\Bigr)\\
 &\wedge\Bigl(
   \operatorname{CMMonotone}(\operatorname{HMinus}(\tau))
   \Longleftrightarrow\operatorname{DMinus}(\tau)\Bigr)\\
 &\wedge\Bigl(
   \operatorname{CMMonotone}(\operatorname{HPlus}(\tau))
   \Longleftrightarrow\tau=\operatorname{diracProb}(0)\Bigr)\\
 &\wedge\Bigl(
   \operatorname{CMMonotone}
    (\operatorname{HZero}(\tau.\operatorname{toFiniteMeasure}))
   \Longleftrightarrow
   \operatorname{suppMeasure}(\operatorname{probMeasure}(\tau))
     \subseteq\operatorname{Set.Icc}(0,1)\Bigr).
\end{align*}
This is a definition of a proposition, not an asserted theorem.  Sufficiency
is proved in Section~4, necessity in Section~5, and the theorem
$\forall\tau,\operatorname{Classifies}(\tau)$ is declared only after both
dependency bundles are available.
\end{definition}

\section{Sufficient conditions for the parameters}
\label{sec:sufficiency}

\subsection{Finite-support measures and the power-mean factorisation}

For $\alpha\in(0,+\infty)$ and $x\in\R^d_{\geq0}$, write
\[
 \norm{x}_\alpha:=
 \operatorname{Real.rpow}\!\left(
   \sum_i\operatorname{Real.rpow}(x_i,\alpha),\frac1\alpha\right),
 \qquad \norm{x}_\infty:=\max_i x_i.
\]
Thus the power mean is total even when some $x_i=0$; every abbreviated
$x_i^\alpha$ in its proofs denotes \texttt{Real.rpow}.  Positivity of
$\alpha$ supplies all zero-base simplification side conditions.

For every finite nonempty type $I$, define the total compactified-order
function
\[
 \operatorname{parameterPowerMean}_I:
 \overline\R_+\to\operatorname{ConeVec}(I)\to\R
\]
by pattern matching on $\alpha:\operatorname{WithTop}(\R_{\geq0})$:
\[
 \operatorname{parameterPowerMean}_I(+\infty,x):=
  \operatorname{finMax}(i\mapsto x(i)),
\]
and, for a finite $r:\R_{\geq0}$,
\[
 \operatorname{parameterPowerMean}_I(r,x):=
 \begin{cases}
  \norm{x}_1,&r=0,\\
  \operatorname{Real.rpow}
   \left(\sum_{i:I}\operatorname{Real.rpow}(x_i,(r:\R)),
         1/(r:\R)\right),&r\neq0.
 \end{cases}
\]
The arbitrary total value at order zero is never used in the norm
factorisation: that argument carries the proof $\alpha\neq0,1$.  This
constructor nevertheless makes every product indexed by
$\alpha:\overline\R_+$ well typed before the endpoint hypotheses are used.

\begin{lemma}[Extension from the punctured cone]\label{lem:origin-extension}
Let $I$ be a finite nonempty type and let
$F:\operatorname{ConeVec}(I)\to\R$ satisfy $F(0)=0$ and
$\operatorname{PosHomOne}(F)$.  Then the following literal conjunction
holds:
\begin{align*}
 &\Bigl[
  \bigl(\forall x,z,\ x\neq0\to z\neq0\to
   \forall(\lambda:\R)(h_\lambda:\lambda\in\operatorname{Set.Icc}(0,1)),\
   F(\operatorname{coneMix}(\lambda,h_\lambda,x,z))
   \geq\lambda F(x)+(1-\lambda)F(z)\bigr)
  \Longrightarrow\operatorname{ConcaveCone}_I(F)\Bigr]\\
 &\quad\wedge\Bigl[
  \bigl(\forall x,z,\ x\neq0\to z\neq0\to
   \forall(\lambda:\R)(h_\lambda:\lambda\in\operatorname{Set.Icc}(0,1)),\
   F(\operatorname{coneMix}(\lambda,h_\lambda,x,z))
   \leq\lambda F(x)+(1-\lambda)F(z)\bigr)
  \Longrightarrow\operatorname{ConvexCone}_I(F)\Bigr].
\end{align*}
\end{lemma}

\begin{proof}
If $x,z\neq0$, there is nothing to prove. If, for example, $x=0$, then for
$0\leq\lambda\leq1$,
\[
 F(\lambda x+(1-\lambda)z)=F((1-\lambda)z)
 =(1-\lambda)F(z)=\lambda F(0)+(1-\lambda)F(z).
\]
The case $z=0$ follows after interchanging $x$ and $z$; if $x=z=0$, both
sides of the curvature inequality equal $F(0)=0$.
\end{proof}

\begin{lemma}[Convexity of the punctured nonnegative cone]
\label{lem:punctured-cone-convex}
For a finite type $I$, the exact conclusion is
\[
 \operatorname{Convex}\!\left(\R,
  \{x:I\to\R\mid \operatorname{Nonneg}(x)\ \wedge\ x\neq0\}\right).
\]
\end{lemma}

\begin{proof}
Coordinatewise nonnegativity is preserved by convex combinations.  If
$x,z$ are nonzero and $0<\lambda<1$, choose a positive coordinate of $x$;
the same coordinate of $\lambda x+(1-\lambda)z$ is positive.  At
$\lambda=0$ or $1$ the combination is one of the two nonzero endpoints.
\end{proof}

\begin{lemma}[Monotone composition rule]\label{lem:composition-rule}
Let $V$ be a real topological vector space, let $C\subseteq V$ be convex,
let $m:\N$, let $U\subseteq(\operatorname{Fin}(m)\to\R)$ be convex, let
$G:V\to(\operatorname{Fin}(m)\to\R)$ satisfy
$\operatorname{MapsTo}(G,C,U)$, and let
$F:(\operatorname{Fin}(m)\to\R)\to\R$.  The exact theorem is the
conjunction
\begin{align*}
&\Bigl[
 \operatorname{ConcaveOn}(\R,U,F)\ 
 \wedge
 \bigl(\forall a\in U,\forall b\in U,
   (\forall j,a(j)\leq b(j))\to F(a)\leq F(b)\bigr)\\[-2pt]
&\hspace{20mm}{}\wedge
 \bigl(\forall j:\operatorname{Fin}(m),
   \operatorname{ConcaveOn}(\R,C,x\mapsto G(x)(j))\bigr)
 \ \to\ \operatorname{ConcaveOn}(\R,C,F\circ G)
 \Bigr]\\
&\quad\wedge\Bigl[
 \forall I_+,I_-:\operatorname{Finset}(\operatorname{Fin}(m)),\quad
 \operatorname{Disjoint}(I_+,I_-)\to
 I_+\cup I_-=\operatorname{Finset.univ}\to\\[-2pt]
&\hspace{8mm}
 \operatorname{ConvexOn}(\R,U,F)\to\\[-2pt]
&\hspace{8mm}
 \bigl(\forall a\in U,\forall b\in U,
   (\forall j\in I_+,a(j)\leq b(j))\to
   (\forall j\in I_-,b(j)\leq a(j))\to F(a)\leq F(b)\bigr)\to\\[-2pt]
&\hspace{8mm}
 \bigl(\forall j\in I_+,
   \operatorname{ConvexOn}(\R,C,x\mapsto G(x)(j))\bigr)\to
 \bigl(\forall j\in I_-,
   \operatorname{ConcaveOn}(\R,C,x\mapsto G(x)(j))\bigr)\to\\[-2pt]
&\hspace{8mm}
 \operatorname{ConvexOn}(\R,C,F\circ G)
 \Bigr].
\end{align*}
An affine coordinate satisfies both of the final two curvature predicates
and may therefore be assigned to either finite set.
\end{lemma}

\begin{proof}
For the first statement,
$G(\lambda x+(1-\lambda)z)\geq\lambda G(x)+(1-\lambda)G(z)$ coordinatewise.
Coordinate monotonicity followed by concavity of $F$ gives the claim. The
second statement is the same argument with the coordinate inequalities
oriented according to $I_+$ and $I_-$.
\end{proof}

\begin{lemma}[Monotone composition on the punctured cone]
\label{lem:cone-composition-rule}
Let $I$ be finite and nonempty, let $m:\N$, let
$G:\operatorname{ConeVec}(I)\to(\operatorname{Fin}(m)\to\R)$, let
$U\subseteq(\operatorname{Fin}(m)\to\R)$, and let
$F:(\operatorname{Fin}(m)\to\R)\to\R$.  Assume
$\operatorname{Convex}(\R,U)$ and
$\forall q:\operatorname{ConeVec}(I),q\neq0\to G(q)\in U$.
Let $x,z:\operatorname{ConeVec}(I)$ have proofs $h_x:x\neq0$ and
$h_z:z\neq0$, let $\lambda:\R$ have proof
$h_\lambda:\lambda\in\operatorname{Set.Icc}(0,1)$, put
$w:=\operatorname{coneMix}(\lambda,h_\lambda,x,z)$, and assume $h_w:w\neq0$.
Then the following two implications form one literal conjunction:
\begin{align*}
 &\Bigl[
  \operatorname{ConcaveOn}(\R,U,F)\ \wedge\\[-2pt]
 &\quad
  \bigl(\forall a\in U,\forall b\in U,
    (\forall j:\operatorname{Fin}(m),a(j)\leq b(j))
    \to F(a)\leq F(b)\bigr)\ \wedge\\[-2pt]
 &\quad
  \bigl(\forall j:\operatorname{Fin}(m),
    \lambda G(x)(j)+(1-\lambda)G(z)(j)\leq G(w)(j)\bigr)\\[-2pt]
 &\hspace{10mm}\Longrightarrow
  \lambda F(G(x))+(1-\lambda)F(G(z))\leq F(G(w))
 \Bigr]\\
 &\quad\wedge\Bigl[
  \forall I_+,I_-:\operatorname{Finset}(\operatorname{Fin}(m)),\
  \operatorname{Disjoint}(I_+,I_-)\to
  I_+\cup I_-=\operatorname{Finset.univ}\to\\[-2pt]
 &\quad
  \operatorname{ConvexOn}(\R,U,F)\to\\[-2pt]
 &\quad
  \bigl(\forall a\in U,\forall b\in U,
   (\forall j\in I_+,a(j)\leq b(j))\to
   (\forall j\in I_-,b(j)\leq a(j))\to F(a)\leq F(b)\bigr)\to\\[-2pt]
 &\quad
  \bigl(\forall j\in I_+,
    G(w)(j)\leq\lambda G(x)(j)+(1-\lambda)G(z)(j)\bigr)\to\\[-2pt]
 &\quad
  \bigl(\forall j\in I_-,
    \lambda G(x)(j)+(1-\lambda)G(z)(j)\leq G(w)(j)\bigr)\to\\[-2pt]
 &\hspace{10mm}
  F(G(w))\leq\lambda F(G(x))+(1-\lambda)F(G(z))
 \Bigr].
\end{align*}
\end{lemma}

\begin{proof}
This is \cref{lem:composition-rule} applied to the underlying ambient
coordinate functions, or directly coordinate monotonicity followed by the
outer curvature inequality.  Every subtype mixture is the explicitly
bundled \texttt{coneMix}; no module instance on \texttt{ConeVec} is used.
\end{proof}

\begin{lemma}[Finite H\"older inequality]\label{lem:finite-holder}
Let $I$ be a finite type, let $p,q:\R$ have proofs $h_p:1<p$,
$h_q:1<q$, and $h_{pq}:1/p+1/q=1$, and let $a,b:I\to\R$ have proofs
$h_a:\forall i,0\leq a(i)$ and $h_b:\forall i,0\leq b(i)$.  Then
\[
 \sum_{i:I}a(i)b(i)
 \leq
 \operatorname{Real.rpow}\!\left(
   \sum_{i:I}\operatorname{Real.rpow}(a(i),p),1/p\right)
 \operatorname{Real.rpow}\!\left(
   \sum_{i:I}\operatorname{Real.rpow}(b(i),q),1/q\right).
\]
\end{lemma}

\begin{proof}
If either power sum vanishes, every corresponding entry vanishes and the
claim follows.  Otherwise divide $a$ and $b$ by their positive power means.
Young's inequality $uv\leq u^p/p+v^q/q$ applied coordinatewise gives a
normalized sum at most $1/p+1/q=1$.  Multiply back by the two normalizing
factors.
\end{proof}

\begin{lemma}[Power-mean curvature]\label{lem:power-curvature}
Let $I$ be a finite nonempty type.  The exact return proposition is
\begin{align*}
 &\Bigl(\forall\alpha:\R,\ 0<\alpha\to\alpha<1\to
   \operatorname{ConcaveCone}_I
    (x\mapsto\operatorname{lpNorm}_I(\alpha,x))\Bigr)\\
 &\quad\wedge
 \Bigl(\forall\alpha:\R,\ 1\leq\alpha\to
   \operatorname{ConvexCone}_I
    (x\mapsto\operatorname{lpNorm}_I(\alpha,x))\Bigr)\\
 &\quad\wedge
 \operatorname{ConvexCone}_I
  \bigl(x\mapsto\operatorname{finMax}(i\mapsto x(i))\bigr).
\end{align*}
No declaration therefore places a real order and $+\infty$ in one
Lean-typed interval.
\end{lemma}

\begin{proof}
First let $\alpha>1$ and put $q=\alpha/(\alpha-1)$.  If
$s=\norm{x+z}_\alpha>0$, then
\begin{align*}
 s^\alpha
 &=\sum_i x_i(x_i+z_i)^{\alpha-1}
   +\sum_i z_i(x_i+z_i)^{\alpha-1}\\
 &\leq(\norm{x}_\alpha+\norm{z}_\alpha)
 \left(\sum_i(x_i+z_i)^{(\alpha-1)q}\right)^{1/q}\\
 &=(\norm{x}_\alpha+\norm{z}_\alpha)s^{\alpha-1},
\end{align*}
where \cref{lem:finite-holder} was used twice.  Division by
$s^{\alpha-1}>0$ gives Minkowski's inequality.  If $s=0$, all vectors vanish.
At $\alpha=1$ the triangle identity holds for nonnegative vectors; at
$+\infty$ it is the inequality
$\max_i(x_i+z_i)\leq\max_i x_i+\max_i z_i$.  Positive homogeneity converts
subadditivity into convexity.

Now let $0<\alpha<1$.  Put $A=\norm{x}_\alpha$ and
$B=\norm{z}_\alpha$.  If one is zero, the corresponding nonnegative vector
is zero.  Otherwise set $p=x/A$, $q=z/B$, and
$\theta=A/(A+B)$.  Concavity of $r^\alpha$ gives
\[
 \sum_i\bigl(\theta p_i+(1-\theta)q_i\bigr)^\alpha
 \geq\theta\sum_ip_i^\alpha+(1-\theta)\sum_iq_i^\alpha=1.
\]
Taking the positive $1/\alpha$ power and multiplying by $A+B$ yields
$\norm{x+z}_\alpha\geq A+B$.  Positive homogeneity converts this
superadditivity into concavity.
\end{proof}

\begin{lemma}[Compactified power-mean bridge]
\label{lem:parameter-power-curvature}
Let $I$ be a finite nonempty type and define
$F_a(x):=\operatorname{parameterPowerMean}_I(a,x)$ for
$a:\overline\R_+$ and $x:\operatorname{ConeVec}(I)$.
The exact proposition is
\begin{align*}
 &\bigl(\forall a:\overline\R_+,\ a<1\to
   \operatorname{ConcaveCone}_I(F_a)\bigr)\\
 &\quad\wedge
 \bigl(\forall a:\overline\R_+,\ 1\leq a\to
   \operatorname{ConvexCone}_I(F_a)\bigr)\\
 &\quad\wedge
 \bigl(\forall(a:\overline\R_+)(x:\operatorname{ConeVec}(I)),\quad
   x\neq0\to0<F_a(x)\bigr).
\end{align*}
\end{lemma}

\begin{proof}
If $a=+\infty$, use the maximum clause of
\cref{lem:power-curvature}.  Otherwise extract the finite nonnegative real
order and use the corresponding real-order clause; at $a=0$ the chosen total
value is the affine mass.  Strict positivity follows from
\cref{lem:fnd-nonzero-mass}: the mass is positive, every positive finite
power sum has a positive term, and the maximum has a positive coordinate.
\end{proof}

\begin{definition}[Hessian quadratic form]
\label{def:hessian-quad}
For $n:\N$, $F:(\operatorname{Fin}(n)\to\R)\to\R$, and
$x,v:\operatorname{Fin}(n)\to\R$, define
\[
 \operatorname{hessianQuad}(F,x,v):=
 \Bigl(
  \operatorname{fderiv}\!\left(\R,
    y\mapsto\operatorname{fderiv}(\R,F,y),x\right)(v)
 \Bigr)(v).
\]
This is a total real term; the Hessian interpretation is used only under the
twice-differentiability hypotheses below.
\end{definition}

\begin{lemma}[Hessian sign criterion]\label{lem:hessian-sign-criterion}
Let $n:\N$, let $U\subseteq(\operatorname{Fin}(n)\to\R)$ satisfy
$\operatorname{IsOpen}(U)$ and $\operatorname{Convex}(\R,U)$, and let
$F:(\operatorname{Fin}(n)\to\R)\to\R$ satisfy
$\operatorname{ContDiffOn}(\R,2,F,U)$.  Then
\begin{align*}
 &\Bigl(
  (\forall x\in U,\forall v:\operatorname{Fin}(n)\to\R,
   \operatorname{hessianQuad}(F,x,v)\leq0)
  \Longrightarrow\operatorname{ConcaveOn}(\R,U,F)\Bigr)\\
 &\quad\wedge\quad
 \Bigl(
  (\forall x\in U,\forall v:\operatorname{Fin}(n)\to\R,
   0\leq\operatorname{hessianQuad}(F,x,v))
  \Longrightarrow\operatorname{ConvexOn}(\R,U,F)\Bigr).
\end{align*}
\end{lemma}

\begin{proof}
For $x,z\in U$, restrict $F$ to the segment
$h(s)=F((1-s)x+sz)$.  The chain rule gives
$h''(s)=\operatorname{hessianQuad}
(F,(1-s)x+sz,z-x)$.  A real twice differentiable function
with nonpositive second derivative is concave; apply this on $[0,1]$.  The
convex case follows by applying the concave result to $-F$.
\end{proof}

For every positive real $y$ and every real exponent $\beta$, the notation
$y^\beta$ in the remainder of this section means
\texttt{Real.rpow}, equivalently
\[
 y^\beta:=\exp(\beta\log y).
\]

\begin{lemma}[Positive-base real-power calculus]
\label{lem:real-rpow-calculus}
For $\beta:\R$, put
$f_\beta(y):=\operatorname{Real.rpow}(y,\beta)$ and
$U:=\operatorname{Set.Ioi}(0)$.  Then the exact theorem is
\begin{align*}
 &\operatorname{ContDiffOn}(\R,2,f_\beta,U)\\
 &\quad\wedge
 \Bigl(\forall y\in U,\quad
  \operatorname{HasDerivAt}
   (f_\beta,\beta\operatorname{Real.rpow}(y,\beta-1),y)\Bigr)\\
 &\quad\wedge
 \Bigl(\forall y\in U,\quad
  \operatorname{secondDeriv}(f_\beta,y)
   =\beta(\beta-1)\operatorname{Real.rpow}(y,\beta-2)\Bigr)\\
 &\quad\wedge
 \bigl(0\leq\beta\to\operatorname{MonotoneOn}(f_\beta,U)\bigr)
 \quad\wedge
 \bigl(\beta\leq0\to\operatorname{AntitoneOn}(f_\beta,U)\bigr)\\
 &\quad\wedge
 \bigl(0<\beta\to\operatorname{StrictMonoOn}(f_\beta,U)\bigr)
 \quad\wedge
 \bigl(\beta<0\to\operatorname{StrictAntiOn}(f_\beta,U)\bigr).
\end{align*}
\end{lemma}

\begin{proof}
Differentiate $\exp(\beta\log y)$ using the chain and product rules and
$\frac{\dd}{\dd y}\log y=1/y$.  The sign of the first derivative proves
the monotonicity assertions by the mean-value theorem.
\end{proof}

\begin{lemma}[Curvature of an affine monomial]\label{lem:monomial-curvature}
Let $N:\N$ and let $\beta:\operatorname{Fin}(N+1)\to\R$ satisfy
$\sum_{j:\operatorname{Fin}(N+1)}\beta_j=1$.  Define the ambient function
\[
 \mathsf M_\beta:(\operatorname{Fin}(N+1)\to\R)\to\R,
 \qquad
 \mathsf M_\beta(y):=\prod_{j:\operatorname{Fin}(N+1)}
   \operatorname{Real.rpow}(y_j,\beta_j),
\]
and let
$U:=\{y:\operatorname{Fin}(N+1)\to\R\mid\forall j,0<y_j\}$.  The following
literal conjunction gives both curvature classifications and all four
coordinate-monotonicity directions:
\begin{align*}
&\Bigl(\operatorname{ConcaveOn}(\R,U,\mathsf M_\beta)
 \Longleftrightarrow\forall j,0\leq\beta(j)\Bigr)\\
&\quad\wedge
\Bigl(\operatorname{ConvexOn}(\R,U,\mathsf M_\beta)
 \Longleftrightarrow
 \exists!k:\operatorname{Fin}(N+1),\quad
 0<\beta(k)\ \wedge\ \forall j,\ j\neq k\to\beta(j)\leq0\Bigr)\\
&\quad\wedge
\Bigl(\forall j,\ 0\leq\beta(j)\to
 \forall a\in U,\forall b\in U,
 (\forall \ell,\ell\neq j\to a(\ell)=b(\ell))\to
 a(j)\leq b(j)\to\mathsf M_\beta(a)\leq\mathsf M_\beta(b)\Bigr)\\
&\quad\wedge
\Bigl(\forall j,\ \beta(j)\leq0\to
 \forall a\in U,\forall b\in U,
 (\forall \ell,\ell\neq j\to a(\ell)=b(\ell))\to
 a(j)\leq b(j)\to\mathsf M_\beta(b)\leq\mathsf M_\beta(a)\Bigr)\\
&\quad\wedge
\Bigl(\forall j,\ 0<\beta(j)\to
 \forall a\in U,\forall b\in U,
 (\forall \ell,\ell\neq j\to a(\ell)=b(\ell))\to
 a(j)<b(j)\to\mathsf M_\beta(a)<\mathsf M_\beta(b)\Bigr)\\
&\quad\wedge
\Bigl(\forall j,\ \beta(j)<0\to
 \forall a\in U,\forall b\in U,
 (\forall \ell,\ell\neq j\to a(\ell)=b(\ell))\to
 a(j)<b(j)\to\mathsf M_\beta(b)<\mathsf M_\beta(a)\Bigr).
\end{align*}
\end{lemma}

\begin{proof}
The positive-base real-power derivative gives
$\partial_j\mathsf M=\beta_j\mathsf M/y_j$.  A second differentiation and
collection of diagonal and off-diagonal terms gives, for $z_j=v_j/y_j$,
\begin{equation}
 \frac{\operatorname{hessianQuad}(\mathsf M_\beta,y,v)}{\mathsf M_\beta(y)}
 =\left(\sum_j\beta_jz_j\right)^2-\sum_j\beta_jz_j^2.
 \label{eq:monomial-hessian}
\end{equation}
If every $\beta_j\geq0$, the right-hand side is minus a weighted variance and
is nonpositive.  Apply \cref{lem:hessian-sign-criterion}.  If some
$\beta_j<0$, taking $z$ supported on that coordinate
gives $\beta_j(\beta_j-1)>0$.  Since the positive orthant $U$ is open and
$y\in U$, continuity of $s\mapsto y+sv$ supplies named witnesses
$\eps>0$ and
$h_U:\forall s\in\operatorname{Set.Ioo}(-\eps,\eps),y+sv\in U$.
Restrict to this interval and apply the concave
half of \cref{lem:fnd-line-curvature}; the positive second derivative at
$s=0$ contradicts concavity.

If two coefficients $\beta_i,\beta_j$ are positive, choose
$z_i=1/\beta_i$, $z_j=-1/\beta_j$, and all other coordinates zero. The first
term in \cref{eq:monomial-hessian} vanishes while the second is strictly
negative.  The same small-line restriction and the convex half of
\cref{lem:fnd-line-curvature} show that convexity fails. Conversely, suppose
$\beta_k>0$ and write
$\beta_j=-\gamma_j\leq0$ for $j\neq k$. Since
$\beta_k=1+\sum_{j\neq k}\gamma_j$, putting
$\delta_j=z_j-z_k$ gives
\[
 \left(\sum_j\beta_jz_j\right)^2-\sum_j\beta_jz_j^2
 =\left(\sum_{j\neq k}\gamma_j\delta_j\right)^2
 +\sum_{j\neq k}\gamma_j\delta_j^2\geq0.
\]
Apply \cref{lem:hessian-sign-criterion}.  The coordinate monotonicity follows
from the displayed first derivative and the one-dimensional mean-value
theorem.
\end{proof}

\begin{lemma}[Enumeration of a finite probability support]
\label{lem:finite-support-enumeration}
Let $\tau:\operatorname{ProbabilityMeasure}(\overline\R_+)$ and assume
$h_{\rm fin}:\operatorname{Set.Finite}
(\operatorname{suppMeasure}(\operatorname{probMeasure}(\tau)))$.  The exact
dependent existential conclusion is
\[
 \exists(N:\N)(h_N:0<N)
   (\alpha:\operatorname{Fin}(N)\to\overline\R_+),\quad
 \operatorname{Function.Injective}(\alpha)\ \wedge\
 \operatorname{Set.range}(\alpha)
 =\operatorname{suppMeasure}(\operatorname{probMeasure}(\tau)).
\]
\end{lemma}

\begin{proof}
The support is nonempty: otherwise its complement is the whole parameter
space and would be null by \cref{lem:fnd-support-complement}, contradicting
the mass-one field of $\tau$.  Take $N$ to be the cardinality of the finite
support and use the standard equivalence between a finite set of cardinality
$N$ and $\operatorname{Fin}(N)$; compose that equivalence with the subtype
inclusion.
\end{proof}

\begin{lemma}[Finite-support measure decomposition]
\label{lem:finite-support-dirac}
Let $N:\N$ satisfy $h_N:0<N$, let
$\alpha:\operatorname{Fin}(N)\to\overline\R_+$ have proof
$h_\alpha:\operatorname{Function.Injective}(\alpha)$, and let
$\tau:\operatorname{ProbabilityMeasure}(\overline\R_+)$ have proof
\[
 h_{\rm supp}:\operatorname{suppMeasure}(\operatorname{probMeasure}(\tau))
 =\operatorname{Set.range}(\alpha).
\]
Then the return proposition begins with the literal local definitions
\begin{align*}
&\texttt{let }m:\operatorname{Fin}(N)\to\operatorname{ENNReal}:=
 \bigl(\ell\mapsto
   \operatorname{probMeasure}(\tau)(\{\alpha(\ell)\})\bigr)\texttt{; }\\[-2pt]
&\texttt{let }w:\operatorname{Fin}(N)\to\R:=
 \bigl(\ell\mapsto\operatorname{ENNReal.toReal}(m(\ell))\bigr)\texttt{; }\\[-2pt]
 &(\forall\ell:\operatorname{Fin}(N),m(\ell)\neq\top)\ \wedge\
  (\forall\ell:\operatorname{Fin}(N),0<m(\ell))\ \wedge\
  (\forall\ell:\operatorname{Fin}(N),0<w(\ell))\\
 &\quad\wedge\quad
  \sum_{\ell:\operatorname{Fin}(N)}w(\ell)=1\\
 &\quad\wedge\quad
  \operatorname{probMeasure}(\tau)
  =\sum_{\ell:\operatorname{Fin}(N)}
    m(\ell)\mathbin{\bullet}\operatorname{diracRaw}(\alpha(\ell))\\
 &\quad\wedge\quad
 \Bigl(\forall f:\overline\R_+\to\R,\
   \operatorname{Measurable}(f)\to
   (\exists C:\R,\forall a,\abs{f(a)}\leq C)\to\\[-2pt]
 &\hspace{30mm}
   \int f\dd\operatorname{probMeasure}(\tau)
   =\sum_{\ell:\operatorname{Fin}(N)}w(\ell) f(\alpha(\ell))\Bigr)\\
 &\quad\wedge\quad
 \Bigl(\forall g:\overline\R_+\to\operatorname{ENNReal},\
   \operatorname{Measurable}(g)\to
   \int^-g(a)\dd\operatorname{probMeasure}(\tau)(a)
   =\sum_{\ell:\operatorname{Fin}(N)}m(\ell) g(\alpha(\ell))\Bigr).
\end{align*}
Here $\bullet$ is the \texttt{ENNReal}-scalar action on raw positive measures.
\end{lemma}

\begin{proof}
The nonemptiness of $\operatorname{Fin}(N)$ is supplied by $h_N$.  For each
$i:\operatorname{Fin}(N)$, the point $\alpha(i)$ belongs to the support.
The Hausdorff property and finiteness give an open set $U_i$ with
$U_i\cap\operatorname{Set.range}(\alpha)=\{\alpha(i)\}$; injectivity is used
here.  Since $\alpha(i)$ lies in the support, $\tau(U_i)>0$.  On the other
hand, $U_i\setminus\{\alpha(i)\}$ is contained in the complement of the
support and is null.  Therefore
$\tau(\{\alpha(i)\})=\tau(U_i)>0$.
The complement of the finite support is null, and finite additivity gives
the sum of the singleton masses.  Since every $m_i\leq1<+\infty$,
\texttt{ENNReal.toReal} preserves positivity and the finite sum, giving
$\sum_iw_i=1$.  For any Borel set $E$,
\[
 \operatorname{probMeasure}(\tau)(E)
 =\sum_{\substack{i:\operatorname{Fin}(N)\\\alpha(i)\in E}}
   \operatorname{probMeasure}(\tau)(\{\alpha(i)\})
 =\left(\sum_i m_i\mathbin{\bullet}
   \operatorname{diracRaw}(\alpha(i))\right)(E),
\]
which proves equality of measures.  The real-integral and
\texttt{lintegral} formulas follow by substitution into this finite measure
sum and the corresponding finite linearity theorems.
\end{proof}

\begin{proposition}[Finite-support sufficiency]\label{prop:finite-sufficiency}
Let $\tau:\operatorname{ProbabilityMeasure}(\overline\R_+)$, let
$h_{\rm fin}:\operatorname{Set.Finite}
(\operatorname{suppMeasure}(\operatorname{probMeasure}(\tau)))$, let
$t:\R$, let $\alpha_*:\overline\R_+$, and let $I$ be a finite nonempty row
type.  Then the following literal conjunction holds:
\begin{align*}
 &\bigl(\operatorname{PosAdm}(t,\tau)\Longrightarrow
        \operatorname{ConcaveCone}_I(\Phi_{t,\tau})\bigr)\\
 &\quad\wedge
 \bigl(\operatorname{NegLowerAdm}(t,\tau)\Longrightarrow
        \operatorname{ConvexCone}_I(\Phi_{t,\tau})\bigr)\\
 &\quad\wedge
 \bigl(\operatorname{NegExcAdm}(t,\tau,\alpha_*)\Longrightarrow
        \operatorname{ConvexCone}_I(\Phi_{t,\tau})\bigr).
\end{align*}
\end{proposition}

\begin{proof}
We first work on the punctured nonnegative cone and assume that $\tau$ has no
atom at $0$ or $1$.  Apply
\cref{lem:finite-support-enumeration} to $\tau,h_{\rm fin}$ and write the
result in Lean as
\[
 \texttt{rcases finiteSupportEnumeration tau h\_fin with}
 \quad\langle N,h_N,\alpha,h_{\rm inj},h_{\rm range}\rangle.
\]
Apply \cref{lem:finite-support-dirac} with $h_{\rm range}.\operatorname{symm}$.
Thus
\[
 \operatorname{probMeasure}(\tau)=
 \sum_{\ell:\operatorname{Fin}(N)}m_\ell\mathbin{\bullet}
 \operatorname{diracRaw}(\alpha(\ell)),
 \qquad w_\ell=\operatorname{ENNReal.toReal}(m_\ell).
\]
For $a:\overline\R_+$ equipped with proofs $h_\top:a\neq+\infty$,
$h_0:a\neq0$, and $h_1:a\neq1$, and for
$x:\operatorname{ConeVec}(I)$ with $h_x:x\neq0$, extract
$r_a:=((a.\operatorname{untop}(h_\top):\R_{\geq0}):\R)$ and apply
\cref{lem:fnd-renyi-norm-formula}.  This gives
\begin{equation}
 H_a(\widehat x)
 =\omega(a)
   \bigl(\log\operatorname{parameterPowerMean}_I(a,x)
          -\log\norm{x}_1\bigr).
 \label{eq:renyi-norm-formula}
\end{equation}
The same identity holds at $+\infty$ with $\omega(+\infty)=-1$. Define
\begin{equation}
 \beta_\ell:=t\omega(\alpha(\ell))w_\ell,
 \qquad
 \beta_0:=1-\sum_{\ell:\operatorname{Fin}(N)}\beta_\ell.
 \label{eq:beta-def}
\end{equation}
Package these values as one Lean family
\[
 \widetilde\beta:\operatorname{Fin}(N+1)\to\R,
 \qquad
 \widetilde\beta:=\operatorname{Fin.cases}(\beta_0)
   (\ell\mapsto\beta_\ell),
\]
and define the inner coordinate map by
\[
 G(x):=\operatorname{Fin.cases}(\norm{x}_1)
  \bigl(\ell\mapsto
    \operatorname{parameterPowerMean}_I(\alpha(\ell),x)\bigr).
\]
Then $\sum_j\widetilde\beta_j=1$, so the outer function is exactly
$\mathsf M_{\widetilde\beta}$ from \cref{lem:monomial-curvature}.
Substitution gives the exact factorisation
\begin{equation}
 \Phi_{t,\tau}(x)
 =\norm{x}_1^{\beta_0}
  \prod_{\ell:\operatorname{Fin}(N)}
  \operatorname{parameterPowerMean}_I(\alpha(\ell),x)^{\beta_\ell}.
 \label{eq:exact-factorisation}
\end{equation}
Every factor is strictly positive when $x\neq0$, even if some coordinates of
$x$ vanish.
The common composition domain is convex by
\cref{lem:punctured-cone-convex}.
All composition steps below invoke the subtype-specialized
\cref{lem:cone-composition-rule} with this $G$; curvature and positivity of
its coordinates are supplied by
\cref{lem:parameter-power-curvature}.  Thus no ambient/subtype conversion is
left implicit.

Suppose $t>0$. Every support point lies below $1$, so
$\beta_\ell>0$. Moreover,
\[
 \beta_0=1-t\operatorname{ENNReal.toReal}(M_{<}(\tau))\geq0.
\]
Here finiteness follows from $\operatorname{PosAdm}$, and applying
\texttt{ENNReal.toReal} to its typed inequality is legitimate because
$1/t\geq0$.
The outer monomial is concave and increasing in every coordinate, while all
inner power means are concave. By \cref{lem:cone-composition-rule}, the
composition is concave.

Now let $t<0$ and suppose all support points lie below $1$. Then every
$\beta_\ell<0$ and $\beta_0>0$, so $\beta_0$ is the unique positive
coefficient. The outer monomial is convex, increasing in the affine
coordinate $\norm{x}_1$, and decreasing in the concave lower-order power
means. The second part of \cref{lem:cone-composition-rule} proves convexity. An
atom at $1$ will be included by approximation below.

Finally, suppose $t<0$ and $\alpha_*>1$ is the unique upper support point.
By \cref{lem:fnd-exceptional-witness-positive},
$\tau(\{\alpha_*\})>0$.  Hence the exponent attached to
$\norm{x}_{\alpha_*}$ is strictly positive, and all nonzero exponents
attached to orders below $1$ are strictly negative. Since
$M(\tau)\leq1/t$ and $t<0$,
\[
 \beta_0=1-tM(\tau)\leq0.
\]
Thus the upper power mean carries the unique positive exponent, while
$\beta_0\leq0$. It is convex;
the lower power means are concave; and $\norm{x}_1=\sum_i x_i$ is affine on
the nonnegative cone. The outer
monomial is increasing in the upper coordinate and decreasing in every other
coordinate, so \cref{lem:cone-composition-rule} proves convexity.

We now include the endpoints through a total construction.  For $n:\N$ put
\[
 r_n:=1/(n+2:\R),\qquad
 \eps_n:=\bigl(\operatorname{ENNReal.ofReal}(r_n):\overline\R_+\bigr),
 \qquad
 \eta_n:=\bigl(\operatorname{ENNReal.ofReal}(1-r_n):\overline\R_+\bigr).
\]
Thus $\eps_n$ and $\eta_n$ have type $\overline\R_+$; the proof records
$0<r_n<1$, so neither is $0$ or $1$.  Let
$a_0:=\operatorname{ENNReal.toReal}(\tau(\{0\}))$, and define the Borel map
\[
 E_n(\alpha):=
 \begin{cases}
  \eps_n,&\alpha=0,\\
  \eta_n,&\alpha=1,\\
  \alpha,&\text{otherwise}.
 \end{cases}
\]
Let $h_{E_n}$ be its explicit Borel, hence $\tau$-a.e.-measurable, proof and
set
\[
 \tau_n:=\operatorname{probMap}(\tau,E_n,h_{E_n}).
\]
The map $E_n$ is Borel but is not globally
finite-range.  Nevertheless, $\tau_n$ is a probability measure with finite
support contained in the finite set $E_n(\supp\tau)$, and it has no mass at
$0$ or $1$.  Formally, if the finite-support decomposition is
$\operatorname{probMeasure}(\tau)=
 \sum_\ell m_\ell\mathbin{\bullet}\operatorname{diracRaw}(\alpha(\ell))$,
then
\[
 \operatorname{probMeasure}(\tau_n)=\sum_\ell
 m_\ell\mathbin{\bullet}\operatorname{diracRaw}(E_n(\alpha(\ell)));
\]
this direct identity proves the support claim without invoking any
continuity theorem for the discontinuous map $E_n$.  In either
moment-constrained branch
one has $\tau(\{1\})=0$, and hence
\[
 M_<(\tau_n)=M_<(\tau)+
 \operatorname{ENNReal.ofReal}\bigl(a_0\omega(\eps_n)\bigr),
\qquad M_>(\tau_n)=M_>(\tau).
\]
Both sides are finite in those branches, and the corresponding real identity
is
\[
 \operatorname{ENNReal.toReal}(M_<(\tau_n))
 =\operatorname{ENNReal.toReal}(M_<(\tau))+a_0\omega(\eps_n).
\]
In the all-lower negative branch an atom at $1$ may be present, but no moment
condition is imposed.

In the all-lower branch put $t_n:=t$.  In either moment-constrained branch,
put
\begin{equation}
 \delta_n:=a_0\omega(\eps_n),
 \qquad
 t_n:=\frac{t}{1+t\delta_n}.
 \label{eq:t-adjustment}
\end{equation}
Since $\delta_n\to0$, there is $n_0:\N$ such that
$1+t\delta_n>0$ for $n\geq n_0$.  In the all-lower branch take $n_0=0$.
Define the actual total Lean sequences, for $k:\N$, by
\[
 \bar t_k:=t_{k+n_0},\qquad \bar\tau_k:=\tau_{k+n_0}.
\]
Every pair $(\bar t_k,\bar\tau_k)$ is therefore admissible.  Whenever
$n\geq n_0$, and hence in particular at $n=k+n_0$, one has
$t_n\neq0$, $t_n$ has the same sign as $t$, and
\[
 \frac1{t_n}=\frac1t+\delta_n.
\]
In the positive branch all terms are finite and nonnegative, so applying
\texttt{ENNReal.toReal} to the typed moment inequality reduces it to
$\operatorname{ENNReal.toReal}(M_<\tau)\leq1/t$.  The preceding real identity,
$1/t_n=1/t+\delta_n$, and the finite-value order embedding give the typed
implication
\[
 M_<\tau\leq\operatorname{ENNReal.ofReal}(1/t)
 \Longrightarrow
 M_<\tau_n\leq\operatorname{ENNReal.ofReal}(1/t_n).
\]
This is obtained in Lean from \texttt{toReal\_add},
\texttt{toReal\_ofReal}, and the two non-top proofs.  In the exceptional
branch, the equality of both one-sided real moments gives
$M_{\R}(\tau_n)=M_{\R}(\tau)+\delta_n$, so
$M_{\R}(\tau_n)\leq1/t_n$ is exactly the original inequality with
$\delta_n$ added to both sides.  Thus every corresponding finite-support
functional $\Phi_{\bar t_k,\bar\tau_k}$ is concave or convex by the
preceding no-endpoint argument.

For every fixed probability vector $p$, if
$a_1:=\operatorname{ENNReal.toReal}(\tau(\{1\}))$, measure linearity and the definition of the
pushforward give the exact identity
\[
 \A_{\tau_n}(p)=\A_\tau(p)
 +a_0\bigl(H_{\eps_n}(p)-H_0(p)\bigr)
 +a_1\bigl(H_{\eta_n}(p)-H_1(p)\bigr).
\]
The two bracketed differences tend to zero by order continuity, and
$\bar t_k\to t$.  The identity also implies
$\A_{\bar\tau_k}(p)\to\A_\tau(p)$.  Therefore, for every fixed nonzero $x$,
$\Phi_{\bar t_k,\bar\tau_k}(x)\to\Phi_{t,\tau}(x)$.
At $x=0$ the same convergence is reflexivity because every displayed
functional is defined to be zero.  Before taking the limit, apply
\cref{lem:origin-extension} to every approximating functional, using its
\texttt{PosHomOne} proof and value zero at the origin.  Thus every
$\Phi_{\bar t_k,\bar\tau_k}$ already satisfies the full custom cone predicate.
Now apply \cref{lem:fnd-pointwise-curvature-limit} on all cone vectors to
obtain the corresponding full predicate for $\Phi_{t,\tau}$.
\end{proof}

\subsection{Arbitrary probability measures}

The finite-support approximation must be independent of the vector: the same
approximating measure must be used at $x$, $z$, and
$\lambda x+(1-\lambda)z$.

\begin{proposition}[General temperate sufficiency]\label{prop:general-sufficiency}
Let $\tau:\operatorname{ProbabilityMeasure}(\overline\R_+)$, let
$t:\R$, let $\alpha_*:\overline\R_+$, and let $I$ be a finite nonempty
row type.  Then
\begin{align*}
 &\bigl(\operatorname{PosAdm}(t,\tau)\Longrightarrow
        \operatorname{ConcaveCone}_I(\Phi_{t,\tau})\bigr)\\
 &\quad\wedge
 \bigl(\operatorname{NegLowerAdm}(t,\tau)\Longrightarrow
        \operatorname{ConvexCone}_I(\Phi_{t,\tau})\bigr)\\
 &\quad\wedge
 \bigl(\operatorname{NegExcAdm}(t,\tau,\alpha_*)\Longrightarrow
        \operatorname{ConvexCone}_I(\Phi_{t,\tau})\bigr).
\end{align*}
\end{proposition}

\begin{proof}
Prove the three conjuncts separately.  For the positive and all-lower
projections of the three-way package, instantiate its irrelevant upper
parameter by the canonical
$a_\top:=\langle+\infty,1<+\infty\rangle:\operatorname{UpperParam}$.

For the positive branch define the total bundled sequence
$\tau_k^+:=\operatorname{discPos}(\tau,k)$.  The first clauses of
\cref{lem:fnd-discretization-basic,lem:fnd-discretization-package} give, for
every $k:\N$, finite support, support in $[0,1)$, zero mass at $1$, and
\[
 M_<\bigl(\operatorname{probMeasure}(\tau_k^+)\bigr)
 \leq M_<\bigl(\operatorname{probMeasure}(\tau)\bigr)
 \leq\operatorname{ENNReal.ofReal}(1/t).
\]
Thus every pair $(t,\tau_k^+)$ satisfies \texttt{PosAdm}, and the first
conjunct of \cref{prop:finite-sufficiency} makes
$\Phi_{t,\tau_k^+}$ concave.

For the all-lower negative branch put
$\tau_k^-:=\operatorname{discLow}(\tau,k)$.  The second package clause gives
finite support and support in $[0,1]$, so \texttt{NegLowerAdm} is preserved
and the second conjunct of \cref{prop:finite-sufficiency} applies.

For the exceptional branch, extract the stored proof $h_*:1<\alpha_*$ from
\texttt{NegExcAdm}, put
$a_*:=\langle\alpha_*,h_*\rangle:\operatorname{UpperParam}$, and define
$\tau_k^{\rm exc}:=\operatorname{discExc}(\tau,a_*,k)$.  The third package
clause gives finite support, preserves the exceptional support and zero atom,
does not increase the lower moment, fixes the upper moment, and hence
preserves the guarded inequality $M_{\R}\leq1/t$.  Apply the third conjunct
of \cref{prop:finite-sufficiency}.

For every $p:\operatorname{ProbVec}(I)$, the three respective conjuncts of
\cref{lem:fnd-discretization-dct}, with
$\nu=\operatorname{probMeasure}(\tau)$, give
\[
 \A_{\tau_k^+}(p)\to\A_\tau(p),\qquad
 \A_{\tau_k^-}(p)\to\A_\tau(p),\qquad
 \A_{\tau_k^{\rm exc}}(p)\to\A_\tau(p)
\]
in the corresponding branch.  Therefore the associated column functions
converge pointwise on every cone vector (at the origin by reflexivity).
Apply \cref{lem:fnd-pointwise-curvature-limit} separately to the three
sequences to obtain the asserted full cone predicates.
\end{proof}

\begin{corollary}[Finite-$t$ sufficiency]\label{cor:finite-t-sufficiency}
Let $\tau:\operatorname{ProbabilityMeasure}(\overline\R_+)$ and let $t:\R$.
If $h:\operatorname{DBulk}(t,\tau)$, the exact dependent-let conclusion is
\[
 \texttt{let }h_{\rm ne}:=\operatorname{DBulk.ne}(h)\texttt{; }\quad
 \operatorname{CMMonotone}
 \bigl(\operatorname{HTemp}(t,h_{\rm ne},\tau)\bigr).
\]
\end{corollary}

\begin{proof}
Split $h:\operatorname{DBulk}(t,\tau)$ into its three defining disjuncts.
Fix arbitrary $X:\texttt{Type u}$ with
$[\operatorname{Fintype}X][\operatorname{Nonempty}X]$.  The corresponding
conjunct of \cref{prop:general-sufficiency} gives respectively the concavity
or convexity of $\Phi_{t,\tau}$ on this row cone.  Apply the positive or
negative rowwise equivalence in
\cref{prop:temperate-shape-reduction}, using the strict-sign field stored in
the disjunct.  Since $X$ was arbitrary, unfold \texttt{CMMonotone}; the proof
argument of \texttt{HTemp} is propositionally equal to
$h_{\rm ne}:=\operatorname{DBulk.ne}(h)$ by proof irrelevance.
\end{proof}

\subsection{Tropical limits}

The already-proved \cref{lem:fnd-log-sum-exp} gives, for probabilities
$p_y$ and real numbers $a_y$,
\begin{align}
 \lim_{t\to-\infty}\frac1t\log\sum_y p_ye^{ta_y}
 &=\min_{y:p_y>0}a_y,
 \label{eq:lse-minus}\\
 \lim_{t\to+\infty}\frac1t\log\sum_y p_ye^{ta_y}
 &=\max_{y:p_y>0}a_y.
 \label{eq:lse-plus}
\end{align}
These are restatement labels only; no duplicate limit lemma is introduced.

\begin{proposition}[Negative tropical sufficiency]
\label{prop:negative-tropical-sufficiency}
Let $\tau:\operatorname{ProbabilityMeasure}(\overline\R_+)$.  Then
\[
 \operatorname{DMinus}(\tau)\Longrightarrow
 \operatorname{CMMonotone}(\operatorname{HMinus}(\tau)).
\]
\end{proposition}

\begin{proof}
Define once, before splitting the two disjuncts of \texttt{DMinus},
\[
 \operatorname{tNeg}:\N\to\R,
 \qquad \operatorname{tNeg}(n):=-(n+1:\R).
\]
If $\supp(\tau)\subseteq[0,1]$, then $(t,\tau)\in\Dbulk$ for every $t<0$.
Apply \cref{cor:finite-t-sufficiency} to the total natural-number sequence
$\operatorname{tNeg}$, and pass each monotonicity inequality to the limit using
\cref{lem:fnd-log-sum-exp} on the finite active conditioning sets.

Suppose the exceptional branch holds. If $M(\tau)<0$, convergence of
$n\mapsto-1/(n+1)$ to $0$ supplies $N_0:\N$ and the named bound
\[
 h_{N_0}:\forall n:\N,\quad N_0\leq n\to
 M(\tau)\leq-1/(n+1)=1/\operatorname{tNeg}(n).
\]
Use the shifted total sequence
$s_n:=\operatorname{tNeg}(n+N_0)$, so every term invoked in the proof is admissible and
$s_n\to-\infty$.  The same sequence argument applies. If $M(\tau)=0$,
define the total sequence
\[
 \eps_n:=1/(n+2:\R),\qquad
 \tau_n^{\rm mix}:=
 \operatorname{probMix}\bigl(
   \tau,\operatorname{diracProb}(\alpha_*),\eps_n,h_n\bigr),
\]
where $h_n:0\leq\eps_n\ \wedge\ \eps_n\leq1$ is the elementary arithmetic proof.
Its underlying measure is the displayed convex combination from the
\texttt{probMix} convention.  It is a probability measure, its support is
contained in the same exceptional set (and now contains $\alpha_*$), it has
no atom at $1$, and its lower moment is finite.  The finite-moment linearity
clause of \cref{lem:fnd-moment-algebra}, together with
$M_{\R}(\operatorname{diracRaw}(\alpha_*))=\omega(\alpha_*)$, gives
\[
 M(\tau_n^{\rm mix})=(1-\eps_n)M(\tau)+\eps_n\omega(\alpha_*)
 =\eps_n\omega(\alpha_*)<0,
\]
so its negative tropical functional is monotone by the preceding case.
Moreover
\[
 \A_{\tau_n^{\rm mix}}(p)=(1-\eps_n)\A_\tau(p)
 +\eps_n H_{\alpha_*}(p),
\]
and the right side tends to $\A_\tau(p)$ because $\eps_n\to0$.  Apply
\cref{lem:fnd-finite-extrema} to the finite active conditioning sets on the
two sides of an arbitrary channel and take $n\to\infty$.
\end{proof}

\begin{lemma}[Positive tropical component formula]
\label{lem:positive-tropical-component}
For all finite nonempty alphabets $X,Y$ and all
$P:\operatorname{JointProb}(X,Y)$,
\[
 \operatorname{HPlus}(\operatorname{diracProb}(0))(P)
 =\left(\texttt{letI := activeNonempty P; }
   \operatorname{finMax}
    (\texttt{fun }y:\operatorname{Active}(P)\texttt{ => }
      H_{(0:\overline\R_+)}(\operatorname{conditional}(P,y)))\right).
\]
\end{lemma}

\begin{proof}
Expand \cref{eq:fnd-H-plus} and use the Dirac integration identity
\cref{lem:fnd-dirac-package}.
\end{proof}

\begin{proposition}[Positive tropical sufficiency]
\label{prop:positive-tropical-sufficiency}
The exact polymorphic conclusion is
\[
 \operatorname{CMMonotone}
   (\operatorname{HPlus}(\operatorname{diracProb}(0))).
\]
\end{proposition}

\begin{proof}
Take $\tau=\operatorname{diracProb}(0)$. It is admissible for every $t>0$.
Apply
\cref{cor:finite-t-sufficiency} to the total sequence
$t_n:=(n+1:\R)$ for $n:\N$ and use
\cref{lem:fnd-log-sum-exp}.  Finally,
$\A_{\operatorname{diracProb}(0)}=H_0$ by
\cref{lem:fnd-dirac-package}, which identifies the
limit with the displayed $H_{+\infty,0}$.
\end{proof}

\subsection{The derivation limit}

\begin{definition}[Derivation polymorphic family]
\label{def:derivation-family}
For a universe $u$ and
$\sigma:\operatorname{FiniteMeasure}(\overline\R_+)$, define
$\operatorname{Delta}(\sigma):\texttt{PolyJointFunctional.\{u\}}$ by
\begin{equation}
 \operatorname{Delta}(\sigma)(P):=
 \sum_{y:\operatorname{Active}(P)}
 \operatorname{colMass}(P,y.1)
 \operatorname{integratedEntropyPos}_X
  (\operatorname{finiteMeasure}(\sigma),\operatorname{conditional}(P,y)).
 \label{eq:derivation-functional}
\end{equation}
The finite row alphabet is inferred from $P$ by the polymorphic signature.
\end{definition}

\begin{proposition}[Derivation sufficiency]\label{prop:derivation-sufficiency}
Let $\sigma:\operatorname{FiniteMeasure}(\overline\R_+)$ satisfy
\[
 \operatorname{suppMeasure}(\operatorname{finiteMeasure}(\sigma))
 \subseteq\operatorname{Set.Icc}(0,1).
\]
Then
\[
 \operatorname{CMMonotone}(\operatorname{Delta}(\sigma))
 \quad\wedge\quad
 \operatorname{Delta}(\sigma)=\operatorname{HZero}(\sigma)
\]
as an equality of polymorphic functions.
\end{proposition}

\begin{proof}
To prove the global predicate, fix an arbitrary row type
$I:\texttt{Type u}$ with
$[\operatorname{Fintype}I][\operatorname{Nonempty}I]$ and prove the
rowwise cone property required by
\cref{cor:global-temperate-shape-reduction}.
For every $\alpha\in[0,1]$, $H_\alpha$ is simplex-concave by
\cref{lem:fnd-renyi-concavity}.  Define the total perspective
\[
 F_\alpha:=\operatorname{perspective}_I(p\mapsto H_\alpha(p)).
\]
It satisfies \texttt{ConcaveCone} by \cref{lem:fnd-perspective}; the zero cases are included in
that lemma.  Finite-sum/integral interchange gives the exact identity
\[
 D_\sigma(x)=\int F_\alpha(x)\dd\operatorname{finiteMeasure}(\sigma)(\alpha).
\]
The integrand is bounded for each fixed finite-dimensional $x$, and positive
integration preserves concavity by \cref{lem:fnd-integration-curvature}.
The finite-positive-measure version of the third clause of
\cref{prop:temperate-shape-reduction} now gives monotonicity.  Expanding
\cref{def:derivation-family,eq:fnd-H-zero} gives the displayed polymorphic
function equality.
\end{proof}

\section{Necessary conditions for the parameters}
\label{sec:necessity}

The necessity proof uses finite-dimensional perturbations whose dimensions
grow with an auxiliary integer $d$. Every differentiation is performed on a
common open interval around $\lambda=0$ on which all coordinates are strictly
positive. A strict curvature obstruction at one sufficiently large finite
$d$ is already a counterexample to monotonicity.

\subsection{Differentiation and elementary analytic lemmas}

\begin{definition}[Signed-measure column functions]
\label{def:signed-column-functions}
For every finite nonempty type $I$, use the exact signatures
\[
 \operatorname{PhiSigned}_I:
 \operatorname{SignedMeasure}(\overline\R_+)
 \to\operatorname{ConeVec}(I)\to\R,
 \qquad
 \operatorname{GSigned}_I:
 \operatorname{SignedMeasure}(\overline\R_+)
 \to\operatorname{PosConeVec}(I)\to\R.
\]
For a finite signed Borel measure $\mu$, define
\begin{equation}
 \operatorname{PhiSigned}_I(\mu,x):=
 \begin{cases}
  \norm{x}_1\exp\!\left(
       \operatorname{signedIntegral}\!\left(\mu,
        \alpha\mapsto H_\alpha\bigl(\operatorname{normalize}
          (\operatorname{toPosCone}(x,h))\bigr)\right)\right),
       &h:x\neq0,\\
  0,&x=0,
 \end{cases}
 \label{eq:Phi-mu}
\end{equation}
and, for $z:\operatorname{PosConeVec}(I)$,
\begin{equation}
 \operatorname{GSigned}_I(\mu,z):=
 \operatorname{signedIntegral}
  \bigl(\mu,\alpha\mapsto H_\alpha(\operatorname{normalize}z)\bigr).
 \label{eq:G-mu}
\end{equation}
The integral is finite because the R\'enyi kernel is bounded in each finite
dimension and $\mu$ has finite total variation.  The notation $\Phi_\mu$
and $G_\mu$ abbreviates these polymorphic functions; the finite carrier $I$
is always an implicit argument.
\end{definition}

\begin{definition}[Positive and negative signed witnesses]
\label{def:signed-witness-predicates}
For $\nu:\operatorname{FiniteMeasure}(\overline\R_+)$, define
\[
 \operatorname{positiveSigned}(\nu):=\operatorname{signedLift}(\nu),
 \qquad
 \operatorname{negativeSigned}(\nu):=-\operatorname{signedLift}(\nu).
\]
For universe $u$, define the three global predicates
with actual Lean headers
\[
 \texttt{def PosPhiConcave.\{u\} (nu : FiniteMeasure Param) : Prop :=},
\quad
 \texttt{def NegPhiConvex.\{u\} (nu : FiniteMeasure Param) : Prop :=},
\quad
 \texttt{def NegGQuasiconvex.\{u\} (nu : FiniteMeasure Param) : Prop :=}.
\]
Their respective bodies are
\begin{align*}
 \texttt{PosPhiConcave.\{u\}}(\nu):\Longleftrightarrow{}
 &\forall\{I:\texttt{Type u}\}
  [\operatorname{Fintype}I][\operatorname{Nonempty}I],\quad
  \operatorname{ConcaveCone}_I
   (\operatorname{PhiSigned}_I(\operatorname{positiveSigned}(\nu))),\\
 \texttt{NegPhiConvex.\{u\}}(\nu):\Longleftrightarrow{}
 &\forall\{I:\texttt{Type u}\}
  [\operatorname{Fintype}I][\operatorname{Nonempty}I],\quad
  \operatorname{ConvexCone}_I
   (\operatorname{PhiSigned}_I(\operatorname{negativeSigned}(\nu))),\\
 \texttt{NegGQuasiconvex.\{u\}}(\nu):\Longleftrightarrow{}
 &\forall\{I:\texttt{Type u}\}
  [\operatorname{Fintype}I][\operatorname{Nonempty}I],\quad
  \operatorname{QCvx}
   (\operatorname{GSigned}_I(\operatorname{negativeSigned}(\nu))).
\end{align*}
There is therefore no informal carrier called a ``negative measure'' in the
formal development: every such object is \texttt{negativeSigned nu} for an
explicit positive finite-measure witness.
\end{definition}

\begin{lemma}[Scaling, total variation, support, and atoms of witnesses]
\label{lem:signed-witness-bridge}
Let $\nu:\operatorname{FiniteMeasure}(\overline\R_+)$ and let
$c:\R$ satisfy $h_c:0\leq c$.  Then the following single conjunction holds:
\begin{align*}
&\operatorname{positiveSigned}(\operatorname{finiteScale}(c,h_c,\nu))
 =c\mathbin{\bullet}\operatorname{positiveSigned}(\nu)\\
&\quad\wedge\quad
 \operatorname{negativeSigned}(\operatorname{finiteScale}(c,h_c,\nu))
 =(-c)\mathbin{\bullet}\operatorname{positiveSigned}(\nu)\\
&\quad\wedge\quad
 \operatorname{signedTV}(\operatorname{positiveSigned}(\nu))
 =\operatorname{finiteMeasure}(\nu)\\
&\quad\wedge\quad
 \operatorname{signedTV}(\operatorname{negativeSigned}(\nu))
 =\operatorname{finiteMeasure}(\nu)\\
&\quad\wedge\quad
 \operatorname{suppSigned}(\operatorname{positiveSigned}(\nu))
 =\operatorname{suppMeasure}(\operatorname{finiteMeasure}(\nu))\\
&\quad\wedge\quad
 \operatorname{suppSigned}(\operatorname{negativeSigned}(\nu))
 =\operatorname{suppMeasure}(\operatorname{finiteMeasure}(\nu))\\
&\quad\wedge\quad
 \forall a:\overline\R_+,\quad
 \operatorname{signedTV}(\operatorname{negativeSigned}(\nu))(\{a\})
 =\operatorname{finiteMeasure}(\nu)(\{a\}).
\end{align*}
\end{lemma}

\begin{proof}
Expand \texttt{finiteScale}, \texttt{positiveSigned}, and
\texttt{negativeSigned}.  The Jordan decompositions are $(\nu,0)$ and
$(0,\nu)$, respectively; the total-variation, support, and atom identities
follow definitionally from the canonical Jordan API.
\end{proof}

\begin{lemma}[Signed-scalar bridge to the candidate column functions]
\label{lem:signed-scalar-column-bridge}
Let $I$ be a finite nonempty type and let
$\tau:\operatorname{ProbabilityMeasure}(\overline\R_+)$.  The exact
conclusion is the following three-way conjunction:
\begin{align*}
&\Bigl(\forall(t:\R)(x:\operatorname{ConeVec}(I)),\quad
 \operatorname{PhiSigned}_I
   (t\mathbin{\bullet}\operatorname{signedLift}
       (\tau.\operatorname{toFiniteMeasure}),x)
 =\Phi_{t,\tau}(x)\Bigr)\\
&\quad\wedge\Bigl(\forall(k:\R)(h_k:0<k)
 (z:\operatorname{PosConeVec}(I)),\quad
 \operatorname{GSigned}_I
   (-k\mathbin{\bullet}\operatorname{signedLift}
       (\tau.\operatorname{toFiniteMeasure}),z)
 =k\,\operatorname{gTrop}_I(\tau,z)\Bigr)\\
&\quad\wedge\Bigl(\forall(k:\R)(h_k:0<k),\quad
 \operatorname{QCvx}(\operatorname{gTrop}_I(\tau))\Longleftrightarrow
 \operatorname{QCvx}\!\left(z\mapsto
  \operatorname{GSigned}_I
   (-k\mathbin{\bullet}\operatorname{signedLift}
    (\tau.\operatorname{toFiniteMeasure}),z)\right)\Bigr).
\end{align*}
\end{lemma}

\begin{proof}
Signed-integral linearity gives
$\int H_\alpha(\widehat x)\dd
 (t\mathbin{\bullet}\operatorname{signedLift}
   (\tau.\operatorname{toFiniteMeasure}))
=t\A_\tau(\widehat x)$ and
$\int H_\alpha(\widehat x)\dd
 (-k\mathbin{\bullet}\operatorname{signedLift}
   (\tau.\operatorname{toFiniteMeasure}))
=-k\A_\tau(\widehat x)$.  Substitute in the definitions.  Multiplication
by the positive scalar $k$ preserves and reflects quasi-convex inequalities.
\end{proof}

\begin{definition}[Embedding a finite real order]
\label{def:finite-param}
Define
\[
 \operatorname{finiteParam}:\R\to\overline\R_+,
 \qquad
 \operatorname{finiteParam}(r):=\operatorname{ENNReal.ofReal}(r).
\]
Every real threshold or exponent used as a compactified parameter below is
inserted through this map.
\end{definition}

\begin{definition}[Positive multiplicative line data]
\label{def:positive-line-data}
For a finite nonempty type $I$, define
\[
 \texttt{structure PositiveLineData (I) where}\qquad
 \texttt{x : I $\to$ $\R$}\qquad
 \texttt{u : I $\to$ $\R$}\qquad
 \texttt{x\_pos : $\forall i$, $0<x(i)$}.
\]
\end{definition}

\begin{definition}[Raw multiplicative line]
\label{def:positive-line-raw}
For finite nonempty $I$, $L:\operatorname{PositiveLineData}(I)$, and
$\lambda:\R$, define
\begin{equation}
 \operatorname{lineRaw}(L,\lambda)(i):=
 L.x(i)(1+L.u(i)\lambda).
 \label{eq:multiplicative-line}
\end{equation}
\end{definition}

\begin{definition}[Positivity predicate for a line]
\label{def:positive-line-predicate}
For finite nonempty $I$, $L:\operatorname{PositiveLineData}(I)$, and
$\lambda:\R$, define
\[
 \operatorname{LinePositive}(L,\lambda):\Longleftrightarrow
 \forall i:I,\ 0<\operatorname{lineRaw}(L,\lambda)(i).
\]
\end{definition}

\begin{lemma}[The base point of a positive line is positive]
\label{lem:line-positive-zero}
For finite nonempty $I$ and $L:\operatorname{PositiveLineData}(I)$,
\[
 \operatorname{linePositiveZero}(L):\operatorname{LinePositive}(L,0).
\]
\end{lemma}

\begin{proof}
For every $i:I$, simplify
$L.x(i)(1+L.u(i)0)$ to $L.x(i)$ and apply $L.x\_pos(i)$.
\end{proof}

\begin{definition}[Proof-carrying cone versions of a positive line]
\label{def:line-cone}
For finite nonempty $I$, $L:\operatorname{PositiveLineData}(I)$,
$\lambda:\R$, and $h:\operatorname{LinePositive}(L,\lambda)$, define
\begin{align*}
 \operatorname{lineCone}(L,\lambda,h)&:=
  \left\langle\operatorname{lineRaw}(L,\lambda),
   i\mapsto\operatorname{le_of_lt}(h(i))\right\rangle
  \in\operatorname{ConeVec}(I),\\
 \operatorname{linePosCone}(L,\lambda,h)&:=
  \left\langle\operatorname{lineCone}(L,\lambda,h),
   \texttt{by
    intro hz
    let i0 : I := Classical.choice (inferInstance : Nonempty I)
    have hi : lineRaw L lambda i0 = 0 := by
      simpa [lineCone] using congrFun hz i0
    exact (ne\_of\_gt (h i0)) hi}
   \right\rangle
  \in\operatorname{PosConeVec}(I).
\end{align*}
Define the total wrappers
\begin{align*}
 \operatorname{lineConeTotal}(L,\lambda)&:=
  \begin{cases}\operatorname{lineCone}(L,\lambda,h),
    &h:\operatorname{LinePositive}(L,\lambda),\\
    0,&\text{otherwise},\end{cases}\\
 \operatorname{linePosConeTotal}(L,\lambda)&:=
  \begin{cases}\operatorname{linePosCone}(L,\lambda,h),
    &h:\operatorname{LinePositive}(L,\lambda),\\
    \operatorname{linePosCone}(L,0,
      \operatorname{linePositiveZero}(L)),&\text{otherwise}.
  \end{cases}
\end{align*}
Proof irrelevance makes both values independent of the particular proof
$h$.  On a positive line their underlying functions are exactly
\texttt{lineRaw}.
\end{definition}

\begin{definition}[Total normalized line]
\label{def:positive-line-prob}
For finite nonempty $I$, $L:\operatorname{PositiveLineData}(I)$, and
$\lambda:\R$, define the literal total function
\[
 \operatorname{lineProb}(L,\lambda):=
 \begin{cases}
  \operatorname{normalize}(\operatorname{linePosCone}(L,\lambda,h)),
    &h:\operatorname{LinePositive}(L,\lambda),\\
  \operatorname{uniformProb}_I,&\text{otherwise}
 \end{cases}
 \in\operatorname{ProbVec}(I).
\]
In Lean this is
\[
 \texttt{if h : LinePositive L lambda then
   normalize (linePosCone L lambda h) else uniformProb\_I}.
\]
It is definitionally the intended normalization in the positive branch.
\end{definition}

\begin{definition}[Scalar quasiconvexity on a set]
\label{def:scalar-qcvx-on}
For $U:\operatorname{Set}(\R)$ and $g:\R\to\R$, define
\begin{align*}
 \operatorname{ScalarQCvxOn}(U,g):\Longleftrightarrow{}
 &\forall a\in U,\forall b\in U,\forall\lambda:\R,
 0\leq\lambda\to\lambda\leq1\to\\[-2pt]
 &\lambda a+(1-\lambda)b\in U\to
 g(\lambda a+(1-\lambda)b)\leq\max\{g(a),g(b)\}.
\end{align*}
\end{definition}

\begin{lemma}[Restricting cone curvature to a positive affine line]
\label{lem:cone-affine-line-bridge}
Let $I$ be finite and nonempty, let
$L:\operatorname{PositiveLineData}(I)$, let $U:\operatorname{Set}(\R)$ have
proof $h_U^{\rm cvx}:\operatorname{Convex}(\R,U)$, and let
$h_U:\forall\lambda\in U,
\operatorname{LinePositive}(L,\lambda)$.  For
$F:\operatorname{ConeVec}(I)\to\R$ and
$g:\operatorname{PosConeVec}(I)\to\R$, the following literal conjunction
holds:
\begin{align*}
&\Bigl(\operatorname{ConcaveCone}_I(F)\to
 \operatorname{ConcaveOn}(\R,U,
  \lambda\mapsto F(\operatorname{lineConeTotal}(L,\lambda)))\Bigr)\\
&\quad\wedge
 \Bigl(\operatorname{ConvexCone}_I(F)\to
 \operatorname{ConvexOn}(\R,U,
  \lambda\mapsto F(\operatorname{lineConeTotal}(L,\lambda)))\Bigr)\\
&\quad\wedge
 \Bigl(\operatorname{QCvx}(g)\to
 \operatorname{ScalarQCvxOn}(U,
  \lambda\mapsto g(\operatorname{linePosConeTotal}(L,\lambda)))\Bigr).
\end{align*}
The proof uses the underlying-function identity
\[
 \operatorname{lineRaw}(L,sa+(1-s)b)
 =s\operatorname{lineRaw}(L,a)+(1-s)\operatorname{lineRaw}(L,b).
\]
\end{lemma}

\begin{proof}
Expand \texttt{lineRaw}; the displayed identity is coordinatewise ring
normalization.  Apply the corresponding cone predicate and erase equality
proof fields by proof irrelevance.
\end{proof}

\begin{definition}[Entropy line]
\label{def:entropy-line}
For finite nonempty $I$, $L:\operatorname{PositiveLineData}(I)$,
$a:\overline\R_+$, and $\lambda:\R$, define
\[
 \operatorname{entropyLine}(L,a,\lambda):=
 H_a(\operatorname{lineProb}(L,\lambda)).
\]
\end{definition}

\begin{definition}[First entropy line derivative]
\label{def:entropy-line-first}
For finite nonempty $I$, $L:\operatorname{PositiveLineData}(I)$,
$a:\overline\R_+$, and $\lambda:\R$, define
\[
 \operatorname{entropyLineFirst}(L,a,\lambda):=
 \operatorname{deriv}(\operatorname{entropyLine}(L,a))(\lambda).
\]
\end{definition}

\begin{definition}[Second entropy line derivative]
\label{def:entropy-line-second}
For finite nonempty $I$, $L:\operatorname{PositiveLineData}(I)$,
$a:\overline\R_+$, and $\lambda:\R$, define
\[
 \operatorname{entropyLineSecond}(L,a,\lambda):=
 \operatorname{secondDeriv}(\operatorname{entropyLine}(L,a),\lambda).
\]
\end{definition}

\begin{definition}[Escort weight]
\label{def:escort-weight}
For finite nonempty $I$, $L:\operatorname{PositiveLineData}(I)$, and
$\alpha,\lambda:\R$, define
\begin{equation}
 \operatorname{escortWeight}(L,\alpha,\lambda,i):=
 \frac{\operatorname{Real.rpow}
       (\operatorname{lineRaw}(L,\lambda)(i),\alpha)}
      {\sum_{j:I}\operatorname{Real.rpow}
       (\operatorname{lineRaw}(L,\lambda)(j),\alpha)}.
 \label{eq:escort}
\end{equation}
\end{definition}

\begin{definition}[Effective line velocity]
\label{def:effective-line-velocity}
For finite nonempty $I$, $L:\operatorname{PositiveLineData}(I)$,
$\lambda:\R$, and $i:I$, define
\[
 \operatorname{effectiveVelocity}(L,\lambda,i):=
 \frac{L.u(i)}{1+L.u(i)\lambda}.
\]
\end{definition}

\begin{definition}[Escort mean]
\label{def:escort-mean}
For finite nonempty $I$, $L:\operatorname{PositiveLineData}(I)$, and
$\alpha,\lambda:\R$, define
\[
 \operatorname{escortMean}(L,\alpha,\lambda):=
 \sum_{i:I}\operatorname{escortWeight}(L,\alpha,\lambda,i)
       \operatorname{effectiveVelocity}(L,\lambda,i).
\]
\end{definition}

\begin{definition}[Escort second moment]
\label{def:escort-second}
For finite nonempty $I$, $L:\operatorname{PositiveLineData}(I)$, and
$\alpha,\lambda:\R$, define
\[
 \operatorname{escortSecond}(L,\alpha,\lambda):=
 \sum_{i:I}\operatorname{escortWeight}(L,\alpha,\lambda,i)
       \operatorname{effectiveVelocity}(L,\lambda,i)^2.
\]
\end{definition}

\begin{definition}[Escort variance]
\label{def:escort-variance}
For finite nonempty $I$, $L:\operatorname{PositiveLineData}(I)$, and
$\alpha,\lambda:\R$, define
\[
 \operatorname{escortVar}(L,\alpha,\lambda):=
 \operatorname{escortSecond}(L,\alpha,\lambda)
 -\operatorname{escortMean}(L,\alpha,\lambda)^2.
\]
\end{definition}

\begin{definition}[Fixed maximal-coordinate condition]
\label{def:fixed-max-coordinate}
For finite nonempty $I$, $L:\operatorname{PositiveLineData}(I)$,
$U:\operatorname{Set}(\R)$, and $i_*:I$, define
\begin{align*}
 \operatorname{FixedMaxCoordinate}(L,U,i_*):\Longleftrightarrow{}
 &\forall\lambda\in U,\quad
   \operatorname{LinePositive}(L,\lambda)\ \wedge\\[-2pt]
 &\quad(\forall i:I,\operatorname{lineRaw}(L,\lambda)(i)
                   \leq\operatorname{lineRaw}(L,\lambda)(i_*))\ \wedge\\[-2pt]
 &\quad(\forall i:I,
   \operatorname{lineRaw}(L,\lambda)(i)
      =\operatorname{lineRaw}(L,\lambda)(i_*)\to
   \operatorname{effectiveVelocity}(L,\lambda,i)
      =\operatorname{effectiveVelocity}(L,\lambda,i_*)).
\end{align*}
\end{definition}

\begin{lemma}[Common positivity and maximal-coordinate intervals]
\label{lem:positivity-intervals}
The following two implications form one literal conjunction:
\begin{align*}
 &\Bigl(\forall m:\N,\forall K:\operatorname{Set}
   (\operatorname{Fin}(m+1)\to\R),\
   \operatorname{IsCompact}(K)\to\\[-2pt]
 &\hspace{12mm}\exists\lambda_0:\R,0<\lambda_0\ \wedge\
   \forall u\in K,\forall j:\operatorname{Fin}(m+1),\forall\lambda:\R,\
   \abs{\lambda}\leq\lambda_0\to
   \tfrac12\leq1+u(j)\lambda\Bigr)\\
 &\quad\wedge\Bigl(\forall p,q:\R,0<p\to0<q\to p+q=1\to
   \forall K:\operatorname{Set}(\R\times\R),\
   \operatorname{IsCompact}(K)\to\\[-2pt]
 &\hspace{12mm}\exists\lambda_0:\R,0<\lambda_0\ \wedge\
   \forall(d:\N),2\leq d\to\forall z\in K,\forall\lambda:\R,\
   \abs{\lambda}\leq\lambda_0\to\\[-2pt]
 &\hspace{22mm}
  q+\frac{z.1\lambda}{\log(d:\R)}\geq\frac q2\ \wedge\
  p-\frac{z.1\lambda}{\log(d:\R)}\geq\frac p2\ \wedge\
  1+z.2\lambda\geq\frac12\Bigr).
\end{align*}
\end{lemma}

\begin{proof}
For the first statement, let $U:=\max_{u\in K,j}\abs{u_j}$ and take
$\lambda_0\leq1/(2U)$, with the evident convention when $U=0$. For the
second, let $H:=\max_K\abs h$ and $V:=\max_K\abs v$. Since
$\log d\geq\log2$, it suffices to take
\[
 \lambda_0\leq
 \min\left\{\frac{\min\{p,q\}\log2}{2H},\frac1{2V}\right\},
\]
again omitting a bound whose denominator vanishes.

\end{proof}

\begin{definition}[Iterated one-variable derivative]
\label{def:iterated-deriv}
For $f:\R\to\R$, $j:\N$, and $x:\R$, define recursively
\[
 \operatorname{iterDeriv}(f,0,x):=f(x),\qquad
 \operatorname{iterDeriv}(f,j+1,x):=
 \operatorname{deriv}(y\mapsto\operatorname{iterDeriv}(f,j,y))(x).
\]
\end{definition}

\begin{definition}[Iterated derivative in the second variable]
\label{def:lambda-deriv}
For $F:\R\times\R\to\R$, define
\[
 \operatorname{lambdaDeriv}(F,j,a,\lambda):=
 \operatorname{iterDeriv}(\lambda'\mapsto F(a,\lambda'),j,\lambda).
\]
\end{definition}

\begin{definition}[Mixed parameter--line derivative]
\label{def:alpha-lambda-deriv}
For $F:\R\times\R\to\R$, define
\[
 \operatorname{alphaLambdaDeriv}(F,j,a,\lambda):=
 \operatorname{deriv}
  (a'\mapsto\operatorname{lambdaDeriv}(F,j,a',\lambda))(a).
\]
\end{definition}

\begin{definition}[Removable Shannon quotient]
\label{def:removable-shannon-quotient}
For $F:\R\times\R\to\R$, define the total function
\[
 \operatorname{removableQuotient}(F,a,\lambda):=
 \begin{cases}
 F(a,\lambda)/(1-a),&a\neq1,\\
 -\operatorname{alphaLambdaDeriv}(F,0,1,\lambda),&a=1.
 \end{cases}
\]
\end{definition}

\begin{lemma}[Parameterized removable quotient at the Shannon point]
\label{lem:parameterized-removable-quotient}
Let $c,d:\R$ have $h_{cd}:c\leq d$, let
$U,V:\operatorname{Set}(\R)$ have proofs
$h_U:\operatorname{IsOpen}(U)$, $h_V:\operatorname{IsOpen}(V)$,
$h_1:1\in U$, and $h_{cd,V}:\operatorname{Set.Icc}(c,d)\subseteq V$.
Let $F:\R\times\R\to\R$ have proofs
$h_F:\operatorname{ContDiffOn}(\R,3,F,U\times^{\mathrm s}V)$ and
$h_{F,1}:\forall\lambda\in V,F(1,\lambda)=0$.  Then the exact conclusion is
\begin{align*}
 \exists\eps:\R,\quad0<\eps\ \wedge\
 \texttt{let }S:=\operatorname{Set.Icc}(1-\eps,1+\eps)
    \times^{\mathrm s}\operatorname{Set.Icc}(c,d)\texttt{; }\quad
 S\subseteq U\times^{\mathrm s}V\ \wedge\\[-2pt]
 &\forall(j:\N),\ j\leq2\to
 \Bigl(\operatorname{ContinuousOn}
   (p\mapsto\operatorname{lambdaDeriv}
      (\operatorname{removableQuotient}(F),j,p.1,p.2),S)\\
 &\quad\wedge\quad
 \forall p:\R\times\R,\quad p\in S\to
 \operatorname{lambdaDeriv}
   (\operatorname{removableQuotient}(F),j,p.1,p.2)
 =-\int_{s\ \mathrm{in}\ \operatorname{Set.Icc}(0,1)}
 \operatorname{alphaLambdaDeriv}
       (F,j,1+s(p.1-1),p.2)\dd s\Bigr).
\end{align*}
\end{lemma}

\begin{proof}
The fundamental theorem of calculus in the $\alpha$ variable gives
$F(\alpha,\lambda)=(\alpha-1)\int_0^1
\partial_\alpha F(1+s(\alpha-1),\lambda)\dd s$.
Divide by $1-\alpha$ when $\alpha\neq1$; the same integral formula has the
declared value when $\alpha=1$.  Continuous differentiation under the
ordinary integral over $[0,1]$ gives the formulas for $j=1,2$.  Uniform
continuity on the relevant compact product proves joint continuity.
\end{proof}

\begin{lemma}[Twice differentiating a finite signed-measure integral]
\label{lem:signed-integral-differentiate-twice}
Let $A:\texttt{Type u}$ carry
$[\operatorname{MeasurableSpace}A]
[\operatorname{TopologicalSpace}A][\operatorname{CompactSpace}A]
[\operatorname{PseudoMetrizableSpace}A][\operatorname{BorelSpace}A]$.
Let $U:\operatorname{Set}(\R)$ have proof $h_U:\operatorname{IsOpen}(U)$, let
$\mu:\operatorname{SignedMeasure}(A)$, and let
$k_0,k_1,k_2:A\to\R\to\R$.  Assume
\begin{align*}
 &\forall\lambda\in U,\quad
   \operatorname{Measurable}(a\mapsto k_0(a,\lambda))\ \wedge\
   \operatorname{Measurable}(a\mapsto k_1(a,\lambda))\ \wedge\
   \operatorname{Measurable}(a\mapsto k_2(a,\lambda)),\\
 &\forall a:A,\forall\lambda\in U,\quad
   \operatorname{HasDerivAt}(k_0(a),k_1(a,\lambda),\lambda)\ \wedge\
   \operatorname{HasDerivAt}(k_1(a),k_2(a,\lambda),\lambda),\\
 &\forall K:\operatorname{Set}(\R),\quad
   \operatorname{IsCompact}(K)\to K\subseteq U\to
   \exists g:A\to\R,\operatorname{Integrable}
     (g,\operatorname{signedTV}(\mu))\ \wedge\\[-2pt]
 &\hspace{23mm}\forall a:A,\forall\lambda\in K,\quad
   \abs{k_0(a,\lambda)}\leq g(a)\ \wedge\
   \abs{k_1(a,\lambda)}\leq g(a)\ \wedge\
   \abs{k_2(a,\lambda)}\leq g(a).
\end{align*}
Then, for every $\lambda\in U$, the literal conjunction
\begin{align*}
 &\operatorname{HasDerivAt}
  \left(s\mapsto\operatorname{signedIntegral}
    (\mu,a\mapsto k_0(a,s)),
    \operatorname{signedIntegral}(\mu,a\mapsto k_1(a,\lambda)),\lambda\right)\\
 &\quad\wedge\quad
 \operatorname{HasDerivAt}
  \left(s\mapsto\operatorname{signedIntegral}
    (\mu,a\mapsto k_1(a,s)),
    \operatorname{signedIntegral}(\mu,a\mapsto k_2(a,\lambda)),\lambda\right)
\end{align*}
holds.
\end{lemma}

\begin{proof}
Choose a compact neighbourhood of $\lambda$ inside $U$.  Apply dominated
differentiation first to $k_0$ and then to $k_1$, separately for
\texttt{signedPos mu} and \texttt{signedNeg mu}; subtract the two identities.
\end{proof}

\begin{definition}[Shannon line slope]
\label{def:shannon-line-slope}
For finite nonempty $I$, $L:\operatorname{PositiveLineData}(I)$, and
$\lambda:\R$, define
\[
 \operatorname{shannonLineSlope}(L,\lambda):=
 -\sum_{i:I}\operatorname{lineProb}(L,\lambda)(i)
 \bigl(\operatorname{effectiveVelocity}(L,\lambda,i)
       -\operatorname{escortMean}(L,1,\lambda)\bigr)
 \log(\operatorname{lineProb}(L,\lambda)(i)).
\]
\end{definition}

\begin{lemma}[Exact entropy derivatives]\label{lem:exact-derivatives}
Let $I$ be finite and nonempty, let
$L:\operatorname{PositiveLineData}(I)$, let $\lambda_0:\R$ satisfy
$0<\lambda_0$, and assume
\[
 \forall\lambda\in\operatorname{Set.Icc}(-\lambda_0,\lambda_0),\quad
 \operatorname{LinePositive}(L,\lambda).
\]
Put $U_0:=\operatorname{Set.Ioo}(-\lambda_0,\lambda_0)$ and
$K_0:=\operatorname{Set.Icc}(-\lambda_0,\lambda_0)$.  Then the following
literal conjunction holds:
\begin{align*}
 &\Bigl(\forall\lambda\in U_0,\forall\alpha:\R,\
  0<\alpha\to\alpha\neq1\to\\[-2pt]
 &\quad
  \operatorname{entropyLineFirst}
   (L,\operatorname{finiteParam}(\alpha),\lambda)
  =\omega(\operatorname{finiteParam}(\alpha))
   \bigl(\operatorname{escortMean}(L,\alpha,\lambda)
         -\operatorname{escortMean}(L,1,\lambda)\bigr)\ \wedge\\[-2pt]
 &\quad
  \operatorname{entropyLineSecond}
   (L,\operatorname{finiteParam}(\alpha),\lambda)
  =-\alpha\operatorname{escortVar}(L,\alpha,\lambda)
   +\omega(\operatorname{finiteParam}(\alpha))
    \bigl(\operatorname{escortMean}(L,1,\lambda)^2
          -\operatorname{escortMean}(L,\alpha,\lambda)^2\bigr)
 \Bigr)\tag{F}\label{eq:H-finite-general}\\
 &\quad\wedge\Bigl(\forall\alpha:\R,0<\alpha\to\alpha\neq1\to\\[-2pt]
 &\quad
  \operatorname{entropyLineFirst}
   (L,\operatorname{finiteParam}(\alpha),0)
  =\omega(\operatorname{finiteParam}(\alpha))
   \bigl(\operatorname{escortMean}(L,\alpha,0)
         -\operatorname{escortMean}(L,1,0)\bigr)\ \wedge\\[-2pt]
 &\quad
  \operatorname{entropyLineSecond}
   (L,\operatorname{finiteParam}(\alpha),0)
  =-\alpha\operatorname{escortVar}(L,\alpha,0)
   +\omega(\operatorname{finiteParam}(\alpha))
    \bigl(\operatorname{escortMean}(L,1,0)^2
          -\operatorname{escortMean}(L,\alpha,0)^2\bigr)\ \wedge\\[-2pt]
 &\quad
  \operatorname{entropyLineSecond}
   (L,\operatorname{finiteParam}(\alpha),0)
  =\frac{\alpha(\alpha-1)\operatorname{escortSecond}(L,\alpha,0)
  -\alpha^2\operatorname{escortMean}(L,\alpha,0)^2
  +\alpha\operatorname{escortMean}(L,1,0)^2}{1-\alpha}
 \Bigr)\tag{F0}\label{eq:H-finite-at-zero}\\
 &\quad\wedge\Bigl(\forall\lambda\in U_0,\\[-2pt]
 &\quad
  \operatorname{entropyLineFirst}(L,0,\lambda)=0\ \wedge\
  \operatorname{entropyLineSecond}(L,0,\lambda)=0\ \wedge\\[-2pt]
 &\quad
  \operatorname{entropyLineFirst}(L,1,\lambda)
   =\operatorname{shannonLineSlope}(L,\lambda)\ \wedge\\[-2pt]
 &\quad
  \operatorname{entropyLineSecond}(L,1,\lambda)
   =-\operatorname{escortVar}(L,1,\lambda)
    -2\operatorname{escortMean}(L,1,\lambda)
       \operatorname{shannonLineSlope}(L,\lambda)
 \Bigr)\tag{E}\label{eq:H-endpoint-derivatives}\\
 &\quad\wedge\Bigl(\forall i_*:I,\
   \operatorname{FixedMaxCoordinate}(L,K_0,i_*)\to
   \forall\lambda\in U_0,\\[-2pt]
 &\quad
  \operatorname{entropyLineFirst}(L,\top,\lambda)
   =\operatorname{escortMean}(L,1,\lambda)
    -\operatorname{effectiveVelocity}(L,\lambda,i_*)\ \wedge\\[-2pt]
 &\quad
  \operatorname{entropyLineSecond}(L,\top,\lambda)
   =\operatorname{effectiveVelocity}(L,\lambda,i_*)^2
    -\operatorname{escortMean}(L,1,\lambda)^2
 \Bigr)\tag{T}\label{eq:H-infty-derivatives}\\
 &\quad\wedge\Bigl(\forall i_*:I,\
   \operatorname{FixedMaxCoordinate}(L,K_0,i_*)\to\\[-2pt]
 &\quad
  \operatorname{ContinuousOn}
   (p\mapsto\operatorname{entropyLine}(L,p.1,p.2),
    \operatorname{Set.univ}\times^{\mathrm s}K_0)\ \wedge\\[-2pt]
 &\quad
  \operatorname{ContinuousOn}
   (p\mapsto\operatorname{entropyLineFirst}(L,p.1,p.2),
    \operatorname{Set.univ}\times^{\mathrm s}K_0)\ \wedge\\[-2pt]
 &\quad
  \operatorname{ContinuousOn}
   (p\mapsto\operatorname{entropyLineSecond}(L,p.1,p.2),
    \operatorname{Set.univ}\times^{\mathrm s}K_0)
 \Bigr).
\end{align*}
\end{lemma}

\begin{proof}
Within this proof only, write $x:=L.x$, $u:=L.u$, and
$x(\lambda):=\operatorname{lineRaw}(L,\lambda)$.
For finite $\alpha\neq1$,
\[
 H_\alpha(\widehat{x(\lambda)})
 =\frac{1}{1-\alpha}
 \left[\log\sum_i x_i^\alpha(1+u_i\lambda)^\alpha
       -\alpha\log\sum_i x_i(1+u_i\lambda)\right].
\]
Two differentiations at an arbitrary $\lambda$ give
\cref{eq:H-finite-general}; equivalently, re-center the
line at $\lambda$ and use effective velocities $w_i(\lambda)$.  Setting
$\lambda=0$ gives all three finite-order identities in
\cref{eq:H-finite-at-zero}.

At order zero the support is constant on the positive interval.  At order
one, differentiation of
$-\sum_ip_i(\lambda)\log p_i(\lambda)$ uses
\[
 p_i'(\lambda)=p_i(\lambda)(w_i(\lambda)-m(\lambda)),
 \qquad m'(\lambda)=-m(\lambda)^2,
\]
and yields the endpoint identities in \cref{eq:H-endpoint-derivatives}.  At infinity,
\[
 H_\infty(\widehat{x(\lambda)})
 =\log\norm{x(\lambda)}_1-\log\norm{x(\lambda)}_\infty.
\]
The total mass and the common maximal coordinate are affine functions of
$\lambda$, so two logarithmic differentiations give
\cref{eq:H-infty-derivatives}.

For continuity through the Shannon point, let the bracketed expression in
the first display be $F(\alpha,\lambda)$.  Finite sums of positive real
powers show that it has all required continuous mixed derivatives near
$\alpha=1$, it satisfies $F(1,\lambda)=0$, and
$H_1=-\partial_\alpha F(1,\lambda)$.  Apply
\cref{lem:parameterized-removable-quotient} to obtain the joint continuous
extension of the entropy and both $\lambda$-derivatives.  Continuity at zero
follows from local constancy.  At infinity, divide every power sum by the
largest coordinate power.  All nonmaximal ratios converge to zero uniformly
on the compact $\lambda$ interval because the maximal set is fixed and
separated from the remaining coordinates; the maximal velocities agree.
The escort expressions then converge uniformly to
\cref{eq:H-infty-derivatives}.  These arguments also prove joint continuity.
\end{proof}
\begin{definition}[Integrated entropy along a line]
\label{def:integrated-entropy-line}
For finite nonempty $I$, $L:\operatorname{PositiveLineData}(I)$,
$\mu:\operatorname{SignedMeasure}(\overline\R_+)$, and $\lambda:\R$, define
\[
 \operatorname{integratedEntropyLine}(\mu,L,\lambda):=
 \operatorname{signedIntegral}
  (\mu,a\mapsto\operatorname{entropyLine}(L,a,\lambda)).
\]
\end{definition}

\begin{definition}[Signed logarithmic column line]
\label{def:signed-log-phi-line}
For finite nonempty $I$, $L:\operatorname{PositiveLineData}(I)$,
$\mu:\operatorname{SignedMeasure}(\overline\R_+)$, and $\lambda:\R$, define
\[
 \operatorname{signedLogPhiLine}(\mu,L,\lambda):=
 \log\!\left(\sum_{i:I}\operatorname{lineRaw}(L,\lambda)(i)\right)
 +\operatorname{integratedEntropyLine}(\mu,L,\lambda).
\]
\end{definition}

\begin{lemma}[Exact bridge from line kernels to signed column functions]
\label{lem:signed-line-column-bridge}
Let $I$ be finite and nonempty, let
$L:\operatorname{PositiveLineData}(I)$,
$\mu:\operatorname{SignedMeasure}(\overline\R_+)$, and let
$\lambda:\R$ have proof $h:\operatorname{LinePositive}(L,\lambda)$.
Then the following literal conjunction holds:
\begin{align*}
&\operatorname{lineProb}(L,\lambda)
 =\operatorname{normalize}(\operatorname{linePosCone}(L,\lambda,h))\\
&\quad\wedge\quad
 \operatorname{integratedEntropyLine}(\mu,L,\lambda)
 =\operatorname{GSigned}_I
   (\mu,\operatorname{linePosCone}(L,\lambda,h))\\
&\quad\wedge\quad
 0<\operatorname{PhiSigned}_I(\mu,\operatorname{lineCone}(L,\lambda,h))\\
&\quad\wedge\quad
 \operatorname{signedLogPhiLine}(\mu,L,\lambda)
 =\log\!\left(\operatorname{PhiSigned}_I
   (\mu,\operatorname{lineCone}(L,\lambda,h))\right)\\
&\quad\wedge\quad
 \operatorname{lineConeTotal}(L,\lambda)
   =\operatorname{lineCone}(L,\lambda,h)\\
&\quad\wedge\quad
 \operatorname{linePosConeTotal}(L,\lambda)
   =\operatorname{linePosCone}(L,\lambda,h).
\end{align*}
\end{lemma}

\begin{proof}
The case distinctions defining \texttt{lineProb}, \texttt{lineConeTotal},
and \texttt{linePosConeTotal} all select their positive branches under $h$;
proof irrelevance identifies the stored witnesses.  Expanding
\cref{eq:Phi-mu,eq:G-mu} gives the two functional identities.  The cone has
strictly positive total mass and the exponential factor is positive, proving
the strict positivity needed by the logarithm.
\end{proof}

\begin{lemma}[Differentiation through the parameter integral]
\label{lem:differentiate-integral}
Let $I$ be finite and nonempty, let
$L:\operatorname{PositiveLineData}(I)$, let $\lambda_0:\R$ satisfy
$0<\lambda_0$, put $U_0:=\operatorname{Set.Ioo}(-\lambda_0,\lambda_0)$ and
$K_0:=\operatorname{Set.Icc}(-\lambda_0,\lambda_0)$, and assume
$\forall\lambda\in K_0,\operatorname{LinePositive}(L,\lambda)$.
Let $i_*:I$ satisfy
$\operatorname{FixedMaxCoordinate}(L,K_0,i_*)$, and let
$\mu:\operatorname{SignedMeasure}(\overline\R_+)$.  Then, for every
$\lambda\in U_0$, the literal conjunction
\begin{align*}
 &\operatorname{HasDerivAt}
  (\operatorname{integratedEntropyLine}(\mu,L),
   \operatorname{signedIntegral}
    (\mu,a\mapsto\operatorname{entropyLineFirst}(L,a,\lambda)),\lambda)\\
 &\quad\wedge\quad
 \operatorname{HasDerivAt}
 \left(s\mapsto\operatorname{signedIntegral}
    (\mu,a\mapsto\operatorname{entropyLineFirst}(L,a,s)),
   \operatorname{signedIntegral}
    (\mu,a\mapsto\operatorname{entropyLineSecond}(L,a,\lambda)),\lambda\right)\\
 &\quad\wedge\quad
 \operatorname{iterDeriv}
   (\operatorname{signedLogPhiLine}(\mu,L),1,\lambda)
 =\operatorname{deriv}
   \left(s\mapsto\log\!\sum_{i:I}\operatorname{lineRaw}(L,s)(i)\right)
   (\lambda)
  +\operatorname{signedIntegral}
    (\mu,a\mapsto\operatorname{entropyLineFirst}(L,a,\lambda))\\
 &\quad\wedge\quad
 \operatorname{iterDeriv}
   (\operatorname{signedLogPhiLine}(\mu,L),2,\lambda)
 =\operatorname{secondDeriv}
   \left(s\mapsto\log\!\sum_{i:I}\operatorname{lineRaw}(L,s)(i),\lambda\right)
  +\operatorname{signedIntegral}
    (\mu,a\mapsto\operatorname{entropyLineSecond}(L,a,\lambda))
\end{align*}
holds.
\end{lemma}

\begin{proof}
By \cref{lem:exact-derivatives}, the entropy and its first two
$\lambda$-derivatives are jointly continuous on the compact set
$[0,+\infty]\times[-\lambda_0,\lambda_0]$. They are Borel in the order and
bounded there by one common constant, which is integrable against the finite
measure $\abs\mu$.  Apply
\cref{lem:signed-integral-differentiate-twice}.  The formulas for
$G_\mu$ and $\log\Phi_\mu$ follow from their definitions.
\end{proof}

\begin{lemma}[Null thresholds]\label{lem:null-thresholds}
Let $\nu:\operatorname{Measure}(\overline\R_+)$ carry an
$[\operatorname{IsFiniteMeasure}(\nu)]$ instance.  Then the following
literal conjunction holds:
\begin{align*}
 &\operatorname{Set.Countable}
   \{a:\overline\R_+\mid\nu(\{a\})\neq0\}\\
 &\quad\wedge\Bigl(\forall U:\operatorname{Set}(\overline\R_+),\
   \operatorname{IsOpen}(U)\to U.\operatorname{Nonempty}\to
   \exists a\in U,\nu(\{a\})=0\Bigr)\\
 &\quad\wedge\Bigl(\forall s:\R,\
   0<\nu(\operatorname{Set.Ioi}(\operatorname{finiteParam}(s)))\to
   \exists r:\R,s<r\ \wedge\nu(\{\operatorname{finiteParam}(r)\})=0\ \wedge\
   0<\nu(\operatorname{Set.Ioi}(\operatorname{finiteParam}(r)))\Bigr)\\
 &\quad\wedge\Bigl(\forall a_1,a_2:\R,\ 0\leq a_1\to a_1<a_2\to
   \operatorname{finiteParam}(a_1)\in\operatorname{suppMeasure}(\nu)\to
   \operatorname{finiteParam}(a_2)\in\operatorname{suppMeasure}(\nu)\to\\[-2pt]
 &\hspace{34mm}
   \exists r:\R,a_1<r\ \wedge\ r<a_2\ \wedge\
   \nu(\{\operatorname{finiteParam}(r)\})=0\Bigr)\\
 &\quad\wedge\Bigl(\exists r:\N\to\R,\operatorname{StrictMono}(r)\ \wedge\
   (\forall n,0<r(n)\ \wedge\ r(n)<1\ \wedge\
    \nu(\{\operatorname{finiteParam}(r(n))\})=0)\ \wedge\
   \operatorname{Tendsto}(r,\texttt{atTop},\mathcal N(1))\Bigr)\\
 &\quad\wedge\Bigl(\exists s:\N\to\R,\operatorname{StrictAnti}(s)\ \wedge\
   (\forall n,1<s(n)\ \wedge\
    \nu(\{\operatorname{finiteParam}(s(n))\})=0)\ \wedge\
   \operatorname{Tendsto}(s,\texttt{atTop},\mathcal N(1))\Bigr).
\end{align*}
For a signed measure, the exact specialization is obtained by passing
$\operatorname{signedTV}(\mu)$ as $\nu$.
\end{lemma}

\begin{proof}
For each $n$, there are only finitely many atoms of mass at least $1/n$;
therefore the set of all atoms is countable.  The complement of a countable
set is dense in every nonempty real interval.  Choose first a null
$r_1\in(0,1)$ and a finite null $s_1\in(1,2)$.  For $n\geq2$, recursively
choose
$r_n\in(\max\{r_{n-1},1-1/n\},1)$ and
$s_n\in(1,\min\{s_{n-1},1+1/n\})$ outside the atom set.  The intervals are
nonempty, which proves strict monotonicity and convergence to $1$. If
$\nu((s,+\infty])>0$, choose
$x\in\supp(\nu)\cap(s,+\infty]$. If $x<+\infty$, choose a null point
$r\in(s,x)$; then a sufficiently small neighbourhood of $x$ is contained in
$(r,+\infty]$ and has positive measure. If $x=+\infty$, every set
$(r,+\infty]$ with finite $r$ is a neighbourhood of $+\infty$, so choose a
finite null $r>s$. The separation statement is proved in the same way.
\end{proof}

\begin{lemma}[Two finite null thresholds between upper parameters]
\label{lem:two-upper-null-thresholds}
Let $\nu:\operatorname{Measure}(\overline\R_+)$ carry an
$[\operatorname{IsFiniteMeasure}(\nu)]$ instance, and let
$\alpha_1,\alpha_2:\overline\R_+$ satisfy
$1<\alpha_1$ and $\alpha_1<\alpha_2$.  Then
\begin{align*}
\exists a,b:\R,\quad&
 1<a\ \wedge\
 \operatorname{finiteParam}(a)<\alpha_1\ \wedge\
 \alpha_1<\operatorname{finiteParam}(b)\ \wedge\
 \operatorname{finiteParam}(b)<\alpha_2\\[-2pt]
&{}\wedge\
 \nu(\{\operatorname{finiteParam}(a)\})=0\ \wedge\
 \nu(\{\operatorname{finiteParam}(b)\})=0.
\end{align*}
\end{lemma}

\begin{proof}
The strict inequality $\alpha_1<\alpha_2$ proves
$h_{\alpha_1}:\alpha_1\neq\top$.  Put
\[
 r_1:=((\alpha_1.\operatorname{untop}(h_{\alpha_1})
       :\R_{\geq0}):\R).
\]
Then $\operatorname{finiteParam}(r_1)=\alpha_1$ and $1<r_1$.
Choose $a$ outside the countable atom set in the nonempty open interval
$(1,r_1)$.  If $\alpha_2=\top$, choose $b$ outside that set in
$(r_1,r_1+1)$; otherwise convert $\alpha_2$ by \texttt{untop} and choose
$b$ in the open interval between the two finite real values.  Rewriting the
two \texttt{finiteParam} identities gives every displayed inequality and
both null-atom clauses.
\end{proof}

\begin{lemma}[Support gives strict weighted mass]
\label{lem:support-strict-integral}
Let $\nu:\operatorname{Measure}(\overline\R_+)$ carry an
$[\operatorname{IsFiniteMeasure}(\nu)]$ instance.  Let
$x:\overline\R_+$, $U,V:\operatorname{Set}(\overline\R_+)$,
$f:\overline\R_+\to\R$, and $c:\R$ satisfy
\begin{align*}
 &x\in\operatorname{suppMeasure}(\nu),\quad
 \operatorname{IsOpen}(U),\quad x\in U,\\
 &\operatorname{IsOpen}(V),\quad x\in V,\quad V\subseteq U,\\
 &\operatorname{Measurable}(f),\quad
 \operatorname{IntegrableOn}(f,U,\nu),\quad0<c.
\end{align*}
Then the literal conjunction
\begin{align*}
 &0<\nu(U)\ \wedge\ 0<\nu(V)\\
 &\quad\wedge\Bigl(
   (\forall a\in U,0\leq f(a))\to
   (\forall a\in V,c\leq f(a))\to
   0<c\operatorname{ENNReal.toReal}(\nu(V))\ \wedge\
   c\operatorname{ENNReal.toReal}(\nu(V))\leq\int_Uf\dd\nu\Bigr)\\
 &\quad\wedge\Bigl(
   (\forall a\in U,f(a)\leq0)\to
   (\forall a\in V,f(a)\leq-c)\to
   \int_Uf\dd\nu\leq-c\operatorname{ENNReal.toReal}(\nu(V))\ \wedge\
   -c\operatorname{ENNReal.toReal}(\nu(V))<0\Bigr)
\end{align*}
holds.
\end{lemma}

\begin{proof}
The first two assertions are the definition of topological support applied
to the open neighbourhoods $U$ and $V$.  Monotonicity of the restricted integral, applied to
$c\mathbf1_V\leq f\mathbf1_U$, gives the displayed bound; finiteness of
$\nu$ makes the real coercion of $\nu(V)$ positive.  Apply the result to
$-f$ for the negative statement.
\end{proof}

\begin{lemma}[Scalar measure support and null points]
\label{lem:scalar-measure-support}
Let $c:\R$ have $h_c:c\neq0$, and let
$\sigma:\operatorname{FiniteMeasure}(\overline\R_+)$.  Put
$\mu:=c\mathbin{\bullet}\operatorname{signedLift}(\sigma)$.  Then the
literal conjunction
\begin{align*}
 &\operatorname{signedTV}(\mu)
 =\operatorname{ENNReal.ofReal}(\abs c)\mathbin{\bullet}
   \operatorname{finiteMeasure}(\sigma)\\
 &\quad\wedge\quad
 \operatorname{suppSigned}(\mu)
 =\operatorname{suppMeasure}(\operatorname{finiteMeasure}(\sigma))\\
 &\quad\wedge\quad
 \forall r:\overline\R_+,\quad
 \operatorname{signedTV}(\mu)(\{r\})=0
 \Longleftrightarrow
 \operatorname{finiteMeasure}(\sigma)(\{r\})=0
\end{align*}
holds.
\end{lemma}

\begin{proof}
The Jordan decomposition is $(c\nu,0)$ for $c>0$ and
$(0,-c\nu)$ for $c<0$.  The support and null claims follow.
\end{proof}

\subsection{Quantified block estimates}

\begin{definition}[Finite dependent block carrier]
\label{def:block-index}
For $J:\N$ and $n:\operatorname{Fin}(J+1)\to\N$, define
\[
 \operatorname{BlockIndex}(n):=
 \mathop{\Sigma}_{j:\operatorname{Fin}(J+1)}\operatorname{Fin}(n(j)).
\]
This sigma type has its canonical \texttt{Fintype} instance.  It has a
\texttt{Nonempty} instance whenever $\forall j,0<n(j)$; the instance is
installed locally from that proof and is never inferred from prose.
\end{definition}

\begin{definition}[Constant-on-block vector]
\label{def:block-vector}
For $J:\N$, $n:\operatorname{Fin}(J+1)\to\N$, and
$b:\operatorname{Fin}(J+1)\to\R$, define
\[
 \operatorname{blockVec}(n,b):\operatorname{BlockIndex}(n)\to\R,
 \qquad
 \operatorname{blockVec}(n,b)(i):=b(i.1).
\]
Thus the first projection of the sigma index selects the block, and its
second projection selects a coordinate inside that block.  The notation
$\col(b_0\one_{n_0},\ldots,b_J\one_{n_J})$ is only a typeset abbreviation
for this definition.
\end{definition}

\begin{definition}[Asymptotic block data]
\label{def:block-data}
For $J:\N$, define the structure
\[
 \operatorname{BlockData}(J):=
 \{m:\operatorname{Fin}(J+1)\to\R,\ 
   m_{\rm pos}:\forall j,0<m(j),\
   a:\operatorname{Fin}(J+1)\to\R\}.
\]
Define the total asymptotic scale
\[
 \operatorname{blockScale}(n):=(n+2:\R),\qquad n:\N.
\]
For $B:\operatorname{BlockData}(J)$, $n:\N$, and
$j:\operatorname{Fin}(J+1)$, define the total natural block count
\begin{align*}
 \operatorname{blockCount}(B,n,j)
 &:=\operatorname{Nat.ceil}
    (\operatorname{Real.rpow}(\operatorname{blockScale}(n),B.m(j))),\\
 \operatorname{blockCounts}(B,n)&:=(j\mapsto
    \operatorname{blockCount}(B,n,j)),\\
 \operatorname{BlockCarrier}(B,n)&:=
    \operatorname{BlockIndex}(\operatorname{blockCounts}(B,n)).
\end{align*}
The count is positive for every $n$ because $2\leq\operatorname{blockScale}(n)$
and $B.m(j)>0$.
\end{definition}

\begin{lemma}[Positive block counts]
\label{lem:block-count-pos}
For $J:\N$, $B:\operatorname{BlockData}(J)$, $n:\N$, and
$j:\operatorname{Fin}(J+1)$, the named proof is
\[
 \operatorname{blockCount\_pos}(B,n,j):
 0<\operatorname{blockCount}(B,n,j).
\]
\end{lemma}

\begin{proof}
The positive real power of the base $n+2\geq2$ is positive, and the natural
ceiling of a positive real number is positive.
\end{proof}

\begin{definition}[Nonempty block-carrier instance]
\label{def:block-carrier-nonempty}
For $J:\N$, $B:\operatorname{BlockData}(J)$, and $n:\N$, define the local
instance-producing term
\[
 \operatorname{blockCarrierNonempty}(B,n):
 \operatorname{Nonempty}(\operatorname{BlockCarrier}(B,n))
 :=\left\langle
   \left\langle0,
    \left\langle0,\operatorname{blockCount\_pos}(B,n,0)\right\rangle
   \right\rangle
 \right\rangle.
\]
Every declaration below whose typeclass search needs this fact begins with
\texttt{letI := blockCarrierNonempty B n}.
\end{definition}

\begin{definition}[Nonempty arbitrary block index]
\label{def:block-index-nonempty}
For $J:\N$, $n:\operatorname{Fin}(J+1)\to\N$,
$h_n:\forall j,1\leq n(j)$, define
\[
 \operatorname{blockIndexNonempty}(n,h_n):
  \operatorname{Nonempty}(\operatorname{BlockIndex}(n))
 :=\left\langle
   \left\langle0,\left\langle0,
    (\operatorname{Nat.succ\_le\_iff}.1(h_n(0)))
   \right\rangle\right\rangle
 \right\rangle.
\]
\end{definition}

\begin{definition}[Maximum of a nonempty block vector]
\label{def:block-max}
For $J:\N$, $n:\operatorname{Fin}(J+1)\to\N$,
$h_n:\forall j,1\leq n(j)$, and $b:\operatorname{Fin}(J+1)\to\R$, define
\[
 \operatorname{blockMax}(n,h_n,b):=
 \left(\texttt{letI := blockIndexNonempty n h\_n; }\quad
 \operatorname{finMax}(i\mapsto\operatorname{blockVec}(n,b)(i))\right).
\]
\end{definition}

\begin{definition}[Typed block line]
\label{def:block-line-data}
Let $J:\N$, $B:\operatorname{BlockData}(J)$, $n:\N$, and
$u:\operatorname{Fin}(J+1)\to\R$.  Define
\begin{align*}
 \operatorname{blockBase}(B,n)(i)
 &:=\operatorname{Real.rpow}(\operatorname{blockScale}(n),B.a(i.1)),\\
 \operatorname{blockVelocity}(u)(i)&:=u(i.1),\\
 \operatorname{blockLineData}(B,n,u)
 &:=\left(\texttt{letI := blockCarrierNonempty B n; }\quad
 \left\langle
   \operatorname{blockBase}(B,n),\operatorname{blockVelocity}(u),
   \left(i\mapsto\operatorname{Real.rpow\_pos\_of\_pos}
    (\texttt{by positivity}:0<\operatorname{blockScale}(n))
    (B.a(i.1))\right)
   \right\rangle
 \right)
 \in\operatorname{PositiveLineData}(\operatorname{BlockCarrier}(B,n)).
\end{align*}
The corresponding raw vector is the declared function
\[
 \operatorname{blockLineRaw}(B,n,u,\lambda):=
 \left(\texttt{letI := blockCarrierNonempty B n; }\quad
 \operatorname{lineRaw}(\operatorname{blockLineData}(B,n,u),\lambda)\right).
\]
It has the coordinate formula
$\operatorname{blockScale}(n)^{B.a(j)}(1+u(j)\lambda)$ on block $j$.
\end{definition}

\begin{definition}[Block exponents, escort weights, and derivative kernels]
\label{def:block-kernels}
For $J:\N$, $B:\operatorname{BlockData}(J)$, $n:\N$, real
$\alpha$, and $j:\operatorname{Fin}(J+1)$, define
\begin{align*}
 \operatorname{blockExponent}(B,j,\alpha)&:=B.m(j)+B.a(j)\alpha,\\
 \operatorname{blockContribution}(B,n,j,\alpha)&:=
  (\operatorname{blockCount}(B,n,j):\R)
  \operatorname{Real.rpow}
   (\operatorname{Real.rpow}(\operatorname{blockScale}(n),B.a(j)),\alpha),\\
 \operatorname{blockEscort}(B,n,j,\alpha)&:=
  \frac{\operatorname{blockContribution}(B,n,j,\alpha)}
  {\sum_{\ell:\operatorname{Fin}(J+1)}
    \operatorname{blockContribution}(B,n,\ell,\alpha)},\\
 \operatorname{blockEscortMean}(B,n,\alpha,u)&:=
   \sum_j\operatorname{blockEscort}(B,n,j,\alpha)u(j),\\
 \operatorname{blockEscortSecond}(B,n,\alpha,u)&:=
   \sum_j\operatorname{blockEscort}(B,n,j,\alpha)u(j)^2,\\
 \operatorname{blockEscortVar}(B,n,\alpha,u)&:=
   \operatorname{blockEscortSecond}(B,n,\alpha,u)
    -\operatorname{blockEscortMean}(B,n,\alpha,u)^2.
\end{align*}
For $u:\operatorname{Fin}(J+1)\to\R$ and
$\beta:\overline\R_+$, define separately
\begin{align*}
 \operatorname{blockKernelFirst}(B,n,u,\beta)&:=
  \left(\texttt{letI := blockCarrierNonempty B n; }\quad
  \operatorname{entropyLineFirst}
   (\operatorname{blockLineData}(B,n,u),\beta,0)\right),\\
 \operatorname{blockKernelSecond}(B,n,u,\beta)&:=
  \left(\texttt{letI := blockCarrierNonempty B n; }\quad
  \operatorname{entropyLineSecond}
   (\operatorname{blockLineData}(B,n,u),\beta,0)\right).
\end{align*}
Both kernels have type $\overline\R_+\to\R$; hence all later signed
integrals bind a function on exactly the measure's carrier.
\end{definition}

\begin{lemma}[Block-vector finite-sum formulas]
\label{lem:block-vector-formulas}
Let $J:\N$, let $n:\operatorname{Fin}(J+1)\to\N$ satisfy
$h_n:\forall j,1\leq n(j)$, and let
$b:\operatorname{Fin}(J+1)\to\R$ satisfy
$h_b:\forall j,0\leq b(j)$.  Then the following literal conjunction holds:
\begin{align*}
 &(\forall i:\operatorname{BlockIndex}(n),\
      0\leq\operatorname{blockVec}(n,b)(i))\\
 &\quad\wedge\quad
 \norm{\operatorname{blockVec}(n,b)}_1
   =\sum_j(n(j):\R)b(j)\\
 &\quad\wedge\quad
 \Bigl(\forall\alpha:\R,0<\alpha\to
  \sum_i\operatorname{Real.rpow}
    (\operatorname{blockVec}(n,b)(i),\alpha)
  =\sum_j(n(j):\R)\operatorname{Real.rpow}(b(j),\alpha)\Bigr)\\
 &\quad\wedge\quad
 \operatorname{blockMax}(n,h_n,b)
   =\operatorname{finMax}(j\mapsto b(j))\\
 &\quad\wedge\quad
 \Bigl(\operatorname{blockVec}(n,b)\neq0
   \Longleftrightarrow\exists j:\operatorname{Fin}(J+1),0<b(j)\Bigr).
\end{align*}
\end{lemma}

\begin{proof}
Reindex a finite sum on the sigma type as an outer sum over $j$ and an inner
sum over $\operatorname{Fin}(n_j)$.  The inner summand is constant, so the
inner sum is $n_j$ times that value.  The maximum formula follows because
each block is nonempty.  The last assertion is
\cref{lem:fnd-nonzero-sum-witness} specialized to disjoint constant blocks.
\end{proof}

For $J:\N$, $B:\operatorname{BlockData}(J)$, $n:\N$,
$u:\operatorname{Fin}(J+1)\to\R$, and $\lambda:\R$, the expression
\begin{equation}
 x_{B,u}^{(n)}(\lambda):=
 \operatorname{blockLineRaw}(B,n,u,\lambda)
 \label{eq:general-block}
\end{equation}
is a function on the declared dependent carrier
$\operatorname{BlockCarrier}(B,n)$.  At $\lambda=0$, its block exponent at
finite order $\alpha:\R$ is
\begin{equation}
 e_{B,j}(\alpha):=\operatorname{blockExponent}(B,j,\alpha),
 \label{eq:block-exponent}
\end{equation}
and the escort mass of block $j$ is
\begin{equation}
 \Pi_{B,j,n}(\alpha):=\operatorname{blockEscort}(B,n,j,\alpha).
 \label{eq:block-escort}
\end{equation}
All limits indexed by $n$ below are ordinary \texttt{Tendsto} statements
along \texttt{atTop}; the displayed scale is always
$d=\operatorname{blockScale}(n)=n+2$.

\begin{lemma}[Rounding estimates]\label{lem:rounding}
The following literal conjunction holds:
\begin{align*}
&\Bigl(\forall(s:\R)(d:\N),\ 0<s\to1\leq d\to
  1\leq
  \frac{(\operatorname{Nat.ceil}
    (\operatorname{Real.rpow}(d:\R,s)):\R)}
       {\operatorname{Real.rpow}(d:\R,s)}\ \wedge\\[-2pt]
&\hspace{28mm}
  \frac{(\operatorname{Nat.ceil}
    (\operatorname{Real.rpow}(d:\R,s)):\R)}
       {\operatorname{Real.rpow}(d:\R,s)}
  \leq1+\operatorname{Real.rpow}(d:\R,-s)\Bigr)
 \tag{R}\label{eq:rounding-estimate}\\
&\quad\wedge\Bigl(\forall(s:\R)(d:\N),\ 0<s\to1\leq d\to
  0\leq\log
  \frac{(\operatorname{Nat.ceil}
    (\operatorname{Real.rpow}(d:\R,s)):\R)}
       {\operatorname{Real.rpow}(d:\R,s)}\ 
  \wedge\\[-4pt]
&\hspace{30mm}
  \log
  \frac{(\operatorname{Nat.ceil}
    (\operatorname{Real.rpow}(d:\R,s)):\R)}
       {\operatorname{Real.rpow}(d:\R,s)}
  \leq\log(1+\operatorname{Real.rpow}(d:\R,-s))\ 
  \wedge\\[-4pt]
&\hspace{30mm}
  \log(1+\operatorname{Real.rpow}(d:\R,-s))
  \leq\operatorname{Real.rpow}(d:\R,-s)\Bigr)
 \tag{RL}\label{eq:rounding-log-estimate}\\
&\quad\wedge\Bigl(\forall(s_0:\R),\ 0<s_0\to
  (\forall(s:\R)(d:\N),\ s_0\leq s\to1\leq d\to
   \operatorname{Real.rpow}(d:\R,-s)
   \leq\operatorname{Real.rpow}(d:\R,-s_0))\ 
  \wedge\\[-4pt]
&\hspace{28mm}
  \operatorname{Tendsto}
   (d\mapsto\operatorname{Real.rpow}(d:\R,-s_0),
    \texttt{atTop},\mathcal N(0))\Bigr).
\end{align*}
\end{lemma}

\begin{proof}
The first inequalities follow from
$d^s\leq\lceil d^s\rceil<d^s+1$. Divide by $d^s>0$, use monotonicity of
$\log$, and apply $\log(1+x)\leq x$ for $x\geq0$. If $s\geq s_0$, then
$d^{-s}\leq d^{-s_0}$ because $d\geq1$. The two limits follow from
$d^{-s_0}\to0$.
\end{proof}

\begin{lemma}[Dominant-block escort estimate]\label{lem:dominant-block}
Let $J:\N$, $B:\operatorname{BlockData}(J)$,
$I:\operatorname{Set}(\R)$, $k:\operatorname{Fin}(J+1)$, and
$\eta:\R$ satisfy $h_\eta:0<\eta$ and
\begin{equation}
 \forall\alpha\in I,\forall j:\operatorname{Fin}(J+1),\quad
 j\neq k\to
 \operatorname{blockExponent}(B,k,\alpha)
 -\operatorname{blockExponent}(B,j,\alpha)\geq\eta.
 \label{eq:dominant-gap}
\end{equation}
The exact return proposition is
\begin{align}
&\Bigl[\exists C:\R,\quad0\leq C\ \wedge\
 \forall(n:\N)(\alpha:\R),\ \alpha\in I\to\\[-2pt]
&\quad
 \sum_{j\in\operatorname{Finset.univ.erase}(k)}
   \operatorname{blockEscort}(B,n,j,\alpha)
 \leq C\operatorname{Real.rpow}(\operatorname{blockScale}(n),-\eta)
 \label{eq:escort-outside}\\[-2pt]
&\quad\wedge\quad
 \sum_{j:\operatorname{Fin}(J+1)}
  \abs{\operatorname{deriv}
   (s\mapsto\operatorname{blockEscort}(B,n,j,s))(\alpha)}
 \leq C\log(\operatorname{blockScale}(n))
       \operatorname{Real.rpow}(\operatorname{blockScale}(n),-\eta)
 \Bigr]\label{eq:escort-alpha-derivative}\\
&\quad\wedge\Bigl[
 \forall(U:\R),\ 0\leq U\to
 \exists C_U:\R,\ 0\leq C_U\ \wedge\
 \forall(u:\operatorname{Fin}(J+1)\to\R),\
 (\forall j,\abs{u(j)}\leq U)\to\\[-2pt]
 \forall(n:\N),\forall\alpha\in I,\quad
 \Bigl[
 &\abs{\operatorname{blockEscortMean}(B,n,\alpha,u)-u(k)}
 +\abs{\operatorname{blockEscortSecond}(B,n,\alpha,u)-u(k)^2}
 +\operatorname{blockEscortVar}(B,n,\alpha,u)\\[-2pt]
 &\hspace{12mm}\leq
 C_U\operatorname{Real.rpow}(\operatorname{blockScale}(n),-\eta)
 \tag{EM}\label{eq:escort-moment-bound}\\[-2pt]
 &\quad\wedge\quad
 \abs{\operatorname{deriv}
  (s\mapsto\operatorname{blockEscortMean}(B,n,s,u))(\alpha)}
 +\abs{\operatorname{deriv}
  (s\mapsto\operatorname{blockEscortSecond}(B,n,s,u))(\alpha)}\\[-2pt]
 &\hspace{12mm}\leq C_U\log(\operatorname{blockScale}(n))
   \operatorname{Real.rpow}(\operatorname{blockScale}(n),-\eta)
 \tag{EMA}\label{eq:escort-moment-alpha}
 \Bigr]\Bigr].
\end{align}
\end{lemma}

\begin{proof}
By \cref{lem:rounding}, the ratio of the contribution of block $j\neq k$ to
that of block $k$ is at most
$2\operatorname{blockScale}(n)^{-\eta}$. Summing over the
finitely many blocks gives \cref{eq:escort-outside}. Differentiating
\cref{eq:block-escort} yields
\[
 \operatorname{deriv}(\alpha\mapsto
   \operatorname{blockEscort}(B,n,j,\alpha))(\alpha)
 =(B.a(j)-\bar a_n)(\log\operatorname{blockScale}(n))
   \operatorname{blockEscort}(B,n,j,\alpha),
 \qquad
 \bar a_n:=\sum_\ell B.a(\ell)
   \operatorname{blockEscort}(B,n,\ell,\alpha).
\]
The dominant block contribution is bounded by the outside mass because
$\abs{B.a(k)-\bar a_n}\leq
\sum_{j\neq k}\abs{B.a(k)-B.a(j)}
 \operatorname{blockEscort}(B,n,j,\alpha)$. This proves
\cref{eq:escort-alpha-derivative}. The moment bounds follow by writing every
mean as its dominant value plus a sum over $j\neq k$. Differentiating the
means and using the preceding identity gives
\cref{eq:escort-moment-alpha}.
\end{proof}

\begin{lemma}[Near-Shannon cancellation for a dominant block]
\label{lem:near-shannon-dominant}
Let $J:\N$, $B:\operatorname{BlockData}(J)$,
$k:\operatorname{Fin}(J+1)$, and $\delta,\eta,U:\R$.  Assume
$0<\delta$, $\delta<1$, $0<\eta$, $0\leq U$, and
\[
 \forall\alpha\in\operatorname{Set.Icc}(1-\delta,1+\delta),
 \forall j\neq k,\quad
 \eta\leq\operatorname{blockExponent}(B,k,\alpha)
       -\operatorname{blockExponent}(B,j,\alpha).
\]
Then the exact conclusion is
\begin{align}
 &\exists C:\R,\quad0\leq C\ \wedge\\[-2pt]
 &\Bigl[\forall(n:\N)(u:\operatorname{Fin}(J+1)\to\R),\
 (\forall j,\abs{u(j)}\leq U)\to
 \forall\alpha\in\operatorname{Set.Icc}(1-\delta,1+\delta),\\[-2pt]
 &\quad
 \abs{\operatorname{blockKernelFirst}
   (B,n,u,\operatorname{finiteParam}(\alpha))}
 +\abs{\operatorname{blockKernelSecond}
   (B,n,u,\operatorname{finiteParam}(\alpha))}\\[-2pt]
 &\qquad\leq
 C(1+\log(\operatorname{blockScale}(n)))
 \operatorname{Real.rpow}(\operatorname{blockScale}(n),-\eta)
 \Bigr]\label{eq:near-shannon-dominant-bound}\\
 &\quad\wedge\quad\operatorname{Tendsto}
 \left(n\mapsto C(1+\log(\operatorname{blockScale}(n)))
  \operatorname{Real.rpow}(\operatorname{blockScale}(n),-\eta),
  \texttt{atTop},\mathcal N(0)\right).
\end{align}
\end{lemma}

\begin{proof}
By \cref{lem:dominant-block}, there is $C_U:\R$ with $0\leq C_U$ such that
$\abs{\operatorname{blockEscortMean}(B,n,\alpha,u)-u(k)}
 \leq C_U\operatorname{Real.rpow}(\operatorname{blockScale}(n),-\eta)$
and the same estimate holds
at order $1$. More importantly,
\cref{eq:escort-moment-alpha} and the mean-value theorem give
\[
 \abs{\operatorname{blockEscortMean}(B,n,\alpha,u)
       -\operatorname{blockEscortMean}(B,n,1,u)}
 \leq C\abs{\alpha-1}(\log\operatorname{blockScale}(n))
       \operatorname{Real.rpow}(\operatorname{blockScale}(n),-\eta).
\]
Multiplication by $\omega(\operatorname{finiteParam}(\alpha))$ therefore gives the stated bound for the
first derivative, including the continuous value at $\alpha=1$.

For the second derivative, use \cref{eq:H-finite-at-zero}. The same lemma
gives
$\operatorname{blockEscortVar}(B,n,\alpha,u)
 \leq C_U\operatorname{Real.rpow}(\operatorname{blockScale}(n),-\eta)$. Moreover,
\[
 \abs{\operatorname{blockEscortMean}(B,n,1,u)^2
       -\operatorname{blockEscortMean}(B,n,\alpha,u)^2}
 \leq C\abs{\alpha-1}(\log\operatorname{blockScale}(n))
       \operatorname{Real.rpow}(\operatorname{blockScale}(n),-\eta).
\]
The same cancellation against $\omega(\operatorname{finiteParam}(\alpha))$ proves
\cref{eq:near-shannon-dominant-bound}.
\end{proof}

\begin{lemma}[Large-order tail bound]\label{lem:large-alpha-tail}
Let $J:\N$, $B:\operatorname{BlockData}(J)$,
$j_\infty:\operatorname{Fin}(J+1)$, and $U:\R$ satisfy $0\leq U$ and
\[
 \forall j:\operatorname{Fin}(J+1),\quad
 j\neq j_\infty\to B.a(j)<B.a(j_\infty).
\]
Then the exact conclusion is
\begin{align}
 &\exists(A,\gamma,C:\R),\quad
 1<A\ \wedge\ 0<\gamma\ \wedge\ 0\leq C\ \wedge\\[-2pt]
 &\forall(n:\N)(\alpha:\R)
  (u:\operatorname{Fin}(J+1)\to\R),\quad
 A\leq\alpha\to(\forall j,\abs{u(j)}\leq U)\to\\[-2pt]
 &\Bigl[
 \sum_{j\in\operatorname{Finset.univ.erase}(j_\infty)}
  \operatorname{blockEscort}(B,n,j,\alpha)
 \leq C\operatorname{Real.rpow}
    (\operatorname{blockScale}(n),-\gamma\alpha)\\[-2pt]
 &\quad\wedge\quad
 \alpha\sum_{j\in\operatorname{Finset.univ.erase}(j_\infty)}
  \operatorname{blockEscort}(B,n,j,\alpha)\leq C
 \label{eq:large-alpha-escort}\\[-2pt]
 &\quad\wedge\quad
 \abs{\operatorname{blockKernelFirst}
   (B,n,u,\operatorname{finiteParam}(\alpha))}\leq C\\[-2pt]
 &\quad\wedge\quad
 \abs{\operatorname{blockKernelSecond}
   (B,n,u,\operatorname{finiteParam}(\alpha))}\leq C\\[-2pt]
 &\quad\wedge\quad
 \abs{\operatorname{blockKernelFirst}(B,n,u,\top)}\leq C
 \quad\wedge\quad
 \abs{\operatorname{blockKernelSecond}(B,n,u,\top)}\leq C
 \Bigr]
\end{align}
\end{lemma}

\begin{proof}
Put $E:=\operatorname{Finset.univ}.\operatorname{erase}(j_\infty)$.  If
$E=\varnothing$, take $A=2$, $\gamma=1$, and discharge both escort sums by
\texttt{simp}.  Otherwise define the finite slope minimum and intercept
bound
\[
 \delta:=\min_{j\in E}\bigl(B.a(j_\infty)-B.a(j)\bigr),
 \qquad C_0:=\max_{j\in E}\abs{B.m(j_\infty)-B.m(j)},
\]
and retain the named proof $h_\delta:0<\delta$ from the strict amplitude
gaps.  Take
\[
 \gamma:=\delta/2,\qquad
 A:=\max\{2,2C_0/\delta\}.
\]
For $j\in E$ and $A\leq\alpha$, expansion of
\texttt{blockExponent} gives the named bound
\[
 \gamma\alpha\leq
 \operatorname{blockExponent}(B,j_\infty,\alpha)
 -\operatorname{blockExponent}(B,j,\alpha).
\]
Insert this finite family of inequalities into the proof of
\cref{lem:dominant-block} to obtain the first estimate.  The continuous
function
$\alpha\mapsto\alpha\operatorname{Real.rpow}(2,-\gamma\alpha)$ tends to
zero at $+\infty$ and is bounded on $[A,+\infty)$, which gives the second.

For the first derivative, $\omega$ is bounded on $[A,+\infty]$ and the escort
means are bounded by the velocities. For the second derivative, use
\cref{eq:H-finite-at-zero}. The variance is bounded by a constant times the
outside escort mass, so the factor $\alpha$ is controlled by
\cref{eq:large-alpha-escort}; the remaining term is bounded because $\omega$
and the means are bounded.  For the endpoint at $+\infty$, use the stated
strict amplitude gaps and finiteness of $\operatorname{Fin}(J+1)$ to choose
$\eps_\infty>0$ and a proof
\[
 \operatorname{FixedMaxCoordinate}
  (\operatorname{blockLineData}(B,n,u),
   \operatorname{Set.Icc}(-\eps_\infty,\eps_\infty),
   \langle j_\infty,
    \langle0,\operatorname{blockCount\_pos}(B,n,j_\infty)\rangle\rangle).
\]
Every coordinate in the maximal block has the common velocity
$u(j_\infty)$.  The endpoint clause \cref{eq:H-infty-derivatives} applied
with this fixed-max proof gives the final two bounds.
\end{proof}

\begin{definition}[Pointwise limit kernels for a block family]
\label{def:block-limit-kernels}
Let $J:\N$, let
$k:\overline\R_+\to\operatorname{Fin}(J+1)$, and let
$u:\operatorname{Fin}(J+1)\to\R$.  Define
\begin{align*}
 \operatorname{blockLimitFirst}(k,u,\beta)&:=
  \omega(\beta)\bigl(u(k(\beta))-u(k(1))\bigr),\\
 \operatorname{blockLimitSecond}(k,u,\beta)&:=
  \omega(\beta)\bigl(u(k(1))^2-u(k(\beta))^2\bigr).
\end{align*}
These definitions are total at $0$, $1$, and $\top$.  In particular both
values at $1$ are zero because $\omega(1)=0$, regardless of the chosen
value of $k(1)$.
\end{definition}

\begin{proposition}[Uniform passage to the block limit]
\label{prop:block-limit-passage}
Let $J:\N$, $B:\operatorname{BlockData}(J)$,
$\mu:\operatorname{SignedMeasure}(\overline\R_+)$,
$T:\operatorname{Finset}(\R)$,
$k:\overline\R_+\to\operatorname{Fin}(J+1)$, and
$j_\infty:\operatorname{Fin}(J+1)$.  Assume the literal hypotheses
\begin{align*}
 &1\notin T,\qquad
 \forall r\in T,\quad
   \operatorname{signedTV}(\mu)(\{\operatorname{finiteParam}(r)\})=0,\\
 &\operatorname{Measurable}(k),\qquad k(\top)=j_\infty,\\
 &\forall\alpha:\R,\quad0\leq\alpha\to\alpha\notin T\to
   \forall j:\operatorname{Fin}(J+1),\quad
   j\neq k(\operatorname{finiteParam}(\alpha))\to\\[-2pt]
 &\hspace{28mm}
   \operatorname{blockExponent}
      (B,j,\alpha)
   <\operatorname{blockExponent}
      (B,k(\operatorname{finiteParam}(\alpha)),\alpha),\\
 &\forall j:\operatorname{Fin}(J+1),\quad
   j\neq j_\infty\to B.a(j)<B.a(j_\infty).
\end{align*}
Let $K:\operatorname{Set}(\operatorname{Fin}(J+1)\to\R)$ satisfy
$h_K:\operatorname{IsCompact}(K)$ and $h_{K0}:K.\operatorname{Nonempty}$,
Then the exact return proposition first binds the compact-parameter subtype
and installs the required instances:
\begin{align*}
&\texttt{let }U_K:=
 \{u:\operatorname{Fin}(J+1)\to\R\mathbin{//}u\in K\}\texttt{; }\\[-2pt]
&\texttt{letI : Nonempty U\_K := }
 \left\langle\left\langle h_{K0}.\operatorname{choose},
 h_{K0}.\operatorname{choose\_spec}\right\rangle\right\rangle\texttt{; }\\[-2pt]
&\texttt{letI : CompactSpace U\_K :=
 isCompact\_iff\_compactSpace.mp h\_K; }\\[-2pt]
&\Bigl(\forall(n:\N)(u:U_K),\quad
 \operatorname{Measurable}
  (\beta\mapsto\operatorname{blockKernelFirst}(B,n,u.1,\beta))\ \wedge\
 \operatorname{Integrable}
  (\beta\mapsto\operatorname{blockKernelFirst}(B,n,u.1,\beta),
    \operatorname{signedTV}(\mu))\Bigr)\\
&\quad\wedge
 \Bigl(\forall(n:\N)(u:U_K),\quad
 \operatorname{Measurable}
  (\beta\mapsto\operatorname{blockKernelSecond}(B,n,u.1,\beta))\ \wedge\
 \operatorname{Integrable}
  (\beta\mapsto\operatorname{blockKernelSecond}(B,n,u.1,\beta),
    \operatorname{signedTV}(\mu))\Bigr)\\
&\quad\wedge
 \Bigl(\forall u:U_K,\quad
 \operatorname{Measurable}
  (\beta\mapsto\operatorname{blockLimitFirst}(k,u.1,\beta))\ \wedge\
 \operatorname{Integrable}
  (\beta\mapsto\operatorname{blockLimitFirst}(k,u.1,\beta),
    \operatorname{signedTV}(\mu))\Bigr)\\
&\quad\wedge
 \Bigl(\forall u:U_K,\quad
 \operatorname{Measurable}
  (\beta\mapsto\operatorname{blockLimitSecond}(k,u.1,\beta))\ \wedge\
 \operatorname{Integrable}
  (\beta\mapsto\operatorname{blockLimitSecond}(k,u.1,\beta),
    \operatorname{signedTV}(\mu))\Bigr)\\
&\quad\wedge
 \operatorname{Tendsto}\Bigl(
 n\mapsto\operatorname{uniformIntegralErrorSigned}
  \bigl(\mu,
   (n\mapsto\beta\mapsto u\mapsto
     \operatorname{blockKernelFirst}(B,n,u.1,\beta)),
   (\beta\mapsto u\mapsto\operatorname{blockLimitFirst}(k,u.1,\beta)),n\bigr),
 \texttt{atTop},\mathcal N(0)\Bigr)\\
&\quad\wedge
 \operatorname{Tendsto}\Bigl(
 n\mapsto\operatorname{uniformIntegralErrorSigned}
  \bigl(\mu,
   (n\mapsto\beta\mapsto u\mapsto
     \operatorname{blockKernelSecond}(B,n,u.1,\beta)),
   (\beta\mapsto u\mapsto\operatorname{blockLimitSecond}(k,u.1,\beta)),n\bigr),
 \texttt{atTop},\mathcal N(0)\Bigr).
 \tag{UB}\label{eq:uniform-block-integral-limit}
\end{align*}
\end{proposition}

\begin{proof}
The exact derivative formulas give measurability of
\texttt{blockKernelFirst} and \texttt{blockKernelSecond}.
The maximizing block $k$ is a measurable finite-valued function, so
\texttt{blockLimitFirst} and \texttt{blockLimitSecond} are measurable.
The arguments below give a
uniform bound on a neighbourhood of $1$, on each bounded cell, and on the
large-order tail; hence both kernels are integrable against the finite
measure $\operatorname{signedTV}(\mu)$.

Fix $K$.  The subtype $U_K$ inherits its topology, pseudometrizability, and
second-countable topology from the ambient finite-dimensional real vector
space; the two displayed local instances supply compactness and nonemptiness
for \cref{lem:fnd-compact-sup-measurable}.  Choose disjoint open
neighbourhoods of the points of $T$, none
containing $1$.  The exact formulas and compactness of $K$ give one bound
$C_K$ for both kernels on these neighbourhoods.  Since every threshold is
$\operatorname{signedTV}(\mu)$-null and there are finitely many, outer regularity permits the
neighbourhoods to be chosen with total $\operatorname{signedTV}(\mu)$-measure below any prescribed
$\varepsilon>0$.  Their contribution to the supremum of the integral error
is at most $2C_K\varepsilon$.

Choose a closed interval around $1$ disjoint from the threshold
neighbourhoods.  The unique maximizing block has a positive uniform exponent
gap there, and \cref{lem:near-shannon-dominant} gives convergence to zero
uniformly in both $\alpha$ and $u\in K$.  On each remaining bounded closed
component, uniqueness and continuity of the finitely many affine differences
give a positive uniform gap; \cref{lem:dominant-block} then gives uniform
convergence to the corresponding named block-limit kernel.  Finally, choose
$A,\gamma,C$ from \cref{lem:large-alpha-tail}.  Uniformly for
$\alpha\geq A$ and $u\in K$, it gives
an error bounded by
\[
 C_K\left(\operatorname{Real.rpow}
   (\operatorname{blockScale}(n),-\eta_1)
 +(1+\alpha)\operatorname{Real.rpow}
   (\operatorname{blockScale}(n),-\gamma\alpha)\right).
\]
When there is another block, choose $\eta_1>0$ as the \texttt{finMin} of
the positive gaps
$\operatorname{blockExponent}(B,k(\operatorname{finiteParam}(1)),1)
-\operatorname{blockExponent}(B,j,1)$ on the finite nonempty subtype
$\{j:\operatorname{Fin}(J+1)\mathbin{//}
  j\neq k(\operatorname{finiteParam}(1))\}$; in the one-block case
the error is identically zero and set $\eta_1=1$.
The first term is the finite-scale error in the order-one escort mean and is
independent of $\alpha$; it cannot be omitted.  The supremum of the displayed
bound over $\alpha\geq A$ tends to zero.  At $+\infty$, the second term is
interpreted as zero and the scale-to-$-\eta_1$ term gives the required control,
including a possible atom there.

The integral error outside the threshold neighbourhoods therefore tends
uniformly to zero.  Take the limit superior, then let $\varepsilon\downarrow
0$, to obtain \cref{eq:uniform-block-integral-limit}.  This is precisely the
signed uniform dominated-convergence interface in
\cref{lem:fnd-uniform-dct-signed}.
\end{proof}

\subsection{Two-block and three-block localization}

\begin{definition}[Specialized block data and signed truncations]
\label{def:special-block-data}
For real $r,R$ with proofs $h_r:0<r$ and $h_R:r<R$, define
$\operatorname{twoBlockData}(r,R,h_r,h_R):\operatorname{BlockData}(1)$ by
\begin{align*}
 m(i)&:=\operatorname{Fin.cases}
       (R,(\_\mapsto R-r),i),&
 \texttt{m\_pos}(i)&:=
       \texttt{by fin\_cases i <;> simp <;> linarith},&
 a(i)&:=\operatorname{Fin.cases}(0,(\_\mapsto1),i).
\end{align*}
For real $a,b,R$ with proof arguments $h_a:0<a$, $h_{ab}:a<b$, and
$h_R:a+b<R$, define
$\operatorname{threeBlockData}(a,b,R,h_a,h_{ab},h_R):
 \operatorname{BlockData}(2)$ by
\begin{align*}
 m(i)&:=\operatorname{Fin.cases}
  \left(R,
   \left(j\mapsto\operatorname{Fin.cases}
     (R-a,(\_\mapsto R-a-b),j)\right),i\right),\\
 \texttt{m\_pos}(i)&:=
   \texttt{by fin\_cases i <;> simp <;> linarith},\\
 a(i)&:=\operatorname{Fin.cases}
  \left(0,
   \left(j\mapsto\operatorname{Fin.cases}(1,(\_\mapsto2),j)\right),i\right).
\end{align*}
The displayed fields are the complete constructor data; the two positivity
proofs use $0<r<R$ and $0<a<b$ together with $a+b<R$, respectively.  For
$\mu:\operatorname{SignedMeasure}(\overline\R_+)$ and $r:\R$, define
\begin{align*}
 \operatorname{lowerMoment}(\mu,r)&:=
  \operatorname{signedIntegral}\left(
   \mu,\bigl(\operatorname{Set.Iio}(\operatorname{finiteParam}(r))\bigr)
          .\operatorname{indicator}(\omega)\right),\\
 \operatorname{upperMoment}(\mu,r)&:=
  \operatorname{signedIntegral}\left(
   \mu,\bigl(\operatorname{Set.Ioi}(\operatorname{finiteParam}(r))\bigr)
          .\operatorname{indicator}(\omega)\right).
\end{align*}
\end{definition}

\begin{definition}[Unique top blocks for the special families]
\label{def:special-block-top-unique}
The unique-largest-amplitude proofs needed at the compactified order
$+\infty$ are named explicitly.  For $r,R:\R$ with $h_r:0<r$ and
$h_R:r<R$, and separately for $a,b,R:\R$ with $h_a:0<a$,
$h_{ab}:a<b$, and $h_R:a+b<R$, define
\begin{align*}
 \operatorname{twoBlockTopUnique}(r,R,h_r,h_R)&:
  \forall j:\operatorname{Fin}(2),\quad j\neq1\to
  (\operatorname{twoBlockData}(r,R,h_r,h_R)).a(j)
   <(\operatorname{twoBlockData}(r,R,h_r,h_R)).a(1)\\[-2pt]
 &:=\texttt{by intro j hj; fin\_cases j <;>
       simp\_all [twoBlockData] <;> norm\_num},\\
 \operatorname{threeBlockTopUnique}(a,b,R,h_a,h_{ab},h_R)&:
  \forall j:\operatorname{Fin}(3),\quad j\neq2\to
  (\operatorname{threeBlockData}(a,b,R,h_a,h_{ab},h_R)).a(j)
   <(\operatorname{threeBlockData}(a,b,R,h_a,h_{ab},h_R)).a(2)\\[-2pt]
 &:=\texttt{by intro j hj; fin\_cases j <;>
       simp\_all [threeBlockData] <;> norm\_num}.
\end{align*}
Thus every later top-order smoothness invocation has a named proof rather
than an inferred maximizer.
\end{definition}

\begin{definition}[Total two-block and three-block dominance maps]
\label{def:block-dominance-maps}
For $a,b:\R$, first define the total order-one norm-dominating block
\[
 \operatorname{threeNormBlock}(a,b):\operatorname{Fin}(3):=
 \texttt{if }1<a\texttt{ then }0
 \texttt{ else if }1<b\texttt{ then }1\texttt{ else }2.
\]
For $r,a,b:\R$, define the total finite-valued dominance functions
\begin{align*}
 \operatorname{twoDominanceMap}(r)
 &: \overline\R_+\to\operatorname{Fin}(2):=
 \bigl(\beta\mapsto\texttt{if }\beta<\operatorname{finiteParam}(r)
      \texttt{ then }0\texttt{ else }1\bigr),\\
 \operatorname{threeDominanceMap}(a,b)
 &: \overline\R_+\to\operatorname{Fin}(3):=
 \bigl(\beta\mapsto\texttt{if }\beta<\operatorname{finiteParam}(a)
      \texttt{ then }0
      \texttt{ else if }\beta<\operatorname{finiteParam}(b)
      \texttt{ then }1\texttt{ else }2\bigr).
\end{align*}
Both definitions include the threshold points and $\top$: equality with a
threshold selects the block on its right, while $\top$ selects the final
block.  No later proof uses an informal piecewise map.
\end{definition}

\begin{lemma}[Exact two-block dominance package]
\label{lem:two-block-dominance-package}
For $r,R:\R$ with $h_r:0<r$ and $h_R:r<R$, the exact return proposition
binds its abbreviations by literal lets:
\begin{align*}
&\texttt{let }B:=\operatorname{twoBlockData}(r,R,h_r,h_R)\texttt{; }\\[-2pt]
&\texttt{let }k:=\operatorname{twoDominanceMap}(r)\texttt{; }\\[-2pt]
&\operatorname{Measurable}(k)\ \wedge\ k(\top)=1\\[-2pt]
&\quad\wedge\Bigl(\forall\alpha:\R,\quad
 0\leq\alpha\to
 \alpha\notin(\{r\}:\operatorname{Finset}(\R))\to
 \forall j:\operatorname{Fin}(2),\quad
 j\neq k(\operatorname{finiteParam}(\alpha))\to\\[-2pt]
&\hspace{25mm}
 \operatorname{blockExponent}(B,j,\alpha)
 <\operatorname{blockExponent}
   (B,k(\operatorname{finiteParam}(\alpha)),\alpha)\Bigr).
\end{align*}
The named theorem term is
$\operatorname{twoBlockDominancePackage}(r,R,h_r,h_R)$.
\end{lemma}

\begin{proof}
The set $\operatorname{Set.Iio}(\operatorname{finiteParam}(r))$ is Borel;
piecewise measurability of two constant maps proves the first conjunct, and
the definition at $\top$ proves the second.  For the last conjunct, use
$0\leq\alpha$ and $0<r$ to rewrite comparison of the two
\texttt{ENNReal.ofReal} values as $\alpha<r$.  Split on that comparison,
unfold \texttt{twoBlockData}, \texttt{blockExponent}, and
\texttt{twoDominanceMap}, then \texttt{fin\_cases j}; the surviving goal is
respectively $\alpha-r<0$ or $0<\alpha-r$.  The excluded singleton gives
strictness.
\end{proof}

\begin{lemma}[Exact three-block dominance package]
\label{lem:three-block-dominance-package}
For $a,b,R:\R$ with $h_a:0<a$, $h_{ab}:a<b$, and $h_R:a+b<R$, the exact
return proposition binds its abbreviations by literal lets:
\begin{align*}
&\texttt{let }B:=\operatorname{threeBlockData}(a,b,R,h_a,h_{ab},h_R)
 \texttt{; }\\[-2pt]
&\texttt{let }k:=\operatorname{threeDominanceMap}(a,b)\texttt{; }\\[-2pt]
&\operatorname{Measurable}(k)\ \wedge\ k(\top)=2\\[-2pt]
&\quad\wedge\Bigl(\forall\alpha:\R,\quad
 0\leq\alpha\to
 \alpha\notin(\{a,b\}:\operatorname{Finset}(\R))\to
 \forall j:\operatorname{Fin}(3),\quad
 j\neq k(\operatorname{finiteParam}(\alpha))\to\\[-2pt]
&\hspace{25mm}
 \operatorname{blockExponent}(B,j,\alpha)
 <\operatorname{blockExponent}
   (B,k(\operatorname{finiteParam}(\alpha)),\alpha)\Bigr).
\end{align*}
The named theorem term is
$\operatorname{threeBlockDominancePackage}(a,b,R,h_a,h_{ab},h_R)$.
\end{lemma}

\begin{proof}
Apply measurable piecewise composition twice to the Borel lower intervals
cut out by $\operatorname{finiteParam}(a)$ and
$\operatorname{finiteParam}(b)$; evaluation at $\top$ is definitional.
For finite nonnegative $\alpha$, the three exponents, after subtracting the
common $R$, are $0$, $\alpha-a$, and $2\alpha-a-b$.  First bind
\texttt{have h\_b : 0 < b := lt\_trans h\_a h\_ab}; this is the proof needed
to reflect both \texttt{ENNReal.ofReal} comparisons back to real
comparisons.  Split successively on
$\alpha<a$ and $\alpha<b$, rewrite the two \texttt{ofReal} comparisons,
and use the excluded two-point finset together with $a<b$.  After
\texttt{fin\_cases j}, every remaining inequality is linear arithmetic.
\end{proof}

\begin{lemma}[Unique order-one block for the two-block family]
\label{lem:two-order-one-dominance}
For $r,R:\R$ with $h_r:0<r$ and $h_R:r<R$, the exact return proposition is
\begin{align*}
&\texttt{let }B_2:=\operatorname{twoBlockData}(r,R,h_r,h_R)\texttt{; }\\[-2pt]
&\Bigl(1<r\to
 \forall j:\operatorname{Fin}(2),\quad j\neq0\to
 \operatorname{blockExponent}(B_2,j,1)
 <\operatorname{blockExponent}(B_2,0,1)\Bigr)\\[-2pt]
&\quad\wedge\Bigl(r<1\to
 \forall j:\operatorname{Fin}(2),\quad j\neq1\to
 \operatorname{blockExponent}(B_2,j,1)
 <\operatorname{blockExponent}(B_2,1,1)\Bigr).
\end{align*}
The named theorem term is
$\operatorname{twoOrderOneDominancePackage}(r,R,h_r,h_R)$; its first and
second projections are applied to proofs of $1<r$ and $r<1$, respectively.
\end{lemma}

\begin{proof}
Introduce the relevant strict comparison, then \texttt{fin\_cases j}; after
unfolding the definitions the nontrivial goal is $1-r<0$ or $0<1-r$.
\end{proof}

\begin{lemma}[Unique order-one block for the three-block family]
\label{lem:three-order-one-dominance}
For $a,b,R:\R$ with $h_a:0<a$, $h_{ab}:a<b$, $h_R:a+b<R$,
$h_{a1}:a\neq1$, and $h_{b1}:b\neq1$, the named proof term has the literal
dependent-let type
\begin{align*}
&\operatorname{threeOrderOneDominant}
 (a,b,R,h_a,h_{ab},h_R,h_{a1},h_{b1}):\\[-2pt]
&\quad\left(\texttt{let }B_3:=\operatorname{threeBlockData}
 (a,b,R,h_a,h_{ab},h_R)\texttt{; }\quad
 \forall j:\operatorname{Fin}(3),\quad
 j\neq\operatorname{threeNormBlock}(a,b)\to
 \operatorname{blockExponent}(B_3,j,1)
 <\operatorname{blockExponent}
   (B_3,\operatorname{threeNormBlock}(a,b),1)\right)
\end{align*}
is total.
\end{lemma}

\begin{proof}
Split the decidable comparisons $1<a$ and $1<b$.  The disequalities
$a\neq1$ and $b\neq1$
turn the negative branches into $a<1$ and $b<1$; now
\texttt{fin\_cases j}, unfold \texttt{threeNormBlock}, and close every
remaining affine inequality by linear arithmetic.
\end{proof}

\begin{definition}[Block logarithmic and norm-free kernels]
\label{def:block-log-g-kernels}
For $J:\N$, $B:\operatorname{BlockData}(J)$,
$\mu:\operatorname{SignedMeasure}(\overline\R_+)$, $n,q:\N$, and
$u:\operatorname{Fin}(J+1)\to\R$, define
\begin{align*}
 \operatorname{blockLogKernel}(\mu,B,n,u,q)&:=
  \left(\texttt{letI := blockCarrierNonempty B n; }\quad
  \operatorname{iterDeriv}
   (\operatorname{signedLogPhiLine}
     (\mu,\operatorname{blockLineData}(B,n,u))),q,0)\right),\\
 \operatorname{blockGKernel}(\mu,B,n,u,q)&:=
  \left(\texttt{letI := blockCarrierNonempty B n; }\quad
  \operatorname{iterDeriv}
   (\operatorname{integratedEntropyLine}
     (\mu,\operatorname{blockLineData}(B,n,u))),q,0)\right).
\end{align*}
Define the corresponding total mass and its logarithmic derivative kernel:
\begin{align*}
 \operatorname{blockMass}(B,n,u,\lambda)&:=
  \left(\texttt{letI := blockCarrierNonempty B n; }\quad
  \sum_{i:\operatorname{BlockCarrier}(B,n)}
   \operatorname{blockLineRaw}(B,n,u,\lambda)(i)\right),\\
 \operatorname{blockLogMassKernel}(B,n,u,q)&:=
  \operatorname{iterDeriv}
   (\lambda\mapsto\log(\operatorname{blockMass}(B,n,u,\lambda)),q,0).
\end{align*}
Define the typed scalar curves used by the curvature arguments:
\begin{align*}
 \operatorname{blockPhiCurve}(\mu,B,n,u,\lambda)&:=
  \left(\texttt{letI := blockCarrierNonempty B n; }\quad
  \operatorname{PhiSigned}_{\operatorname{BlockCarrier}(B,n)}
   \left(\mu,\operatorname{lineConeTotal}
    (\operatorname{blockLineData}(B,n,u),\lambda)\right)\right),\\
 \operatorname{blockGCurve}(\mu,B,n,u,\lambda)&:=
  \left(\texttt{letI := blockCarrierNonempty B n; }\quad
  \operatorname{GSigned}_{\operatorname{BlockCarrier}(B,n)}
   \left(\mu,\operatorname{linePosConeTotal}
    (\operatorname{blockLineData}(B,n,u),\lambda)\right)\right).
\end{align*}
For a type $U$, $K:\operatorname{Set}(U)$,
$f:\N\to U\to\R$, $g:U\to\R$, and $n:\N$, define
\[
 \operatorname{compactUniformError}(K,f,g,n):=
 \operatorname{sSup}\{z:\R\mid\exists u\in K,\ z=\abs{f(n,u)-g(u)}\}.
\]
For the same $U,K,f,g$, define the exact convergence predicate
\[
 \operatorname{CompactUniformConverges}(K,f,g):\Longleftrightarrow
 \operatorname{Tendsto}
  (n\mapsto\operatorname{compactUniformError}(K,f,g,n),
   \texttt{atTop},\mathcal N(0)).
\]
\end{definition}

\begin{lemma}[Evaluation of compact-uniform convergence]
\label{lem:compact-uniform-eval}
Let $U$ be a finite-dimensional real normed space, let
$K:\operatorname{Set}(U)$ be nonempty and compact, let
$f:\N\to U\to\R$ and $g:U\to\R$ satisfy
$\forall n,\operatorname{ContinuousOn}(f(n),K)$ and
$\operatorname{ContinuousOn}(g,K)$, and let $u:U$ satisfy $h_u:u\in K$.
Then
\begin{align*}
&\operatorname{Tendsto}
 (n\mapsto\operatorname{compactUniformError}(K,f,g,n),
  \texttt{atTop},\mathcal N(0))\\[-2pt]
&\hspace{15mm}\Longrightarrow
 \operatorname{Tendsto}(n\mapsto f(n,u),
  \texttt{atTop},\mathcal N(g(u))).
\end{align*}
Equivalently, \texttt{CompactUniformConverges K f g} may be evaluated at
any named member of $K$.
\end{lemma}

\begin{proof}
Continuity on the compact set makes the defining error range bounded above.
It contains $\abs{f(n,u)-g(u)}$, so this value is bounded above by
\texttt{compactUniformError}; the latter is nonnegative because $K$ is
nonempty.  Squeeze the absolute difference to zero and use the metric
characterization of real convergence.
\end{proof}

\begin{lemma}[Addition closure for compact-uniform convergence]
\label{lem:compact-uniform-add}
Let $U:\texttt{Type u}$ carry
$[\operatorname{NormedAddCommGroup}(U)]$,
$[\operatorname{NormedSpace}(\R,U)]$, and
$[\operatorname{FiniteDimensional}(\R,U)]$.  Let
$K:\operatorname{Set}(U)$ have proofs $h_{K0}:K.\operatorname{Nonempty}$
and $h_K:\operatorname{IsCompact}(K)$.  Let
$f_1,f_2:\N\to U\to\R$ and $g_1,g_2:U\to\R$ have proofs
\begin{align*}
h_{f_1}&:\forall n:\N,\ \operatorname{ContinuousOn}(f_1(n),K),&
h_{f_2}&:\forall n:\N,\ \operatorname{ContinuousOn}(f_2(n),K),\\
h_{g_1}&:\operatorname{ContinuousOn}(g_1,K),&
h_{g_2}&:\operatorname{ContinuousOn}(g_2,K).
\end{align*}
If
\[
 H_1:\operatorname{CompactUniformConverges}(K,f_1,g_1)
 \qquad\text{and}\qquad
 H_2:\operatorname{CompactUniformConverges}(K,f_2,g_2),
\]
then the exact conclusion is
\[
 \operatorname{CompactUniformConverges}\!\left(
 K,\bigl(n\mapsto u\mapsto f_1(n,u)+f_2(n,u)\bigr),
   u\mapsto g_1(u)+g_2(u)\right).
\]
The named theorem term is
\texttt{compactUniformConverges\_add
 K f1 f2 g1 g2 h\_K0 h\_K h\_f1 h\_f2 h\_g1 h\_g2 H1 H2}.
\end{lemma}

\begin{proof}
For every $n$ and $u\in K$, the triangle inequality gives
\[
 \abs{(f_1(n,u)+f_2(n,u))-(g_1(u)+g_2(u))}
 \leq \abs{f_1(n,u)-g_1(u)}+\abs{f_2(n,u)-g_2(u)}.
\]
All three error ranges are nonempty by $h_{K0}$ and bounded above because
the displayed functions are continuous on the compact set $K$.  Apply
\texttt{le\_csSup} to each summand and \texttt{csSup\_le} to the left-hand
range to obtain, for every $n$,
\[
 0\leq
 \operatorname{compactUniformError}
  (K,n\mapsto u\mapsto f_1(n,u)+f_2(n,u),u\mapsto g_1(u)+g_2(u),n)
 \leq
 \operatorname{compactUniformError}(K,f_1,g_1,n)
 +\operatorname{compactUniformError}(K,f_2,g_2,n).
\]
The right-hand side tends to zero by $H_1$, $H_2$, and
\texttt{Tendsto.add}; squeeze.
\end{proof}

\begin{lemma}[Evaluation on a singleton]\label{lem:compact-uniform-singleton}
For every type $U$, $f:\N\to U\to\R$, $g:U\to\R$, and $u:U$,
the following literal conjunction holds:
\begin{align*}
&\Bigl(\forall n:\N,\quad
 \operatorname{compactUniformError}(\{u\},f,g,n)
 =\abs{f(n,u)-g(u)}\Bigr)\\
&\quad\wedge\Bigl(\operatorname{Tendsto}
 (n\mapsto\operatorname{compactUniformError}(\{u\},f,g,n),
  \texttt{atTop},\mathcal N(0))
 \Longrightarrow
 \operatorname{Tendsto}(n\mapsto f(n,u),
  \texttt{atTop},\mathcal N(g(u)))\Bigr).
\end{align*}
\end{lemma}

\begin{proof}
Use \texttt{constructor}.  For the first clause, introduce $n$ and apply
\texttt{simp [compactUniformError]}; the set under \texttt{sSup} is the
singleton containing the displayed absolute error.  The second clause is
the metric characterization of convergence together with the first clause.
\end{proof}

\begin{lemma}[Subtype bridge for uniform signed-integral errors]
\label{lem:subtype-uniform-error-bridge}
Let $E$ carry a measurable space, let $\mu:\operatorname{SignedMeasure}(E)$,
let $U$ be a type, let $K:\operatorname{Set}(U)$, let
$f_n:\N\to E\to U\to\R$, $f:E\to U\to\R$, and $n:\N$.  The exact return
proposition binds the subtype and both auxiliary functions by literal local
lets:
\begin{align*}
&\texttt{let }U_K:=\{u:U\mathbin{//}u\in K\}\texttt{; }\\[-2pt]
&\texttt{let }F_n:=
 (m\mapsto u\mapsto\operatorname{signedIntegral}
   (\mu,e\mapsto f_n(m,e,u)))\texttt{; }\\[-2pt]
&\texttt{let }F:=
 (u\mapsto\operatorname{signedIntegral}(\mu,e\mapsto f(e,u)))\texttt{; }\\[-2pt]
&\operatorname{uniformIntegralErrorSigned}\!\left(
 \mu,(m\mapsto e\mapsto u:U_K\mapsto f_n(m,e,u.1)),
      (e\mapsto u:U_K\mapsto f(e,u.1)),n\right)\\
&\hspace{25mm}
 =\operatorname{compactUniformError}(K,F_n,F,n)
\end{align*}
holds.
\end{lemma}

\begin{proof}
Unfold both definitions.  Extensionality identifies the two sets under
\texttt{sSup}: a witness $u:U_K$ is exactly an ambient witness $u:U$
together with $u\in K$.  The signed-integral expressions are then
definitionally equal.
\end{proof}

\begin{lemma}[Block curve--kernel bridge]
\label{lem:block-curve-kernel-bridge}
For $J:\N$, $B:\operatorname{BlockData}(J)$,
$\mu:\operatorname{SignedMeasure}(\overline\R_+)$, $n:\N$, and
$u:\operatorname{Fin}(J+1)\to\R$, let
$j_\infty:\operatorname{Fin}(J+1)$ have the explicit proof
\[
 h_\infty:\forall j:\operatorname{Fin}(J+1),\quad j\neq j_\infty\to
 B.a(j)<B.a(j_\infty).
\]
Then there exists $\eps:\R$ such that
\begin{align*}
&0<\eps\ \wedge\
 \bigl(\forall\lambda\in\operatorname{Set.Ioo}(-\eps,\eps),\quad
   \operatorname{LinePositive}
    (\operatorname{blockLineData}(B,n,u),\lambda)\bigr)\\
&\quad\wedge
 0<\operatorname{blockPhiCurve}(\mu,B,n,u,0)\\
&\quad\wedge
 \operatorname{ContDiffOn}\!\left(\R,2,
  \lambda\mapsto\operatorname{blockPhiCurve}(\mu,B,n,u,\lambda),
  \operatorname{Set.Ioo}(-\eps,\eps)\right)\\
&\quad\wedge
 \operatorname{ContDiffOn}\!\left(\R,2,
  \lambda\mapsto\operatorname{blockGCurve}(\mu,B,n,u,\lambda),
  \operatorname{Set.Ioo}(-\eps,\eps)\right)\\
&\quad\wedge\Bigl(\forall q:\N,\ q\leq2\to
 \operatorname{iterDeriv}
  (\lambda\mapsto\log(\operatorname{blockPhiCurve}
    (\mu,B,n,u,\lambda)),q,0)
 =\operatorname{blockLogKernel}(\mu,B,n,u,q)\Bigr)\\
&\quad\wedge\Bigl(\forall q:\N,\ q\leq2\to
 \operatorname{iterDeriv}
  (\operatorname{blockGCurve}(\mu,B,n,u),q,0)
 =\operatorname{blockGKernel}(\mu,B,n,u,q)\Bigr).
\end{align*}
\end{lemma}

\begin{proof}
Choose $\eps$ from the finite positive base coordinates and the finitely many
velocities.  On this interval the two total line wrappers select their
positive branches.  At $\lambda=0$, apply
\cref{lem:signed-line-column-bridge} with
\texttt{linePositiveZero (blockLineData B n u)} and rewrite its
\texttt{lineConeTotal} equality; its strict-positivity conjunct is exactly
$0<\operatorname{blockPhiCurve}(\mu,B,n,u,0)$.  At order $+\infty$,
$h_\infty$ makes $j_\infty$ the
unique maximal base coordinate throughout a possibly smaller common
neighbourhood, so the fixed-max-coordinate endpoint clause applies.  The exact entropy-derivative package and
\cref{lem:signed-integral-differentiate-twice} give the two displayed
\texttt{ContDiffOn} clauses.  Now use
\cref{lem:signed-line-column-bridge}; differentiating its logarithmic and
norm-free identities zero, one, and two times gives the two kernel equalities.
\end{proof}

\begin{lemma}[Block (G)-kernels are signed kernel integrals]
\label{lem:block-g-kernel-integrals}
For $J:\N$, $B:\operatorname{BlockData}(J)$,
$\mu:\operatorname{SignedMeasure}(\overline\R_+)$, $n:\N$,
$u:\operatorname{Fin}(J+1)\to\R$, and
$j_\infty:\operatorname{Fin}(J+1)$, assume
\[
 h_\infty:\forall j:\operatorname{Fin}(J+1),\quad j\neq j_\infty\to
 B.a(j)<B.a(j_\infty).
\]
Then the two derivative--integral identities form one literal conjunction:
\begin{align*}
&\operatorname{blockGKernel}(\mu,B,n,u,1)
 =\operatorname{signedIntegral}\!\left(
   \mu,\beta\mapsto\operatorname{blockKernelFirst}(B,n,u,\beta)\right)\\
&\quad\wedge\quad
 \operatorname{blockGKernel}(\mu,B,n,u,2)
 =\operatorname{signedIntegral}\!\left(
   \mu,\beta\mapsto\operatorname{blockKernelSecond}(B,n,u,\beta)\right).
\end{align*}
The named theorem term is
$\operatorname{blockGKernelIntegralBridge}
 (\mu,B,n,u,j_\infty,h_\infty)$.
\end{lemma}

\begin{proof}
Install \texttt{blockCarrierNonempty B n}.  Shrink the positive interval
from \cref{lem:block-curve-kernel-bridge} so that the closed interval is
still positive.  Use the coordinate in block $j_\infty$ with inner index
zero.  The strict amplitude gaps in $h_\infty$, together with the common
velocity on that block, give the required \texttt{FixedMaxCoordinate} on
this closed interval.  Apply \cref{lem:differentiate-integral} at
$\lambda=0$ once and twice, unfold \texttt{iterDeriv},
\texttt{blockGKernel}, \texttt{blockKernelFirst}, and
\texttt{blockKernelSecond}, and take the two resulting equalities as the
displayed conjunction.
\end{proof}

\begin{lemma}[Block logarithmic-mass splitting and limit]
\label{lem:block-log-mass-limit}
Let $J:\N$, $B:\operatorname{BlockData}(J)$,
$\mu:\operatorname{SignedMeasure}(\overline\R_+)$, and
$k:\operatorname{Fin}(J+1)$.  Assume
\[
 h_k:\forall j:\operatorname{Fin}(J+1),\quad j\neq k\to
 \operatorname{blockExponent}(B,j,1)
 <\operatorname{blockExponent}(B,k,1).
\]
Also let $j_\infty:\operatorname{Fin}(J+1)$ have proof
\[
 h_\infty:\forall j:\operatorname{Fin}(J+1),\quad j\neq j_\infty\to
 B.a(j)<B.a(j_\infty).
\]
Then the following two clauses form one literal conjunction:
\begin{align*}
&\Bigl(\forall(n:\N)(u:\operatorname{Fin}(J+1)\to\R)(q:\N),\quad
 q\leq2\to
 \operatorname{blockLogKernel}(\mu,B,n,u,q)
 =\operatorname{blockLogMassKernel}(B,n,u,q)
  +\operatorname{blockGKernel}(\mu,B,n,u,q)\Bigr)\\
&\quad\wedge\Bigl(\forall K:\operatorname{Set}
   (\operatorname{Fin}(J+1)\to\R),\quad
 K.\operatorname{Nonempty}\to\operatorname{IsCompact}(K)\to\\[-2pt]
&\hspace{8mm}
 \operatorname{CompactUniformConverges}
  (K,(n\mapsto u\mapsto\operatorname{blockLogMassKernel}(B,n,u,1)),
     u\mapsto u(k))\ \wedge\\[-2pt]
&\hspace{8mm}
 \operatorname{CompactUniformConverges}
  (K,(n\mapsto u\mapsto\operatorname{blockLogMassKernel}(B,n,u,2)),
     u\mapsto-u(k)^2)\Bigr).
\end{align*}
\end{lemma}

\begin{proof}
The defining identity
\texttt{signedLogPhiLine = log blockMass + integratedEntropyLine} holds on
a common open neighbourhood of zero; its order-$+\infty$ branch uses the
unique maximal coordinate supplied by $h_\infty$.  Apply
\cref{lem:differentiate-integral} and differentiate the finite sum to obtain
the first clause for orders zero, one, and two.  For the second clause,
divide the block mass by the contribution of block $k$.  The strict exponent
gap in $h_k$, the ceiling estimates from \cref{lem:rounding}, and compactness
of $K$ make the sum of all other ratios tend uniformly to zero.  The first
two logarithmic derivatives consequently tend uniformly to $u(k)$ and
$-u(k)^2$.
\end{proof}

\begin{lemma}[Continuity of logarithmic block-mass kernels]
\label{lem:block-log-mass-continuity}
For $J:\N$, $B:\operatorname{BlockData}(J)$, $n,q:\N$, and
$h_q:q\leq2$, the named proof term has the literal type
\[
 \operatorname{blockLogMassKernelContinuous}(B,n,q,h_q):
 \operatorname{Continuous}\!\left(
  u:\bigl(\operatorname{Fin}(J+1)\to\R\bigr)\mapsto
  \operatorname{blockLogMassKernel}(B,n,u,q)\right).
\]
\end{lemma}

\begin{proof}
At $\lambda=0$ the block mass is a finite sum of strictly positive
coordinates, hence is positive.  Differentiate that finite sum explicitly
for $q=0,1,2$.  The resulting expressions are quotients of finite sums of
polynomials in the coordinates of $u$, with the positive mass (or its
square) as denominator.  They are continuous.  The proof splits
$q=0,1,2$ using $h_q$ and closes each branch with the standard continuity
rules for finite sums, products, and division by a nowhere-zero function.
\end{proof}

\begin{lemma}[Linearity and continuity of the block kernels]
\label{lem:block-kernel-regularity}
For $J:\N$, $B:\operatorname{BlockData}(J)$,
$\mu:\operatorname{SignedMeasure}(\overline\R_+)$,
$j_\infty:\operatorname{Fin}(J+1)$ with proof
\[
 h_\infty:\forall j:\operatorname{Fin}(J+1),\quad j\neq j_\infty\to
 B.a(j)<B.a(j_\infty),
\]
and $n:\N$, the exact named theorem term
$\operatorname{blockKernelRegularity}(\mu,B,j_\infty,h_\infty,n)$ has type
\begin{align*}
&\Bigl(\forall u,v:\operatorname{Fin}(J+1)\to\R,\quad
 \operatorname{blockGKernel}(\mu,B,n,u+v,1)
 =\operatorname{blockGKernel}(\mu,B,n,u,1)
  +\operatorname{blockGKernel}(\mu,B,n,v,1)\Bigr)\\
&\quad\wedge\Bigl(\forall(c:\R)(u:\operatorname{Fin}(J+1)\to\R),\quad
 \operatorname{blockGKernel}(\mu,B,n,c\mathbin{\bullet}u,1)
 =c\operatorname{blockGKernel}(\mu,B,n,u,1)\Bigr)\\
&\quad\wedge\operatorname{Continuous}
  (u\mapsto\operatorname{blockGKernel}(\mu,B,n,u,1))\\
&\quad\wedge\operatorname{Continuous}
  (u\mapsto\operatorname{blockGKernel}(\mu,B,n,u,2)).
\end{align*}
\end{lemma}

\begin{proof}
Use the first two formulas of \cref{lem:exact-derivatives}, signed-integral
linearity, and \cref{lem:differentiate-integral}; the compactified top-order
branch uses the fixed maximizer $j_\infty$ and its strict-gap proof
$h_\infty$.  The first derivative is a
finite-dimensional linear form in $u$; the second is a quadratic polynomial,
so both are continuous.
\end{proof}

\begin{definition}[Two-block limit polynomials]
\label{def:two-block-limits}
For $\mu:\operatorname{SignedMeasure}(\overline\R_+)$, $r:\R$, and
$u:\operatorname{Fin}(2)\to\R$, define the following total terms.  Each
abbreviation is introduced by a literal \texttt{let} inside its own body:
\begin{align}
 \operatorname{twoUpperLogFirst}(\mu,r,u)&:=
  \left(\texttt{let A := }\operatorname{upperMoment}(\mu,r)\texttt{; }
   (1-A)u(0)+Au(1)\right),
 \label{eq:two-r-greater-first}\\
 \operatorname{twoUpperLogSecond}(\mu,r,u)&:=
  \left(\texttt{let A := }\operatorname{upperMoment}(\mu,r)\texttt{; }
   -(1-A)u(0)^2-Au(1)^2\right),
 \label{eq:two-r-greater-second}\\
 \operatorname{twoLowerLogFirst}(\mu,r,u)&:=
  \left(\texttt{let B := }\operatorname{lowerMoment}(\mu,r)\texttt{; }
   Bu(0)+(1-B)u(1)\right),
 \label{eq:two-r-less-first}\\
 \operatorname{twoLowerLogSecond}(\mu,r,u)&:=
  \left(\texttt{let B := }\operatorname{lowerMoment}(\mu,r)\texttt{; }
   -Bu(0)^2-(1-B)u(1)^2\right),
 \label{eq:two-r-less-second}\\
 \operatorname{twoUpperGFirst}(\mu,r,u)&:=
  \left(\texttt{let A := }\operatorname{upperMoment}(\mu,r)\texttt{; }
   A(u(1)-u(0))\right),
 \label{eq:two-G-r-greater-first}\\
 \operatorname{twoUpperGSecond}(\mu,r,u)&:=
  \left(\texttt{let A := }\operatorname{upperMoment}(\mu,r)\texttt{; }
   A(u(0)^2-u(1)^2)\right),
 \label{eq:two-G-r-greater-second}\\
 \operatorname{twoLowerGFirst}(\mu,r,u)&:=
  \left(\texttt{let B := }\operatorname{lowerMoment}(\mu,r)\texttt{; }
   B(u(0)-u(1))\right),
 \label{eq:two-G-r-less-first}\\
 \operatorname{twoLowerGSecond}(\mu,r,u)&:=
  \left(\texttt{let B := }\operatorname{lowerMoment}(\mu,r)\texttt{; }
   B(u(1)^2-u(0)^2)\right).
 \label{eq:two-G-r-less-second}
\end{align}
\end{definition}

\begin{definition}[Three-block cells and limit polynomials]
\label{def:three-block-limits}
For $\mu:\operatorname{SignedMeasure}(\overline\R_+)$ and $a,b:\R$, define
one total function on $\operatorname{Fin}(3)$:
\begin{align*}
 \operatorname{threeCellMoment}(\mu,a,b):\operatorname{Fin}(3)\to\R:=
 \operatorname{Fin.cases}&\left(
  \operatorname{signedIntegral}\!\left(\mu,
   \bigl(\operatorname{Set.Iio}(\operatorname{finiteParam}(a))\bigr)
    .\operatorname{indicator}(\omega)\right),\right.\\[-2pt]
 &\left.j\mapsto\operatorname{Fin.cases}\!\left(
  \operatorname{signedIntegral}\!\left(\mu,
   \bigl(\operatorname{Set.Ioo}(\operatorname{finiteParam}(a),
          \operatorname{finiteParam}(b))\bigr)
    .\operatorname{indicator}(\omega)\right),\right.\right.\\[-2pt]
 &\left.\left.\_\mapsto
  \operatorname{signedIntegral}\!\left(\mu,
   \bigl(\operatorname{Set.Ioi}(\operatorname{finiteParam}(b))\bigr)
    .\operatorname{indicator}(\omega)\right),j\right)\right).
\end{align*}
For $a,b:\R$ and $u:\operatorname{Fin}(3)\to\R$, define each polynomial
with its dominating index bound by a literal local \texttt{let}:
\begin{align}
 \operatorname{threeLogFirst}(\mu,a,b,u)&:=
  \left(\texttt{let k := threeNormBlock a b; }u(k)+
  \sum_{j\in\operatorname{Finset.univ.erase}(k)}
   \operatorname{threeCellMoment}(\mu,a,b,j)(u(j)-u(k))\right),
 \label{eq:three-first}\\
 \operatorname{threeLogSecond}(\mu,a,b,u)&:=
  \left(\texttt{let k := threeNormBlock a b; }-u(k)^2+
  \sum_{j\in\operatorname{Finset.univ.erase}(k)}
   \operatorname{threeCellMoment}(\mu,a,b,j)(u(k)^2-u(j)^2)\right),
 \label{eq:three-second}\\
 \operatorname{threeGFirst}(\mu,a,b,u)&:=
  \left(\texttt{let k := threeNormBlock a b; }
  \sum_{j\in\operatorname{Finset.univ.erase}(k)}
   \operatorname{threeCellMoment}(\mu,a,b,j)(u(j)-u(k))\right),
 \label{eq:three-G-first}\\
 \operatorname{threeGSecond}(\mu,a,b,u)&:=
  \left(\texttt{let k := threeNormBlock a b; }
  \sum_{j\in\operatorname{Finset.univ.erase}(k)}
   \operatorname{threeCellMoment}(\mu,a,b,j)(u(k)^2-u(j)^2)\right).
 \label{eq:three-G-second}
\end{align}
\end{definition}

\begin{lemma}[Evaluation of the two-block limit integrals]
\label{lem:two-block-limit-integral-eval}
Let $\mu:\operatorname{SignedMeasure}(\overline\R_+)$, let $r:\R$ have
proof $h_r:0<r$, let $u:\operatorname{Fin}(2)\to\R$, and assume
\[
 h_{\mu,r}:\operatorname{signedTV}(\mu)
  (\{\operatorname{finiteParam}(r)\})=0.
\]
Then the exact return proposition is the single conjunction
\begin{align*}
&\Bigl(1<r\to
 \operatorname{signedIntegral}\!\left(
  \mu,\beta\mapsto\operatorname{blockLimitFirst}
   (\operatorname{twoDominanceMap}(r),u,\beta)\right)
 =\operatorname{twoUpperGFirst}(\mu,r,u)\\[-2pt]
&\hspace{19mm}\wedge\quad
 \operatorname{signedIntegral}\!\left(
  \mu,\beta\mapsto\operatorname{blockLimitSecond}
   (\operatorname{twoDominanceMap}(r),u,\beta)\right)
 =\operatorname{twoUpperGSecond}(\mu,r,u)\Bigr)\\
&\quad\wedge\Bigl(r<1\to
 \operatorname{signedIntegral}\!\left(
  \mu,\beta\mapsto\operatorname{blockLimitFirst}
   (\operatorname{twoDominanceMap}(r),u,\beta)\right)
 =\operatorname{twoLowerGFirst}(\mu,r,u)\\[-2pt]
&\hspace{19mm}\wedge\quad
 \operatorname{signedIntegral}\!\left(
  \mu,\beta\mapsto\operatorname{blockLimitSecond}
   (\operatorname{twoDominanceMap}(r),u,\beta)\right)
 =\operatorname{twoLowerGSecond}(\mu,r,u)\Bigr).
\end{align*}
The named theorem term is
$\operatorname{twoBlockLimitIntegralEval}(\mu,r,h_r,u,h_{\mu,r})$.
\end{lemma}

\begin{proof}
First split the implications $1<r$ and $r<1$.  In the first branch,
$\operatorname{twoDominanceMap}(r,\operatorname{finiteParam}(1))=0$; in
the second it equals $1$.  In either branch, simplify the coefficient of
that order-one cell to zero \emph{before} invoking signed-integral
linearity.  The only remaining nonzero cell is separated from $1$, so its
$\omega$-indicator is integrable against both Jordan parts: use
\cref{lem:fnd-omega-measurable}, boundedness of $\omega$ on that closed
side of the threshold, and finiteness of \texttt{signedTV}.  The map is
constant on the two strict cells cut out by
$\operatorname{finiteParam}(r)$.  Use the null hypothesis to replace the
kernel at the threshold by either adjacent value via
\texttt{signedIntegral\_congr\_ae}.  Now unfold
\texttt{blockLimitFirst}, \texttt{blockLimitSecond},
\texttt{twoUpperGFirst}, \texttt{twoUpperGSecond},
\texttt{twoLowerGFirst}, and \texttt{twoLowerGSecond}; signed-integral
linearity leaves exactly \texttt{upperMoment} or \texttt{lowerMoment}
times the displayed velocity difference.
\end{proof}

\begin{lemma}[Evaluation of the three-block limit integrals]
\label{lem:three-block-limit-integral-eval}
Let $\mu:\operatorname{SignedMeasure}(\overline\R_+)$, let $a,b:\R$ have
proofs $h_a:0<a$, $h_{ab}:a<b$, $h_{a1}:a\neq1$, and $h_{b1}:b\neq1$,
and let $u:\operatorname{Fin}(3)\to\R$.  Assume
\begin{align*}
 h_{\mu,a}&:\operatorname{signedTV}(\mu)
   (\{\operatorname{finiteParam}(a)\})=0,\\
 h_{\mu,b}&:\operatorname{signedTV}(\mu)
   (\{\operatorname{finiteParam}(b)\})=0.
\end{align*}
Then the exact return proposition is
\begin{align*}
&\operatorname{signedIntegral}\!\left(
  \mu,\beta\mapsto\operatorname{blockLimitFirst}
   (\operatorname{threeDominanceMap}(a,b),u,\beta)\right)
 =\operatorname{threeGFirst}(\mu,a,b,u)\\[-2pt]
&\quad\wedge\quad
 \operatorname{signedIntegral}\!\left(
  \mu,\beta\mapsto\operatorname{blockLimitSecond}
   (\operatorname{threeDominanceMap}(a,b),u,\beta)\right)
 =\operatorname{threeGSecond}(\mu,a,b,u).
\end{align*}
The named theorem term is
$\operatorname{threeBlockLimitIntegralEval}
(\mu,a,b,h_a,h_{ab},h_{a1},h_{b1},u,h_{\mu,a},h_{\mu,b})$.
\end{lemma}

\begin{proof}
First split the comparisons $1<a$ and $1<b$ to prove
\[
 \operatorname{threeDominanceMap}
  (a,b,\operatorname{finiteParam}(1))
 =\operatorname{threeNormBlock}(a,b).
\]
Off the two null thresholds, the dominance map is respectively $0$, $1$,
and $2$ on the three strict cells in \texttt{threeCellMoment}.  In each
comparison branch, first simplify the coefficient of the unique cell
containing $1$ to zero.  Each remaining nonzero cell is separated from
$1$; its $\omega$-indicator is integrable by
\cref{lem:fnd-omega-measurable}, the corresponding bounded-away-from-$1$
estimate, and finiteness of \texttt{signedTV}.  Apply
\texttt{signedIntegral\_congr\_ae} using $h_{\mu,a},h_{\mu,b}$, unfold the
two block-limit kernels, and distribute the signed integral only over those
integrable nonzero cells (the order-one cell is discharged by
\texttt{simp}).  The three scalar integrals are definitionally
\texttt{threeCellMoment}; collecting the velocity differences gives
\texttt{threeGFirst} and \texttt{threeGSecond}.
\end{proof}

\begin{proposition}[Two-block and three-block localisation]
\label{prop:block-localisation}
Let $\mu:\operatorname{SignedMeasure}(\overline\R_+)$; let
$r,R:\R$ satisfy $h_r:0<r$, $h_{r1}:r\neq1$, and $h_R:r<R$; and let
$K_2:\operatorname{Set}(\operatorname{Fin}(2)\to\R)$ have proofs
$h_{K_2,0}:K_2.\operatorname{Nonempty}$ and
$h_{K_2}:\operatorname{IsCompact}(K_2)$.  Also let $a,b,R_3:\R$ have proofs
$h_a:0<a$, $h_{ab}:a<b$, $h_{a1}:a\neq1$, $h_{b1}:b\neq1$, and
$h_{R_3}:a+b<R_3$, and let
$K_3:\operatorname{Set}(\operatorname{Fin}(3)\to\R)$ have proofs
$h_{K_3,0}:K_3.\operatorname{Nonempty}$ and
$h_{K_3}:\operatorname{IsCompact}(K_3)$.  Assume
\[
 h_{\mu,r}:\operatorname{signedTV}(\mu)
   (\{\operatorname{finiteParam}(r)\})=0,
 \quad
 h_{\mu,a}:\operatorname{signedTV}(\mu)
   (\{\operatorname{finiteParam}(a)\})=0,
 \quad
 h_{\mu,b}:\operatorname{signedTV}(\mu)
   (\{\operatorname{finiteParam}(b)\})=0.
\]
Then the following single literal dependent-let proposition of convergence
statements holds:
\begin{align*}
&\texttt{let }B_2:=\operatorname{twoBlockData}(r,R,h_r,h_R)\texttt{; }\\[-2pt]
&\texttt{let }B_3:=\operatorname{threeBlockData}
 (a,b,R_3,h_a,h_{ab},h_{R_3})\texttt{; }\\[-2pt]
&\Bigl(1<r\to
 \operatorname{Tendsto}
  (n\mapsto\operatorname{compactUniformError}
   (K_2,(m\mapsto u\mapsto\operatorname{blockLogKernel}(\mu,B_2,m,u,1)),
    u\mapsto\operatorname{twoUpperLogFirst}(\mu,r,u),n),
   \texttt{atTop},\mathcal N(0))\\[-2pt]
&\hspace{11mm}\wedge
 \operatorname{Tendsto}
  (n\mapsto\operatorname{compactUniformError}
   (K_2,(m\mapsto u\mapsto\operatorname{blockLogKernel}(\mu,B_2,m,u,2)),
    u\mapsto\operatorname{twoUpperLogSecond}(\mu,r,u),n),
   \texttt{atTop},\mathcal N(0))\\[-2pt]
&\hspace{11mm}\wedge
 \operatorname{Tendsto}
  (n\mapsto\operatorname{compactUniformError}
   (K_2,(m\mapsto u\mapsto\operatorname{blockGKernel}(\mu,B_2,m,u,1)),
    u\mapsto\operatorname{twoUpperGFirst}(\mu,r,u),n),
   \texttt{atTop},\mathcal N(0))\\[-2pt]
&\hspace{11mm}\wedge
 \operatorname{Tendsto}
  (n\mapsto\operatorname{compactUniformError}
   (K_2,(m\mapsto u\mapsto\operatorname{blockGKernel}(\mu,B_2,m,u,2)),
    u\mapsto\operatorname{twoUpperGSecond}(\mu,r,u),n),
   \texttt{atTop},\mathcal N(0))\Bigr)\\
&\quad\wedge\Bigl(r<1\to
 \operatorname{Tendsto}
  (n\mapsto\operatorname{compactUniformError}
   (K_2,(m\mapsto u\mapsto\operatorname{blockLogKernel}(\mu,B_2,m,u,1)),
    u\mapsto\operatorname{twoLowerLogFirst}(\mu,r,u),n),
   \texttt{atTop},\mathcal N(0))\\[-2pt]
&\hspace{11mm}\wedge
 \operatorname{Tendsto}
  (n\mapsto\operatorname{compactUniformError}
   (K_2,(m\mapsto u\mapsto\operatorname{blockLogKernel}(\mu,B_2,m,u,2)),
    u\mapsto\operatorname{twoLowerLogSecond}(\mu,r,u),n),
   \texttt{atTop},\mathcal N(0))\\[-2pt]
&\hspace{11mm}\wedge
 \operatorname{Tendsto}
  (n\mapsto\operatorname{compactUniformError}
   (K_2,(m\mapsto u\mapsto\operatorname{blockGKernel}(\mu,B_2,m,u,1)),
    u\mapsto\operatorname{twoLowerGFirst}(\mu,r,u),n),
   \texttt{atTop},\mathcal N(0))\\[-2pt]
&\hspace{11mm}\wedge
 \operatorname{Tendsto}
  (n\mapsto\operatorname{compactUniformError}
   (K_2,(m\mapsto u\mapsto\operatorname{blockGKernel}(\mu,B_2,m,u,2)),
    u\mapsto\operatorname{twoLowerGSecond}(\mu,r,u),n),
   \texttt{atTop},\mathcal N(0))\Bigr)\\
&\quad\wedge
 \operatorname{Tendsto}
  (n\mapsto\operatorname{compactUniformError}
   (K_3,(m\mapsto u\mapsto\operatorname{blockLogKernel}(\mu,B_3,m,u,1)),
    u\mapsto\operatorname{threeLogFirst}(\mu,a,b,u),n),
   \texttt{atTop},\mathcal N(0))\\[-2pt]
&\quad\wedge
 \operatorname{Tendsto}
  (n\mapsto\operatorname{compactUniformError}
   (K_3,(m\mapsto u\mapsto\operatorname{blockLogKernel}(\mu,B_3,m,u,2)),
    u\mapsto\operatorname{threeLogSecond}(\mu,a,b,u),n),
   \texttt{atTop},\mathcal N(0))\\[-2pt]
&\quad\wedge
 \operatorname{Tendsto}
  (n\mapsto\operatorname{compactUniformError}
   (K_3,(m\mapsto u\mapsto\operatorname{blockGKernel}(\mu,B_3,m,u,1)),
    u\mapsto\operatorname{threeGFirst}(\mu,a,b,u),n),
   \texttt{atTop},\mathcal N(0))\\[-2pt]
&\quad\wedge
 \operatorname{Tendsto}
  (n\mapsto\operatorname{compactUniformError}
   (K_3,(m\mapsto u\mapsto\operatorname{blockGKernel}(\mu,B_3,m,u,2)),
    u\mapsto\operatorname{threeGSecond}(\mu,a,b,u),n),
   \texttt{atTop},\mathcal N(0)).
\end{align*}
\end{proposition}

\begin{proof}
Apply \texttt{dsimp only} to expose the two leading local lets in the
conclusion.  Let $H_2$ be the result of applying
\cref{prop:block-limit-passage} to $B_2$ with the exact terms
\[
 T=(\{r\}:\operatorname{Finset}(\R)),\qquad
 k=\operatorname{twoDominanceMap}(r),\qquad j_\infty=1.
\]
Let $P_2:=\operatorname{twoBlockDominancePackage}(r,R,h_r,h_R)$.
Its measurable, top-value, and strict finite-dominance arguments are
$P_2.1$, $P_2.2.1$, and $P_2.2.2$, and its final
unique-top argument is
$\operatorname{twoBlockTopUnique}(r,R,h_r,h_R)$.  Pass
\[
 \texttt{by simpa [eq\_comm] using h\_r1}
\]
for $1\notin T$.  For the finite-set null family pass
\[
 \texttt{by intro s hs; simpa only [Finset.mem\_singleton] at hs;
 subst s; exact h\_mu\_r}.
\]
Here \texttt{h\_mu\_r} is the Lean binder displayed as $h_{\mu,r}$.
Also pass $K_2,h_{K_2},h_{K_2,0}$ explicitly.
After \texttt{dsimp only at H\_2}, the two raw uniform signed-integral
convergences are exactly
\[
 \texttt{H\_2.2.2.2.2.1}
 \qquad\text{and}\qquad
 \texttt{H\_2.2.2.2.2.2}.
\]
Let $H_3$ be the result of applying the same proposition to $B_3$ with
\[
 T=(\{a,b\}:\operatorname{Finset}(\R)),\qquad
 k=\operatorname{threeDominanceMap}(a,b),\qquad j_\infty=2,
\]
Let $P_3:=\operatorname{threeBlockDominancePackage}
(a,b,R_3,h_a,h_{ab},h_{R_3})$.  Pass $P_3.1$, $P_3.2.1$, and $P_3.2.2$,
and
$\operatorname{threeBlockTopUnique}(a,b,R_3,h_a,h_{ab},h_{R_3})$; here
$1\notin T$ is
\texttt{by simpa [eq\_comm] using And.intro h\_a1 h\_b1}.  The finite-set
null family is
\[
 \texttt{by intro s hs; simp only [Finset.mem\_insert,
 Finset.mem\_singleton] at hs; rcases hs with rfl | rfl;
 all\_goals assumption}.
\]
The assumptions used by \texttt{all\_goals assumption} are the Lean binders
\texttt{h\_mu\_a} and \texttt{h\_mu\_b} displayed as
$h_{\mu,a}$ and $h_{\mu,b}$.
Pass $K_3,h_{K_3},h_{K_3,0}$ explicitly.
After \texttt{dsimp only at H\_3}, the two corresponding raw convergence
fields are
\[
 \texttt{H\_3.2.2.2.2.1}
 \qquad\text{and}\qquad
 \texttt{H\_3.2.2.2.2.2}.
\]
For every invocation of
\cref{lem:block-curve-kernel-bridge,lem:block-log-mass-limit,lem:block-kernel-regularity}
with $B_2$, take $j_\infty=1$ and
$h_\infty:=\operatorname{twoBlockTopUnique}(r,R,h_r,h_R)$; with $B_3$, take
$j_\infty=2$ and
$h_\infty:=\operatorname{threeBlockTopUnique}
 (a,b,R_3,h_a,h_{ab},h_{R_3})$.
For every $n$ and $u$, instantiate
\cref{lem:block-g-kernel-integrals} with those same two pairs
$(j_\infty,h_\infty)$.  Rewrite the two signed integrals returned by
\cref{prop:block-limit-passage} as
the $q=1$ and $q=2$ values of
\texttt{blockGKernel $\mu$ B\_2 n u q} and
\texttt{blockGKernel $\mu$ B\_3 n u q}, respectively; then apply
\cref{lem:subtype-uniform-error-bridge} to identify the two uniform integral
errors with the four \texttt{compactUniformError} terms in the statement.
In the $1<r$ and $r<1$ branches, the order-one proof supplied to
\cref{lem:block-log-mass-limit} is respectively
$(\operatorname{twoOrderOneDominancePackage}(r,R,h_r,h_R)).1(h_{1r})$
and
$(\operatorname{twoOrderOneDominancePackage}(r,R,h_r,h_R)).2(h_{r1}')$,
where $h_{1r}:1<r$ and $h_{r1}':r<1$ are the implication binders of those
branches.  For $B_3$, use
$\operatorname{threeOrderOneDominant}
(a,b,R_3,h_a,h_{ab},h_{R_3},h_{a1},h_{b1})$ and
$k=\operatorname{threeNormBlock}(a,b)$.
For each family and each $q\in\{1,2\}$, name the two convergence proofs
$H_{\rm mass}$ and $H_G$.  The former is the appropriate projection of
the second conjunct of \cref{lem:block-log-mass-limit}; the latter is the
rewritten projection
\texttt{H\_2.2.2.2.2.1}/\texttt{H\_2.2.2.2.2.2}, or its $H_3$
analogue.  Apply \cref{lem:compact-uniform-add} to these two proofs.
For its source-continuity hypotheses use
\cref{lem:block-log-mass-continuity} and, respectively,
\[
 \bigl(\operatorname{blockKernelRegularity}
  (\mu,B,j_\infty,h_\infty,n)\bigr).2.2.1
 \quad\hbox{or}\quad
 \bigl(\operatorname{blockKernelRegularity}
  (\mu,B,j_\infty,h_\infty,n)\bigr).2.2.2.
\]
The coordinate and negative-square mass targets are continuous by
\texttt{fun\_prop}; the two- and three-block $G$ targets are finite
polynomials and are also discharged by \texttt{fun\_prop}.  Finally rewrite
the source sum with the first conjunct of
\cref{lem:block-log-mass-limit} and simplify the target sums using
\cref{def:two-block-limits,def:three-block-limits}.  This is the exact
compact-uniform addition step producing all logarithmic-kernel conclusions;
pointwise \texttt{Tendsto.add} is not being used in its place.
The positivity and maximal-coordinate intervals follow from
\cref{lem:positivity-intervals}. After factoring out
$\operatorname{blockScale}(n)^R$, the two
power-sum exponents are
\[
 0\quad\text{and}\quad\alpha-r.
\]
Thus block $0$ dominates below $r$ and block $1$ above $r$. For the
three-block family, the exponents are
\[
 0,\qquad\alpha-a,\qquad2\alpha-a-b,
\]
so blocks $0,1,2$ dominate on the three sets used by
\texttt{threeCellMoment}, respectively. The block that
dominates at order $1$ also dominates the $\ell^1$ norm.

At every nonthreshold order, \cref{lem:exact-derivatives} and
\cref{lem:dominant-block} give, when block $j$ dominates at $\alpha$ and block
$k$ dominates at $1$,
\begin{align*}
 H_{\operatorname{finiteParam}(\alpha)}'(0)
   &\longrightarrow\omega(\operatorname{finiteParam}(\alpha))(u_j-u_k),\\
 H_{\operatorname{finiteParam}(\alpha)}''(0)
   &\longrightarrow\omega(\operatorname{finiteParam}(\alpha))(u_k^2-u_j^2).
\end{align*}
The first two derivatives of the logarithm of the norm converge to $u_k$ and
$-u_k^2$.  For each $u$, apply
\cref{lem:two-block-limit-integral-eval} with $h_{\mu,r}$ in the two
implication branches and apply
\cref{lem:three-block-limit-integral-eval} with $h_{\mu,a},h_{\mu,b}$.
These are the exact signed-integral equalities giving the displayed
norm-free formulas; the threshold values disappear by the named null
hypotheses, not by definitional reduction.  The
interchange of limit and integral, including the neighbourhood of $1$, the
active crossing points, and the large-order tail, is justified by
\cref{prop:block-limit-passage}.
\end{proof}

The combined proposition is convenient when both perturbations are already
available.  The necessity proof also needs the following two dependency-clean
interfaces; neither asks for data belonging to the other perturbation.

\begin{proposition}[Two-block localisation interface]
\label{prop:two-block-localisation}
Let $\mu:\operatorname{SignedMeasure}(\overline\R_+)$, let $r,R:\R$ have
proofs $h_r:0<r$, $h_{r1}:r\neq1$, and $h_R:r<R$, and let
$K:\operatorname{Set}(\operatorname{Fin}(2)\to\R)$ have proofs
$h_{K0}:K.\operatorname{Nonempty}$ and
$h_K:\operatorname{IsCompact}(K)$.  Assume
\[
 h_{\mu,r}:\operatorname{signedTV}(\mu)
   (\{\operatorname{finiteParam}(r)\})=0.
\]
Then the exact conclusion is
\begin{align*}
&\texttt{let }B:=\operatorname{twoBlockData}(r,R,h_r,h_R)\texttt{; }\\[-2pt]
&\Bigl(1<r\to
 \operatorname{CompactUniformConverges}
   (K,(n\mapsto u\mapsto\operatorname{blockLogKernel}(\mu,B,n,u,1)),
      u\mapsto\operatorname{twoUpperLogFirst}(\mu,r,u))\\[-2pt]
&\hspace{17mm}\wedge
 \operatorname{CompactUniformConverges}
   (K,(n\mapsto u\mapsto\operatorname{blockLogKernel}(\mu,B,n,u,2)),
      u\mapsto\operatorname{twoUpperLogSecond}(\mu,r,u))\\[-2pt]
&\hspace{17mm}\wedge
 \operatorname{CompactUniformConverges}
   (K,(n\mapsto u\mapsto\operatorname{blockGKernel}(\mu,B,n,u,1)),
      u\mapsto\operatorname{twoUpperGFirst}(\mu,r,u))\\[-2pt]
&\hspace{17mm}\wedge
 \operatorname{CompactUniformConverges}
   (K,(n\mapsto u\mapsto\operatorname{blockGKernel}(\mu,B,n,u,2)),
      u\mapsto\operatorname{twoUpperGSecond}(\mu,r,u))\Bigr)\\
&\quad\wedge\Bigl(r<1\to
 \operatorname{CompactUniformConverges}
   (K,(n\mapsto u\mapsto\operatorname{blockLogKernel}(\mu,B,n,u,1)),
      u\mapsto\operatorname{twoLowerLogFirst}(\mu,r,u))\\[-2pt]
&\hspace{17mm}\wedge
 \operatorname{CompactUniformConverges}
   (K,(n\mapsto u\mapsto\operatorname{blockLogKernel}(\mu,B,n,u,2)),
      u\mapsto\operatorname{twoLowerLogSecond}(\mu,r,u))\\[-2pt]
&\hspace{17mm}\wedge
 \operatorname{CompactUniformConverges}
   (K,(n\mapsto u\mapsto\operatorname{blockGKernel}(\mu,B,n,u,1)),
      u\mapsto\operatorname{twoLowerGFirst}(\mu,r,u))\\[-2pt]
&\hspace{17mm}\wedge
 \operatorname{CompactUniformConverges}
   (K,(n\mapsto u\mapsto\operatorname{blockGKernel}(\mu,B,n,u,2)),
      u\mapsto\operatorname{twoLowerGSecond}(\mu,r,u))\Bigr).
\end{align*}
\end{proposition}

\begin{proof}
Apply \texttt{dsimp only} to expose the leading local definition of $B$.
Let
$P_2:=\operatorname{twoBlockDominancePackage}(r,R,h_r,h_R)$, and let
$H_2$ be \cref{prop:block-limit-passage} instantiated with
$T=(\{r\}:\operatorname{Finset}(\R))$,
$k=\operatorname{twoDominanceMap}(r)$, $j_\infty=1$, and
$K,h_K,h_{K0}$.  Pass $P_2.1$, $P_2.2.1$, and $P_2.2.2$ as the measurable,
top-value, and strict-dominance hypotheses, and pass
$\operatorname{twoBlockTopUnique}(r,R,h_r,h_R)$ as the unique-top
hypothesis.  The remaining two proof terms are literally
\[
 \texttt{by simpa [eq\_comm] using h\_r1}
\]
for $1\notin T$, and
\[
 \texttt{by intro s hs; simpa only [Finset.mem\_singleton] at hs;
 subst s; exact h\_mu\_r}
\]
for the finite-set null family, where \texttt{h\_mu\_r} is the Lean binder
displayed as $h_{\mu,r}$.  After \texttt{dsimp only at H\_2}, its two raw
uniform signed-integral convergence proofs are exactly
\[
 \texttt{H\_2.2.2.2.2.1}
 \qquad\text{and}\qquad
 \texttt{H\_2.2.2.2.2.2}.
\]
The only crossing is the supplied null point
$r$.  In each implication branch,
\cref{lem:two-block-limit-integral-eval}, instantiated with $h_{\mu,r}$,
identifies the two resulting limit integrals with the appropriate two
norm-free polynomials in \cref{def:two-block-limits}.  Before passing from
the two uniform signed-integral conclusions to the displayed
\texttt{blockGKernel} conclusions, rewrite at every $n,u$ with
\cref{lem:block-g-kernel-integrals} instantiated by $j_\infty=1$ and
$h_\infty:=\operatorname{twoBlockTopUnique}(r,R,h_r,h_R)$; then
\cref{lem:subtype-uniform-error-bridge} gives the required compact-uniform
errors.  Obtain the two logarithmic-mass limits from
\cref{lem:block-log-mass-limit}, passing
$j_\infty=1$ and
$h_\infty:=\operatorname{twoBlockTopUnique}(r,R,h_r,h_R)$; the block that
dominates at order $1$ is $0$ when $1<r$ and $1$ when $r<1$.  The splitting
identity in that lemma gives a pointwise sum, not yet a
compact-uniform conclusion.  Apply \cref{lem:compact-uniform-add} separately
for $q=1$ and $q=2$, using
\cref{lem:block-log-mass-continuity} for the mass source,
\texttt{(blockKernelRegularity $\mu$ B 1
 (twoBlockTopUnique r R h\_r h\_R) n).2.2.1} or
\texttt{.2.2.2} for the $G$ source, and \texttt{fun\_prop} for the
polynomial targets.  Rewrite by the splitting identity and simplify
\cref{def:two-block-limits}; these are the two logarithmic-column limits.
More
exactly, its $h_k$ argument in the two implication branches is
$(\operatorname{twoOrderOneDominancePackage}(r,R,h_r,h_R)).1(h_{1r})$ and
$(\operatorname{twoOrderOneDominancePackage}(r,R,h_r,h_R)).2(h_{r1}')$,
where $h_{1r}:1<r$ and $h_{r1}':r<1$ are the corresponding local binders.
\end{proof}

\begin{proposition}[Three-block localisation interface]
\label{prop:three-block-localisation}
Let $\mu:\operatorname{SignedMeasure}(\overline\R_+)$ and let $a,b,R:\R$
have proofs $h_a:0<a$, $h_{ab}:a<b$, $h_{a1}:a\neq1$,
$h_{b1}:b\neq1$, and $h_R:a+b<R$.  Let
$K:\operatorname{Set}(\operatorname{Fin}(3)\to\R)$ have proofs
$h_{K0}:K.\operatorname{Nonempty}$ and
$h_K:\operatorname{IsCompact}(K)$.  Assume
\[
 h_\mu:\operatorname{signedTV}(\mu)
   (\{\operatorname{finiteParam}(a)\})=0
 \quad\wedge\quad
 \operatorname{signedTV}(\mu)(\{\operatorname{finiteParam}(b)\})=0.
\]
Then the following literal conjunction holds:
\begin{align*}
&\texttt{let }B:=\operatorname{threeBlockData}
 (a,b,R,h_a,h_{ab},h_R)\texttt{; }\\[-2pt]
&\operatorname{CompactUniformConverges}
 (K,(n\mapsto u\mapsto\operatorname{blockLogKernel}(\mu,B,n,u,1)),
    u\mapsto\operatorname{threeLogFirst}(\mu,a,b,u))\\
&\quad\wedge\operatorname{CompactUniformConverges}
 (K,(n\mapsto u\mapsto\operatorname{blockLogKernel}(\mu,B,n,u,2)),
    u\mapsto\operatorname{threeLogSecond}(\mu,a,b,u))\\
&\quad\wedge\operatorname{CompactUniformConverges}
 (K,(n\mapsto u\mapsto\operatorname{blockGKernel}(\mu,B,n,u,1)),
    u\mapsto\operatorname{threeGFirst}(\mu,a,b,u))\\
&\quad\wedge\operatorname{CompactUniformConverges}
 (K,(n\mapsto u\mapsto\operatorname{blockGKernel}(\mu,B,n,u,2)),
    u\mapsto\operatorname{threeGSecond}(\mu,a,b,u)).
\end{align*}
\end{proposition}

\begin{proof}
Obtain
$\langle h_{\mu,a},h_{\mu,b}\rangle:=h_\mu$, and apply
\texttt{dsimp only} to expose the leading definition of $B$.  Let
$P_3:=\operatorname{threeBlockDominancePackage}
(a,b,R,h_a,h_{ab},h_R)$, and let $H_3$ be
\cref{prop:block-limit-passage} instantiated with
$T=(\{a,b\}:\operatorname{Finset}(\R))$,
$k=\operatorname{threeDominanceMap}(a,b)$, $j_\infty=2$, and
$K,h_K,h_{K0}$.  Supply $P_3.1$, $P_3.2.1$, $P_3.2.2$, and
$\operatorname{threeBlockTopUnique}(a,b,R,h_a,h_{ab},h_R)$.  The exact
$1\notin T$ argument is
\[
 \texttt{by simpa [eq\_comm] using And.intro h\_a1 h\_b1},
\]
and the exact finite-set null-family argument is
\[
 \texttt{by intro s hs; simp only [Finset.mem\_insert,
 Finset.mem\_singleton] at hs; rcases hs with rfl | rfl;
 all\_goals assumption}.
\]
After \texttt{dsimp only at H\_3}, its two raw uniform signed-integral
convergence proofs are
\[
 \texttt{H\_3.2.2.2.2.1}
 \qquad\text{and}\qquad
 \texttt{H\_3.2.2.2.2.2}.
\]
Its only
crossings are the two assumed total-variation-null points.  Apply
\cref{lem:three-block-limit-integral-eval} with these two named proofs to
identify the integrals of the cellwise limits with the targets in
\cref{def:three-block-limits}; this is an almost-everywhere equality, not a
definitional one.  Rewrite the two uniform signed-integral
conclusions as the two \texttt{blockGKernel} conclusions by
\cref{lem:block-g-kernel-integrals}, at every $n,u$, with $j_\infty=2$ and
$h_\infty:=\operatorname{threeBlockTopUnique}
(a,b,R,h_a,h_{ab},h_R)$; \cref{lem:subtype-uniform-error-bridge} gives the
displayed compact-uniform errors.  Add the order-$1$ logarithmic-mass limits
from \cref{lem:block-log-mass-limit}, passing $j_\infty=2$ and
$h_\infty:=\operatorname{threeBlockTopUnique}(a,b,R,h_a,h_{ab},h_R)$, and
passing
$\operatorname{threeOrderOneDominant}
(a,b,R,h_a,h_{ab},h_R,h_{a1},h_{b1})$ as $h_k$ for
$k=\operatorname{threeNormBlock}(a,b)$.  Use its exact splitting identity
and apply \cref{lem:compact-uniform-add} for $q=1$ and $q=2$.
The required source continuity is supplied by
\cref{lem:block-log-mass-continuity} and the projections
\texttt{.2.2.1}, \texttt{.2.2.2} of
\texttt{blockKernelRegularity $\mu$ B 2
 (threeBlockTopUnique a b R h\_a h\_ab h\_R) n}; all target functions are
finite polynomials, so \texttt{fun\_prop} supplies their continuity.
Simplification with \cref{def:three-block-limits} identifies the target sum
and gives the two logarithmic-column limits.
\end{proof}

\subsection{Shannon-point localization}

The next perturbation detects an atom at order $1$.  Its entire dependent
carrier and every derivative kernel are fixed by the following declarations.

\begin{definition}[Shannon block data]
\label{def:shannon-data}
Define the structure
\[
 \texttt{structure ShannonData where}\qquad
 \texttt{c R p : $\R$}\qquad
 \texttt{c\_gt\_one : 1 < c}\qquad
 \texttt{R\_gt : c + 1 < R}\qquad
 \texttt{p\_pos : 0 < p}\qquad
 \texttt{p\_lt\_one : p < 1}.
\]
For $\theta:\operatorname{ShannonData}$, define the namespace declaration
\[
 \texttt{def ShannonData.q (theta : ShannonData) : $\R$ :=
   1 - theta.p}.
\]
Thus $\theta.q$ is an actual identifier rather than an informal field.
Define also the named total scale functions
\[
 \operatorname{shannonScale}(n):=\operatorname{blockScale}(n),\qquad
 \operatorname{shannonLogScale}(n):=
   \log(\operatorname{shannonScale}(n):\R).
\]
The paper symbols $d_n$ and $L_n$ are pure notation for these two functions;
they are not scoped local variables.  Define one exhaustive function on
$\operatorname{Fin}(3)$:
\begin{align*}
 \operatorname{shannonCount}(\theta,n)&:=
  \operatorname{Fin.cases}\!\left(
   \operatorname{Nat.ceil}
    (\operatorname{Real.rpow}(\operatorname{shannonScale}(n),\theta.R)),
   j\mapsto\operatorname{Fin.cases}\!\left(
    \operatorname{Nat.ceil}
     (\operatorname{Real.rpow}
       (\operatorname{shannonScale}(n),\theta.R-1)),
    \_\mapsto\operatorname{Nat.ceil}
     (\operatorname{Real.rpow}
       (\operatorname{shannonScale}(n),\theta.R-1-\theta.c)),j\right)
   \right),\\
 \operatorname{ShannonIndex}(\theta,n)&:=
  \mathop{\Sigma}_{j:\operatorname{Fin}(3)}
    \operatorname{Fin}(\operatorname{shannonCount}(\theta,n,j)).
\end{align*}
\end{definition}

\begin{lemma}[Positivity of the complementary Shannon weight]
\label{lem:shannon-q-pos}
For $\theta:\operatorname{ShannonData}$, the exported proof term is
\[
 \texttt{ShannonData.q\_pos} (\theta):0<\theta.q.
\]
\end{lemma}

\begin{proof}
This is
\texttt{sub\_pos.mpr theta.p\_lt\_one} after unfolding
\texttt{ShannonData.q}.
\end{proof}

\begin{lemma}[Positivity of the Shannon logarithmic scale]
\label{lem:shannon-log-scale-pos}
For $n:\N$, the exported proof term is
\[
 \operatorname{shannonLogScale\_pos}(n):
 0<\operatorname{shannonLogScale}(n).
\]
\end{lemma}

\begin{proof}
Since $1<(n+2:\R)$, this is \texttt{Real.log\_pos} after unfolding
\texttt{shannonLogScale}, \texttt{shannonScale}, and
\texttt{blockScale}.
\end{proof}

\begin{lemma}[Positive Shannon block counts]
\label{lem:shannon-count-pos}
For $\theta:\operatorname{ShannonData}$, $n:\N$, and
$j:\operatorname{Fin}(3)$, the named proof is
\[
 \operatorname{shannonCount\_pos}(\theta,n,j):
 0<\operatorname{shannonCount}(\theta,n,j).
\]
\end{lemma}

\begin{proof}
Split on $j$.  The three real exponents are positive by the fields
$1<\theta.c$ and $\theta.c+1<\theta.R$; the base is at least two, so its
real power and then its natural ceiling are positive.
\end{proof}

\begin{definition}[Nonempty Shannon index and block representative]
\label{def:shannon-index-nonempty}
For $\theta:\operatorname{ShannonData}$ and $n:\N$, define
\[
 \operatorname{shannonIndexNonempty}(\theta,n):
 \operatorname{Nonempty}(\operatorname{ShannonIndex}(\theta,n))
 :=\left\langle
  \left\langle0,
   \left\langle0,\operatorname{shannonCount\_pos}(\theta,n,0)\right\rangle
  \right\rangle
 \right\rangle.
\]
For $j:\operatorname{Fin}(3)$, define the genuine coordinate
\[
 \operatorname{shannonRepresentative}(\theta,n,j):=
 \left\langle j,
  \left\langle0,\operatorname{shannonCount\_pos}(\theta,n,j)\right\rangle
 \right\rangle
 \in\operatorname{ShannonIndex}(\theta,n).
\]
\end{definition}

\begin{definition}[Shannon line and its escort]
\label{def:shannon-line-data}
For $\theta:\operatorname{ShannonData}$, $n:\N$, and $z:\R\times\R$,
define on $\operatorname{ShannonIndex}(\theta,n)$
\begin{align*}
 \operatorname{shannonBase}(\theta,n,z)(i)&:=
  \operatorname{Fin.cases}\left(
   \theta.q,
   \left(j\mapsto\operatorname{Fin.cases}
    (\operatorname{shannonScale}(n)\theta.p,
     (\_\mapsto\operatorname{shannonScale}(n)^2),j)\right),i.1\right),\\
 \operatorname{shannonVelocity}(\theta,n,z)(i)&:=
  \operatorname{Fin.cases}\left(
   \frac{z.1}{\theta.q\operatorname{shannonLogScale}(n)},
   \left(j\mapsto\operatorname{Fin.cases}
    \left(-\frac{z.1}{\theta.p\operatorname{shannonLogScale}(n)},
      (\_\mapsto z.2),j\right)\right),i.1\right),\\
 \operatorname{shannonBase\_pos}(\theta,n,z,i)&:
   0<\operatorname{shannonBase}(\theta,n,z)(i)
   :=\texttt{by
    rcases i with $\langle$j, hj$\rangle$
    fin\_cases j
    $\cdot$ exact theta.q\_pos
    $\cdot$ exact mul\_pos (by positivity) theta.p\_pos
    $\cdot$ exact pow\_pos (by positivity) 2},\\
 \operatorname{shannonLineData}(\theta,n,z)&:=
  \left(\texttt{letI := shannonIndexNonempty theta n; }\quad
  \left\langle\operatorname{shannonBase}(\theta,n,z),
   \operatorname{shannonVelocity}(\theta,n,z),
   \operatorname{shannonBase\_pos}(\theta,n,z)\right\rangle\right)
  \in\operatorname{PositiveLineData}(\operatorname{ShannonIndex}(\theta,n)),\\
 \operatorname{shannonLineRaw}(\theta,n,z,\lambda)&:=
  \left(\texttt{letI := shannonIndexNonempty theta n; }\quad
  \operatorname{lineRaw}(\operatorname{shannonLineData}(\theta,n,z),\lambda)\right).
 \tag{SB}\label{eq:shannon-block}
\end{align*}
Define the total mass and the normalized mass of an entire block by
\begin{align*}
 \operatorname{shannonMass}(\theta,n,z,\lambda)&:=
  \left(\texttt{letI := shannonIndexNonempty theta n; }\quad
  \sum_{i:\operatorname{ShannonIndex}(\theta,n)}
   \operatorname{shannonLineRaw}(\theta,n,z,\lambda)(i)\right),\\
 \operatorname{shannonBlockMass}(\theta,n,z,j,\lambda)&:=
  \frac{(\operatorname{shannonCount}(\theta,n,j):\R)\,
    \operatorname{shannonLineRaw}(\theta,n,z,\lambda)
      (\operatorname{shannonRepresentative}(\theta,n,j))}
   {\operatorname{shannonMass}(\theta,n,z,\lambda)}.
\end{align*}
For $j:\operatorname{Fin}(3)$ and $\alpha:\R$, define
\begin{align*}
 \operatorname{shannonContribution}(\theta,n,z,j,\alpha)&:=
 (\operatorname{shannonCount}(\theta,n,j):\R)
 \operatorname{Real.rpow}
  (\operatorname{shannonBase}(\theta,n,z)
    (\operatorname{shannonRepresentative}(\theta,n,j)),\alpha),\\
 \operatorname{shannonEscort}(\theta,n,z,j,\alpha)&:=
 \frac{\operatorname{shannonContribution}(\theta,n,z,j,\alpha)}
 {\sum_{\ell:\operatorname{Fin}(3)}
  \operatorname{shannonContribution}(\theta,n,z,\ell,\alpha)},\\
 \operatorname{shannonBlockVelocity}(\theta,n,z,j)&:=
  \operatorname{shannonVelocity}(\theta,n,z)
   (\operatorname{shannonRepresentative}(\theta,n,j)),\\
 \operatorname{shannonLogBase}(\theta,n,z,j)&:=
  \log\!\left(\operatorname{shannonBase}(\theta,n,z)
   (\operatorname{shannonRepresentative}(\theta,n,j))\right),\\
 \operatorname{shannonEscortLogMean}(\theta,n,z,\alpha)&:=
  \sum_j\operatorname{shannonEscort}(\theta,n,z,j,\alpha)
    \operatorname{shannonLogBase}(\theta,n,z,j),\\
 \operatorname{shannonMean}(\theta,n,z,\alpha)&:=
  \sum_j\operatorname{shannonEscort}(\theta,n,z,j,\alpha)
   \operatorname{shannonBlockVelocity}(\theta,n,z,j),\\
 \operatorname{shannonSecond}(\theta,n,z,\alpha)&:=
  \sum_j\operatorname{shannonEscort}(\theta,n,z,j,\alpha)
   \operatorname{shannonBlockVelocity}(\theta,n,z,j)^2,\\
 \operatorname{shannonVar}(\theta,n,z,\alpha)&:=
  \operatorname{shannonSecond}(\theta,n,z,\alpha)
   -\operatorname{shannonMean}(\theta,n,z,\alpha)^2.
\end{align*}
\end{definition}

\begin{lemma}[The third Shannon block is strictly maximal at the base point]
\label{lem:shannon-top-unique-at-zero}
For $\theta:\operatorname{ShannonData}$, $n:\N$, and
$z:\R\times\R$, define the named proof
\[
 \operatorname{shannonTopUniqueAtZero}(\theta,n,z):
 \forall i:\operatorname{ShannonIndex}(\theta,n),\quad
 i.1\neq2\to
 \operatorname{shannonLineRaw}(\theta,n,z,0)(i)
 <\operatorname{shannonLineRaw}(\theta,n,z,0)
   (\operatorname{shannonRepresentative}(\theta,n,2)).
\]
Thus the assertion is strict only between blocks; all coordinates inside
block $2$ may tie, and they have the same velocity $z.2$.
\end{lemma}

\begin{proof}
At $\lambda=0$ the three coordinate values are
$\theta.q$, $\operatorname{shannonScale}(n)\theta.p$, and
$\operatorname{shannonScale}(n)^2$.  The fields
$0<\theta.p<1$ give $0<\theta.q<1$, while
$2\leq\operatorname{shannonScale}(n)$.  Introduce $i$, run
\texttt{rcases i with $\langle$j,hj$\rangle$}, and then
\texttt{fin\_cases j}; the first two cases follow by multiplication
monotonicity and the third is excluded by the transported proof $j\neq2$.
In Lean this finite case split is followed by
\texttt{simp [shannonLineRaw, shannonBase, shannonScale]} and
\texttt{nlinarith}.
\end{proof}

\begin{definition}[Named Shannon kernels]
\label{def:shannon-kernels}
For $\theta:\operatorname{ShannonData}$, $n:\N$, $z:\R\times\R$, and
$\beta:\overline\R_+$, define
\begin{align*}
 \operatorname{shannonKOne}(\theta,n,z,\beta)&:=
  \left(\texttt{letI := shannonIndexNonempty theta n; }\quad
  \operatorname{entropyLineFirst}
    (\operatorname{shannonLineData}(\theta,n,z),\beta,0)\right),\\
 \operatorname{shannonKTwo}(\theta,n,z,\beta)&:=
  \left(\texttt{letI := shannonIndexNonempty theta n; }\quad
  \operatorname{entropyLineSecond}
    (\operatorname{shannonLineData}(\theta,n,z),\beta,0)\right).
\end{align*}
For $\mu:\operatorname{SignedMeasure}(\overline\R_+)$ and $q:\N$, define
\begin{align*}
 \operatorname{shannonLogKernel}(\mu,\theta,n,z,q)&:=
  \left(\texttt{letI := shannonIndexNonempty theta n; }\quad
  \operatorname{iterDeriv}
   (\operatorname{signedLogPhiLine}
     (\mu,\operatorname{shannonLineData}(\theta,n,z))),q,0)\right),\\
 \operatorname{shannonGKernel}(\mu,\theta,n,z,q)&:=
  \left(\texttt{letI := shannonIndexNonempty theta n; }\quad
  \operatorname{iterDeriv}
   (\operatorname{integratedEntropyLine}
     (\mu,\operatorname{shannonLineData}(\theta,n,z))),q,0)\right),\\
 \operatorname{shannonLogMassKernel}(\theta,n,z,q)&:=
  \operatorname{iterDeriv}
   (\lambda\mapsto\log(\operatorname{shannonMass}
      (\theta,n,z,\lambda))),q,0).
\end{align*}
For the same arguments define the fully typed scalar curves
\begin{align*}
 \operatorname{shannonPhiCurve}(\mu,\theta,n,z,\lambda):=
  \left(\texttt{letI := shannonIndexNonempty theta n; }\quad
  \operatorname{PhiSigned}_{\operatorname{ShannonIndex}(\theta,n)}
   \left(\mu,\operatorname{lineConeTotal}
    (\operatorname{shannonLineData}(\theta,n,z),\lambda)\right)\right),\\
 \operatorname{shannonGCurve}(\mu,\theta,n,z,\lambda):=
  \left(\texttt{letI := shannonIndexNonempty theta n; }\quad
  \operatorname{GSigned}_{\operatorname{ShannonIndex}(\theta,n)}
   \left(\mu,\operatorname{linePosConeTotal}
    (\operatorname{shannonLineData}(\theta,n,z),\lambda)\right)\right).
\end{align*}
\end{definition}

\begin{lemma}[Shannon curve--kernel bridge]
\label{lem:shannon-curve-kernel-bridge}
For $\mu:\operatorname{SignedMeasure}(\overline\R_+)$,
$\theta:\operatorname{ShannonData}$, $n:\N$, and $z:\R\times\R$, there
exists $\eps:\R$ such that the following literal conjunction holds:
\begin{align*}
&0<\eps\\
&\quad\wedge\Bigl(\forall\lambda\in\operatorname{Set.Ioo}(-\eps,\eps),
 \operatorname{LinePositive}
  (\operatorname{shannonLineData}(\theta,n,z),\lambda)\Bigr)\\
&\quad\wedge\quad
 0<\operatorname{shannonPhiCurve}(\mu,\theta,n,z,0)\\
&\quad\wedge
 \operatorname{ContDiffOn}\!\left(\R,2,
  \lambda\mapsto\operatorname{shannonPhiCurve}(\mu,\theta,n,z,\lambda),
  \operatorname{Set.Ioo}(-\eps,\eps)\right)\\
&\quad\wedge
 \operatorname{ContDiffOn}\!\left(\R,2,
  \lambda\mapsto\operatorname{shannonGCurve}(\mu,\theta,n,z,\lambda),
  \operatorname{Set.Ioo}(-\eps,\eps)\right)\\
&\quad\wedge\Bigl(\forall q:\N,\ q\leq2\to
 \operatorname{iterDeriv}
  (\lambda\mapsto\log(\operatorname{shannonPhiCurve}
    (\mu,\theta,n,z,\lambda)),q,0)
 =\operatorname{shannonLogKernel}(\mu,\theta,n,z,q)\Bigr)\\
&\quad\wedge\Bigl(\forall q:\N,\ q\leq2\to
 \operatorname{iterDeriv}
  (\operatorname{shannonGCurve}(\mu,\theta,n,z),q,0)
 =\operatorname{shannonGKernel}(\mu,\theta,n,z,q)\Bigr).
\end{align*}
\end{lemma}

\begin{proof}
The carrier $\operatorname{ShannonIndex}(\theta,n)$ is finite and every
coordinate of \texttt{shannonBase} is strictly positive.  Put
\[
 i_\infty:=\operatorname{shannonRepresentative}(\theta,n,2).
\]
By \cref{lem:shannon-top-unique-at-zero}, every coordinate outside block
$2$ is strictly smaller than $i_\infty$ at zero, while every coordinate in
block $2$ has the same base value and velocity as $i_\infty$.  Finiteness
and continuity therefore give named witnesses
$\eps_{\rm top}:\R$, $h_{\rm top}:0<\eps_{\rm top}$, and
\[
 \operatorname{FixedMaxCoordinate}
  (\operatorname{shannonLineData}(\theta,n,z),
   \operatorname{Set.Icc}(-\eps_{\rm top},\eps_{\rm top}),i_\infty).
\]
Choose $\eps>0$ with $\eps\leq\eps_{\rm top}$ and so that
\[
 \eps\,
 \abs{\operatorname{shannonVelocity}(\theta,n,z)(i)}<1
 \qquad\text{for every }i:\operatorname{ShannonIndex}(\theta,n).
\]
Restriction from the $\eps_{\rm top}$ interval supplies the same
\texttt{FixedMaxCoordinate} proof on
$\operatorname{Set.Icc}(-\eps,\eps)$.
Then every raw coordinate is positive on
$\operatorname{Set.Ioo}(-\eps,\eps)$, with no appeal to a later uniform
expansion.  On this interval the total line wrappers select their positive
branches.  Applying \cref{lem:signed-line-column-bridge} at zero gives
$0<\operatorname{shannonPhiCurve}(\mu,\theta,n,z,0)$.  The fixed-max proof
above supplies the compactified-order hypothesis in the exact
entropy-derivative package, and
\cref{lem:signed-integral-differentiate-twice} give the two
\texttt{ContDiffOn} clauses.  Finally use
\cref{lem:signed-line-column-bridge} and differentiate its logarithmic and
norm-free identities through order two.
\end{proof}

\begin{lemma}[Shannon derivative--integral and logarithmic splitting bridge]
\label{lem:shannon-kernel-integral-bridge}
For $\mu:\operatorname{SignedMeasure}(\overline\R_+)$,
$\theta:\operatorname{ShannonData}$, $n:\N$, and $z:\R\times\R$, the
following four exact equalities form one literal conjunction:
\begin{align*}
&\operatorname{shannonGKernel}(\mu,\theta,n,z,1)
 =\operatorname{signedIntegral}
   (\mu,\beta\mapsto\operatorname{shannonKOne}(\theta,n,z,\beta))\\
&\quad\wedge\quad
 \operatorname{shannonGKernel}(\mu,\theta,n,z,2)
 =\operatorname{signedIntegral}
   (\mu,\beta\mapsto\operatorname{shannonKTwo}(\theta,n,z,\beta))\\
&\quad\wedge\quad
 \operatorname{shannonLogKernel}(\mu,\theta,n,z,1)
 =\operatorname{shannonLogMassKernel}(\theta,n,z,1)
  +\operatorname{shannonGKernel}(\mu,\theta,n,z,1)\\
&\quad\wedge\quad
 \operatorname{shannonLogKernel}(\mu,\theta,n,z,2)
 =\operatorname{shannonLogMassKernel}(\theta,n,z,2)
  +\operatorname{shannonGKernel}(\mu,\theta,n,z,2).
\end{align*}
\end{lemma}

\begin{proof}
Put
$i_\infty:=\operatorname{shannonRepresentative}(\theta,n,2)$.
Use \cref{lem:shannon-top-unique-at-zero} and finiteness to choose
$\lambda_0>0$ such that the Shannon line is positive on
$\operatorname{Set.Icc}(-\lambda_0,\lambda_0)$ and
\[
 \operatorname{FixedMaxCoordinate}
  (\operatorname{shannonLineData}(\theta,n,z),
   \operatorname{Set.Icc}(-\lambda_0,\lambda_0),i_\infty).
\]
Instantiate \cref{lem:differentiate-integral} at $\lambda=0$.  Its first
two \texttt{HasDerivAt} clauses, together with the fact that the first
derivative identity holds on a neighbourhood of zero, give the first two
displayed equalities.  Its last two clauses give the corresponding
derivatives of \texttt{signedLogPhiLine}; unfold
\texttt{shannonLogKernel}, \texttt{shannonLogMassKernel}, and the first two
equalities to obtain the two splitting identities.
\end{proof}

\begin{lemma}[Linearity and continuity of the Shannon kernels]
\label{lem:shannon-kernel-regularity}
For $\mu:\operatorname{SignedMeasure}(\overline\R_+)$,
$\theta:\operatorname{ShannonData}$, and $n:\N$, the named theorem term
$\operatorname{shannonKernelRegularity}(\mu,\theta,n)$ has the exact type
\begin{align*}
&\Bigl(\forall z,w:\R\times\R,\quad
 \operatorname{shannonGKernel}(\mu,\theta,n,z+w,1)
 =\operatorname{shannonGKernel}(\mu,\theta,n,z,1)
  +\operatorname{shannonGKernel}(\mu,\theta,n,w,1)\Bigr)\\
&\quad\wedge\Bigl(\forall(c:\R)(z:\R\times\R),\quad
 \operatorname{shannonGKernel}(\mu,\theta,n,c\mathbin{\bullet}z,1)
 =c\operatorname{shannonGKernel}(\mu,\theta,n,z,1)\Bigr)\\
&\quad\wedge\operatorname{Continuous}
  (z\mapsto\operatorname{shannonGKernel}(\mu,\theta,n,z,1))\\
&\quad\wedge\operatorname{Continuous}
  (z\mapsto\operatorname{shannonGKernel}(\mu,\theta,n,z,2)).
\end{align*}
\end{lemma}

\begin{proof}
For each $z$, use \cref{lem:shannon-top-unique-at-zero} and shrink a
neighbourhood of zero exactly as in
\cref{lem:shannon-curve-kernel-bridge}; this supplies the
\texttt{FixedMaxCoordinate} hypothesis needed by the order-$\top$ clause of
\cref{lem:exact-derivatives}.  The first Shannon derivative is then a
finite-dimensional linear form in the two velocity coordinates and the
second is a quadratic polynomial.  Signed-integral linearity and
\cref{lem:differentiate-integral} give the displayed identities and
continuity.
\end{proof}

\begin{lemma}[Dedicated Shannon escort, compact-region, and tail estimates]
\label{lem:shannon-dedicated-escort}
Let $\theta:\operatorname{ShannonData}$ and let
$K:\operatorname{Set}(\R\times\R)$ have proofs
$h_K:\operatorname{IsCompact}(K)$ and
$h_{K0}:K.\operatorname{Nonempty}$.  Then the
following three clauses form one literal conjunction:
\begin{align*}
&\Bigl[\forall\delta:\R,\quad
 0<\delta\to
 \delta<\min\{1/2,(\theta.c-1)/4\}\to
 \exists C:\R,\ 0\leq C\ \wedge\\[-2pt]
&\quad\forall(n:\N)(z:\R\times\R),\ z\in K\to
 \forall\alpha\in\operatorname{Set.Icc}(1-\delta,1+\delta),\quad
 \operatorname{shannonEscort}(\theta,n,z,2,\alpha)
 \leq C\operatorname{Real.rpow}
  (\operatorname{blockScale}(n),-(\theta.c-1)/2)\ \wedge\\[-2pt]
&\quad
 \abs{\operatorname{deriv}
  (s\mapsto\operatorname{shannonEscort}(\theta,n,z,2,s))(\alpha)}
 \leq C\log(\operatorname{blockScale}(n))
  \operatorname{Real.rpow}
   (\operatorname{blockScale}(n),-(\theta.c-1)/2)
 \Bigr]\\
&\quad\wedge\Bigl[\forall I:\operatorname{Set}(\R),\quad
 \operatorname{IsCompact}(I)\to I.\operatorname{Nonempty}\to
 \operatorname{OrdConnected}(I)\to
 I\subseteq\operatorname{Set.Ici}(0)\to
 I\cap\{1,\theta.c\}=\varnothing\to\\[-2pt]
&\quad\exists(k:\operatorname{Fin}(3))(\eta,C_I,C_K:\R),\quad
 0<\eta\ \wedge\ 0\leq C_I\ \wedge\ 0\leq C_K\ \wedge\\[-2pt]
&\quad\forall(n:\N)(z:\R\times\R),\ z\in K\to\forall\alpha\in I,\quad
 \sum_{j\in\operatorname{Finset.univ.erase}(k)}
  \operatorname{shannonEscort}(\theta,n,z,j,\alpha)
 \leq C_I\operatorname{Real.rpow}(\operatorname{blockScale}(n),-\eta)\ 
 \wedge\\[-2pt]
&\quad
 \abs{\operatorname{shannonMean}(\theta,n,z,\alpha)
       -\operatorname{shannonBlockVelocity}(\theta,n,z,k)}
 +\abs{\operatorname{shannonSecond}(\theta,n,z,\alpha)
       -\operatorname{shannonBlockVelocity}(\theta,n,z,k)^2}
 +\operatorname{shannonVar}(\theta,n,z,\alpha)\\[-2pt]
&\hspace{24mm}\leq
 C_K\operatorname{Real.rpow}(\operatorname{blockScale}(n),-\eta)\ 
 \wedge\\[-2pt]
&\quad
 \abs{\operatorname{deriv}
  (s\mapsto\operatorname{shannonMean}(\theta,n,z,s))(\alpha)}
 +\abs{\operatorname{deriv}
  (s\mapsto\operatorname{shannonSecond}(\theta,n,z,s))(\alpha)}\\[-2pt]
&\hspace{24mm}\leq C_K\log(\operatorname{blockScale}(n))
  \operatorname{Real.rpow}(\operatorname{blockScale}(n),-\eta)
 \Bigr]\\
&\quad\wedge\Bigl[
 \exists(A,\gamma,C:\R),\quad
 \theta.c<A\ \wedge\ 0<\gamma\ \wedge\ 0\leq C\ \wedge\\[-2pt]
&\quad\forall(n:\N)(z:\R\times\R),\ z\in K\to
 \forall\alpha:\R,\ A\leq\alpha\to\\[-2pt]
&\quad
 \operatorname{shannonEscort}(\theta,n,z,0,\alpha)
 +\operatorname{shannonEscort}(\theta,n,z,1,\alpha)
 \leq C\operatorname{Real.rpow}
  (\operatorname{blockScale}(n),-\gamma\alpha)\ 
 \wedge\\[-2pt]
&\quad
 \alpha\bigl(\operatorname{shannonEscort}(\theta,n,z,0,\alpha)
 +\operatorname{shannonEscort}(\theta,n,z,1,\alpha)\bigr)\leq C\ 
 \wedge\\[-2pt]
&\quad
 \abs{\operatorname{shannonKOne}
  (\theta,n,z,\operatorname{finiteParam}(\alpha))}\leq C\ 
 \wedge\
 \abs{\operatorname{shannonKTwo}
  (\theta,n,z,\operatorname{finiteParam}(\alpha))}\leq C\ 
 \wedge\\[-2pt]
&\quad
 \abs{\operatorname{shannonKOne}(\theta,n,z,\top)}\leq C\ 
 \wedge\
 \abs{\operatorname{shannonKTwo}(\theta,n,z,\top)}\leq C
 \Bigr].
\end{align*}
\end{lemma}

\begin{proof}
For arbitrary $n:\N$ and $z:\R\times\R$, introduce the following local
\texttt{let} bindings
\[
 d:=\operatorname{shannonScale}(n),\quad
 L_d:=\operatorname{shannonLogScale}(n),\quad
 p:=\theta.p,\quad q:=\theta.q,\quad c:=\theta.c,\quad R:=\theta.R,
\]
\[
 C_{j,n,z}(\alpha):=\operatorname{shannonContribution}(\theta,n,z,j,\alpha),
 \qquad
 \Pi^{\rm Sh}_{j,n,z}(\alpha):=
   \operatorname{shannonEscort}(\theta,n,z,j,\alpha),
 \qquad
 u_{j,n,z}:=\operatorname{shannonBlockVelocity}(\theta,n,z,j).
\]
These abbreviations are universally quantified by the preceding sentence;
all constants below are independent of $n$ and $z$.  By
\cref{lem:rounding}, each $C_{j,n,z}$ differs by a factor between two fixed
positive constants from
\[
 \operatorname{Real.rpow}(q,\alpha)\operatorname{Real.rpow}(d,R),\qquad
 \operatorname{Real.rpow}(p,\alpha)
   \operatorname{Real.rpow}(d,R-1+\alpha),\qquad
 \operatorname{Real.rpow}(d,R-1-c+2\alpha),
\]
respectively, on every bounded order interval.  Thus their affine exponents
are $R$, $R-1+\alpha$, and $R-1-c+2\alpha$, with crossings exactly at
$1$ and $c$.  On the interval around $1$, the third exponent has deficit at
least $(c-1)/2$, proving the first escort bound.  Differentiation gives
\[
 \partial_\alpha\Pi^{\rm Sh}_{j,n,z}(\alpha)
 =\bigl(\operatorname{shannonLogBase}(\theta,n,z,j)
  -\operatorname{shannonEscortLogMean}(\theta,n,z,\alpha)\bigr)
   \Pi^{\rm Sh}_{j,n,z}(\alpha),
\]
which proves its derivative bound.

On a compact set separated from the crossings, the smallest difference of
the three affine exponents is positive.  The constant factors
$\operatorname{Real.rpow}(p,\alpha)$ and
$\operatorname{Real.rpow}(q,\alpha)$ are bounded above and away from zero
there, so the proof of
\cref{lem:dominant-block} can be repeated: divide each nondominant
$C_{j,n,z}$ by the dominant $C_{k,n,z}$, use the positive exponent gap to bound
the ratio by $C_K\operatorname{Real.rpow}(d,-\eta)$, sum the two ratios, and
differentiate the
softmax identity displayed above.  This gives the three claimed escort,
moment, and derivative estimates with adjusted constants.  The
$d$-dependent velocities remain in a common compact set because
$L_d\geq\log2$.

For $\alpha\geq A>c$, the third block dominates the second with exponent gap
$\alpha-c$ and the first with gap $2\alpha-1-c$.  Increasing $A$ absorbs
the factors $\operatorname{Real.rpow}(p,\alpha)$ and
$\operatorname{Real.rpow}(q,\alpha)$ into a smaller positive gap
$\gamma\alpha$.  The tail bounds and the derivative-kernel bound then follow
from the variance argument in \cref{lem:large-alpha-tail}.

For the two endpoint bounds, put
\[
 i_\infty:=\operatorname{shannonRepresentative}(\theta,n,2).
\]
Invoke $\operatorname{shannonTopUniqueAtZero}(\theta,n,z)$ and the same
finite-coordinate continuity argument as in
\cref{lem:shannon-curve-kernel-bridge} to choose
$\eps_\infty:\R$, $h_{\eps_\infty}:0<\eps_\infty$, and the named proof
\[
 h_{\rm fixed}:\operatorname{FixedMaxCoordinate}
  (\operatorname{shannonLineData}(\theta,n,z),
   \operatorname{Set.Icc}(-\eps_\infty,\eps_\infty),i_\infty).
\]
Every coordinate tied with $i_\infty$ lies in block $2$ and has the literal
velocity $z.2$.  Instantiate the order-$\top$ clause of
\cref{eq:H-infty-derivatives} with $i_\infty$ and $h_{\rm fixed}$; this is
the required fixed-max invocation and yields the last two endpoint estimates.
\end{proof}

\begin{definition}[Shannon main term and uniform remainder]
\label{def:shannon-remainder}
For $\theta:\operatorname{ShannonData}$, $n:\N$, $z:\R\times\R$,
and $\lambda:\R$, define
\begin{align*}
 \operatorname{shannonMain}(\theta,n,z,\lambda)&:=
 \left(\theta.R-\theta.p+
   \frac{z.1\lambda}{\operatorname{shannonLogScale}(n)}\right)
   \operatorname{shannonLogScale}(n)\\[-2pt]
 &\quad-\left(\theta.q+
   \frac{z.1\lambda}{\operatorname{shannonLogScale}(n)}\right)
   \log\left(\theta.q+
   \frac{z.1\lambda}{\operatorname{shannonLogScale}(n)}\right)\\[-2pt]
 &\quad-\left(\theta.p-
   \frac{z.1\lambda}{\operatorname{shannonLogScale}(n)}\right)
   \log\left(\theta.p-
   \frac{z.1\lambda}{\operatorname{shannonLogScale}(n)}\right),\\
 \operatorname{shannonError}(\theta,n,z,\lambda)&:=
  \left(\texttt{letI := shannonIndexNonempty theta n; }\quad
  \operatorname{entropyLine}
   (\operatorname{shannonLineData}(\theta,n,z),1,\lambda)
  -\operatorname{shannonMain}(\theta,n,z,\lambda)\right).
 \tag{SE}\label{eq:shannon-expansion}
\end{align*}
For $K:\operatorname{Set}(\R\times\R)$, $\lambda_0:\R$, and $j,n:\N$,
define
\[
 \operatorname{shannonCtwoError}(\theta,K,\lambda_0,j,n):=
 \operatorname{sSup}\left\{r:\R\ \middle|\
 \exists z\in K,\exists\lambda\in\operatorname{Set.Icc}(-\lambda_0,\lambda_0),
 r=\abs{\operatorname{iterDeriv}
  (\operatorname{shannonError}(\theta,n,z),j,\lambda)}\right\}.
\]
\end{definition}

\begin{lemma}[Uniform Shannon expansion]\label{lem:uniform-shannon-expansion}
Let $\theta:\operatorname{ShannonData}$ and let
$K:\operatorname{Set}(\R\times\R)$ satisfy
$h_K:\operatorname{IsCompact}(K)$ and $h_{K0}:K.\operatorname{Nonempty}$.
Then there exists $\lambda_0:\R$ with $0<\lambda_0$ such that the following
single conjunction holds:
\begin{align*}
&\Bigl(\forall(n:\N)(z:\R\times\R),\ z\in K\to
 \bigl(\texttt{letI := shannonIndexNonempty theta n; }\quad
 \forall\lambda\in\operatorname{Set.Icc}(-\lambda_0,\lambda_0),\quad
 \operatorname{LinePositive}
   (\operatorname{shannonLineData}(\theta,n,z),\lambda)\bigr)\Bigr)\\
&\quad\wedge
 \Bigl(\forall j:\N,\ j\leq2\to
  \operatorname{Tendsto}
   (n\mapsto\operatorname{shannonCtwoError}
      (\theta,K,\lambda_0,j,n),\texttt{atTop},\mathcal N(0))\Bigr)
 \tag{C2}\label{eq:shannon-explicit-C2}\\
&\quad\wedge
 \operatorname{Tendsto}
  \left(n\mapsto\operatorname{compactUniformError}
   \left(K,(m\mapsto z\mapsto
    \operatorname{shannonKOne}(\theta,m,z,1)),
    (z\mapsto z.1),n\right),\texttt{atTop},\mathcal N(0)\right)
 \tag{S1}\label{eq:shannon-H1-first}\\
&\quad\wedge
 \operatorname{Tendsto}
  \left(n\mapsto\operatorname{compactUniformError}
   \left(K,(m\mapsto z\mapsto
    \operatorname{shannonKTwo}(\theta,m,z,1)),
    (z\mapsto0),n\right),\texttt{atTop},\mathcal N(0)\right).
 \tag{S2}\label{eq:shannon-H1-second}
\end{align*}
\end{lemma}

\begin{proof}
Apply the second conjunct of \cref{lem:positivity-intervals} with
\[
 p:=\theta.p,\qquad q:=\theta.q,\qquad
 h_p:=\theta.\texttt{p\_pos},\qquad
 h_q:=\theta.\texttt{q\_pos},\qquad
 h_{pq}:=\texttt{by simp [ShannonData.q]},
\]
and the compact set $K$.  Choose its witnesses
$\lambda_0:\R$, $h_{\lambda_0}:0<\lambda_0$, and $h_{\rm pos}$.
For each $n$, instantiate $h_{\rm pos}$ at the natural number $n+2$;
the proof $2\leq n+2$ is \texttt{Nat.le\_add\_left 2 n} after
normalization.  Using \cref{lem:shannon-log-scale-pos}, the three resulting
lower bounds unfold exactly to
\[
 \forall z\in K,\ \forall\lambda\in
  \operatorname{Set.Icc}(-\lambda_0,\lambda_0),\quad
 \operatorname{LinePositive}
  (\operatorname{shannonLineData}(\theta,n,z),\lambda).
\]
Fix this named $\lambda_0$ for the rest of the proof; it is the witness of
the existential in the lemma statement.

Use the global sequences
\[
 d_n:=\operatorname{shannonScale}(n),\qquad
 L_n:=\operatorname{shannonLogScale}(n),
\]
and, for all $n,z,j,\lambda$, abbreviate only the already declared term
\[
 w_{j,n,z}(\lambda):=
 \operatorname{shannonBlockMass}(\theta,n,z,j,\lambda),
 \qquad N_{j,n}:=\operatorname{shannonCount}(\theta,n,j).
\]
Because entries are equal within each block, the Shannon entropy has the exact
decomposition
\begin{equation}
 \left(\texttt{letI := shannonIndexNonempty theta n; }\quad
 H_1\!\left(\operatorname{lineProb}
   (\operatorname{shannonLineData}(\theta,n,z),\lambda)\right)\right)
 =-\sum_{j:\operatorname{Fin}(3)}w_{j,n,z}(\lambda)
      \log w_{j,n,z}(\lambda)
  +\sum_{j:\operatorname{Fin}(3)}w_{j,n,z}(\lambda)
      \log(N_{j,n}:\R),
 \label{eq:block-Shannon-exact}
\end{equation}
on the common positivity interval.  By \cref{lem:rounding} and the quotient rule, there is
$C_K:\R$ with $0\leq C_K$ such that the explicit sequence
\[
 \eta(n):=C_K\bigl(
  \operatorname{Real.rpow}(d_n,1-\theta.R)
  +\operatorname{Real.rpow}(d_n,1-\theta.c)\bigr)
\]
satisfies, for $j=0,1,2$,
\begin{align*}
 \sup_{\substack{z\in K\\\abs\lambda\leq\lambda_0}}
 \bigg(&\abs{\partial_\lambda^j
   \bigl(w_{0,n,z}(\lambda)-\theta.q-z.1\lambda/L_n\bigr)}
 +\abs{\partial_\lambda^j
   \bigl(w_{1,n,z}(\lambda)-\theta.p+z.1\lambda/L_n\bigr)}
 +\abs{\partial_\lambda^jw_{2,n,z}(\lambda)}\bigg)
 \leq\eta(n).
\end{align*}
Indeed, after division by $\operatorname{Real.rpow}(d_n,\theta.R)$, the
rounding errors in the first two block masses are at most
$\operatorname{Real.rpow}(d_n,-\theta.R)$ and
$\operatorname{Real.rpow}(d_n,1-\theta.R)$, while the entire third-block mass
and its first two $\lambda$-derivatives are at most a compact-set constant
times $\operatorname{Real.rpow}(d_n,1-\theta.c)
+\operatorname{Real.rpow}(d_n,2-\theta.R)$.  The denominator is uniformly bounded below by
$\min\{\theta.p,\theta.q\}/2$ on the common positivity interval.  Applying the quotient
rule zero, one, and two times gives the displayed $\eta(n)$ bound.  Since
$\theta.R>\theta.c+1$ and $\theta.c>1$, both $\eta(n)\to0$ and
$L_n\eta(n)\to0$.
Moreover, \cref{eq:rounding-log-estimate} gives the explicit bounds
\begin{align*}
 0\leq\log(N_{0,n}:\R)-\theta.RL_n
   &\leq\operatorname{Real.rpow}(d_n,-\theta.R),\\
 0\leq\log(N_{1,n}:\R)-(\theta.R-1)L_n
   &\leq\operatorname{Real.rpow}(d_n,1-\theta.R),\\
 0\leq\log(N_{2,n}:\R)-(\theta.R-1-\theta.c)L_n
   &\leq\operatorname{Real.rpow}(d_n,1+\theta.c-\theta.R).
\end{align*}
The product of each right-hand side with $L_n$ tends to zero. After enlarging
$C_K$, there is $n_0$ such that, for $n_0\leq n$ and $j=0,1,2$,
\[
 \sup_{\substack{z\in K\\\abs\lambda\leq\lambda_0}}
 \abs{\partial_\lambda^j\bigl(w_{2,n,z}(\lambda)
   \abs{\log w_{2,n,z}(\lambda)}\bigr)}
 \leq C_K\operatorname{Real.rpow}(d_n,1-\theta.c)L_n,
\]
whose right-hand side tends to zero because $\theta.c>1$. Substituting these
estimates into
\cref{eq:block-Shannon-exact} gives \cref{eq:shannon-expansion}.  The stated
error bounds hold after zero, one, or two $\lambda$-derivatives, uniformly on
the compact parameter set.  Taking the three suprema gives
\cref{eq:shannon-explicit-C2}. Differentiating at zero yields
\cref{eq:shannon-H1-first,eq:shannon-H1-second}.
\end{proof}

\begin{lemma}[Uniform control near the Shannon point]
\label{lem:shannon-neighbourhood}
Let $\theta:\operatorname{ShannonData}$,
$K:\operatorname{Set}(\R\times\R)$, and $\delta:\R$.  Assume
$h_K:\operatorname{IsCompact}(K)$,
$h_{K0}:K.\operatorname{Nonempty}$,
$h_{\delta0}:0<\delta$, and
$h_\delta:\delta<\min\{1/2,(\theta.c-1)/4\}$.  Then there exists $C:\R$ with
$0\leq C$ such that the following literal conjunction holds:
\begin{align*}
&\Bigl(\forall(n:\N)(z:\R\times\R),\ z\in K\to
 \forall\alpha\in\operatorname{Set.Icc}(1-\delta,1+\delta),\quad
 \abs{\operatorname{shannonKOne}
  (\theta,n,z,\operatorname{finiteParam}(\alpha))}\leq C\ 
 \wedge\
 \abs{\operatorname{shannonKTwo}
  (\theta,n,z,\operatorname{finiteParam}(\alpha))}\leq C\Bigr)\\
&\quad\wedge\Bigl(\forall\alpha:\R,\quad
 \alpha\in\operatorname{Set.Icc}(1-\delta,1+\delta)\to\alpha\neq1\to\\[-2pt]
&\hspace{12mm}
 \operatorname{Tendsto}
  \left(n\mapsto\operatorname{compactUniformError}
   \left(K,(m\mapsto z\mapsto\operatorname{shannonKOne}
     (\theta,m,z,\operatorname{finiteParam}(\alpha))),
    (z\mapsto0),n\right),\texttt{atTop},\mathcal N(0)\right)\ 
 \wedge\\[-2pt]
&\hspace{12mm}
 \operatorname{Tendsto}
  \left(n\mapsto\operatorname{compactUniformError}
   \left(K,(m\mapsto z\mapsto\operatorname{shannonKTwo}
     (\theta,m,z,\operatorname{finiteParam}(\alpha))),
    (z\mapsto0),n\right),\texttt{atTop},\mathcal N(0)\right)\Bigr)\\
&\quad\wedge
 \operatorname{Tendsto}
  \left(n\mapsto\operatorname{compactUniformError}
   \left(K,(m\mapsto z\mapsto\operatorname{shannonKOne}(\theta,m,z,1)),
    (z\mapsto z.1),n\right),\texttt{atTop},\mathcal N(0)\right)\\
&\quad\wedge
 \operatorname{Tendsto}
  \left(n\mapsto\operatorname{compactUniformError}
   \left(K,(m\mapsto z\mapsto\operatorname{shannonKTwo}(\theta,m,z,1)),
    (z\mapsto0),n\right),\texttt{atTop},\mathcal N(0)\right).
\end{align*}
\end{lemma}

\begin{proof}
For arbitrary $n:\N$ and $\zeta:\R\times\R$ with $\zeta\in K$, write
\[
 d_n:=\operatorname{shannonScale}(n),\quad
 L_n:=\operatorname{shannonLogScale}(n),\quad
 h:=\zeta.1,\quad v:=\zeta.2,
\]
\[
 \Pi_{j,n,\zeta}(\alpha):=
  \operatorname{shannonEscort}(\theta,n,\zeta,j,\alpha),\qquad
 \bar u_{n,\zeta}(\alpha):=
  \operatorname{shannonMean}(\theta,n,\zeta,\alpha),
\]
and define the exact escort after deletion of block $2$ by
\[
 \widetilde\Pi_{1,n,\zeta}(\alpha):=
 \frac{\Pi_{1,n,\zeta}(\alpha)}
      {\Pi_{0,n,\zeta}(\alpha)+\Pi_{1,n,\zeta}(\alpha)}.
\]
Define also the rounding ratio and deleted-block mean by
\begin{align*}
 \rho(n)&:=
  \frac{d_n\,(\operatorname{shannonCount}(\theta,n,1):\R)}
       {(\operatorname{shannonCount}(\theta,n,0):\R)}-1,\\
 \widetilde{\operatorname{shannonMean}}(n,\zeta,\alpha)&:=
 \frac{\Pi_{0,n,\zeta}(\alpha)
          \operatorname{shannonBlockVelocity}(\theta,n,\zeta,0)
       +\Pi_{1,n,\zeta}(\alpha)
          \operatorname{shannonBlockVelocity}(\theta,n,\zeta,1)}
      {\Pi_{0,n,\zeta}(\alpha)+\Pi_{1,n,\zeta}(\alpha)}.
\end{align*}
All constants below are uniform in the quantified $n,\zeta$.  Project the
first conjunct of \cref{lem:shannon-dedicated-escort} at
$\delta,h_{\delta0},h_\delta$ and choose its typed witnesses
\[
 C_{\rm esc}:\R,\qquad h_{C_{\rm esc}}:0\leq C_{\rm esc},\qquad h_{\rm esc}.
\]
Thus, throughout the chosen interval, the third-block escort mass is at most
\[
 C_{\rm esc}\operatorname{Real.rpow}
   (d_n,-(\theta.c-1)/2),
\]
and the absolute value of its first $\alpha$-derivative is at most
\[
 C_{\rm esc}L_n\operatorname{Real.rpow}
   (d_n,-(\theta.c-1)/2).
\]
It therefore suffices to analyse the first two blocks.

At $\lambda=0$, their escort-mass ratio is
\begin{equation}
 \frac{\Pi_{1,n,\zeta}(\alpha)}{\Pi_{0,n,\zeta}(\alpha)}
 =\operatorname{Real.rpow}(d_n,\alpha-1)
   \operatorname{Real.rpow}\left(\frac{\theta.p}{\theta.q},\alpha\right)
   (1+\rho(n)),
 \label{eq:shannon-logistic-ratio}
\end{equation}
The two applications of \cref{lem:rounding} give typed witnesses
\[
 C_{\rm round}:\R,\qquad h_{C_{\rm round}}:0\leq C_{\rm round},\qquad
 h_{\rm round}:\forall n:\N,\quad
 \abs{\rho(n)}\leq C_{\rm round}
   \operatorname{Real.rpow}(d_n,1-\theta.R).
\]
Put
$s=(\alpha-1)L_n$. The
second-block escort mass after deleting the third block is therefore a
logistic function of
$s+\alpha\log(\theta.p/\theta.q)+\log(1+\rho(n))$. Its derivative with respect to this
argument is at most $1/4$. The block velocities are
\[
 u_0=\frac{h}{\theta.qL_n},\qquad
 u_1=-\frac{h}{\theta.pL_n},\qquad u_2=v.
\]
The first two blocks give the exact identity
\begin{equation}
 \widetilde{\operatorname{shannonMean}}(n,\zeta,\alpha)
 =\frac{h}{\theta.p\theta.qL_n}
   \bigl(\theta.p-\widetilde\Pi_{1,n,\zeta}(\alpha)\bigr).
 \label{eq:shannon-escort-mean}
\end{equation}
By compactness, choose a typed nonnegative coordinate bound
\[
 C_{\rm coord}:\R,\qquad h_{C_{\rm coord}}:0\leq C_{\rm coord},\qquad
 h_{\rm coord}:\forall\zeta\in K,\quad
   \abs{\zeta.1}+\abs{\zeta.2}\leq C_{\rm coord}.
\]
Insert $h_{\rm esc}$, $h_{\rm round}$, and $h_{\rm coord}$ into
\cref{eq:shannon-logistic-ratio,eq:shannon-escort-mean}.  Positivity of
$\theta.p$, $\theta.q$, and $L_n$ and the logistic Lipschitz constant $1/4$
then give a typed witness
$C_{\rm mean}:\R$ with $h_{C_{\rm mean}}:0\leq C_{\rm mean}$ for all the
following mean estimates.  In particular, the third block changes the mean
by at most
$C_{\rm mean}\operatorname{Real.rpow}(d_n,-(\theta.c-1)/2)$ and its
$\alpha$-derivative by at most
$C_{\rm mean}L_n\operatorname{Real.rpow}(d_n,-(\theta.c-1)/2)$.  Consequently,
\[
 \abs{\bar u_{n,\zeta}(1)}\leq
 C_{\rm mean}\operatorname{Real.rpow}(d_n,1-\theta.R)/L_n
 +C_{\rm mean}\operatorname{Real.rpow}(d_n,-(\theta.c-1)/2).
\]
The Lipschitz bound for the
logistic function, followed by the mean-value theorem for the rounding and
third-block remainders, gives
\begin{equation}
 \abs{\bar u_{n,\zeta}(\alpha)-\bar u_{n,\zeta}(1)}
 \leq \frac{C_{\rm mean}\abs h}{L_n}
       \min\{1,\abs{\alpha-1}L_n\}
 +C_{\rm mean}\abs{\alpha-1}L_n
   \operatorname{Real.rpow}(d_n,-(\theta.c-1)/2).
 \label{eq:shannon-mean-bound}
\end{equation}
Multiplication by
$\abs{\omega(\operatorname{finiteParam}(\alpha))}$ gives a uniform bound for the first
derivative in \cref{eq:H-finite-at-zero}. It also shows pointwise convergence to zero
for every fixed $\alpha\neq1$. The value at $1$ is supplied by
\cref{lem:uniform-shannon-expansion}.

For the second derivative, use \cref{eq:H-finite-at-zero}.  The same finite
maximum construction gives a typed witness
$C_{\rm second}:\R$ with $h_{C_{\rm second}}:0\leq C_{\rm second}$ such that
the first two blocks contribute a variance bounded by
$C_{\rm second}h^2/L_n^2$, and the third block contributes at most
$C_{\rm second}\operatorname{Real.rpow}(d_n,-(\theta.c-1)/2)$.  The difference
$\bar u_{n,\zeta}(1)^2-\bar u_{n,\zeta}(\alpha)^2$ is controlled by
\cref{eq:shannon-mean-bound} and the bound
$\abs{\bar u_{n,\zeta}(\alpha)}+\abs{\bar u_{n,\zeta}(1)}
\leq C_{\rm second}/L_n+C_{\rm second}
\operatorname{Real.rpow}(d_n,-(\theta.c-1)/2)$. The singular factor
$\omega(\operatorname{finiteParam}(\alpha))$ therefore
cancels exactly as in the first-derivative bound. This proves uniform
boundedness and pointwise convergence to zero. The value at $1$ again follows
from \cref{lem:uniform-shannon-expansion}.

More explicitly, the preceding bounds and compactness of $K$ yield typed
bounds
\[
 C_1,C_2:\R,\qquad h_{C_1}:0\leq C_1,\qquad h_{C_2}:0\leq C_2,
\]
for the absolute values of
\texttt{shannonKOne} and \texttt{shannonKTwo}, respectively, on every
quantified $n,z,\alpha$.  Define the single final witness
\[
 C_{\rm final}:=\max C_1 C_2,
 \qquad h_{C_{\rm final}}:0\leq C_{\rm final}
   :=\texttt{le\_trans h\_C1 (le\_max\_left C1 C2)}.
\]
Use $C_{\rm final}$ for the existential $C$ in the lemma statement.  The two
pointwise estimates above prove its first conjunct, while the convergence
arguments and the two order-one limits prove the remaining three conjuncts.
\end{proof}

\begin{definition}[Signed atom and its removal]
\label{def:remove-signed-atom}
For $\mu:\operatorname{SignedMeasure}(\overline\R_+)$ and
$a:\overline\R_+$, define
\begin{align*}
 \operatorname{signedAtom}(\mu,a)&:=
  \operatorname{signedIntegral}
   \left(\mu,\{a\}.\operatorname{indicator}(x\mapsto1)\right),\\
 \operatorname{removeSignedAtom}(\mu,a)&:=
  \mu-\operatorname{signedAtom}(\mu,a)\mathbin{\bullet}
       \operatorname{signedLift}(\operatorname{diracFinite}(a)).
\end{align*}
Both terms have type \texttt{SignedMeasure Param}; the scalar in the second
line is real.
\end{definition}

\begin{lemma}[Removing a signed atom]
\label{lem:remove-signed-atom}
Let $\mu:\operatorname{SignedMeasure}(\overline\R_+)$ and
$a:\overline\R_+$.  Put $m:=\operatorname{signedAtom}(\mu,a)$ and
$\mu_0:=\operatorname{removeSignedAtom}(\mu,a)$.  Then the following single
conjunction holds:
\begin{align*}
 &\operatorname{signedTV}(\mu_0)(\{a\})=0\\
 &\quad\wedge\quad
 \forall f:\overline\R_+\to\R,\quad
 \operatorname{Measurable}(f)\to
 (\exists C:\R,0\leq C\ \wedge\ \forall x,\abs{f(x)}\leq C)\to\\[-2pt]
 &\hspace{25mm}
 \operatorname{signedIntegral}(\mu,f)
 =m f(a)+\operatorname{signedIntegral}(\mu_0,f).
\end{align*}
\end{lemma}

\begin{proof}
Evaluation on the singleton gives $\mu_0(\{a\})=0$.  A signed measure
restricted to a singleton is its mass times the Dirac measure, so its total
variation there is the absolute value of that mass; hence the total variation
is also zero.  The integral identity follows by signed-measure linearity.
\end{proof}

\begin{lemma}[Witness atoms and the positive lower truncation]
\label{lem:signed-witness-integral-bridge}
Let $\nu:\operatorname{FiniteMeasure}(\overline\R_+)$ and put
$\rho:=\operatorname{finiteMeasure}(\nu)$.  Then the following literal
conjunction holds:
\begin{align*}
&\Bigl(\forall a:\overline\R_+,\quad
 \operatorname{signedAtom}(\operatorname{positiveSigned}(\nu),a)
 =\operatorname{ENNReal.toReal}(\rho(\{a\}))\Bigr)\\
&\quad\wedge\Bigl(\forall a:\overline\R_+,\quad
 \operatorname{signedAtom}(\operatorname{negativeSigned}(\nu),a)
 =-\operatorname{ENNReal.toReal}(\rho(\{a\}))\Bigr)\\
&\quad\wedge\Bigl(\forall(r:\R),\quad0<r\to r<1\to
 \operatorname{lowerMoment}(\operatorname{positiveSigned}(\nu),r)
 =\operatorname{ENNReal.toReal}\!\left(
   \int^-_{\operatorname{Set.Iio}(\operatorname{finiteParam}(r))}
    \operatorname{omegaLower}\dd\rho\right)\Bigr).
\end{align*}
\end{lemma}

\begin{proof}
For the first two clauses, evaluate the signed lifts on the singleton and
use finiteness of $\rho(\{a\})$ to pass from the extended nonnegative value
to its real value; signed integration against the negative lift introduces
the displayed minus sign.  For $0<r<1$, $\omega$ is nonnegative on
\texttt{Set.Iio (finiteParam r)} and equals the real coercion of
\texttt{omegaLower} there.  The finite-set-separated kernel is bounded, so
the standard \texttt{ofReal}-lintegral/Bochner-integral agreement gives the
third clause.
\end{proof}

\begin{lemma}[Shannon-block logarithmic mass derivatives]
\label{lem:shannon-log-mass}
Let $\theta:\operatorname{ShannonData}$ and let
$K:\operatorname{Set}(\R\times\R)$ have proofs
$h_{K0}:K.\operatorname{Nonempty}$ and
$h_K:\operatorname{IsCompact}(K)$.  Then the
following literal conjunction holds:
\begin{align*}
&\operatorname{Tendsto}
 \left(n\mapsto\operatorname{compactUniformError}
  \left(K,(m\mapsto z\mapsto
    \operatorname{shannonLogMassKernel}(\theta,m,z,1)),
    (z\mapsto0),n\right),\texttt{atTop},\mathcal N(0)\right)\\
&\quad\wedge
 \operatorname{Tendsto}
 \left(n\mapsto\operatorname{compactUniformError}
  \left(K,(m\mapsto z\mapsto
    \operatorname{shannonLogMassKernel}(\theta,m,z,2)),
    (z\mapsto0),n\right),\texttt{atTop},\mathcal N(0)\right)\\
&\quad\wedge
 \exists C:\R,\ 0\leq C\ \wedge\
 \forall(n:\N)(z:\R\times\R),\ z\in K\to\\[-2pt]
&\hspace{9mm}
 \abs{\operatorname{shannonLogMassKernel}(\theta,n,z,1)}
 \leq C\left(
  \frac{\operatorname{Real.rpow}
    (\operatorname{blockScale}(n),1-\theta.R)}
       {\log(\operatorname{blockScale}(n))}
  +\operatorname{Real.rpow}
    (\operatorname{blockScale}(n),1-\theta.c)\right)\ 
 \wedge\\[-2pt]
&\hspace{9mm}
 \operatorname{shannonLogMassKernel}(\theta,n,z,2)
 =-\operatorname{shannonLogMassKernel}(\theta,n,z,1)^2.
\end{align*}
\end{lemma}

\begin{proof}
For $n:\N$ and $z:\R\times\R$, put
\[
 d_n:=\operatorname{shannonScale}(n),\quad
 L_n:=\operatorname{shannonLogScale}(n),\quad
 N_{j,n}:=\operatorname{shannonCount}(\theta,n,j),\quad
 S_{n,z}(\lambda):=\operatorname{shannonMass}(\theta,n,z,\lambda).
\]
Direct finite-sum differentiation gives the exact identity
\[
 \operatorname{deriv}(S_{n,z})(0)
 =\frac {z.1}{L_n}(N_{0,n}-d_nN_{1,n})
   +z.2\,d_n^2N_{2,n}.
\]
The rounding inequalities give
$\abs{N_{0,n}-d_nN_{1,n}}\leq1+d_n$,
$d_n^2N_{2,n}\leq
 \operatorname{Real.rpow}(d_n,\theta.R+1-\theta.c)+d_n^2$, and a constant
$C_\theta:\R$ with $0<C_\theta$ such that
$C_\theta\operatorname{Real.rpow}(d_n,\theta.R)\leq S_{n,z}(0)$ for every
$n$.  Compactness of
$K$ uniformly bounds $\abs{z.1}$ and $\abs{z.2}$.  Division gives the stated
bound, which tends to zero because
$\theta.R>\theta.c+1>2$ and $\theta.c>1$.  The mass is affine
in $\lambda$, so $\operatorname{secondDeriv}(S_{n,z},0)=0$ and two
differentiations of its logarithm give
the square identity.
\end{proof}

\begin{definition}[Shannon localization targets]
\label{def:shannon-localization-targets}
For $\mu:\operatorname{SignedMeasure}(\overline\R_+)$,
$\theta:\operatorname{ShannonData}$, and $z:\R\times\R$, define
\begin{align*}
 \operatorname{shannonTailMoment}(\mu,\theta)&:=
  \operatorname{upperMoment}(\mu,\theta.c),\\
 \operatorname{shannonLimitFirst}(\mu,\theta,z)&:=
  \operatorname{signedAtom}(\mu,1)z.1
  +\operatorname{shannonTailMoment}(\mu,\theta)z.2,\\
 \operatorname{shannonLimitSecond}(\mu,\theta,z)&:=
  -\operatorname{shannonTailMoment}(\mu,\theta)z.2^2.
\end{align*}
\end{definition}

\begin{proposition}[Shannon-point localisation]\label{prop:shannon-localisation}
Let $\mu:\operatorname{SignedMeasure}(\overline\R_+)$,
$\theta:\operatorname{ShannonData}$, and
$K:\operatorname{Set}(\R\times\R)$.  Assume
$h_{\mu,c}:\operatorname{signedTV}(\mu)
 (\{\operatorname{finiteParam}(\theta.c)\})=0$,
$h_K:\operatorname{IsCompact}(K)$, and
$h_{K0}:K.\operatorname{Nonempty}$.  Then the
following four convergence statements form one literal conjunction:
\begin{align}
&\operatorname{Tendsto}
 \left(n\mapsto\operatorname{compactUniformError}
  \left(K,(m\mapsto z\mapsto
   \operatorname{shannonLogKernel}(\mu,\theta,m,z,1)),
   (z\mapsto\operatorname{shannonLimitFirst}(\mu,\theta,z)),n\right),
  \texttt{atTop},\mathcal N(0)\right)
 \label{eq:shannon-first}\\
&\quad\wedge
 \operatorname{Tendsto}
 \left(n\mapsto\operatorname{compactUniformError}
  \left(K,(m\mapsto z\mapsto
   \operatorname{shannonLogKernel}(\mu,\theta,m,z,2)),
   (z\mapsto\operatorname{shannonLimitSecond}(\mu,\theta,z)),n\right),
  \texttt{atTop},\mathcal N(0)\right)
 \label{eq:shannon-second}\\
&\quad\wedge
 \operatorname{Tendsto}
 \left(n\mapsto\operatorname{compactUniformError}
  \left(K,(m\mapsto z\mapsto
   \operatorname{shannonGKernel}(\mu,\theta,m,z,1)),
   (z\mapsto\operatorname{shannonLimitFirst}(\mu,\theta,z)),n\right),
  \texttt{atTop},\mathcal N(0)\right)\\
&\quad\wedge
 \operatorname{Tendsto}
 \left(n\mapsto\operatorname{compactUniformError}
  \left(K,(m\mapsto z\mapsto
   \operatorname{shannonGKernel}(\mu,\theta,m,z,2)),
   (z\mapsto\operatorname{shannonLimitSecond}(\mu,\theta,z)),n\right),
  \texttt{atTop},\mathcal N(0)\right).
\end{align}
\end{proposition}

\begin{proof}
For every $n:\N$ put
$d_n:=\operatorname{shannonScale}(n)$ and
$L_n:=\operatorname{shannonLogScale}(n)$.  For every
$z:\R\times\R$, its two perturbation coordinates are literally $z.1$ and
$z.2$.  After factoring out
$\operatorname{Real.rpow}(d_n,\theta.R)$, the power-sum exponents are
\[
 0,\qquad\alpha-1,\qquad2\alpha-1-\theta.c.
\]
Thus the first block dominates below $1$, the second on
$(1,\theta.c)$, and the third above $\theta.c$. The first two blocks tie at
$1$, while the third has an exponent deficit $\theta.c-1$ there.

For every fixed $\alpha\neq1,\theta.c$, the first two block velocities tend
to zero.  The exact derivative formulas therefore give pointwise limits zero
below $\theta.c$,
and
\[
 \operatorname{shannonKOne}
   (\theta,n,z,\operatorname{finiteParam}(\alpha))
   \longrightarrow\omega(\operatorname{finiteParam}(\alpha))z.2,
 \qquad
 \operatorname{shannonKTwo}
   (\theta,n,z,\operatorname{finiteParam}(\alpha))
   \longrightarrow-\omega(\operatorname{finiteParam}(\alpha))z.2^2
\]
above $\theta.c$. At $\alpha=1$,
\cref{lem:uniform-shannon-expansion} gives the limits $z.1$ and $0$.

To pass through the measure integral, use
\cref{lem:shannon-neighbourhood} near $1$. Separate the atom from the
nonatomic remainder using \texttt{removeSignedAtom} and
\cref{lem:remove-signed-atom}. Dominated convergence
shows that the nonatomic contribution near $1$ vanishes, while the atom
contributes $\operatorname{signedAtom}(\mu,1)z.1$ to the first derivative and
zero to the second. A small
neighbourhood of $\theta.c$ is separated from $1$, so the exact formulas are
uniformly bounded there; the point $\theta.c$ itself is
$\operatorname{signedTV}(\mu)$-null by $h_{\mu,c}$. On the remaining
bounded regions, the second clause of
\cref{lem:shannon-dedicated-escort} applies, and its third clause controls the
large-order tail.  At every finite $n$ and every $z$, the first two
equalities of \cref{lem:shannon-kernel-integral-bridge} identify these two
signed parameter integrals literally with
\texttt{shannonGKernel} at orders one and two.  This proves the displayed
limits for the two $G_\mu$ kernels.

To obtain the stated local uniformity, take the supremum of the kernel error
over the arbitrary compact set $K$ of vectors $z=(z.1,z.2)$.  The estimates in
\cref{lem:uniform-shannon-expansion,lem:shannon-neighbourhood,lem:shannon-dedicated-escort}
give pointwise convergence of this supremum and
a common integrable dominator on each region.  More literally, put the
subtype $U_K:=\{z:\R\times\R\mathbin{//}z\in K\}$ and install
\texttt{letI : Nonempty U\_K} using $h_{K0}$ together with
\texttt{letI : CompactSpace U\_K :=
isCompact\_iff\_compactSpace.mp h\_K}.  The subtype inherits the required
topology, pseudometrizability, and second-countable topology.  Apply the exact
identity \cref{lem:subtype-uniform-error-bridge} to identify
\texttt{uniformIntegralErrorSigned} on $U_K$ with
\texttt{compactUniformError K}.  Apply
\cref{lem:fnd-compact-sup-measurable} to obtain measurability of the error
supremum, and then apply \cref{lem:fnd-uniform-dct-signed} to the signed
nonatomic remainder.

\Cref{lem:shannon-log-mass} shows that the first two derivatives of
$\log(\operatorname{shannonMass}(\theta,n,z,\lambda))$ tend uniformly to
zero.  The last two equalities of
\cref{lem:shannon-kernel-integral-bridge} therefore transfer the two
$G_\mu$ limits to \texttt{shannonLogKernel}; these are exactly the first two
conjuncts of the proposition.  No differentiation of a value-only bridge is
used here.
\end{proof}

\begin{definition}[Normalized scalar-line curvature]
\label{def:normalized-curvature}
For $\phi:\R\to\R$, define
\[
 \operatorname{normalizedCurvature}(\phi):=
 \frac{\operatorname{secondDeriv}(\phi,0)}{\phi(0)}.
\]
\end{definition}

\begin{lemma}[Logarithmic curvature identity]
\label{lem:log-curvature-identity}
Let $\phi:\R\to\R$ satisfy $0<\phi(0)$ and
$\operatorname{ContDiffAt}(\R,2,\phi,0)$.  Then the exact identity
\begin{equation}
 \operatorname{normalizedCurvature}(\phi)
 =\operatorname{secondDeriv}(\lambda\mapsto\log(\phi(\lambda)),0)
  +\operatorname{deriv}(\lambda\mapsto\log(\phi(\lambda)))(0)^2
 \label{eq:log-curvature}
\end{equation}
holds.
\end{lemma}

\begin{proof}
Continuity and $0<\phi(0)$ give a neighbourhood of zero on which
$0<\phi$.  On that neighbourhood,
$\phi=\exp(\log\circ\phi)$.  Differentiate this identity twice at zero and
divide by $\phi(0)$.
\end{proof}

\begin{lemma}[Logarithmic curvature limit obstruction]
\label{lem:log-curvature-obstruction}
Let $\phi_n:\N\to\R\to\R$ and $a,b:\R$.  Assume the single conjunction
\begin{align*}
 &\Bigl(\forall n,\quad0<\phi_n(n)(0)\ \wedge\
   \operatorname{ContDiffAt}(\R,2,\phi_n(n),0)\Bigr)\\
 &\quad\wedge\quad\operatorname{Tendsto}
  (n\mapsto\operatorname{deriv}
    (\lambda\mapsto\log(\phi_n(n)(\lambda)))(0),
   \texttt{atTop},\mathcal N(a))\\
 &\quad\wedge\quad\operatorname{Tendsto}
  (n\mapsto\operatorname{secondDeriv}
    (\lambda\mapsto\log(\phi_n(n)(\lambda)),0),
   \texttt{atTop},\mathcal N(b)).
\end{align*}
Then the following literal conjunction holds:
\begin{align*}
 &\operatorname{Tendsto}
  (n\mapsto\operatorname{normalizedCurvature}(\phi_n(n)),
   \texttt{atTop},\mathcal N(b+a^2))\\
 &\quad\wedge\quad
 \bigl(0<b+a^2\to
   \operatorname{Filter.Eventually}
    (n\mapsto 0<\operatorname{normalizedCurvature}(\phi_n(n)))
    (\texttt{Filter.atTop})\bigr)\\
 &\quad\wedge\quad
 \bigl(b+a^2<0\to
   \operatorname{Filter.Eventually}
    (n\mapsto\operatorname{normalizedCurvature}(\phi_n(n))<0)
    (\texttt{Filter.atTop})\bigr).
\end{align*}
Together with \cref{lem:cone-affine-line-bridge,lem:fnd-line-curvature},
the last two conjuncts are the precise concavity and convexity obstructions.
\end{lemma}

\begin{proof}
Apply \cref{lem:log-curvature-identity} pointwise to $\phi_n(n)$.  The two
assumed limits, continuity of squaring,
and continuity of addition give the first conjunct.  Convergence to a
strictly positive or strictly negative real gives the corresponding
eventual strict sign.
\end{proof}

The tropical argument requires exact stationarity at the finite dimension
where quasi-convexity is invoked.  Package the repeated hypotheses once:

\begin{definition}[Stationarity convergence package]
\label{def:stationarity-package}
Let $V:\texttt{Type u}$ carry
$[\operatorname{NormedAddCommGroup}(V)]
 [\operatorname{NormedSpace}(\R,V)]$.
For $a_n,b_n:\N\to V\to\R$ and $a,b:V\to\R$, define
$\operatorname{StationarityPackage}(a_n,a,b_n,b):\operatorname{Prop}$ to
be the following literal conjunction:
\begin{align*}
&\Bigl(\forall(n:\N)(x,y:V),\quad
  a_n(n)(x+y)=a_n(n)(x)+a_n(n)(y)\Bigr)\\
&\quad\wedge\Bigl(\forall(n:\N)(c:\R)(x:V),\quad
  a_n(n)(c\mathbin{\bullet}x)=c\,a_n(n)(x)\Bigr)\\
&\quad\wedge\Bigl(\forall x,y:V,\quad a(x+y)=a(x)+a(y)\Bigr)\\
&\quad\wedge\Bigl(\forall(c:\R)(x:V),\quad
  a(c\mathbin{\bullet}x)=c\,a(x)\Bigr)\\
&\quad\wedge\Bigl(\forall n,\operatorname{Continuous}(a_n(n))\Bigr)
 \quad\wedge\quad\operatorname{Continuous}(a)\\
&\quad\wedge\Bigl(\forall n,\operatorname{Continuous}(b_n(n))\Bigr)
 \quad\wedge\quad\operatorname{Continuous}(b)\\
&\quad\wedge\Bigl(\forall K:\operatorname{Set}(V),\quad
 \operatorname{IsCompact}(K)\to K.\operatorname{Nonempty}\to\\[-2pt]
&\hspace{10mm}
 \operatorname{Tendsto}
  (n\mapsto\operatorname{compactUniformError}
    (K,(m\mapsto x\mapsto a_n(m)(x)),a,n),
   \texttt{atTop},\mathcal N(0))\ \wedge\\[-2pt]
&\hspace{10mm}
 \operatorname{Tendsto}
  (n\mapsto\operatorname{compactUniformError}
    (K,(m\mapsto x\mapsto b_n(m)(x)),b,n),
   \texttt{atTop},\mathcal N(0))\Bigr).
\end{align*}
\end{definition}

\begin{lemma}[Exact stationarity correction]\label{lem:stationarity-correction}
Let $V:\texttt{Type u}$ carry
$[\operatorname{NormedAddCommGroup}(V)]
 [\operatorname{NormedSpace}(\R,V)]
 [\operatorname{FiniteDimensional}(\R,V)]$.
Let
$a_n:\N\to V\to\R$, $a:V\to\R$,
$b_n:\N\to V\to\R$, and $b:V\to\R$.  Assume
$h_{\rm stat}:\operatorname{StationarityPackage}(a_n,a,b_n,b)$.
Let $u,e:V$ satisfy $a(u)=0$, $b(u)<0$, and $a(e)\neq0$.  Then there exists
$u_n:\N\to V$ such that the following literal conjunction holds:
\begin{align*}
 &\operatorname{Tendsto}(u_n,\texttt{atTop},\mathcal N(u))\\
 &\quad\wedge\quad
 \operatorname{Filter.Eventually}
  (n\mapsto a_n(n)(u_n(n))=0)(\texttt{Filter.atTop})\\
 &\quad\wedge\quad
 \operatorname{Tendsto}(n\mapsto b_n(n)(u_n(n)),
    \texttt{atTop},\mathcal N(b(u)))\\
 &\quad\wedge\quad
 \operatorname{Filter.Eventually}
  (n\mapsto b_n(n)(u_n(n))<0)(\texttt{Filter.atTop}).
\end{align*}
\end{lemma}

\begin{proof}
Local uniform convergence at the singletons $\{e\}$ and $\{u\}$ gives
$a_n(n)(e)\to a(e)\neq0$ and $a_n(n)(u)\to a(u)=0$.  Hence
$a_n(n)(e)\neq0$ eventually.  Define the total sequence
\[
 u_n(n):=
 \begin{cases}
  u,&h:a_n(n)(e)=0,\\
  u-\dfrac{a_n(n)(u)}{a_n(n)(e)}\mathbin{\bullet}e,&
    h:a_n(n)(e)\neq0.
 \end{cases}
\]
Its literal Lean body is
\texttt{if h : a\_n n e = 0 then u else
u - (a\_n n u / a\_n n e) • e}.  On the eventual nonvanishing set the
second branch is selected, and the quotient tends to zero.  Hence
$u_n\to u$, and linearity gives exact eventual
stationarity.  A finite-dimensional closed ball containing this now-known
convergent sequence is compact.  The second compact-uniform-convergence
component of $h_{\rm stat}$, instantiated at that nonempty compact closed
ball, gives
$b_n(n)(u_n(n))\to b(u)<0$; eventual negativity is immediate.
\end{proof}

\begin{lemma}[Stationary curvature of a quasi-convex function]
\label{lem:quasiconvex-second}
Let $g:\R\to\R$ and $\eps:\R$ satisfy $0<\eps$,
$\operatorname{ContDiffOn}(\R,2,g,\operatorname{Set.Ioo}(-\eps,\eps))$,
and
$\operatorname{ScalarQCvxOn}(\operatorname{Set.Ioo}(-\eps,\eps),g)$.
Then
\[
 \operatorname{deriv}(g)(0)=0\quad\Longrightarrow\quad
 0\leq\operatorname{secondDeriv}(g,0).
\]
\end{lemma}

\begin{proof}
If $g''(0)<0$, Taylor's theorem gives
$g(0)>\max\{g(-\lambda),g(\lambda)\}$ for every sufficiently small
$\lambda>0$. This contradicts quasi-convexity at the midpoint.
\end{proof}

\begin{lemma}[Typed curvature contradiction from logarithmic kernels]
\label{lem:typed-log-curvature-contradiction}
Let $\phi:\N\to\R\to\R$ and $a,b:\R$.  Assume
\begin{align*}
&\forall n:\N,\quad
 \exists\eps_n:\R,\quad 0<\eps_n\ \wedge\
 0<\phi(n)(0)\ \wedge\\[-2pt]
&\hspace{24mm}
 \operatorname{ContDiffOn}
  (\R,2,\phi(n),\operatorname{Set.Ioo}(-\eps_n,\eps_n)),
 \tag{TC0}\label{eq:typed-curvature-smooth}\\
&\operatorname{Tendsto}
 (n\mapsto\operatorname{deriv}
   (\lambda\mapsto\log(\phi(n)(\lambda)))(0),
  \texttt{atTop},\mathcal N(a)),\tag{TC1}\label{eq:typed-curvature-first}\\
&\operatorname{Tendsto}
 (n\mapsto\operatorname{secondDeriv}
   (\lambda\mapsto\log(\phi(n)(\lambda)),0),
  \texttt{atTop},\mathcal N(b)).\tag{TC2}\label{eq:typed-curvature-second}
\end{align*}
Then the following two implications form one literal conjunction:
\begin{align*}
&\Bigl(0<b+a^2\to
  (\forall n:\N,\ \exists\eps_n:\R,\ 0<\eps_n\ \wedge
   \operatorname{ConcaveOn}
    (\R,\operatorname{Set.Ioo}(-\eps_n,\eps_n),\phi(n)))
  \to\operatorname{False}\Bigr)\\
&\quad\wedge
\Bigl(b+a^2<0\to
  (\forall n:\N,\ \exists\eps_n:\R,\ 0<\eps_n\ \wedge
   \operatorname{ConvexOn}
    (\R,\operatorname{Set.Ioo}(-\eps_n,\eps_n),\phi(n)))
  \to\operatorname{False}\Bigr).
\end{align*}
\end{lemma}

\begin{proof}
Apply \cref{lem:log-curvature-obstruction} with the named argument
\texttt{phi\_n := phi}.  Clause \cref{eq:typed-curvature-smooth} gives the
required \texttt{ContDiffAt} statement because $0$ lies in every displayed
open interval.  In the first implication choose $n$ in the eventual set on
which the normalized curvature is positive.  Destructure
\cref{eq:typed-curvature-smooth} at $n$ as
$(\eps_{\rm sm},h_{\rm sm},h_{\phi0},h_{C^2})$ and the asserted concavity
package at $n$ as $(\eps_{\rm sh},h_{\rm sh},h_{\rm conc})$.  Put
\[
 \delta:=\min\{\eps_{\rm sm},\eps_{\rm sh}\}.
\]
The named proofs $h_{\rm sm}$ and $h_{\rm sh}$ give $0<\delta$.
Restrict $h_{C^2}$ and $h_{\rm conc}$ to
$\operatorname{Set.Ioo}(-\delta,\delta)$ by \texttt{ContDiffOn.mono} and
\texttt{ConcaveOn.mono}.  Since $h_{\phi0}$ is positive, the positive
normalized curvature makes
$0<\operatorname{secondDeriv}(\phi(n),0)$, while
\cref{lem:fnd-line-curvature} on this common interval gives the opposite
weak inequality.  For the second implication use the same construction
with the convexity radius and the eventual negative-curvature set.  This
lemma deliberately allows both radii to depend on $n$.
\end{proof}

\begin{lemma}[Corrected stationary-line obstruction]
\label{lem:corrected-stationary-line-obstruction}
Let $V:\texttt{Type u}$ carry
$[\operatorname{NormedAddCommGroup}(V)]
 [\operatorname{NormedSpace}(\R,V)]
 [\operatorname{FiniteDimensional}(\R,V)]$.
Let $\gamma:\N\to V\to\R\to\R$,
$a_n,b_n:\N\to V\to\R$, and $a,b:V\to\R$.
Assume
$h_{\rm stat}:\operatorname{StationarityPackage}(a_n,a,b_n,b)$ and the
exact line package
\begin{align*}
\forall(n:\N)(v:V),\quad\exists\eps:\R,\quad
 &0<\eps\ \wedge
 \operatorname{ContDiffOn}
  (\R,2,\gamma(n,v),\operatorname{Set.Ioo}(-\eps,\eps))\\[-2pt]
 &{}\wedge
 \operatorname{ScalarQCvxOn}
  (\operatorname{Set.Ioo}(-\eps,\eps),\gamma(n,v))\\[-2pt]
 &{}\wedge
 \operatorname{deriv}(\gamma(n,v))(0)=a_n(n)(v)\\[-2pt]
 &{}\wedge
 \operatorname{secondDeriv}(\gamma(n,v),0)=b_n(n)(v).
 \tag{CSL}\label{eq:corrected-stationary-line-package}
\end{align*}
If there are $v,e:V$ with
$a(v)=0$, $b(v)<0$, and $a(e)\neq0$, then
$\operatorname{False}$.
\end{lemma}

\begin{proof}
Apply \cref{lem:stationarity-correction} and choose the resulting sequence
$v_n$.  Eventually its exact first derivative is zero and its exact second
derivative is negative.  For one index in that eventual set, instantiate
\cref{eq:corrected-stationary-line-package} at $(n,v_n(n))$ and apply
\cref{lem:quasiconvex-second}.  The two derivative identities give the
contradiction.  Thus every use of stationarity below has a named scalar curve,
a positivity interval, a smoothness proof, and a scalar quasi-convexity proof.
\end{proof}

\subsection{Necessity in the positive temperate branch}

We begin with the ordinary-order case. Here the representing measure is
positive, and monotonicity requires concavity of the corresponding column
functional.

\begin{proposition}[Positive-measure obstruction]
\label{prop:positive-necessity}
Let $\nu:\operatorname{FiniteMeasure}(\overline\R_+)$.  If
$\texttt{PosPhiConcave.\{u\}}(\nu)$, then the following literal
conjunction holds:
\begin{equation}
 \operatorname{suppMeasure}(\operatorname{finiteMeasure}(\nu))
   \subseteq\operatorname{Set.Icc}(0,1)
 \quad\wedge\quad
 \operatorname{finiteMeasure}(\nu)(\{1\})=0
 \quad\wedge\quad
 \operatorname{MLower}(\operatorname{finiteMeasure}(\nu))
   \leq\operatorname{ENNReal.ofReal}(1).
 \label{eq:positive-mu-condition}
\end{equation}
The last quantity is an extended nonnegative integral and is therefore
well-defined before its finiteness is established.
\end{proposition}

\begin{proof}
Put
\[
 \rho:=\operatorname{finiteMeasure}(\nu),\qquad
 \mu_s:=\operatorname{positiveSigned}(\nu).
\]
The witness identities in
\cref{lem:signed-witness-bridge,lem:signed-witness-integral-bridge} give
\[
 \operatorname{signedTV}(\mu_s)=\rho,\qquad
 \operatorname{signedAtom}(\mu_s,a)
 =\operatorname{ENNReal.toReal}(\rho(\{a\}))
 \quad(a:\overline\R_+).
\]
Suppose $h_1:\rho(\{1\})\neq0$ and put
$m:=\operatorname{ENNReal.toReal}(\rho(\{1\}))$.
Finiteness of $\rho$ and $h_1$ give $h_m:0<m$.
Use \cref{lem:null-thresholds} to choose $c:\R$ and proofs
\[
 h_c:1<c,\qquad
 h_{\rho,c}:\rho(\{\operatorname{finiteParam}(c)\})=0,
\]
then choose $R:\R$ and $h_R:c+1<R$.  Define
\[
 \theta:=
 \left\{c:=c,\ R:=R,\ p:=\tfrac12,\
 \texttt{c\_gt\_one}:=h_c,\ \texttt{R\_gt}:=h_R,\
 \texttt{p\_pos}:=\texttt{by norm\_num},\
 \texttt{p\_lt\_one}:=\texttt{by norm\_num}\right\}
 \in\operatorname{ShannonData},
\]
and set
\[
 z_S:=(1,0),\qquad K_S:=\{z_S\},\qquad
 \phi_S(n,\lambda):=
  \operatorname{shannonPhiCurve}(\mu_s,\theta,n,z_S,\lambda).
\]
The singleton $K_S$ is nonempty and compact.  Apply
\cref{prop:shannon-localisation} with the null-atom proof obtained by
rewriting $h_{\rho,c}$ through \cref{lem:signed-witness-bridge}, and evaluate
its first two conjuncts at $z_S$ using
\cref{lem:compact-uniform-singleton}.  Unfolding
\cref{def:shannon-localization-targets} yields
\begin{align*}
 &\operatorname{Tendsto}
  (n\mapsto\operatorname{shannonLogKernel}
    (\mu_s,\theta,n,z_S,1),\texttt{atTop},\mathcal N(m)),\\
 &\operatorname{Tendsto}
  (n\mapsto\operatorname{shannonLogKernel}
    (\mu_s,\theta,n,z_S,2),\texttt{atTop},\mathcal N(0)).
\end{align*}
For every $n$, first install the local instance
\texttt{letI := shannonIndexNonempty theta n}.  Then
\cref{lem:shannon-curve-kernel-bridge} supplies one
$\eps_n>0$, positivity, \texttt{ContDiffOn}, and both logarithmic derivative
identities for $\phi_S(n)$.  The global \texttt{PosPhiConcave} hypothesis
and \cref{lem:cone-affine-line-bridge} supply
\[
 \operatorname{ConcaveOn}
  (\R,\operatorname{Set.Ioo}(-\eps_n,\eps_n),\phi_S(n))
\]
on that same interval.  Now
\cref{lem:typed-log-curvature-contradiction}, with $a=m$ and $b=0$,
contradicts $0<m^2$.  Consequently $\rho(\{1\})=0$.

Suppose next that $\operatorname{suppMeasure}(\rho)$ meets
$\operatorname{Set.Ioi}(1)$. By
\cref{lem:null-thresholds}, there is $r:\R$ with $1<r$ such that
$\rho(\{\operatorname{finiteParam}(r)\})=0$ and
$0<\rho(\operatorname{Set.Ioi}(\operatorname{finiteParam}(r)))$. Define
\[
 A_r:=\operatorname{upperMoment}(\mu_s,r).
\]
The function $\omega$ is bounded and strictly negative on
$\operatorname{Set.Ioi}(\operatorname{finiteParam}(r))$.  Since the tail
has positive $\rho$-mass,
\cref{lem:support-strict-integral} applied to $-\omega$ shows that $A_r$ is
a finite strictly negative number.  Bind
$h_{r0}:0<r$ and $h_{r1}:r\neq1$, choose $R_r:\R$ with
$h_{R_r}:r<R_r$, and put
\begin{align*}
 \texttt{B\_r\_data}&:=\operatorname{twoBlockData}(r,R_r,h_{r0},h_{R_r}),\\
 u_U&:=\operatorname{Fin.cases}(0,(\_\mapsto1)):
       \operatorname{Fin}(2)\to\R,\\
 K_U&:=\{u_U\},\\
 \phi_U(n,\lambda)&:=
  \operatorname{blockPhiCurve}(\mu_s,\texttt{B\_r\_data},n,u_U,\lambda).
\end{align*}
Apply \cref{prop:two-block-localisation} to the nonempty compact singleton
$K_U$.  Its upper logarithmic targets simplify at $u_U$ to
$A_r$ and $-A_r$.  Thus \cref{lem:compact-uniform-singleton} and
\cref{lem:block-curve-kernel-bridge}, instantiated with $j_\infty=1$ and
$h_\infty:=\operatorname{twoBlockTopUnique}(r,R_r,h_{r0},h_{R_r})$, give
the two log-derivative limits for
$\phi_U$ and their exact derivative identities.  For each $n$, install
\texttt{letI := blockCarrierNonempty B\_r\_data n} before applying the
global shape predicate and the cone-affine bridge.  The same bridge gives
positivity and smoothness; \texttt{PosPhiConcave} together with
\cref{lem:cone-affine-line-bridge} gives concavity on its interval.  Apply
\cref{lem:typed-log-curvature-contradiction} with
$a=A_r$, $b=-A_r$.  Since
$0<A_r^2-A_r=A_r(A_r-1)$, this is a contradiction.  Therefore
\[
 \operatorname{suppMeasure}(\rho)
 \subseteq\operatorname{Set.Icc}(0,1).
\]

Finally, let $r:\R$ have named proofs $h_{r0}:0<r$ and
$h_{r<1}:r<1$, and
$\rho(\{\operatorname{finiteParam}(r)\})=0$, and put
\[
 \ell_r:=\int^-_{\operatorname{Set.Iio}(\operatorname{finiteParam}(r))}
   \omega_<\dd\rho,
 \qquad B_r:=\operatorname{ENNReal.toReal}(\ell_r).
\]
The extended integral $\ell_r$ is finite because $\omega$ is bounded on
$\operatorname{Set.Iic}(\operatorname{finiteParam}(r))$, and $0\leq B_r$.
The witness-integral identities from
\cref{lem:signed-witness-integral-bridge} give the typed equality
\[
 B_r=\operatorname{lowerMoment}(\mu_s,r).
\]
Bind $h_{r\neq1}:r\neq1$, choose $R_L$ with $h_{R_L}:r<R_L$, and put
\[
 B_L:=\operatorname{twoBlockData}(r,R_L,h_{r0},h_{R_L}),\qquad
 u_L:=\operatorname{Fin.cases}(1,(\_\mapsto0)):
       \operatorname{Fin}(2)\to\R,\qquad K_L:=\{u_L\},
\]
\[
 \phi_L(n,\lambda):=
  \operatorname{blockPhiCurve}(\mu_s,B_L,n,u_L,\lambda).
\]
The lower logarithmic pair in \cref{prop:two-block-localisation}, evaluated
at $u_L$ by \cref{lem:compact-uniform-singleton}, is exactly $(B_r,-B_r)$.
If $1<B_r$, then for each $n$ install
\texttt{letI := blockCarrierNonempty B\_L n}.  The block curve bridge
supplies its logarithmic derivative
identities and smoothness, instantiated with $j_\infty=1$ and
$h_\infty:=\operatorname{twoBlockTopUnique}
 (r,R_L,h_{r0},h_{R_L})$, while the cone-affine bridge supplies concavity.
\Cref{lem:typed-log-curvature-contradiction} then contradicts
$0<B_r(B_r-1)$.  Hence $B_r\leq1$.  Since $\ell_r<\top$, the
order-embedding property of
\texttt{ENNReal.toReal} gives
$\ell_r\leq\operatorname{ENNReal.ofReal}(1)$.

Choose, using \cref{lem:null-thresholds}, an increasing sequence of
$\rho$-null points $r_n\uparrow1$. The nonnegative functions
$\omega_<\mathbf 1_{\operatorname{Set.Iio}
 (\operatorname{finiteParam}(r_n))}$ increase pointwise to $\omega_<$. The
monotone convergence theorem in $[0,+\infty]$ for the positive measure
$\rho$ therefore yields
\[
 \operatorname{MLower}(\rho)
 =\lim_{n\to\infty}\ell_{r_n}
 \leq\operatorname{ENNReal.ofReal}(1).
\]
This proves \cref{eq:positive-mu-condition}.
\end{proof}

\begin{corollary}[Positive temperate necessity for probability measures]
\label{cor:positive-temperate-probability-necessity}
At universe $u$, let
$\tau:\operatorname{ProbabilityMeasure}(\overline\R_+)$ and $t:\R$.
For every proof $h_t:0<t$,
\[
 \operatorname{CMMonotone}
  (\operatorname{HTemp}(t,\operatorname{ne\_of\_gt}(h_t),\tau))
 \Longrightarrow \operatorname{PosAdm}(t,\tau).
\]
\end{corollary}

\begin{proof}
Assume the displayed monotonicity and bind
\[
 h_{t0}:0\leq t:=\operatorname{le\_of\_lt}(h_t),\qquad
 \nu_t:=\operatorname{finiteScale}
  (t,h_{t0},\tau.\operatorname{toFiniteMeasure}).
\]
The first conjunct of
\cref{cor:global-temperate-shape-reduction} supplies
$\operatorname{ConcaveCone}_I(\Phi_{t,\tau})$ for every finite nonempty
$I:\texttt{Type u}$.  The positive-scaling identity in
\cref{lem:signed-witness-bridge}, followed by the first conjunct of
\cref{lem:signed-scalar-column-bridge}, rewrites this family of statements
as $\texttt{PosPhiConcave.\{u\}}(\nu_t)$.  Apply
\cref{prop:positive-necessity}.

Unfolding \texttt{finiteScale}, the scalar is
$\operatorname{ENNReal.ofReal}(t)$, which is finite and nonzero because
$h_t:0<t$.  Hence \cref{lem:scalar-measure-support} and the elementary
atom formula for a nonzero scalar multiple transfer the support inclusion
and the null atom at $1$ from $\nu_t$ to $\tau$.  The first scaling identity
of \cref{lem:fnd-moment-algebra} gives
\[
 \operatorname{MLower}(\operatorname{finiteMeasure}(\nu_t))
 =\operatorname{ENNReal.ofReal}(t)\,
  \operatorname{MLower}(\operatorname{probMeasure}(\tau)).
\]
Cancel the positive finite scalar in the bound supplied by
\cref{prop:positive-necessity}, using
$\operatorname{ENNReal.ofReal}(1/t)$ as its reciprocal.  The result is
\[
 \operatorname{MLower}(\operatorname{probMeasure}(\tau))
 \leq\operatorname{ENNReal.ofReal}(1/t).
\]
Together with $h_t$ and the transferred support and atom clauses, this is
exactly $\operatorname{PosAdm}(t,\tau)$.  Notice that the condition still
permits $1$ to be a limit point of the support.
\end{proof}

\subsection{Necessity in the negative temperate branch}

We now bind a positive finite-measure witness $\nu$ and use the typed signed
measure $\operatorname{negativeSigned}(\nu)$.  Monotonicity requires the
global predicate $\texttt{NegPhiConvex.\{u\}}(\nu)$.

\begin{proposition}[A Shannon atom excludes upper support]
\label{prop:negative-shannon-obstruction}
Let $\nu:\operatorname{FiniteMeasure}(\overline\R_+)$ satisfy
$\operatorname{finiteMeasure}(\nu)(\{1\})\neq0$.  If
$\texttt{NegPhiConvex.\{u\}}(\nu)$, then
\[
 \operatorname{suppMeasure}(\operatorname{finiteMeasure}(\nu))
 \subseteq\operatorname{Set.Icc}(0,1).
\]
\end{proposition}

\begin{proof}
Put $\rho:=\operatorname{finiteMeasure}(\nu)$ and
$\mu:=\operatorname{negativeSigned}(\nu)$.
Assume that there is support above $1$. By \cref{lem:null-thresholds},
choose $c:\R$ with named proofs
\[
 h_c:1<c,\qquad
 h_{\rho,c}:\rho(\{\operatorname{finiteParam}(c)\})=0,\qquad
 h_{\rm tail}:0<\rho(\operatorname{Set.Ioi}(\operatorname{finiteParam}(c))).
\]
Set $A_c:=\operatorname{upperMoment}(\mu,c)$.  Obtain the named proof
$h_{A_c}:0<A_c$ from \cref{lem:support-strict-integral}, applied to
the positive measure $\rho$ and the positive function $-\omega$ on a
positive-mass neighbourhood inside the tail.
Set $m:=\operatorname{signedAtom}(\mu,1)$.  By
\cref{lem:signed-witness-integral-bridge}, obtain the separate named proofs
\[
 h_m^{\rm eq}:m=-\operatorname{ENNReal.toReal}(\rho(\{1\})),
 \qquad h_m^{\rm neg}:m<0.
\]
Choose $R$ and $h_R:c+1<R$, and define
\[
 \theta:=
 \left\{c:=c,\ R:=R,\ p:=\tfrac12,\
 \texttt{c\_gt\_one}:=h_c,\ \texttt{R\_gt}:=h_R,\
 \texttt{p\_pos}:=\texttt{by norm\_num},\
 \texttt{p\_lt\_one}:=\texttt{by norm\_num}\right\}
 \in\operatorname{ShannonData},
\]
\[
 z_N:=\left(-\frac{A_c}{m},1\right),\qquad
 K_N:=\{z_N\},\qquad
 \phi_N(n,\lambda):=
  \operatorname{shannonPhiCurve}(\mu,\theta,n,z_N,\lambda).
\]
Apply \cref{prop:shannon-localisation} to $K_N$; the required null proof is
$h_{\rho,c}$ rewritten through \cref{lem:signed-witness-bridge}.  Evaluate
the first two conjuncts with \cref{lem:compact-uniform-singleton}.  Since
$m\neq0$, the two targets simplify exactly to
\[
 \operatorname{shannonLimitFirst}(\mu,\theta,z_N)=0,\qquad
 \operatorname{shannonLimitSecond}(\mu,\theta,z_N)=-A_c.
\]
For every $n$, first install
\texttt{letI := shannonIndexNonempty theta n}.  Then
\cref{lem:shannon-curve-kernel-bridge} supplies positivity,
smoothness, and the logarithmic derivative identities for $\phi_N(n)$.
The global \texttt{NegPhiConvex} hypothesis and
\cref{lem:cone-affine-line-bridge} supply convexity on the same interval.
Apply \cref{lem:typed-log-curvature-contradiction} with $a=0$ and
$b=-A_c$.  Since $-A_c<0$, this is a contradiction.
\end{proof}

When the witness has no atom at $1$, its lower weighted moment need not
initially be known to be finite. We therefore work with truncated moments.

\begin{definition}[Positive lower and upper truncations]
\label{def:negative-truncations}
For $\nu:\operatorname{FiniteMeasure}(\overline\R_+)$ and $r:\R$, define
\begin{align*}
 \operatorname{lowerTrunc}(\nu,r)&:=
  \int_{\operatorname{Set.Iio}(\operatorname{finiteParam}(r))}
    \omega\dd\operatorname{finiteMeasure}(\nu),\\
 \operatorname{upperTrunc}(\nu,r)&:=
  \int_{\operatorname{Set.Ioi}(\operatorname{finiteParam}(r))}
    (-\omega)\dd\operatorname{finiteMeasure}(\nu).
\end{align*}
Both are ordinary real integrals over sets on the \texttt{Param} carrier.
\end{definition}

\begin{lemma}[Negative-witness truncation identities]
\label{lem:negative-truncation-bridge}
Let $\nu:\operatorname{FiniteMeasure}(\overline\R_+)$ and put
$\mu:=\operatorname{negativeSigned}(\nu)$.  Then
\begin{align*}
&\Bigl(\forall(a:\R),\quad0<a\to a<1\to
 \operatorname{lowerMoment}(\mu,a)
 =-\operatorname{lowerTrunc}(\nu,a)\Bigr)\\
&\quad\wedge\Bigl(\forall(b:\R),\quad1<b\to
 \operatorname{upperMoment}(\mu,b)
 =\operatorname{upperTrunc}(\nu,b)\Bigr).
\end{align*}
\end{lemma}

\begin{proof}
On the lower truncated set $\omega\geq0$ and is bounded; on the upper
truncated set $\omega<0$ and $-\omega$ is bounded.  Expand
\texttt{negativeSigned} and use signed-integral negation in the first case
and negation of the integrand in the second.  The separation from $1$
supplies the required integrability.
\end{proof}

\begin{proposition}[Truncated-moment obstruction]
\label{prop:truncated-moment}
Let $\nu:\operatorname{FiniteMeasure}(\overline\R_+)$ satisfy
$\operatorname{finiteMeasure}(\nu)(\{1\})=0$ and
$\texttt{NegPhiConvex.\{u\}}(\nu)$.  Let $a,b:\R$ satisfy
$0<a$, $a<1$, $1<b$,
\begin{align*}
 &\operatorname{finiteMeasure}(\nu)(\{\operatorname{finiteParam}(a)\})=0,\\
 &\operatorname{finiteMeasure}(\nu)(\{\operatorname{finiteParam}(b)\})=0,\\
 &0<\operatorname{finiteMeasure}(\nu)
   (\operatorname{Set.Ioi}(\operatorname{finiteParam}(b))).
\end{align*}
Then the following single proposition holds:
\begin{align}
 &\left(1\leq\operatorname{upperTrunc}(\nu,b)
       -\operatorname{lowerTrunc}(\nu,a)\right)
 \label{eq:truncated-condition}\\
 &\quad\wedge\quad
 \operatorname{IntegrableOn}
  (\omega,\operatorname{Set.Iio}(\operatorname{finiteParam}(a)),
   \operatorname{finiteMeasure}(\nu))\notag\\
 &\quad\wedge\quad
 \operatorname{IntegrableOn}
  (-\omega,\operatorname{Set.Ioi}(\operatorname{finiteParam}(b)),
   \operatorname{finiteMeasure}(\nu)).
 \label{eq:A-B-truncated}
\end{align}
\end{proposition}

\begin{proof}
Put $\rho:=\operatorname{finiteMeasure}(\nu)$,
$\mu:=\operatorname{negativeSigned}(\nu)$,
$L_a:=\operatorname{lowerTrunc}(\nu,a)$,
$A_b:=\operatorname{upperTrunc}(\nu,b)$, and $B_a:=-L_a$.
On $\operatorname{Set.Iio}(\operatorname{finiteParam}(a))$ the weight is
bounded and nonnegative, while on
$\operatorname{Set.Ioi}(\operatorname{finiteParam}(b))$ its negative is
bounded, strictly positive, and bounded away from zero.  Thus both integrals
are finite; retain the named proofs
\[
 h_{B_a}:B_a\leq0,\qquad h_{A_b}^{\rm nn}:0\leq A_b.
\]
The positive tail-mass hypothesis and
\cref{lem:support-strict-integral} applied to $\rho$ and $-\omega$ give the
strict witness
\[
 h_{A_b}:0<A_b.
\]
To prove \cref{eq:truncated-condition}, suppose for contradiction that it
fails.  After rewriting with $B_a=-L_a$, retain the named strict inequality
\[
 h_{\rm fail}:A_b+B_a<1.
\]
Put
\[
 C:=1-A_b-B_a.
\]
Then bind
\[
 h_C:0<C
 \quad\text{by unfolding \(C\) and applying linear arithmetic to
 \(h_{\rm fail}\)}.
\]
Together, $h_C$ and $h_{A_b}$ justify every division and the final strict
inequality below.  Bind the proofs
\[
 h_{a0}:0<a,\quad h_{ab}:a<b,\quad
 h_{a1}:a\neq1,\quad h_{b1}:b\neq1,
\]
choose $R_3$ with $h_{R_3}:a+b<R_3$, and define
\[
 B_3:=\operatorname{threeBlockData}
  (a,b,R_3,h_{a0},h_{ab},h_{R_3}),
\]
\[
 u_T:=\operatorname{Fin.cases}\!\left(
  0,j\mapsto\operatorname{Fin.cases}
   \left(\frac1C,(\_\mapsto-\frac1{A_b}),j\right)\right):
  \operatorname{Fin}(3)\to\R,\qquad K_T:=\{u_T\},
\]
\[
 \phi_T(n,\lambda):=
  \operatorname{blockPhiCurve}(\mu,B_3,n,u_T,\lambda).
\]
The truncation bridge gives
\[
 \operatorname{threeCellMoment}(\mu,a,b,0)=B_a,\qquad
 \operatorname{threeCellMoment}(\mu,a,b,2)=A_b,
\]
and $a<1<b$ gives
$\operatorname{threeNormBlock}(a,b)=1$.  Hence the two targets in
\cref{def:three-block-limits} simplify at $u_T$ to
\[
 \operatorname{threeLogFirst}(\mu,a,b,u_T)=0,\qquad
 \operatorname{threeLogSecond}(\mu,a,b,u_T)
 =-\frac1C-\frac1{A_b}<0.
\]
Apply \cref{prop:three-block-localisation} to the nonempty compact singleton
$K_T$; rewrite its two null assumptions with
\cref{lem:signed-witness-bridge}, and evaluate with
\cref{lem:compact-uniform-singleton}.  Instantiate
\cref{lem:block-curve-kernel-bridge} with $j_\infty=2$ and
$h_\infty:=\operatorname{threeBlockTopUnique}
 (a,b,R_3,h_{a0},h_{ab},h_{R_3})$.
It supplies positivity, \texttt{ContDiffOn}, and both log-kernel identities.
For each $n$, install
\texttt{letI := blockCarrierNonempty B\_3 n}; the global
\texttt{NegPhiConvex} hypothesis and the cone-affine bridge then give
convexity on that same interval.  The negative half of
\cref{lem:typed-log-curvature-contradiction}, with $a=0$ and
$b=-1/C-1/A_b$, is impossible.  Thus
\cref{eq:truncated-condition} holds, and the two already-established
integrability statements complete the required conjunction.
\end{proof}

\begin{proposition}[At most one upper support point]
\label{prop:one-upper-point}
Let $\nu:\operatorname{FiniteMeasure}(\overline\R_+)$.  If
$\texttt{NegPhiConvex.\{u\}}(\nu)$, then
\[
 \bigl(\operatorname{suppMeasure}(\operatorname{finiteMeasure}(\nu))
       \cap\operatorname{Set.Ioi}(1)\bigr).\operatorname{Subsingleton}.
\]
\end{proposition}

\begin{proof}
Put $\rho:=\operatorname{finiteMeasure}(\nu)$ and
$\mu:=\operatorname{negativeSigned}(\nu)$, and abbreviate
\[
 S:=\operatorname{suppMeasure}(\rho)\cap\operatorname{Set.Ioi}(1).
\]
First prove the following ordered-pair contradiction:
\[
 h_{\rm ord}:\ 
 \forall\alpha_1\,\alpha_2:\texttt{Param},\
 \alpha_1\in\operatorname{suppMeasure}(\rho)\to
 \alpha_2\in\operatorname{suppMeasure}(\rho)\to
 1<\alpha_1\to\alpha_1<\alpha_2\to\operatorname{False}.
\]
Fix $\alpha_1,\alpha_2$ and name the four hypotheses
$h_{\alpha,1}$, $h_{\alpha,2}$,
$h_{1,\alpha}:1<\alpha_1$, and
$h_{\alpha,12}:\alpha_1<\alpha_2$.
By \cref{lem:two-upper-null-thresholds}, choose $a,b:\R$ with named proofs
\[
 h_a:1<a,\qquad
 h_{a,\alpha}:\operatorname{finiteParam}(a)<\alpha_1,\qquad
 h_{\alpha,b}:\alpha_1<\operatorname{finiteParam}(b),\qquad
 h_{b,\alpha}:\operatorname{finiteParam}(b)<\alpha_2,
\]
and retain the two separate null proofs
\[
 h_{\rho,a}:\rho(\{\operatorname{finiteParam}(a)\})=0,\qquad
 h_{\rho,b}:\rho(\{\operatorname{finiteParam}(b)\})=0.
\]
Set
\[
 A_1:=\operatorname{threeCellMoment}(\mu,a,b,1),
 \qquad A_2:=\operatorname{upperMoment}(\mu,b).
\]
Both intervals contain an assumed support point.  Their strict positivity
gives named proofs $h_{A_1}:0<A_1$ and $h_{A_2}:0<A_2$ by
\cref{lem:support-strict-integral} applied to $\rho$ and $-\omega$.
Bind
\[
 h_{a0}:0<a,\qquad h_{ab}:a<b,\qquad
 h_{a1}:a\neq1,\qquad h_{b1}:b\neq1,
\]
choose $R_O$ with $h_{R_O}:a+b<R_O$, and define
\[
 B_O:=\operatorname{threeBlockData}
  (a,b,R_O,h_{a0},h_{ab},h_{R_O}),
\]
\[
 u_O:=\operatorname{Fin.cases}\!\left(
  0,j\mapsto\operatorname{Fin.cases}
    \left(\frac1{A_1},(\_\mapsto-\frac1{A_2}),j\right)\right):
    \operatorname{Fin}(3)\to\R,\qquad K_O:=\{u_O\},
\]
\[
 \phi_O(n,\lambda):=
  \operatorname{blockPhiCurve}(\mu,B_O,n,u_O,\lambda).
\]
Since $1<a$, \texttt{threeNormBlock a b = 0}.  Expanding the two targets
therefore gives the exact identities
\[
 \operatorname{threeLogFirst}(\mu,a,b,u_O)=0,\qquad
 \operatorname{threeLogSecond}(\mu,a,b,u_O)
 =-\frac1{A_1}-\frac1{A_2}<0.
\]
Apply \cref{prop:three-block-localisation} to $K_O$, rewrite its null
assumptions through \cref{lem:signed-witness-bridge}, and use
\cref{lem:compact-uniform-singleton}.  Instantiate the block curve bridge
with $j_\infty=2$ and
$\operatorname{threeBlockTopUnique}
 (a,b,R_O,h_{a0},h_{ab},h_{R_O})$.
For each $n$, install
\texttt{letI := blockCarrierNonempty B\_O n}.  The global
\texttt{NegPhiConvex} hypothesis gives interval convexity through
\cref{lem:cone-affine-line-bridge}; the bridge gives the matching
smoothness and logarithmic derivative identities.  The negative half of
\cref{lem:typed-log-curvature-contradiction}, with limiting pair
$(0,-1/A_1-1/A_2)$, is the required contradiction.
This proves $h_{\rm ord}$.

It remains to prove $S.\operatorname{Subsingleton}$.  Take
$\alpha_1,\alpha_2:\texttt{Param}$ with
$h_1:\alpha_1\in S$ and $h_2:\alpha_2\in S$.
Apply trichotomy to $\alpha_1$ and $\alpha_2$.  In the first branch retain
$h_{12}:\alpha_1<\alpha_2$ and eliminate
$h_{\rm ord}(\alpha_1,\alpha_2,h_1.1,h_2.1,h_1.2,h_{12})$.
The equality branch is the desired conclusion.  In the final branch retain
$h_{21}:\alpha_2<\alpha_1$ and eliminate
$h_{\rm ord}(\alpha_2,\alpha_1,h_2.1,h_1.1,h_2.2,h_{21})$.
Thus $\alpha_1=\alpha_2$, and unfolding $S$ proves the displayed
\texttt{Subsingleton} proposition.
\end{proof}

\begin{proposition}[Complete negative temperate necessity]
\label{prop:negative-temperate-necessity}
Let $\nu:\operatorname{FiniteMeasure}(\overline\R_+)$.  If
$\texttt{NegPhiConvex.\{u\}}(\nu)$, then the following single disjunction
holds:
\begin{align}
&\operatorname{suppMeasure}(\operatorname{finiteMeasure}(\nu))
  \subseteq\operatorname{Set.Icc}(0,1)\nonumber\\[-2pt]
&\quad\lor\quad
 \exists\alpha_*:\overline\R_+,\quad
 1<\alpha_*\ \wedge\
 \operatorname{suppMeasure}(\operatorname{finiteMeasure}(\nu))
  \subseteq\operatorname{Set.Icc}(0,1)\cup\{\alpha_*\}\ 
 \wedge\nonumber\\[-2pt]
&\hspace{17mm}
 \operatorname{finiteMeasure}(\nu)(\{1\})=0\ 
 \wedge\
 \operatorname{MLower}(\operatorname{finiteMeasure}(\nu))<\top\ 
 \wedge\
 \operatorname{MomFin}(\operatorname{finiteMeasure}(\nu))\ 
 \wedge\nonumber\\[-2pt]
&\hspace{17mm}
 \operatorname{MReal}(\operatorname{finiteMeasure}(\nu))\leq-1.
 \label{eq:negative-mu-moment}
\end{align}
\end{proposition}

\begin{proof}
Put $\rho:=\operatorname{finiteMeasure}(\nu)$ and
$\mu:=\operatorname{negativeSigned}(\nu)$.
If there is no upper support, the first alternative holds. Suppose that there
is upper support. By \cref{prop:one-upper-point}, its intersection with
$(1,+\infty]$ consists of a single point $\alpha_*$. By
\cref{prop:negative-shannon-obstruction}, $\rho(\{1\})=0$.

Use \texttt{by\_cases h\_top : alpha\_* = top}.  In the equality branch,
apply the nonempty-open-set
clause of \cref{lem:null-thresholds} to the finite interval $(1,2)$; in the
non-equality branch, retain the propositional proof
$h_{\rm top}:\alpha_*\neq\top$, put
$r_*:=((\alpha_*.\operatorname{untop}(h_{\rm top}):\R_{\geq0}):\R)$, and apply it to
$(1,r_*)$.  In either branch retain one $b:\R$ and the uniform named proofs
\[
 h_{1b}:1<b,\qquad
 h_{b*}:\operatorname{finiteParam}(b)<\alpha_*,\qquad
 h_{\rho,b}:\rho(\{\operatorname{finiteParam}(b)\})=0.
\]
Since there is no support in
$(1,+\infty]\setminus\{\alpha_*\}$, the quantity
\[
 A_b:=\operatorname{upperMoment}(\mu,b)
\]
is finite and equals the whole upper contribution.  More precisely,
\cref{lem:negative-truncation-bridge,lem:fnd-support-complement,lem:fnd-moment-algebra}
give the named equality
\[
 h_A^{\rm whole}:A_b=
  \operatorname{ENNReal.toReal}(\operatorname{MUpper}(\rho)).
\]
The isolated support point $\alpha_*$ has positive mass, so
\cref{lem:support-strict-integral} gives $h_{A_b}:0<A_b$.

For each real $\rho$-null $a$ with $0<a<1$, set
$B_a:=-\operatorname{lowerTrunc}(\nu,a)$.  The truncated obstruction gives
the named proof $h_{B_a}:1-A_b\leq B_a$.
Since also $B_a\leq0$, these two inequalities imply $A_b\geq1$.  Thus
$A_b-1$ is a legitimate nonnegative finite upper bound.
Using the positive raw measure $\rho$, define
\[
 \ell_a:=\int^-_{\operatorname{Set.Iio}(\operatorname{finiteParam}(a))}
   \omega_<\dd\rho.
\]
Boundedness on the truncated set gives the separate named proofs
\[
 h_{\ell_a}:\ell_a<\top,\qquad
 h_{\ell_a}^{\rm real}:\operatorname{ENNReal.toReal}(\ell_a)=-B_a,
 \qquad h_{\ell_a}^{\rm bd}:\operatorname{ENNReal.toReal}(\ell_a)\leq A_b-1.
\]
The equality with $-B_a$ uses boundedness of $\omega$ on the truncated set
and the signed identity $\mu=\operatorname{negativeSigned}(\nu)$.  Since
$A_b-1\geq0$, the finite-value order
embedding gives
$\ell_a\leq\operatorname{ENNReal.ofReal}(A_b-1)$.
Choose an increasing sequence of $\rho$-null points $a_n\uparrow1$. The
functions $\omega_<\mathbf1_{\operatorname{Set.Iio}
 (\operatorname{finiteParam}(a_n))}$ are nonnegative and increase to
$\omega_<$. Applying monotone convergence in $[0,+\infty]$ to the
\emph{positive} measure $\rho$ yields
\[
 \operatorname{MLower}(\rho)=\lim_{n\to\infty}\ell_{a_n}
 \leq\operatorname{ENNReal.ofReal}(A_b-1)<+\infty.
\]
Thus the lower one-sided moment is finite.  The support condition makes
$\operatorname{MUpper}(\rho)$ finite, so $\operatorname{MomFin}(\rho)$.
Put $L:=\operatorname{ENNReal.toReal}(\operatorname{MLower}(\rho))$.  By
\cref{lem:fnd-moment-algebra}, this is the ordinary lower integral.  Since
the upper contribution is $A_b$ by $h_A^{\rm whole}$, unfolding
\texttt{MReal} gives the named equality
\[
 h_M:-\operatorname{MReal}(\rho)=A_b-L.
\]
The preceding bound on $\operatorname{MLower}(\rho)$ gives $L\leq A_b-1$,
so $h_M$ yields $1\leq-\operatorname{MReal}(\rho)$.
Thus $\operatorname{MReal}(\rho)\leq-1$; all the stored exceptional-branch
conjuncts now follow.
\end{proof}

\begin{corollary}[Negative temperate necessity for probability measures]
\label{cor:negative-temperate-probability-necessity}
At universe $u$, let
$\tau:\operatorname{ProbabilityMeasure}(\overline\R_+)$ and $t:\R$.
For every proof $h_t:t<0$,
\[
 \operatorname{CMMonotone}
  (\operatorname{HTemp}(t,\operatorname{ne\_of\_lt}(h_t),\tau))
 \Longrightarrow
 \operatorname{NegLowerAdm}(t,\tau)\ \lor\
 \exists\alpha_*:\overline\R_+,\,
  \operatorname{NegExcAdm}(t,\tau,\alpha_*).
\]
\end{corollary}

\begin{proof}
Assume the displayed monotonicity.  Bind
\[
 h_{-t}:0<-t:=\operatorname{neg\_pos.mpr}(h_t),\qquad
 h_{-t}^{\rm nn}:0\leq-t:=\operatorname{le\_of\_lt}(h_{-t}),
\]
and put
\[
 \nu_t:=\operatorname{finiteScale}
  (-t,h_{-t}^{\rm nn},\tau.\operatorname{toFiniteMeasure}).
\]
The second conjunct of
\cref{cor:global-temperate-shape-reduction} supplies
$\operatorname{ConvexCone}_I(\Phi_{t,\tau})$ for every finite nonempty
$I:\texttt{Type u}$.  By the negative-scaling identity in
\cref{lem:signed-witness-bridge},
\[
 \operatorname{negativeSigned}(\nu_t)
 =-(-t)\mathbin{\bullet}
   \operatorname{signedLift}(\tau.\operatorname{toFiniteMeasure})
 =t\mathbin{\bullet}
   \operatorname{signedLift}(\tau.\operatorname{toFiniteMeasure}).
\]
The first conjunct of \cref{lem:signed-scalar-column-bridge} therefore
rewrites the rowwise convexity family as
$\texttt{NegPhiConvex.\{u\}}(\nu_t)$.  Apply
\cref{prop:negative-temperate-necessity}.

In the lower-support branch, support invariance under the nonzero positive
scalar $-t$ transfers the inclusion to $\tau$, and $h_t$ completes
$\operatorname{NegLowerAdm}(t,\tau)$.  In the exceptional branch,
destructure the witness $\alpha_*$ and its named conjuncts.
The support and atom fields transfer by
\cref{lem:scalar-measure-support}.  The positive-scalar clauses of
\cref{lem:fnd-moment-algebra} reflect
$\operatorname{MomFin}$ and give the guarded identity
\[
 \operatorname{MReal}(\operatorname{finiteMeasure}(\nu_t))
 =(-t)\operatorname{MReal}(\operatorname{probMeasure}(\tau)).
\]
Consequently the bound
$\operatorname{MReal}(\operatorname{finiteMeasure}(\nu_t))\leq-1$
is equivalent, after division by the positive number $-t$, to
\[
 \operatorname{MReal}(\operatorname{probMeasure}(\tau))\leq1/t.
\]
Moment finiteness supplies the required lower-moment finiteness, and all
remaining fields are invariant under the same nonzero scalar.  Reassemble
them to obtain
$\operatorname{NegExcAdm}(t,\tau,\alpha_*)$.  No unguarded use of
\texttt{MReal} occurs.
\end{proof}

\subsection{Necessity in the negative tropical branch}

\begin{proposition}[Complete negative tropical necessity]
\label{prop:negative-tropical-necessity}
Let $\tau:\operatorname{ProbabilityMeasure}(\overline\R_+)$.  Then
\[
 \operatorname{CMMonotone}(\operatorname{HMinus}(\tau))
 \Longrightarrow\operatorname{DMinus}(\tau).
\]
\end{proposition}

\begin{proof}
Assume the displayed monotonicity, put
\[
 \nu:=\tau.\operatorname{toFiniteMeasure},
 \qquad
 \rho:=\operatorname{probMeasure}(\tau),
 \qquad
 \mu:=\operatorname{negativeSigned}(\nu).
\]
By
\cref{lem:signed-scalar-column-bridge,cor:global-tropical-shape-reduction},
the hypothesis gives $\texttt{NegGQuasiconvex.\{u\}}(\nu)$.

Assume first that $\rho(\{1\})\neq0$ and that there is upper support.  Set
$m:=\operatorname{signedAtom}(\mu,1)$.  By
\cref{lem:signed-witness-integral-bridge}, obtain the separate named proofs
\[
 h_m^{\rm eq}:m=-\operatorname{ENNReal.toReal}(\rho(\{1\})),
 \qquad h_m^{\rm neg}:m<0.
\]
By \cref{lem:null-thresholds}, choose $c:\R$ with named proofs
\[
 h_c:1<c,\qquad
 h_{\rho,c}:\rho(\{\operatorname{finiteParam}(c)\})=0,\qquad
 0<\rho(\operatorname{Set.Ioi}(\operatorname{finiteParam}(c))),
\]
and set $A_c:=\operatorname{upperMoment}(\mu,c)$.  The theorem
\cref{lem:support-strict-integral} gives the named proof $h_{A_c}:0<A_c$.
Choose $R$ and $h_R:c+1<R$, and set
\[
\theta:=
 \left\{c:=c,\ R:=R,\ p:=\tfrac12,\
 \texttt{c\_gt\_one}:=h_c,\ \texttt{R\_gt}:=h_R,\
 \texttt{p\_pos}:=\texttt{by norm\_num},\
 \texttt{p\_lt\_one}:=\texttt{by norm\_num}\right\}
 \in\operatorname{ShannonData}.
\]
Define, with literal types,
\begin{align*}
 a_S(n,z)&:=\operatorname{shannonGKernel}(\mu,\theta,n,z,1),&
 b_S(n,z)&:=\operatorname{shannonGKernel}(\mu,\theta,n,z,2),\\
 a_S^\infty(z)&:=\operatorname{shannonLimitFirst}(\mu,\theta,z),&
 b_S^\infty(z)&:=\operatorname{shannonLimitSecond}(\mu,\theta,z),\\
 \gamma_S(n,z,\lambda)&:=
   \operatorname{shannonGCurve}(\mu,\theta,n,z,\lambda).
\end{align*}
The exact conjunction
\[
 \operatorname{StationarityPackage}
  (a_S,a_S^\infty,b_S,b_S^\infty)
\]
follows by unfolding \cref{def:stationarity-package}: additivity,
homogeneity, and continuity are
\cref{lem:shannon-kernel-regularity}, the limit targets are polynomial and
continuous, and for each nonempty compact $K$ the last two conjuncts are the
third and fourth clauses of \cref{prop:shannon-localisation}, using
$h_{\rho,c}$ and \cref{lem:signed-witness-bridge}.
For every $n,z$, first install
\texttt{letI := shannonIndexNonempty theta n}.  The exact line package
\cref{eq:corrected-stationary-line-package} for $\gamma_S$ follows from
\cref{lem:shannon-curve-kernel-bridge}; its scalar quasi-convexity clause is
\texttt{NegGQuasiconvex} restricted by
\cref{lem:cone-affine-line-bridge}.

Finally put
\[
 z_0:=\left(-\frac{A_c}{m},1\right),\qquad e_S:=(0,1).
\]
Unfolding the targets gives
\[
 a_S^\infty(z_0)=0,\qquad
 b_S^\infty(z_0)=-A_c<0,\qquad
 a_S^\infty(e_S)=A_c\neq0.
\]
Thus \cref{lem:corrected-stationary-line-obstruction} is a contradiction.
Hence upper support forces $\rho(\{1\})=0$.

Next suppose that $\alpha_1,\alpha_2:\texttt{Param}$ are support points with
named proofs $h_{1,\alpha}:1<\alpha_1$ and
$h_{\alpha,12}:\alpha_1<\alpha_2$.  Use
\cref{lem:two-upper-null-thresholds} to choose $a,b:\R$ with
\[
 h_a:1<a,\qquad
 h_{a,\alpha}:\operatorname{finiteParam}(a)<\alpha_1,\qquad
 h_{\alpha,b}:\alpha_1<\operatorname{finiteParam}(b),\qquad
 h_{b,\alpha}:\operatorname{finiteParam}(b)<\alpha_2,
\]
and retain the two null proofs
$h_{\rho,a}:\rho(\{\operatorname{finiteParam}(a)\})=0$ and
$h_{\rho,b}:\rho(\{\operatorname{finiteParam}(b)\})=0$.  Put
\[
 A_1:=\operatorname{threeCellMoment}(\mu,a,b,1),
 \qquad A_2:=\operatorname{upperMoment}(\mu,b).
\]
The theorem \cref{lem:support-strict-integral} gives named proofs
$h_{A_1}:0<A_1$ and $h_{A_2}:0<A_2$.  Bind
$h_{a0}:0<a$, $h_{ab}:a<b$, $h_{a1}:a\neq1$, and $h_{b1}:b\neq1$; choose
$R_D$ with $h_{R_D}:a+b<R_D$; and put
\[
 B_D:=\operatorname{threeBlockData}
  (a,b,R_D,h_{a0},h_{ab},h_{R_D}).
\]
Define
\begin{align*}
 a_D(n,u)&:=\operatorname{blockGKernel}(\mu,B_D,n,u,1),&
 b_D(n,u)&:=\operatorname{blockGKernel}(\mu,B_D,n,u,2),\\
 a_D^\infty(u)&:=\operatorname{threeGFirst}(\mu,a,b,u),&
 b_D^\infty(u)&:=\operatorname{threeGSecond}(\mu,a,b,u),\\
 \gamma_D(n,u,\lambda)&:=
   \operatorname{blockGCurve}(\mu,B_D,n,u,\lambda).
\end{align*}
Unfolding \cref{def:stationarity-package}, the unique-top instance
\[
 h_\infty:=\operatorname{threeBlockTopUnique}
  (a,b,R_D,h_{a0},h_{ab},h_{R_D})
\]
and \cref{lem:block-kernel-regularity,prop:three-block-localisation} give
\[
 \operatorname{StationarityPackage}
  (a_D,a_D^\infty,b_D,b_D^\infty).
\]
Here the localization proposition is applied separately to every arbitrary
nonempty compact $K$; its null assumptions are rewritten with
\cref{lem:signed-witness-bridge}.  For each $n$, install
\texttt{letI := blockCarrierNonempty B\_D n}.  The block curve bridge with
$(j_\infty,h_\infty)=(2,h_\infty)$ and the cone-affine restriction of
\texttt{NegGQuasiconvex} give
\cref{eq:corrected-stationary-line-package} for $\gamma_D$.

Let
\[
 u_D:=\operatorname{Fin.cases}\!\left(
  0,j\mapsto\operatorname{Fin.cases}
   \left(\frac1{A_1},(\_\mapsto-\frac1{A_2}),j\right)\right):
   \operatorname{Fin}(3)\to\R,
\]
\[
 e_D:=\operatorname{Fin.cases}
  (0,(j\mapsto\operatorname{Fin.cases}(1,(\_\mapsto0),j))):
  \operatorname{Fin}(3)\to\R.
\]
Because $1<a$, the norm block is $0$, and direct finite-case simplification
gives
\[
 a_D^\infty(u_D)=0,\qquad
 b_D^\infty(u_D)=-\frac1{A_1}-\frac1{A_2}<0,\qquad
 a_D^\infty(e_D)=A_1\neq0.
\]
\Cref{lem:corrected-stationary-line-obstruction} is a contradiction.
Therefore at most one upper support point is possible.

Suppose now that such a point $\alpha_*>1$ exists.  It is isolated from
$[0,1]$; because it belongs to the support it has strictly positive
$\rho$-mass.  The preceding atom argument gives $\rho(\{1\})=0$.
Choose $\rho$-null thresholds $a,b:\R$ by the convergent-null-sequence and
nonempty-open-set clauses of \cref{lem:null-thresholds}, splitting on
$\alpha_*=\top$ for the choice of $b$.  Retain the named proofs
\[
 h_{a0}:0<a,\quad h_{a1}:a<1,\quad h_{b1}:1<b,\quad
 h_{b*}:\operatorname{finiteParam}(b)<\alpha_*,\quad
 h_{\rho,a}:\rho(\{\operatorname{finiteParam}(a)\})=0,\quad
 h_{\rho,b}:\rho(\{\operatorname{finiteParam}(b)\})=0.
\]
Define, without referring to an equation label,
\[
 B_a:=-\operatorname{lowerTrunc}(\nu,a),\qquad
 A_b:=\operatorname{upperTrunc}(\nu,b).
\]
Then $B_a\leq0$, $0<A_b$, and
\cref{lem:negative-truncation-bridge} gives the exact identities
\[
 B_a=\operatorname{lowerMoment}(\mu,a),\qquad
 A_b=\operatorname{upperMoment}(\mu,b).
\]
The exceptional support inclusion, $h_{b*}$, and
\cref{lem:fnd-support-complement,lem:fnd-moment-algebra} also give the named
whole-upper identity
\[
 h_A^{\rm whole}:A_b=
  \operatorname{ENNReal.toReal}(\operatorname{MUpper}(\rho)).
\]
Assume for contradiction the named inequality
\[
 h_{\rm fail}:A_b+B_a<0.
\]
Put $C:=-A_b-B_a$ and retain
\[
 h_C:0<C
 \quad\text{by unfolding \(C\) and applying linear arithmetic to
 \(h_{\rm fail}\)}.
\]
Bind $h_{ab}:a<b$, $h_{a\ne1}:a\neq1$, and $h_{b\ne1}:b\neq1$, choose
$R_T$ with $h_{R_T}:a+b<R_T$, and set
\[
 B_T:=\operatorname{threeBlockData}
  (a,b,R_T,h_{a0},h_{ab},h_{R_T}).
\]
Define
\begin{align*}
 a_T(n,u)&:=\operatorname{blockGKernel}(\mu,B_T,n,u,1),&
 b_T(n,u)&:=\operatorname{blockGKernel}(\mu,B_T,n,u,2),\\
 a_T^\infty(u)&:=\operatorname{threeGFirst}(\mu,a,b,u),&
 b_T^\infty(u)&:=\operatorname{threeGSecond}(\mu,a,b,u),\\
 \gamma_T(n,u,\lambda)&:=
  \operatorname{blockGCurve}(\mu,B_T,n,u,\lambda).
\end{align*}
Unfold \cref{def:stationarity-package}.  The additivity, homogeneity, and
continuity clauses are \cref{lem:block-kernel-regularity}; the two limit
targets are finite polynomials and hence continuous; and, for every
arbitrary nonempty compact $K$, the final two clauses are the norm-free
parts of \cref{prop:three-block-localisation}.  Therefore
\[
 \operatorname{StationarityPackage}
  (a_T,a_T^\infty,b_T,b_T^\infty);
\]
the application uses the two endpoint-null proofs, every arbitrary nonempty
compact $K$, and
$h_\infty:=\operatorname{threeBlockTopUnique}
 (a,b,R_T,h_{a0},h_{ab},h_{R_T})$.
For each $n$, install
\texttt{letI := blockCarrierNonempty B\_T n}.  The block curve bridge with
$(2,h_\infty)$ and
\texttt{NegGQuasiconvex} restricted by the cone-affine bridge give the exact
line package for $\gamma_T$.

Put
\[
 u_T:=\operatorname{Fin.cases}\!\left(
  0,j\mapsto\operatorname{Fin.cases}
   \left(\frac1C,(\_\mapsto-\frac1{A_b}),j\right)\right):
   \operatorname{Fin}(3)\to\R,\qquad
 e_T:=\operatorname{Fin.cases}
  (0,(j\mapsto\operatorname{Fin.cases}(0,(\_\mapsto1),j))):
  \operatorname{Fin}(3)\to\R.
\]
Since $a<1<b$, \texttt{threeNormBlock a b = 1}; finite-case simplification
therefore gives
\[
 a_T^\infty(u_T)=0,\qquad
 b_T^\infty(u_T)=-\frac1C-\frac1{A_b}<0,\qquad
 a_T^\infty(e_T)=A_b\neq0.
\]
This contradicts \cref{lem:corrected-stationary-line-obstruction}.
Consequently,
\begin{equation}
 A_b+B_a\geq0
 \label{eq:tropical-truncated}
\end{equation}
for every such null threshold $a<1$.

To pass to $a\uparrow1$, we again avoid any monotone-convergence argument for
a signed measure. Since $\rho$ is positive, define
\[
 \ell_a:=\int^-_{\operatorname{Set.Iio}(\operatorname{finiteParam}(a))}
   \omega_<\dd\rho.
\]
Boundedness on the truncated interval and \cref{eq:tropical-truncated} give
the separate named proofs
\[
 h_{\ell_a}:\ell_a<\top,\qquad
 h_{\ell_a}^{\rm real}:\operatorname{ENNReal.toReal}(\ell_a)=-B_a,
 \qquad h_{\ell_a}^{\rm bd}:\operatorname{ENNReal.toReal}(\ell_a)\leq A_b.
\]
Since $A_b>0$, the finite-value order embedding gives
$\ell_a\leq\operatorname{ENNReal.ofReal}(A_b)$.  Along an increasing
sequence of $\rho$-null points $a_n\uparrow1$, monotone convergence for
$\rho$ in $[0,+\infty]$ gives
\[
 \operatorname{MLower}(\rho)=\lim_{n\to\infty}\ell_{a_n}
 \leq\operatorname{ENNReal.ofReal}(A_b)<+\infty.
\]
The unique-upper-support conclusion makes $\operatorname{MUpper}(\rho)$
finite, hence $\operatorname{MomFin}(\rho)$.  Put
$L:=\operatorname{ENNReal.toReal}(\operatorname{MLower}(\rho))$ and retain
the named proof $h_L:L\leq A_b$ from the preceding bound; the guarded
moment-agreement clause identifies $L$ with the ordinary lower integral.
Using $h_A^{\rm whole}$ and unfolding \texttt{MReal} gives
\[
 h_M:-\operatorname{MReal}(\rho)=A_b-L.
\]
Together $h_M$ and $h_L$ give $0\leq-\operatorname{MReal}(\rho)$.
Since $\rho=\operatorname{probMeasure}(\tau)$, this is exactly
\[
 \operatorname{MReal}(\operatorname{probMeasure}(\tau))\leq0.
\]
This proves the necessity of the two branches in \cref{def:Dminus}.
\end{proof}

Together with \cref{prop:negative-tropical-sufficiency}, this completes the
negative tropical classification.

\subsection{Necessity in the positive tropical branch}

\begin{definition}[Probability support as a positive cone vector]
\label{def:prob-as-pos-cone}
For a finite nonempty type $I$ and $p:\operatorname{ProbVec}(I)$, define
\begin{align*}
 \operatorname{probSupport}(p)&:\operatorname{Set}(I):=
   \{i:I\mid0<p.1(i)\},\\
 \operatorname{ProbSupport}(p)&:=
   \{i:I\mathbin{//}i\in\operatorname{probSupport}(p)\},\\
 \operatorname{probAsPosCone}(p)&:=
  \left\langle\left\langle p.1,p.2.1\right\rangle,
   \texttt{by
    intro hp
    have hm :=
      congrArg (fun x : ConeVec I $\Rightarrow$ l1Mass x.1) hp
    have hz : l1Mass p.1 = 0 := by
      simpa [l1Mass] using hm
    exact one\_ne\_zero ((p.2.2).symm.trans hz)}
   \right\rangle
  \in\operatorname{PosConeVec}(I).
\end{align*}
\end{definition}

\begin{lemma}[Probability-to-cone identities]
\label{lem:prob-as-pos-cone-identities}
For finite nonempty $I$, $p,q:\operatorname{ProbVec}(I)$,
$\lambda:\R$, and $h_\lambda:0<\lambda\ \wedge\ \lambda<1$, put
$h_\lambda^{\rm cl}:=
\langle\operatorname{le_of_lt}(h_\lambda.1),
       \operatorname{le_of_lt}(h_\lambda.2)\rangle$.
The following literal conjunction holds:
\begin{align*}
 &\operatorname{normalize}(\operatorname{probAsPosCone}(p))=p\\
 &\quad\wedge\quad
 \operatorname{probAsPosCone}
   (\operatorname{mixProbVec}(\lambda,h_\lambda^{\rm cl},p,q))
  =\operatorname{posMix}(\lambda,h_\lambda,
    \operatorname{probAsPosCone}(p),\operatorname{probAsPosCone}(q)).
\end{align*}
The second equality uses proof irrelevance for the bundled positivity fields.
\end{lemma}

\begin{proof}
The first identity follows from the mass-one field of $p$.  The second is
coordinatewise reflexivity after expanding the two mixture constructors;
proof irrelevance closes the subtype fields.
\end{proof}

\begin{lemma}[A probability vector has nonempty positive support]
\label{lem:prob-support-nonempty}
For every finite nonempty $I$ and $p:\operatorname{ProbVec}(I)$,
\[
 \operatorname{probSupportNonempty}(p):
  \operatorname{Nonempty}\!\left(\operatorname{ProbSupport}(p)\right).
\]
\end{lemma}

\begin{proof}
If the support were empty, every nonnegative coordinate would be zero.
The finite sum would then be zero, contradicting the stored mass identity
$\operatorname{l1Mass}(p.1)=1$.
\end{proof}

\begin{lemma}[Decomposition inside a simplex face]
\label{lem:same-support-decomposition}
Let $I$ be finite and nonempty and let
$p,q:\operatorname{ProbVec}(I)$ satisfy
$\operatorname{probSupport}(p)=\operatorname{probSupport}(q)$.  Then
\begin{align*}
 \exists\eps:\R,\quad0<\eps\ \wedge\ \eps\leq1\ \wedge\
 \forall\lambda:\R,\quad0<\lambda\to\lambda<\eps\to
 \exists(h_\lambda:0\leq\lambda\ \wedge\ \lambda\leq1)
 (r:\operatorname{ProbVec}(I)),\\[-2pt]
 \operatorname{probSupport}(r)=\operatorname{probSupport}(p)
 \quad\wedge\quad
 p=\operatorname{mixProbVec}(\lambda,h_\lambda,q,r).
\end{align*}
\end{lemma}

\begin{proof}
Install
\texttt{letI : Nonempty (ProbSupport p) :=
  probSupportNonempty p}.
On this finite subtype, the ratios
$p(i)/q(i)$ are strictly positive.  Their \texttt{finMin}, and its minimum
with $1$, is the required $\eps$.
The chosen inequality gives $0<\lambda<1$ and makes every
supported coordinate of $p-\lambda q$ strictly positive, while coordinates
outside the common support remain zero.  Its total mass is $1-\lambda$.
Define $r(i):=(p(i)-\lambda q(i))/(1-\lambda)$ and bundle the verified
nonnegative coordinates and mass-one equality.  Division by the positive
scalar $1-\lambda$ proves the support and mixture equalities.
\end{proof}

\begin{lemma}[Probability supported at zero is the zero Dirac measure]
\label{lem:probability-eq-dirac-zero}
For $\tau:\operatorname{ProbabilityMeasure}(\overline\R_+)$,
\[
 \operatorname{probMeasure}(\tau)(\operatorname{Set.Ioi}(0))=0
 \Longleftrightarrow
 \tau=\operatorname{diracProb}(0).
\]
\end{lemma}

\begin{proof}
The parameter space is the disjoint union of $\{0\}$ and
\texttt{Set.Ioi 0}.  The forward implication gives mass one to the singleton;
measure extensionality then identifies the underlying measure with the Dirac
measure.  The reverse implication is evaluation of a Dirac measure.
\end{proof}

\begin{proposition}[Only the support-rank measure survives]
\label{prop:positive-tropical-necessity}
Let $\tau:\operatorname{ProbabilityMeasure}(\overline\R_+)$.  Then
\[
 \operatorname{CMMonotone}(\operatorname{HPlus}(\tau))
 \Longrightarrow\tau=\operatorname{diracProb}(0).
\]
\end{proposition}

\begin{proof}
Assume monotonicity and put
$\A_I(p):=\operatorname{aTrop}_I
 (\tau,\operatorname{probAsPosCone}(p))$. By
\cref{cor:global-tropical-shape-reduction,lem:prob-as-pos-cone-identities},
monotonicity implies
\begin{equation}
 \A_I(\operatorname{mixProbVec}
   (\lambda,\langle\operatorname{le_of_lt}(h_\lambda.1),
     \operatorname{le_of_lt}(h_\lambda.2)\rangle,p,q))
 \geq\max\{\A_I(p),\A_I(q)\},
 \label{eq:strong-qc}
\end{equation}
for every finite nonempty $I$, $p,q:\operatorname{ProbVec}(I)$,
$\lambda:\R$, and $h_\lambda:0<\lambda\ \wedge\ \lambda<1$.

Suppose that $p$ and $q$ have the same support.  Destructure
\cref{lem:same-support-decomposition} as
$(\eps,h_{\eps0},h_{\eps1},h_{\rm dec})$, put $\lambda:=\eps/2$, and retain
the arithmetic proofs
\[
 h_{\lambda0}:0<\lambda,\qquad h_{\lambda\eps}:\lambda<\eps,\qquad
 h_{\lambda1}:\lambda<1.
\]
Destructure $h_{\rm dec}(\lambda,h_{\lambda0},h_{\lambda\eps})$ as
$(h_\lambda^{\rm cl},r,h_{\rm supp},h_{\rm mix})$.  Rewrite
\cref{eq:strong-qc} with $h_{\rm mix}$ and the strict pair
$\langle h_{\lambda0},h_{\lambda1}\rangle$; its first maximum inequality
gives $\A_I(p)\geq\A_I(q)$.  Apply the same destructuring with $p$ and $q$
interchanged to obtain the reverse inequality and hence equality.
Therefore $\A$ depends only on the support.

Take
$p:=\operatorname{binaryProb}(1/3,
 \texttt{(by norm\_num)},\texttt{(by norm\_num)})$. It has the same
support as the uniform binary vector, so
\[
 \A_{\operatorname{Fin}(2)}(p)=\log2.
\]
On the other hand,
\[
 H_{(0:\overline\R_+)}(p)=\log2,
 \qquad
 H_\alpha(p)<\log2
 \quad\text{for every }\alpha:\overline\R_+\text{ with }0<\alpha.
\]
Because $\tau$ is positive and $H_\alpha(p)\leq\log2$ at every order, the
identity
\[
 \operatorname{integratedEntropyPos}_{\operatorname{Fin}(2)}
  (\operatorname{probMeasure}(\tau),p)=\log2
\]
implies
\[
 \int\bigl(\log2-H_\alpha(p)\bigr)
   \dd\operatorname{probMeasure}(\tau)(\alpha)=0.
\]
The integrand is Borel, nonnegative everywhere, and strictly positive on
$\operatorname{Set.Ioi}(0)$ by \cref{lem:fnd-strict-binary-gap}.  It is bounded by
$\log2$ and therefore integrable against the finite measure
$\operatorname{probMeasure}(\tau)$.  Hence
\cref{lem:fnd-zero-integral-positive} gives
$\operatorname{probMeasure}(\tau)(\operatorname{Set.Ioi}(0))=0$.
The probability-measure extensionality lemma for the partition
$\{0\}\cup\operatorname{Set.Ioi}(0)$ now gives
$\tau=\operatorname{diracProb}(0)$.
\end{proof}

Together with \cref{prop:positive-tropical-sufficiency}, this proves the
positive tropical part of the final classification theorem.

\subsection{Necessity for derivations}

\begin{definition}[Derivation column function]
\label{def:derivation-column-function}
For finite nonempty $I$ and
$\nu:\operatorname{FiniteMeasure}(\overline\R_+)$, define
\[
 \texttt{abbrev derivationColumn\_I :=
   columnDeriv\_I}.
\]
Thus
\texttt{derivationColumn\_I nu x = columnDeriv\_I nu x}
definitionally; no second case split or proof-irrelevance obligation is
created.  This is the column function whose conditional finite sum is
$\operatorname{Delta}(\nu)$.
\end{definition}

\begin{proposition}[Derivation necessity]
\label{prop:derivation-necessity}
Let $\nu:\operatorname{FiniteMeasure}(\overline\R_+)$.  If
\[
 \forall\{I:\texttt{Type u}\}
 [\operatorname{Fintype}I][\operatorname{Nonempty}I],\quad
 \operatorname{ConcaveCone}_I(\operatorname{derivationColumn}_I(\nu)),
\]
then
\[
 \operatorname{suppMeasure}(\operatorname{finiteMeasure}(\nu))
 \subseteq\operatorname{Set.Icc}(0,1).
\]
\end{proposition}

\begin{proof}
Put $\rho:=\operatorname{finiteMeasure}(\nu)$ and
$\mu_s:=\operatorname{positiveSigned}(\nu)$. Suppose that
$\operatorname{suppMeasure}(\rho)$ meets $\operatorname{Set.Ioi}(1)$. By
\cref{lem:null-thresholds}, choose a real $r>1$ with
$\rho(\{\operatorname{finiteParam}(r)\})=0$ and
$0<\rho(\operatorname{Set.Ioi}(\operatorname{finiteParam}(r)))$, and set
$A_r:=\operatorname{upperMoment}(\mu_s,r)$.
The integral is finite because the tail is separated from $1$, and it is
strictly negative by \cref{lem:support-strict-integral} since the tail has
positive $\rho$-mass; retain this as $h_{A_r}:A_r<0$.  Bind $h_{r0}:0<r$
and $h_{r1}:r\neq1$, choose
$R:\R$ with $h_R:r<R$, and put
\[
 B_2:=\operatorname{twoBlockData}(r,R,h_{r0},h_R),\qquad
 u:=\operatorname{Fin.cases}(0,(\_\mapsto1)):
   \operatorname{Fin}(2)\to\R.
\]
For each $n:\N$ first install
\[
 \texttt{letI : Nonempty (BlockCarrier B\_2 n) :=
   blockCarrierNonempty B\_2 n},
\]
and then write
\[
 L_n:=\operatorname{blockLineData}(B_2,n,u),\qquad
 S_n(\lambda):=\sum_i\operatorname{lineRaw}(L_n,\lambda)(i),
 \qquad
 K_n(\lambda):=\operatorname{integratedEntropyLine}(\mu_s,L_n,\lambda).
\]
Define the total curve
\[
 F_n(\lambda):=
 \operatorname{derivationColumn}_{\operatorname{BlockCarrier}(B_2,n)}
  (\nu,\operatorname{lineConeTotal}(L_n,\lambda)).
\]
By \cref{lem:block-curve-kernel-bridge}, instantiated with
$j_\infty=1$ and
$\operatorname{twoBlockTopUnique}(r,R,h_{r0},h_R)$, choose
$\eps_n>0$ together with its named line-positivity proof $h_{\rm pos}$ and
the named proof $h_G$ that
$\operatorname{blockGCurve}(\mu_s,B_2,n,u)$ is \texttt{ContDiffOn} of order
two.  Put $U_n:=\operatorname{Set.Ioo}(-\eps_n,\eps_n)$.  For every
$\lambda\in U_n$, the second and sixth conjuncts of
\cref{lem:signed-line-column-bridge}, instantiated with
$h_{\rm pos}(\lambda)$, give the named equality
\[
 h_{KG}(\lambda):K_n(\lambda)
  =\operatorname{blockGCurve}(\mu_s,B_2,n,u,\lambda).
\]
Transfer $h_G$ along $h_{KG}$ to obtain
$h_K:\operatorname{ContDiffOn}(\R,2,K_n,U_n)$.
On this interval, and only on this interval, unfolding
\texttt{derivationColumn}, \texttt{lineConeTotal}, and normalization gives
\[
 \forall\lambda\in\operatorname{Set.Ioo}(-\eps_n,\eps_n),\qquad
 F_n(\lambda)=S_n(\lambda)K_n(\lambda).
\]
Name this pointwise identity $h_{FK}$.  Since $S_n$ is affine, it is
\texttt{ContDiffOn} of every order; combine this with $h_K$ using the
product rule to obtain
\texttt{ContDiffOn}$(\R,2,\lambda\mapsto S_n(\lambda)K_n(\lambda),U_n)$,
and transfer that proof along $h_{FK}$ to obtain
$h_F:\operatorname{ContDiffOn}(\R,2,F_n,U_n)$.  The identity $h_{FK}$ also
identifies the first two derivatives at zero.  Since $0<S_n(0)$, the
product rule gives
\begin{equation}
 \frac{\operatorname{secondDeriv}(F_n,0)}{S_n(0)}
 =2\frac{\operatorname{deriv}(S_n)(0)}{S_n(0)}
   \operatorname{deriv}(K_n)(0)+\operatorname{secondDeriv}(K_n,0).
 \label{eq:derivation-second}
\end{equation}
Put $d_n:=\operatorname{blockScale}(n)$,
$N_{0,n}:=\operatorname{blockCount}(B_2,n,0)$, and
$N_{1,n}:=\operatorname{blockCount}(B_2,n,1)$. Only the second block moves,
and therefore
\[
 S_n(0)=N_{0,n}+d_nN_{1,n},\qquad
 \operatorname{deriv}(S_n)(0)=d_nN_{1,n}.
\]
Using the rounding bounds gives
\[
 0\leq\frac{\operatorname{deriv}(S_n)(0)}{S_n(0)}
 \leq
 \operatorname{Real.rpow}(d_n:\R,1-r)
 +\operatorname{Real.rpow}(d_n:\R,1-R),
\]
and the right-hand side tends to zero because $1<r<R$.

Apply \cref{prop:two-block-localisation} to $\mu_s$, the endpoint-null proof
rewritten through \cref{lem:signed-witness-bridge}, and an arbitrary
nonempty compact set (or the singleton $\{u\}$ for this proof).  Its upper
norm-free pair, together with
\cref{lem:compact-uniform-singleton}, gives
\[
 \operatorname{Tendsto}
  (n\mapsto\operatorname{deriv}(K_n)(0),
    \texttt{atTop},\mathcal N(A_r)),
\qquad
 \operatorname{Tendsto}
  (n\mapsto\operatorname{secondDeriv}(K_n,0),
    \texttt{atTop},\mathcal N(-A_r)).
\]
The kernel identities needed here are the $q=1,2$ clauses of the block
G-kernel integral bridge; the upper targets simplify at $u$ to
$A_r$ and $-A_r$.  The convergent first-derivative sequence is bounded.
Passing to the limit in \cref{eq:derivation-second} yields
\[
 \operatorname{Tendsto}\!\left(
  n\mapsto\frac{\operatorname{secondDeriv}(F_n,0)}{S_n(0)},
  \texttt{atTop},\mathcal N(-A_r)\right),
\qquad 0<-A_r.
\]
Consequently $\operatorname{secondDeriv}(F_n,0)>0$ eventually, because
$S_n(0)>0$.  On the other hand, the assumed cone concavity and
\cref{lem:cone-affine-line-bridge} make $F_n$ concave on $U_n$.  Apply
\cref{lem:fnd-line-curvature} using this concavity and the explicitly named
smoothness proof $h_F$ to make the same derivative nonpositive.  Choose
one index in the eventual set to obtain the contradiction.
\end{proof}

Combining \cref{prop:derivation-necessity} with
\cref{prop:derivation-sufficiency}, the monotone derivations are precisely the
positive barycentres of $H_{0,\alpha}$ over $\alpha\in[0,1]$. Under the
normalisation used for conditional entropies, the representing measure is a
probability measure.

\subsection{Final assembly and consistency checks}

The positive and negative temperate necessity results, together with
\cref{cor:finite-t-sufficiency}, prove the finite-$t$ part of the final
classification theorem. The tropical and derivation arguments prove
the remaining parts. The following point-mass specialisations provide useful
checks on every endpoint and inequality direction.

\begin{center}
\begin{tabular}{>{\raggedright\arraybackslash}p{0.17\textwidth}
                >{\raggedright\arraybackslash}p{0.28\textwidth}
                >{\raggedright\arraybackslash}p{0.38\textwidth}}
\toprule
Order $\alpha$ & Finite-$t$ condition for $\tau=\delta_\alpha$
& Limiting regimes \\
\midrule
$\alpha=0$
& every $t\neq0$
& $t=0$, $t=-\infty$, and $t=+\infty$ are all admissible \\
$0<\alpha<1$
& every $t<0$, and $0<t\leq(1-\alpha)/\alpha$
& $t=0$ and $t=-\infty$ are admissible; $t=+\infty$ is not \\
$\alpha=1$
& every $t<0$; no $t>0$
& $t=0$ and $t=-\infty$ are admissible \\
$1<\alpha<+\infty$
& $t\leq(1-\alpha)/\alpha$
& $t=-\infty$ is admissible; $t=0$ and $t=+\infty$ are not \\
$\alpha=+\infty$
& $t\leq-1$
& $t=-\infty$ is admissible \\
\bottomrule
\end{tabular}
\end{center}

For $0<\alpha<1$, the positive-branch inequality
\[
 \frac{\alpha}{1-\alpha}\leq\frac1t
\]
is equivalent to $0<t\leq(1-\alpha)/\alpha$. For
$1<\alpha<+\infty$, the exceptional negative-branch inequality
\[
 \frac{\alpha}{1-\alpha}\leq\frac1t,
 \qquad t<0,
\]
is equivalent to $t\leq(1-\alpha)/\alpha$. At $\alpha=+\infty$, the
convention $\omega(+\infty)=-1$ gives $t\leq-1$. The Shannon point belongs
to the all-lower negative branch, which explains why every $t<0$ is allowed
there, whereas the positive branch excludes an atom at $1$.

All parameter ranges in the final classification theorem have now been covered,
and the sufficient and necessary conditions agree at every finite and
limiting boundary.

\subsection{Necessity bundle in final Lean theorem form}

\begin{theorem}[All necessary parameter conditions]
\label{thm:necessity-bundle}
For universe $u$ and
$\tau:\operatorname{ProbabilityMeasure}(\overline\R_+)$, the following
single four-way conjunction holds:
\begin{align*}
&\Bigl(\forall(t:\R)(h_t:t\neq0),\quad
 \operatorname{CMMonotone}(\operatorname{HTemp}(t,h_t,\tau))
 \Longrightarrow\operatorname{DBulk}(t,\tau)\Bigr)\\
&\quad\wedge\Bigl(
 \operatorname{CMMonotone}(\operatorname{HMinus}(\tau))
 \Longrightarrow\operatorname{DMinus}(\tau)\Bigr)\\
&\quad\wedge\Bigl(
 \operatorname{CMMonotone}(\operatorname{HPlus}(\tau))
 \Longrightarrow\tau=\operatorname{diracProb}(0)\Bigr)\\
&\quad\wedge\Bigl(
 \operatorname{CMMonotone}
  (\operatorname{HZero}(\tau.\operatorname{toFiniteMeasure}))
 \Longrightarrow
 \operatorname{suppMeasure}(\operatorname{probMeasure}(\tau))
  \subseteq\operatorname{Set.Icc}(0,1)\Bigr).
\end{align*}
\end{theorem}

The corresponding Lean header is
\[
 \texttt{theorem necessityBundle.\{u\}
  (tau : ProbabilityMeasure Param) :}
\]
followed by this conjunction.

\begin{proof}
Fix $t$ and split on its sign using $h_t$.  If
$h_{\rm pos}:0<t$, apply
\cref{cor:positive-temperate-probability-necessity} and inject
\texttt{PosAdm t tau} into the first disjunct of \texttt{DBulk}.  If
$h_{\rm neg}:t<0$, apply
\cref{cor:negative-temperate-probability-necessity}; its lower-support
alternative injects into the second disjunct, and its exceptional witness
injects into the third.  Proof irrelevance identifies the sign-derived
nonzero proof with $h_t$ in \texttt{HTemp}.

The second and third conjuncts are
\cref{prop:negative-tropical-necessity,prop:positive-tropical-necessity}.
For the fourth, global temperate shape reduction at the derivation branch
makes \texttt{derivationColumn} concave, and
\cref{prop:derivation-necessity} applies.
\end{proof}

\begin{theorem}[Complete parameter classification]
\label{thm:main-classification}
At universe $u$, for every
$\tau:\operatorname{ProbabilityMeasure}(\overline\R_+)$,
\[
 \texttt{Classifies.\{u\}}(\tau).
\]
\end{theorem}

The actual header is
\[
 \texttt{theorem mainClassification.\{u\}
  (tau : ProbabilityMeasure Param) : Classifies.\{u\} tau :=}.
\]

\begin{proof}
Finite-$t$ sufficiency is \cref{cor:finite-t-sufficiency}; its necessity is
the first conjunct of \cref{thm:necessity-bundle}.  The two tropical
sufficiency directions are
\cref{prop:negative-tropical-sufficiency,prop:positive-tropical-sufficiency},
and the converse directions are the second and third conjuncts of
\cref{thm:necessity-bundle}.  Derivation sufficiency is
\cref{prop:derivation-sufficiency}, and the fourth conjunct of
\cref{thm:necessity-bundle} is its converse.  These are exactly the four
fields of \texttt{Classifies.\{u\}}.
\end{proof}

\subsection{Lift to the embedded relation and the exact axiom bundle}

\begin{definition}[Conditional-entropy axiom bundle]
\label{def:conditional-entropy-axioms}
For $F:\texttt{PolyJointFunctional.\{u\}}$, define separately
\begin{align*}
 \operatorname{NonnegativeJointFunctional}(F):\Longleftrightarrow{}
 &\forall\{X,Y:\texttt{Type u}\}
 [\operatorname{Fintype}X][\operatorname{Nonempty}X]
 [\operatorname{Fintype}Y][\operatorname{Nonempty}Y],\
 \forall P:\operatorname{JointProb}(X,Y),\quad0\leq F(P),\\
 \operatorname{FairBitNormalized}(F):\Longleftrightarrow{}
 &\forall\{Y:\texttt{Type u}\}
 [\operatorname{Fintype}Y][\operatorname{Nonempty}Y],\
 \forall q:\operatorname{ProbVec}(Y),\\[-2pt]
 &F(\operatorname{independentJoint}
    (\operatorname{uniformProb}_{\operatorname{Fin}(2)},q))=\log2,\\
 \operatorname{CEmbedsMonotone}(F):\Longleftrightarrow{}
 &\forall\{X,Y,X',Y':\texttt{Type u}\}
 [\operatorname{Fintype}X][\operatorname{Nonempty}X]
 [\operatorname{Fintype}Y][\operatorname{Nonempty}Y]\\[-2pt]
 &\hspace{5mm}
 [\operatorname{Fintype}X'][\operatorname{Nonempty}X']
 [\operatorname{Fintype}Y'][\operatorname{Nonempty}Y'],\
 \forall P:\operatorname{JointProb}(X,Y),
 \forall Q:\operatorname{JointProb}(X',Y'),\\[-2pt]
 &\hspace{35mm}\operatorname{CEmbeds}(P,Q)\to F(P)\leq F(Q).
\end{align*}
Finally define the single proposition
\begin{align*}
 \operatorname{ConditionalEntropyAxioms}(F):\Longleftrightarrow{}
 &\operatorname{NonnegativeJointFunctional}(F)\ 
 \wedge\operatorname{JointEmbeddingInvariant}(F)\ 
 \wedge\operatorname{JointTensorAdditive}(F)\\[-2pt]
 &\wedge\operatorname{FairBitNormalized}(F)\ 
 \wedge\operatorname{CEmbedsMonotone}(F).
\end{align*}
\end{definition}

\begin{lemma}[Embedding lift of fixed-row monotonicity]
\label{lem:embedding-lift}
For $F:\texttt{PolyJointFunctional.\{u\}}$,
\[
 \operatorname{JointEmbeddingInvariant}(F)
 \Longrightarrow
 \operatorname{CMMonotone}(F)
 \Longrightarrow
 \operatorname{CEmbedsMonotone}(F).
\]
\end{lemma}

\begin{proof}
Unpack \cref{eq:foundation-embedded-preorder}.  It supplies finite common
alphabets, four injections, and a conditional mixing datum from
$\widetilde P:=E_{e_X,e_Y}P$ to
$\widetilde Q:=E_{e'_X,e'_Y}Q$.  Zero-extension invariance gives
$F(P)=F(\widetilde P)$ and $F(Q)=F(\widetilde Q)$; fixed-row monotonicity
gives $F(\widetilde P)\leq F(\widetilde Q)$.
\end{proof}

\begin{corollary}[Admissible candidates satisfy the complete axiom bundle]
\label{cor:all-conditional-entropy-axioms}
At universe $u$, the following single four-way conjunction holds:
\begin{align*}
&\Bigl(\forall(t:\R)
 (\tau:\operatorname{ProbabilityMeasure}(\overline\R_+)),\quad
 \forall h:\operatorname{DBulk}(t,\tau),\quad
 \operatorname{ConditionalEntropyAxioms}
  (\operatorname{HTemp}(t,\operatorname{DBulk.ne}(h),\tau))\Bigr)\\
&\quad\wedge\Bigl(
 \forall\tau:\operatorname{ProbabilityMeasure}(\overline\R_+),\quad
 \operatorname{DMinus}(\tau)\to
 \operatorname{ConditionalEntropyAxioms}(\operatorname{HMinus}(\tau))
 \Bigr)\\
&\quad\wedge
 \operatorname{ConditionalEntropyAxioms}
  (\operatorname{HPlus}(\operatorname{diracProb}(0)))\\
&\quad\wedge\Bigl(
 \forall\tau:\operatorname{ProbabilityMeasure}(\overline\R_+),\quad
 \operatorname{suppMeasure}(\operatorname{probMeasure}(\tau))
   \subseteq\operatorname{Set.Icc}(0,1)\to\\[-2pt]
&\hspace{32mm}
 \operatorname{ConditionalEntropyAxioms}
  (\operatorname{HZero}(\tau.\operatorname{toFiniteMeasure}))\Bigr).
\end{align*}
\end{corollary}

The corresponding declaration begins
\[
 \texttt{corollary allConditionalEntropyAxioms.\{u\} :}
\]
and continues with the displayed conjunction.

\begin{proof}
Well-definedness and nonnegativity are
\cref{lem:fnd-conditional-well-defined}.  Invariance, tensor additivity, and
normalization are \cref{prop:fnd-easy-axioms}.  Fixed-row monotonicity follows
from
\cref{cor:finite-t-sufficiency,prop:negative-tropical-sufficiency,prop:positive-tropical-sufficiency,prop:derivation-sufficiency};
apply \cref{lem:embedding-lift}.
\end{proof}

By \cref{thm:main-classification}, no parameter outside these sets satisfies
the required \texttt{CMMonotone} conjunct.

\subsection{Lean module manifest}

The following modules form an acyclic implementation plan.  The labels in
parentheses are the principal exported declarations.
\begin{enumerate}[label=\texttt{M\arabic*.}]
 \item Finite carriers, cones, masses, normalization, columns, and channel
 algebra
 (\cref{lem:fnd-nonzero-mass,lem:fnd-normalization-package,lem:fnd-conditional-channel-stochastic}).
 \item The compactified order, finite measures, Jordan integration, support,
 pushforwards, and one-sided moments
 (\cref{lem:fnd-support-complement,lem:fnd-pushforward-integration,lem:fnd-moment-algebra}).
 \item Total R\'enyi entropy, endpoint continuity, bounds, tensor and
 embedding identities, Schur concavity, and lower-order concavity
 (\cref{prop:fnd-renyi-order-continuity,lem:fnd-renyi-schur,lem:fnd-renyi-concavity}).
 \item Conditional candidates, easy axioms, column functions, finite mixing,
 perspectives, curvature closure, and shape reductions
 (\cref{prop:fnd-easy-axioms,lem:fnd-perspective,prop:temperate-shape-reduction,prop:tropical-shape-reduction}).
 \item Power means, monomials, finite-support factorization, endpoint
 approximation, and finite-range pushforwards
 (\cref{lem:power-curvature,lem:monomial-curvature,prop:finite-sufficiency,prop:general-sufficiency}).
 \item Tropical and derivation sufficiency
 (\cref{prop:negative-tropical-sufficiency,prop:positive-tropical-sufficiency,prop:derivation-sufficiency}).
 \item Entropy derivatives, differentiation under the parameter integral,
 atoms, and null thresholds
 (\cref{lem:exact-derivatives,lem:differentiate-integral,lem:null-thresholds}).
 \item Block counts, escort bounds, uniform integration, and concrete
 two-/three-block localizations
 (\cref{lem:dominant-block,prop:block-limit-passage,prop:two-block-localisation,prop:three-block-localisation}).
 \item Dedicated Shannon estimates and localization
 (\cref{lem:shannon-dedicated-escort,lem:uniform-shannon-expansion,prop:shannon-localisation}).
 \item Curvature obstruction logic and the five necessity branches
 (\cref{lem:log-curvature-obstruction,lem:stationarity-correction,thm:necessity-bundle}).
 \item Final equivalence and lift to the main body
 (\cref{thm:main-classification,lem:embedding-lift,cor:all-conditional-entropy-axioms}).
\end{enumerate}

\end{document}
